\documentclass[11pt,a4paper]{article}
\usepackage[utf8]{inputenc}
\usepackage[T1]{fontenc}
\usepackage{lmodern,amsmath,amssymb,textcomp}
\usepackage[margin=24mm]{geometry}
\usepackage{longtable,booktabs,array,enumitem}
\usepackage[hidelinks]{hyperref}
\setlength{\parindent}{0pt}
\setlength{\parskip}{6pt}
\emergencystretch=3em
\begin{document}
{\LARGE\bfseries Exact constructions, quantitative bounds and formal proofs from OpenMath 2026\par}\medskip
{\small Original-result working anthology with source-faithful scope and verification records
Editorial compilation by Alejandro Zarzuelo Urdiales | 5 October 2026\par}\medskip
{\small Judging and evaluation record: the submissions discussed in this document have been judged and evaluated by Alejandro Zarzuelo Urdiales for OpenMath 2026. This dated review records the stated verification, classification and unresolved holds; it is a preliminary evaluation rather than a signed award decision.\par}

{\small Event provenance: OpenMath 2026 ran 27 September-2 October 2026 with worldwide online participation. Detailed organizer listings specify Harvard Innovation Labs for the 27 September in-person event; the OpenMath programme advertises a hybrid kickoff at MIT. Sundai identifies its Harvard/MIT community. Organizer sources: \url{https://rsihouse.ai/openmath} ; \url{https://luma.com/yzp9abvr} ; \url{https://www.sundai.club/events/boston/recursive-learning-hack-with-harvard-innovation-labs}\par}

The mathematical authors are credited per chapter. This anthology is a working manuscript for contributor and specialist review. It records 19 proposed original-result families; unresolved numerical-priority candidates have separate manuscripts and are excluded from its theorem chapters.

Central preserved source archive: \url{https://github.com/alejandrozu/openmath-2026-judging}

\clearpage
\section{Contents}
\tableofcontents
\section{Abstract}
This working anthology explains the proposed original mathematical contributions in the OpenMath 2026 submission archive. It covers exact answers to sequence and subsidiary extremal questions, quantitative number-theory bounds, structural combinatorics, a sharp singular-likelihood classification, an explicit sparse matrix-multiplication scheme and certified graph/geometry constructions. For every family it records the exact mathematical scope, construction or proof mechanism, relation to prior work and formal evidence. Correctness, originality, competition placement and manuscript readiness are evaluated separately. The two older weighted Ramsey manuscripts, the restricted Grothendieck manuscript and the 61-line priority manuscript are stored outside the original-results anthology. The latest Luke/Madhan structural Ramsey packet and Matt generic pivot packet are included as separately dated original-result candidates, with their acceptance and verification conditions retained. The collection is prepared for author and specialist review; it is not evidence of a refereed publication or a signed final jury decision.

\section{Editorial provenance and author attribution}
The frozen judging archive is alejandrozu/openmath-2026-judging at revision d0ed487f487fa7ec814ff705ab55249983fe8707. It preserves source packets, contribution descriptions, manifests and dated evidence. The Sakana bundle SHA 256 is f22e456eecf1a5c071aafa0666fa517dc9ebbd102f483a963f694ee64748f769. HTPeo is pinned to76c63b8e425e551a930060a6871b84fe3b4780db, Chandragupt to479c2416b6261ee841052eef4881df0e40054f3f, and Rohith to f462d8e18aea2a458376c523c9b6c2980237071f.

The mathematical contribution authors are recorded per chapter. Rachel Teo, Wenyi Wang and Hamdi Abderrahmene share the Sakana packet; Woohyuk Kang, Hyunjin Lee, Sanghyeon Lee and Youchan Oh share HTPeo; Chandragupt Sharma and Rohith Poola have individual packets. These packet authorship records do not establish a contribution-by-contribution percentage split. Verified affiliations appear in the shared author register; missing details remain unconfirmed. Organizer/editor status must be described separately from mathematical authorship, and an anthology-wide author list requires the contributors' approval.

This text is a source-grounded editorial synthesis. It does not silently attribute reviewer-derived observations, post-event repairs or older source mathematics to entrants. Shared foundational lemmas and known methods are credited in the chapters. The final release should carry a contribution statement and an AI-assistance disclosure reflecting the teams' verified workflow, without inventing an allocation of human versus automated proof work.

\section{What is meant by an original result}
A result may be a complete answer to a registered question, a new infinite family, a stronger quantitative bound, a finite obstruction to a specific strategy, or a useful new explicit construction. Those are distinct scientific scopes. A result that settles an explicit subsidiary question is not thereby a solution of the entire parent conjecture. A formal proof of known mathematics can be an important contribution to verification infrastructure, but it belongs in the accompanying known-mathematics account when no original mathematical delta is established.

Originality here is a bounded source comparison, not a certificate that all historical literature has been exhausted. Named unresolved prior sources are preserved in the chapters. The January 2026 Technion Ramsey announcement is numerically stronger than the two older submitted weighted constructions, which therefore cannot be advertised as numerical records. The latest Luke/Madhan exact bound is below the announced threshold, but its comparison with the unpublished exact earlier bound remains unresolved; its structural contribution must be assessed separately. The 61-line Kobon candidate faces an unresolved straight-line versus pseudoline recursion comparison. The three-row Grothendieck certificate is valid but does not improve the universal constant and is weaker than an available restricted comparison.

The competition's approved October 4 amendments govern placement: late delivery is allowed, the frozen difficulty proposals are retained, bonuses are removed, independently achieved genuinely new contributions may count even if superseded, and Company/Non-company are the two classes. The user retains progress p=.15 for HTPeo's DMS family and requests a mention for Leanification. None of those administrative decisions proves historical novelty or substitutes for exact theorem verification.

\section{Verification methodology and its limits}
For each family, verification has three layers. First, source fidelity asks whether the formal statement represents the actual mathematical problem, with the same count, graph model, positivity, coordinate ordering and boundary convention. Second, proof closure asks whether the selected endpoint and every load-bearing own dependency compile in the specified environment without unresolved proof assumptions. Third, independent exact or ordinary checks exercise finite certificates and analytic arguments from a separate implementation.

Archived receipts are preserved evidence of earlier runs. Fresh execution in this publication pass must be identified by its own date, command, compiler, dependency pins, source hashes and endpoint axiom list. An alternate-version core compilation is useful corroboration but must not be labelled a full pinned-project replay. A no-sorry text scan is insufficient: an upstream conjecture may be irrelevant to the selected endpoint, while a hidden trusted dependency can matter even when the final file looks clean. The common verification register records these distinctions.

The independent audit reproduced all recorded score formulas and several substantial exact witnesses: weighted Ramsey counts, rational matrix tensor identities and its parameter polynomial, Kobon face counts, paired matrix ranks, graph-colouring instances and selected finite certificates. Sample computations do not prove their corresponding infinite theorems. Conversely a closed Lean theorem can remain historically known or mathematically narrow. Mathematical reviewers should inspect definitions, proof scope, prior sources and attribution in addition to compilation.

Human-readable proofs and source-derived expansions explain the mechanisms actually present in the submissions. Across the eighteen base families and the separately incorporated Matt generic-pivot family, all nineteen main original candidates now have complete ordinary arguments at their exact displayed scopes. Nine family arguments use explicit finite certificates. The quartic chapter includes quantitative analytic appendices proving the complex-log remainder, positive Gaussian core, modulus-deficit estimates, finite windows and uniform exponent assembly. The DMS GP(n,k) construction for 1\ensuremath{\leq}k\ensuremath{\leq}15 includes exact period/seam/exception tables, the finite-window reduction and an independent checker. Its broader matching-preserving two-/three-cut seam lemmas, conditional inductions and universal host/marking census bridge are also explained from the delivered proof. BASE12, SIMPLE14 and B14D are discharged by that delivered chain; the five infinite structural hypotheses remain open. Matt's ordinary proof closes only the selected generic matrix/eigenphase theorem, while the circuit-node construction and universal CNOT claim remain unproved. Ordinary-proof completeness does not certify historical originality, fresh formal verification or contributor approval; those are reported separately.

\section{Scientific organization of the portfolio}
The collection naturally divides into four mathematical strands. The arithmetic strand contains subset-product interpolation gaps, prime-reversal nonperiodicity, the three-prime residue divisibility exclusion, weighted fibre counts, cube-sum representation bounds and an explicit reciprocal-sum construction. Its reusable themes are exact integer encoding, finite certificates after global reductions, and weighted counting.

The combinatorial strand contains the additive-three hypergraph cover bound, Ferrers factor-preserving hierarchy, finite critical-graph order obstructions, half-turn defect and the diagonal-pullback obstruction. These contributions depend on quantifiers that should be visible in ordinary prose: every supplied factor, every order, every diagonal pullback, or exactly the stated finite graph class.

The analytic/statistical strand contains the Hermitian quartic's exact exponent semigroup and logarithmic boundary scale, and the parity RBM's sharp collapsed-fibre likelihood threshold. Both require higher-order or boundary analysis; a finite numerical fit or a second-derivative calculation would not establish their strongest endpoints.

The construction strand contains the sparse 23-product matrix scheme, HTPeo's star-colouring families and exact 39-line Kobon arrangement. Their machine-readable finite data are scientifically useful only together with the bridges to the intended tensor, graph or geometry. The one-leg rigidity note complements the matrix scheme by delimiting one search strategy without contradicting a joint factor change.

\section{Publication and reuse}
This anthology is intended as a comprehensive technical reference, while individual family manuscripts provide natural units for specialist review. Closely related strengthenings remain in the same paper: the exact eight-level subset-product fibre belongs with the interpolation refutation; all quartic asymptotics and exceptional exponents form one family; the graph order inequalities belong together; gauges and the parameter family support the 138 scheme.

The supplementary judging repository should provide one stable index of source revisions, exact statements, proof/checking commands and data. Public release should follow author approval, affiliation confirmation, appropriate source/licence attribution and the explicit novelty holds. A DOI or archival version should identify which results have closed proofs and which are editorial candidates. The short Wolfram-facing document can then point to this collection and the individual papers without requiring him to reconstruct the competition correspondence.

\clearpage
\section{Arbitrarily long prime-exponent gaps in complete-interval subset products}
{\small Provisional mathematical authors: Rachel Teo, Wenyi Wang, Hamdi Abderrahmene\par}\medskip
{\small Rachel Teo  -  Sakana AI.\par}

{\small Wenyi Wang  -  Sakana AI.\par}

{\small Hamdi Abderrahmene  -  Sakana AI.\par}

{\small Affiliations follow the recorded author declarations or public institutional/profile sources. Preferred author names, author order and publication affiliations await author review. This editorial compilation does not establish anthology coauthorship.\par}

\subsection{Abstract}
The distinct-subset-product interpolation question associated with OEIS A060957 is refuted. More strongly, for every positive gap length and every prescribed lower bound on a prime, the construction realizes an exact fibre consisting of two blocks of four attainable exponents separated by the prescribed missing block.

\subsection{The original interpolation problem}
Let P(n) be the set of products of subsets of \{1,...,n\}. A factor may be chosen once or omitted: this is not the unrestricted multiplicative monoid generated by that interval. The question asks whether, when m and p\textasciicircum{}a m are attainable for a prime p\ensuremath{\leq}n, all the intermediate p-exponent levels are attainable. Its complete-interval condition matters. A hole in a hand-picked collection of factors would not answer the question, because the omitted interval factors could fill that hole.

The submitted result handles all factors of the interval simultaneously. Its basic endpoint supplies m, p and n with m and p\textasciicircum{}5m in P(n), but p\textasciicircum{}4m outside P(n). The stronger endpoint specifies the entire exponent fibre, rather than just proving that one level is missing. For every g\ensuremath{\geq}1 and lower bound B, there are n,p,d,h with p prime, B<p\ensuremath{\leq}n, d>0 and p not dividing d, for which p\textasciicircum{}e d is attainable precisely at e=h,h+1,h+2,h+3,h+g+4,h+g+5,h+g+6,h+g+7. The g omitted levels are h+4 through h+g+3.

\subsection{Encoding and the exposed-face construction}
The proof represents integer products by their prime-exponent vectors, separating the distinguished p coordinate from all other coordinates. A graph-based atom construction controls the possible exponent increments. In the generalized construction, a circulant system decomposes the complete graph on an odd number of vertices into oriented spanning two-factors. Every vertex sees every other vertex once through the shifts and their inverses. This identity is proved from the circulant model; it is not left as an existence hypothesis in the final theorem.

Integral feasibility conditions in the height-one slice classify the allowable atoms. Summing inequalities by colour gives a degree-two budget; summing by vertex gives the complementary bound. These estimates constrain integral solutions to the two desired exponent blocks. The argument is algebraic and combinatorial rather than a search through arbitrary integer factors.

To transport that local construction to the complete interval, the proof chooses prime coordinates and an isolating rational linear functional. The positive-weight factors form a forced baseline W. Any subset product with the same p-free exponent vector must contain these positive-weight factors and can differ only through zero-weight atoms in the prescribed subspace. Consequently unrelated factors from \{1,...,n\} cannot silently create extra attainable levels. The complete-face theorem proves both directions of this classification. Finally, dividing the baseline by its p-primary factor produces d not divisible by p and an absolute exponent offset h.

\subsection{Formal target and proof obligations}
The important proof boundary is the passage from a model fibre to products of actual distinct interval factors. The selected Main and Generic.Final endpoints close that boundary through the complete-face, product-support and normalization modules. This publication pass independently replayed all 45 exact frozen authored source modules and the selected audit under Lean 4.33.1 and Mathlib 0df444a360eaa60ab8c11dca51a86af692955474, finishing on 5 October 2026 at 03:10:13 UTC. All four requested prints passed: the explicit subset-product counterexample, falsity of interpolation, exact arbitrary-gap eight-level fibre, and the consecutive-above corollary each depend only on propext, Classical.choice and Quot.sound. No selected endpoint uses admitted, native-evaluation or unknown axioms. Archived dual-version checks remain distinct from this fresh 45-source replay. Receipt: Verification/sakana-a060957-fresh-build.json.

The ordinary-proof expansion now gives the complete atom classification, the choice of prime parameters and separating linear functional, and the bridge preventing the remaining factors of the complete interval from filling the gap. The formal proof supplies a precise supplementary specification alongside that human-readable argument. No minimum counterexample size is asserted, and the theorem does not provide one fixed prime that works for every gap length.

\subsection{Prior work and significance}
Prime-coordinate encoding, exposed faces and semigroup holes are established methods. The proposed new contribution is their use to refute the complete-interval interpolation statement and to prescribe gaps of arbitrary length with an exact eight-level fibre. The closest recorded sources did not supply that theorem. This is a bounded priority comparison, so the paper should avoid an unqualified first-discovery claim until the author and referee checks are complete.

This family is a strong standalone number-theory/combinatorics preprint candidate. The qualitative counterexample answers the original question, while the arbitrary-gap theorem supplies the more durable mathematical story. The two statements should appear in one paper and receive one contribution-family treatment.

\subsection{An explicit finite-dimensional gadget}
Fix m\ensuremath{\geq}2 and put v=2m+1. Take the vertex set Z/vZ and the m shifts \ensuremath{\sigma}\_i(w)=w+i for i=1,...,m. The neighbours \ensuremath{\sigma}\_i(w) and \ensuremath{\sigma}\_i\ensuremath{^{-1}}(w), as i varies, are exactly the v\ensuremath{-}1 vertices different from w. Thus, for every function y on the vertices with sum S, the sum of y over these successors and predecessors is S\ensuremath{-}y\_w. This elementary circulant identity is the only graph existence input.

Use one coordinate for each oriented edge (i,w), one link coordinate for each vertex w, and one height coordinate. Define the linear map L by edge coordinate x\_i+y\_w+y\_\{\ensuremath{\sigma}\_i(w)\}, link coordinate y\_w, and height X+S, where X=\ensuremath{\Sigma}\_i x\_i and S=\ensuremath{\Sigma}\_w y\_w. The atom space is the rational range of L. Its integer points will be constrained by coordinate nonnegativity. The sum of all edge and link coordinates is v(X+S): each x\_i occurs v times, each y\_w occurs 2m times on edges and once as a link.

There are m cycle atoms A\_i=L(e\_i,0), v vertex atoms B\_w=L(0,e\_w), and one links atom C=L(\ensuremath{-}2\ensuremath{\cdot}1,1). Every atom has height one and nonnegative integer coordinates. The cycle atom has ones on the edges of one colour; the vertex atom has ones on its incident edges and its own link; C has ones on all links and zeros on the edges. The identity \ensuremath{\Sigma}\_w B\_w=2\ensuremath{\Sigma}\_i A\_i+C follows coordinate by coordinate. This is the balance relation from which the exponent gap will arise.

\subsection{Lemma: classification of nonnegative integral height-one points}
Suppose L(x,y) is integral, nonnegative and has height one. The link coordinates show that all y\_w are integers. An edge coordinate then shows that every x\_i is an integer. Thus y\_w\ensuremath{\geq}0, X+S=1, and x\_i+y\_w+y\_\{\ensuremath{\sigma}\_i(w)\}\ensuremath{\geq}0 for every i,w. Summing the last inequalities over w gives v x\_i+2S\ensuremath{\geq}0. Summing these colour inequalities gives vX+2mS=v\ensuremath{-}S\ensuremath{\geq}0, so S\ensuremath{\leq}v.

There is a second useful estimate. Fix w and sum the inequalities on both edges incident to w in every colour. Their total is 2X+2m y\_w+(S\ensuremath{-}y\_w)=2\ensuremath{-}S+(2m\ensuremath{-}1)y\_w\ensuremath{\geq}0. Consequently S\ensuremath{\leq}2+(2m\ensuremath{-}1)y\_w. If some y\_w=0, then S\ensuremath{\leq}2. If an x\_i were negative, integrality would give x\_i\ensuremath{\leq}\ensuremath{-}1, and hence v x\_i+2S\ensuremath{\leq}\ensuremath{-}v+4<0 because v\ensuremath{\geq}5. Therefore all x\_i are nonnegative. Since X+S=1 and all variables are now nonnegative integers, exactly one of them is one and all the others are zero. This gives a cycle or vertex atom.

If no y\_w is zero, then y\_w\ensuremath{\geq}1 for all w, hence S\ensuremath{\geq}v. Together with S\ensuremath{\leq}v this yields y\_w=1 for every w. The edge inequalities give x\_i\ensuremath{\geq}\ensuremath{-}2, while X=1\ensuremath{-}v=\ensuremath{-}2m. There are m such variables, so each is \ensuremath{-}2. This gives C. We have therefore proved that the only integral, nonnegative, height-one points in the atom space are the m+v+1 atoms just described. At height zero, nonnegativity and the coordinate-sum identity force every coordinate to vanish. These are complete classifications, not assumptions about a chosen generating set.

\subsection{Lemma: realizing exactly the allowed integer factors}
Assign pairwise distinct small primes to the edge coordinates. Let L\ensuremath{_0} be their product. Assign distinct link primes, each larger than L\ensuremath{_0}, and let C\ensuremath{_0} be their product. In particular L\ensuremath{_0}<C\ensuremath{_0}. Choose a prime p above both the prescribed lower bound and 2L\ensuremath{_0}\ensuremath{^2}+C\ensuremath{_0}L\ensuremath{_0}. Bertrand's postulate supplies a prime q with pL\ensuremath{_0}<q\ensuremath{\leq}2pL\ensuremath{_0}. Use q as the height prime and set n=pqL\ensuremath{_0}. The inequalities imply p\ensuremath{^2}<n<p\ensuremath{^3} and n<q\ensuremath{^2}; also all coordinate primes differ from p and from q.

For a cycle atom let a\_i be the product of its edge primes. Then 1\ensuremath{\leq}a\_i\ensuremath{\leq}L\ensuremath{_0}, so both qa\_i and pqa\_i lie in \{1,...,n\}, while p\ensuremath{^2}qa\_i>n. For a vertex atom let b\_w be its link prime times the product of its incident edge primes. Then L\ensuremath{_0}<b\_w\ensuremath{\leq}C\ensuremath{_0}L\ensuremath{_0}<p. Hence qb\_w\ensuremath{\leq}n but pqb\_w>n. The links atom similarly corresponds to qC\ensuremath{_0}\ensuremath{\leq}n, and pqC\ensuremath{_0}>n. Finally, p\ensuremath{^2}<n<p\ensuremath{^3} means that the allowable pure p-powers in the interval are 1,p,p\ensuremath{^2}.

We claim that these are every interval factor whose p-free exponent vector lies in the atom space. For such a factor j\ensuremath{\leq}n, its exponent of q is at most one because n<q\ensuremath{^2}. If that exponent is zero, the height-zero classification forces all its p-free coordinates to vanish; unique prime factorization makes j a pure p-power, hence 1,p or p\ensuremath{^2}. If the q exponent is one, the height-one classification makes its p-free part one of the cycle, vertex or links bases. The preceding upper bounds allow exactly the two p-variants of a cycle and the single variant of each vertex and the links atom. An exponent vector in the embedded atom space has no prime coordinates outside the chosen ones, so no extra prime factor has been overlooked. Unique factorization also proves that every displayed integer is distinct: equal numbers would have the same atom and p exponent.

The full allowed set is therefore \{1,p,p\ensuremath{^2}\}, the pairs qa\_i,pqa\_i, the v vertex numbers qb\_w, and qC\ensuremath{_0}. This is the arithmetic bridge from a rational linear construction to actual distinct factors in the complete interval.

\subsection{Lemma: the exact two-block exponent support}
Let T=\ensuremath{\Sigma}\_w B\_w be the target p-free vector, and let t\ensuremath{_0} be its corresponding positive integer with no p factor. To make a selected product have p-free vector T, write d\_w\ensuremath{\in}\{0,1\} for the selection of the vertex atom B\_w, c\ensuremath{\in}\{0,1\} for C, and a\_i,b\_i\ensuremath{\in}\{0,1\} for the two cycle numbers qa\_i and pqa\_i. The link coordinates give d\_w+c=1 for every w. The edge coordinates give a\_i+b\_i+d\_w+d\_\{\ensuremath{\sigma}\_i(w)\}=2.

If c=0, then every d\_w=1 and every a\_i=b\_i=0. If c=1, then every d\_w=0 and every a\_i=b\_i=1. Thus there are exactly two structural ways to make the p-free target: choose all vertices, or choose C and both numbers of each cycle. The latter contributes exactly m extra p factors. Independently choosing p and p\ensuremath{^2} contributes any exponent in \{0,1,2,3\}; choosing 1 has no effect on the product. Hence a subset of the allowed factors has product p\textasciicircum{}e t\ensuremath{_0} exactly when e\ensuremath{\in}\{0,1,2,3\}\ensuremath{\cup}\{m,m+1,m+2,m+3\}. Every exponent on the displayed right side is realized by the indicated choices. The equations exclude every other exponent, even when the subset omits or includes the factor 1.

\subsection{Lemma: embedding the gadget in all of \{1,...,n\}}
The preceding calculation still concerns a particular collection of factors. To prevent the other interval integers from filling a gap, let F(j) be the p-free prime-exponent vector of j. Work in the finite-dimensional rational vector space spanned by F(1),...,F(n), and the atom space A. There is a rational linear functional \ensuremath{\ell} vanishing on A and nonzero on every F(j) outside A. Indeed, each outside vector excludes one proper hyperplane of the annihilator of A; a finite union of proper rational hyperplanes cannot fill that annihilator. Choosing outside their union gives \ensuremath{\ell}.

Give interval factor j the rational weight w\_j=\ensuremath{\ell}(F(j)), and let W be the set of factors with positive weight. All gadget factors have weight zero. Define the baseline M=(\ensuremath{\prod}\_\{j\ensuremath{\in}W\}j)t\ensuremath{_0}. Suppose an arbitrary interval subset U has product with the same p-free vector as M. Applying \ensuremath{\ell} gives \ensuremath{\Sigma}\_\{j\ensuremath{\in}U\}w\_j=\ensuremath{\Sigma}\_\{j\ensuremath{\in}W\}w\_j. The right side is the largest possible sum over an interval subset: it includes every positive term and no negative term. Equality therefore forces every positive factor into U and excludes every negative factor. The remaining factors have weight zero, hence have p-free vector in A by the choice of \ensuremath{\ell}, and have already been classified as gadget factors.

Consequently U is W together with a subset of the gadget. Cancelling the baseline factors reduces its product to the exact two-block support proved above. Conversely each supported gadget subset is disjoint from W and can be adjoined to W. Thus, for every positive integer x with F(x)=F(M), membership in the complete-interval subset-product set is equivalent to x=p\textasciicircum{}e M for e\ensuremath{\in}\{0,1,2,3\}\ensuremath{\cup}\{m,m+1,m+2,m+3\}. This is an exact fibre statement over the whole interval.

Finally write M=p\textasciicircum{}h d with p\ensuremath{\not\mid}d. Unique factorization gives an arbitrary x in this fibre the form p\textasciicircum{}k d. The exact attainable exponents are h through h+3 and h+m through h+m+3. Given a prescribed positive gap length g, choose m=g+4. The g exponents h+4,...,h+g+3 are missing. The construction allows p above any prescribed bound. In particular, taking g=1 supplies two attainable endpoints surrounding a missing interpolation level. No assertion about a smallest counterexample or a uniform fixed prime is needed.

\subsection{Verification and reproducibility}
Fresh execution: sakana-a 060957: PASS; 45/45 custom modules compiled successfully; receipt Verification/sakana-a060957-fresh-build.json. Compiler pins, source hashes and endpoint axiom outputs are in Verification. Historical receipts and finite checks do not replace fresh replay.

\begin{samepage}
\subsection{SubsetProductGap.subset\_product\_counterexample}
\begin{quote}\small\ttfamily theorem subset\_product\_counterexample : \ensuremath{\exists} n p m : \ensuremath{\mathbb{N}}, p.Prime \ensuremath{\wedge} p \ensuremath{\leq} n \ensuremath{\wedge} 0 < m \ensuremath{\wedge}\newline     m \ensuremath{\in} productsOfSubsets n \ensuremath{\wedge} p\textasciicircum{}5*m \ensuremath{\in} productsOfSubsets n \ensuremath{\wedge} p\textasciicircum{}4*m \ensuremath{\notin} productsOfSubsets n\end{quote}
{\small Source: proofs/a060957/OpHack/SubsetProducts/Main.lean:89.\par}

{\small Source SHA-256: 5b9d85abd44f036bbbfa11f241b62ebe8dce01a2ffd89305caa965669c72b0fc\par}

\end{samepage}
\begin{samepage}
\subsection{SubsetProductGap.interpolation\_conjecture\_false}
\begin{quote}\small\ttfamily theorem interpolation\_conjecture\_false : \ensuremath{\neg} (\ensuremath{\forall} (n p : \ensuremath{\mathbb{N}}), p.Prime \ensuremath{\to} p \ensuremath{\leq} n \ensuremath{\to}\newline     \ensuremath{\forall} (m a : \ensuremath{\mathbb{N}}), m \ensuremath{\in} productsOfSubsets n \ensuremath{\to} p\textasciicircum{}a*m \ensuremath{\in} productsOfSubsets n \ensuremath{\to}\newline     \ensuremath{\forall} k : \ensuremath{\mathbb{N}}, 0 < k \ensuremath{\to} k < a \ensuremath{\to} p\textasciicircum{}k*m \ensuremath{\in} productsOfSubsets n)\end{quote}
{\small Source: proofs/a060957/OpHack/SubsetProducts/Main.lean:97.\par}

{\small Source SHA-256: 5b9d85abd44f036bbbfa11f241b62ebe8dce01a2ffd89305caa965669c72b0fc\par}

\end{samepage}
\begin{samepage}
\subsection{SubsetProductGap.Generic.arbitrary\_gap\_eight\_levels\_above}
\begin{quote}\small\ttfamily theorem arbitrary\_gap\_eight\_levels\_above (g : \ensuremath{\mathbb{N}}) (hg : 1 \ensuremath{\leq} g) (lower : \ensuremath{\mathbb{N}}) :\newline     \ensuremath{\exists} n p d h : \ensuremath{\mathbb{N}}, p.Prime \ensuremath{\wedge} lower < p \ensuremath{\wedge} p \ensuremath{\leq} n \ensuremath{\wedge} 0 < d \ensuremath{\wedge} \ensuremath{\neg} p \ensuremath{\mid} d \ensuremath{\wedge}\newline       \ensuremath{\forall} e : \ensuremath{\mathbb{N}}, p\textasciicircum{}e * d \ensuremath{\in} productsOfSubsets n \ensuremath{\leftrightarrow}\newline         e = h \ensuremath{\vee} e = h + 1 \ensuremath{\vee} e = h + 2 \ensuremath{\vee} e = h + 3 \ensuremath{\vee}\newline           e = h + g + 4 \ensuremath{\vee} e = h + g + 5 \ensuremath{\vee} e = h + g + 6 \ensuremath{\vee} e = h + g + 7\end{quote}
{\small Source: proofs/a060957/OpHack/SubsetProducts/Generic/Final.lean:90.\par}

{\small Source SHA-256: a2e26af97d551980ecca1eb62f5db1362f75ada87f28795f99fe526651007d53\par}

\end{samepage}
\begin{samepage}
\subsection{SubsetProductGap.Generic.arbitrary\_gap\_consecutive\_above}
\begin{quote}\small\ttfamily theorem arbitrary\_gap\_consecutive\_above (g : \ensuremath{\mathbb{N}}) (hg : 1 \ensuremath{\leq} g) (lower : \ensuremath{\mathbb{N}}) :\newline     \ensuremath{\exists} n p M : \ensuremath{\mathbb{N}}, p.Prime \ensuremath{\wedge} lower < p \ensuremath{\wedge} p \ensuremath{\leq} n \ensuremath{\wedge} 0 < M \ensuremath{\wedge}\newline       M \ensuremath{\in} productsOfSubsets n \ensuremath{\wedge} p\textasciicircum{}(g + 1) * M \ensuremath{\in} productsOfSubsets n \ensuremath{\wedge}\newline         \ensuremath{\forall} j : \ensuremath{\mathbb{N}}, 1 \ensuremath{\leq} j \ensuremath{\to} j \ensuremath{\leq} g \ensuremath{\to} p\textasciicircum{}j * M \ensuremath{\notin} productsOfSubsets n\end{quote}
{\small Source: proofs/a060957/OpHack/SubsetProducts/Generic/Final.lean:113.\par}

{\small Source SHA-256: a2e26af97d551980ecca1eb62f5db1362f75ada87f28795f99fe526651007d53\par}

{\small\textbf{Sources and attribution:} \href{https://oeis.org/A060957}{1. oeis.org / A060957}; \href{https://arxiv.org/abs/2501.02695}{2. arXiv source: 2501.02695}; \href{https://github.com/alejandrozu/openmath-2026-judging/blob/d0ed487f487fa7ec814ff705ab55249983fe8707/OpenMath-Judging/review/entries/E02-a060957.txt}{3. alejandrozu/openmath-2026-judging / E02-a 060957.txt}; \href{https://github.com/alejandrozu/openmath-2026-judging}{Central source and adjudication archive}.\par}

\end{samepage}
\clearpage
\section{No positive-start periodic orbit for the prime-index reversal map}
{\small Provisional mathematical authors: Rachel Teo, Wenyi Wang, Hamdi Abderrahmene\par}\medskip
{\small Rachel Teo  -  Sakana AI.\par}

{\small Wenyi Wang  -  Sakana AI.\par}

{\small Hamdi Abderrahmene  -  Sakana AI.\par}

{\small Affiliations follow the recorded author declarations or public institutional/profile sources. Preferred author names, author order and publication affiliations await author review. This editorial compilation does not establish anthology coauthorship.\par}

\subsection{Abstract}
For the map sending n to the decimal reversal of its nth prime, every orbit beginning at a positive integer is not ultimately periodic. The proof combines a uniform large-index growth bound with an exact finite exclusion of the remaining cycles.

\subsection{Definition and exact result}
Write p\_n for the nth prime with p\_1=2, and let T(n) be the integer obtained by reversing the usual decimal digits of p\_n. Start at x>0 and iterate T. The theorem states that this sequence is not ultimately periodic for every positive x. This includes x=1. Any convention assigning a value to the zeroth prime index is outside the original positive-start problem.

An eventually periodic orbit would eventually enter a directed cycle in the functional graph of T. It therefore suffices to exclude every positive cycle. The proof is not just an empirical assertion that tested trajectories grow: it separates all sufficiently large indices, where a uniform inequality forbids a cycle, from a finite domain on which the entire graph can be checked.

\subsection{The source proof: uniform three-step escape}
The selected source proves that every n>0 has n<T(n), n<T\ensuremath{^2}(n), or n<T\ensuremath{^3}(n). Its elementary Chebyshev and decimal-digit estimates handle every prime at least 10,000. An exact small-prime prefix certificate identifies precisely 17 nonincreasing indices below that threshold; 16 escape on their second step and index 1118 escapes on its third.

The complete ordinary argument and every exceptional transition are given in the proof sections below. The independent exact recount also verifies the finite prime-counting bounds and source exception list. A periodic tail has a maximum, whereas the three-step lemma produces a later larger value, so no positive orbit can be ultimately periodic.

\subsection{Formalization and reproducibility}
The endpoint OeisA100475.not\_periodic takes the explicit assumption 0<x. Its canonical conjecture bridge excludes exactly the positive-start periodicity assertion. This publication pass independently replayed all 84 exact frozen authored source modules and the selected audit under Lean 4.33.1 and Mathlib 0df444a360eaa60ab8c11dca51a86af692955474, finishing on 5 October 2026 at 03:53:25 UTC. Both requested endpoint prints depend only on propext, Classical.choice and Quot.sound, with no admitted, native-evaluation or unknown axioms. The archived 1 October replay remains a separate historical record. The prime and digit modules support a single result and are not additional original theorems to be scored separately. Receipt: Verification/sakana-a100475-fresh-build.json.

The distributed supplement includes the exact-integer finite checker check\_final\_four\_certificates.py and its recorded prime-counting, exception-list and transition checks in final\_four\_independent\_certificate\_audit.json under Verification/Independent checkers. The ordinary proof below states the threshold bounds and every exceptional transition. A reader can reproduce the finite rejection independently of the Lean sieve. The decimal base and ordinary digit convention remain part of the statement.

\subsection{Historical comparison and publication scope}
The ingredients are classical and the final argument is short once its definitions are fixed. That creates a real historical priority question. The recorded review identified the Ball-Coxeter discussion and the 2004 SeqFan archive as specific unresolved leads. The mathematical result can be sound while its originality remains contingent on those sources.

If no earlier positive-start proof is found, this is a complete answer to the sequence question and a suitable concise preprint with formal certificates. If an earlier solution is found, move it to the known-mathematics formalization account without presenting zero-index totalization as a replacement novelty claim. Its value then lies in faithful formal treatment and an auditable finite/analytic proof chain.

A focused 5 October 2026 recheck found that the current OEIS record still asks the loop question and records computational extensions from December 2004. The September 2026 formal-conjectures correction handles the positive-start statement and the artificial zero-start solution; it does not exclude positive cycles. The linked Wolfram Reversal exposition cites Ball-Coxeter pp14-15 for general digit-reversal identities and products. The exact 13th-edition pages and original November/December 2004 SeqFan message bodies could not be retrieved in this bounded check. Thus neither an earlier positive-start nonperiodicity proof nor its absence has been established. Historical originality remains conditional, without a score change based on non-discovery.

\subsection{The three-step escape lemma and its dynamical consequence}
Let p\_n be the nth prime for n\ensuremath{\geq}1, and let R(p) reverse its ordinary decimal digits. Put T(n)=R(p\_n). The submitted proof establishes the stronger local assertion that, for every n>0, at least one of T(n), T\ensuremath{^2}(n), T\ensuremath{^3}(n) is greater than n. This assertion allows small decreases; it is stronger and more precise than an observation that sufficiently large indices grow. We give its actual prime-counting proof, rather than making the printed argument depend on an unstated nth-prime asymptotic.

Lemma 1 (finite-tail maximum). If a positive integer sequence is ultimately periodic, its periodic tail takes only finitely many values and therefore has a largest value m. Select an occurrence of m in that tail. Every later value is at most m. If the sequence is an orbit of T, the three-step assertion at m produces a later value greater than m, a contradiction. It follows that no positive-start orbit is ultimately periodic. Positivity causes no difficulty: a prime has a nonzero final decimal digit, so its reversal is positive and induction keeps every orbit term positive.

The theorem concerns ultimate periodicity, including an arbitrary preperiod, and includes the start 1. It therefore excludes every positive directed cycle. For a deterministic map on the positive integers, it also implies that every orbit is unbounded: any bounded orbit lies in a finite set, repeats a state by the pigeonhole principle, and hence is ultimately periodic. This last observation is a consequence of the same result and is not counted as a separate original family.

\subsection{Digit reversal and a sufficient explicit prime-counting bound}
Lemma 2 (digit size). If p is prime and has d decimal digits, its last digit is nonzero: otherwise 10 divides p, which contradicts primality. The reversal therefore has d digits too, and R(p)\ensuremath{\geq}10\textasciicircum{}(d\ensuremath{-}1), while p<10\textasciicircum{}d. Consequently p<10R(p). More generally, p\ensuremath{\geq}10\textasciicircum{}k implies R(p)\ensuremath{\geq}10\textasciicircum{}k. These statements also cover the small primes 2 and 5; no exception to the nonzero-last-digit statement is needed.

Write \ensuremath{\pi}(x) for the number of primes at most x, and \ensuremath{\theta}(x)=\ensuremath{\Sigma}\_\{p\ensuremath{\leq}x\}log p. We use the classical elementary estimate \ensuremath{\theta}(x)\ensuremath{\leq}(log 4)x and the exact partial-summation identity \ensuremath{\pi}(x)=\ensuremath{\theta}(x)/log x+\ensuremath{\int}\ensuremath{_2}\ensuremath{^{x}} \ensuremath{\theta}(t)/(t(log t)\ensuremath{^2}) dt, valid for x\ensuremath{\geq}2. These are background facts formalized in Mathlib, not new contributions of this chapter. Bounding \ensuremath{\theta} in both terms reduces the estimate to the integral I(x)=\ensuremath{\int}\ensuremath{_2}\ensuremath{^{x}}(log t)\textasciicircum{}(\ensuremath{-}2)dt.

For 1<a\ensuremath{\leq}b and 0<L\ensuremath{\leq}log a, monotonicity of log gives \ensuremath{\int}\ensuremath{_{a}}\ensuremath{^{b}}(log t)\textasciicircum{}(\ensuremath{-}2)dt\ensuremath{\leq}(b\ensuremath{-}a)/L\ensuremath{^2}. On [2,10] use log 2>1/2, so the integral is at most 32. On each decimal interval [10\textasciicircum{}k,10\textasciicircum{}(k+1)] for 1\ensuremath{\leq}k\ensuremath{\leq}6, use log(10\textasciicircum{}k)\ensuremath{\geq}2k. Adding the seven rational bounds gives I(10\ensuremath{^7})\ensuremath{\leq}80,000. For x\ensuremath{\geq}10\ensuremath{^7} we have log x\ensuremath{\geq}16, because log 10>23/10. Thus I(x)\ensuremath{\leq}80,000+(x\ensuremath{-}10\ensuremath{^7})/256\ensuremath{\leq}x/125; the last inequality is an elementary comparison of linear functions whose difference is zero at 10\ensuremath{^7} and has positive slope.

Since log 4\ensuremath{\leq}7/5 and log x\ensuremath{\geq}16, the \ensuremath{\theta}(x)/log x term is at most 7x/80. The integral term is at most (7/5)I(x)\ensuremath{\leq}7x/625. Their sum is (987/10,000)x=.0987x. Hence \ensuremath{\pi}(x)\ensuremath{\leq}.0987x for every x\ensuremath{\geq}10\ensuremath{^7}. Applying this at a prime p, and using p<10R(p), gives \ensuremath{\pi}(p)<.987R(p)<R(p). The strict inequality follows from R(p)>0.

\subsection{The three intermediate decades and the finite exceptional table}
Lemma 3 (all primes at least 10,000 grow). For 10\ensuremath{^4}\ensuremath{\leq}p<10\ensuremath{^5}, monotonicity and the finite certificate give \ensuremath{\pi}(p)\ensuremath{\leq}\ensuremath{\pi}(10\ensuremath{^5})\ensuremath{\leq}9,595<10\ensuremath{^4}\ensuremath{\leq}R(p). For 10\ensuremath{^5}\ensuremath{\leq}p<10\ensuremath{^6}, use \ensuremath{\pi}(10\ensuremath{^6})\ensuremath{\leq}94,067<10\ensuremath{^5}\ensuremath{\leq}R(p). For 10\ensuremath{^6}\ensuremath{\leq}p<10\ensuremath{^7}, the analytic bound at 10\ensuremath{^7} gives \ensuremath{\pi}(p)\ensuremath{\leq}987,000<10\ensuremath{^6}\ensuremath{\leq}R(p). For p\ensuremath{\geq}10\ensuremath{^7}, use Lemma 2 and the preceding .0987p bound. Since \ensuremath{\pi}(p\_n)=n, this proves T(n)>n whenever p\_n\ensuremath{\geq}10,000.

The submitted finite bounds are intentionally loose. In the independent ordinary Eratosthenes recount, \ensuremath{\pi}(100,000)=9,592 and \ensuremath{\pi}(1,000,000)=78,498, which imply the submitted inequalities. The formal million-range certificate instead counts survivors of a smaller-prime divisibility screen: every actual prime survives, and the complete survivor census is 94,067. Confusing that survivor total with the exact prime count would misdescribe the certificate. Neither loose constant is a conjectural prime estimate.

Lemma 4 (small exceptions). Enumerate all 1,229 primes below 10,000 with their one-based prime indices. Every pair (p,n) satisfies n<R(p) except precisely the following 17 pairs: (61,18), (71,20), (601,110), (701,126), (811,141), (821,142), (911,156), (941,160), (8101,1019), (8501,1060), (9001,1118), (9011,1120), (9311,1152), (9511,1178), (9601,1185), (9811,1210), (9901,1221). This is a finite certificate, not an extrapolation from sampled trajectories. Each integer is checked for primality by trial divisors up to its square root, and a cumulative prefix count gives the exact index.

For each exceptional index, the source certifies the nth primes needed for the next transitions. The complete escape table, stopping as soon as a value exceeds the initial index, is: 18 \ensuremath{\to} 16 \ensuremath{\to} 35; 20 \ensuremath{\to} 17 \ensuremath{\to} 95; 110 \ensuremath{\to} 106 \ensuremath{\to} 775; 126 \ensuremath{\to} 107 \ensuremath{\to} 785; 141 \ensuremath{\to} 118 \ensuremath{\to} 746; 142 \ensuremath{\to} 128 \ensuremath{\to} 917; 156 \ensuremath{\to} 119 \ensuremath{\to} 356; 160 \ensuremath{\to} 149 \ensuremath{\to} 958; 1019 \ensuremath{\to} 1018 \ensuremath{\to} 3908; 1060 \ensuremath{\to} 1058 \ensuremath{\to} 1648; 1118 \ensuremath{\to} 1009 \ensuremath{\to} 1108 \ensuremath{\to} 3988; 1120 \ensuremath{\to} 1109 \ensuremath{\to} 3298; 1152 \ensuremath{\to} 1139 \ensuremath{\to} 7819; 1178 \ensuremath{\to} 1159 \ensuremath{\to} 1739; 1185 \ensuremath{\to} 1069 \ensuremath{\to} 1858; 1210 \ensuremath{\to} 1189 \ensuremath{\to} 9269; 1221 \ensuremath{\to} 1099 \ensuremath{\to} 1288. All but the index 1118 escape on their second step. That exceptional case has 1118\ensuremath{\to}1009\ensuremath{\to}1108\ensuremath{\to}3988, and escapes on the third. All relevant nth primes are explicitly listed in the certificate and independently recovered by the ordinary prime sieve.

For any n>0, p\_n is either at least 10,000, when Lemma 3 gives immediate growth, or is below 10,000, when Lemma 4 gives immediate growth or one of the displayed two/three-step escapes. This proves the three-step assertion for every positive n. Lemma 1 then proves the desired nonperiodicity theorem without any assumption that an individual trajectory is monotone.

\subsection{Finite reproduction and the source bridge}
A minimal ordinary checker creates a Boolean Eratosthenes sieve through 1,000,000, extracts its ordered prime list, reverses the digits of each prime below 10,000, identifies the nonincreasing pairs, and computes only the next two transitions for those 17 indices. It checks the two loose prime-counting inequalities and compares the complete exception list with the printed table. The verification uses exact integer arithmetic. The independent checker for this editorial draft records every exceptional orbit prefix in final\_four\_independent\_certificate\_audit.json; it found exactly the source list and every required escape.

The source assembly is A100475Bounds\ensuremath{\to}A100475Large and A100475Certificate\ensuremath{\to}A100475Transitions, joined in A100475Conclusion.three\_step\_growth and not\_periodic\_of\_counts. The final OeisA100475.not\_periodic supplies both finite bounds. OeisA100475.conjecture is the canonical positive-start bridge. This ordinary exposition explains every infinite step of that assembly; the large finite computations remain explicit reproducible certificates attached to the proof. A successful fresh replay establishes the selected formal statements, while historical priority must still be assessed on its own evidence.

\subsection{Verification and reproducibility}
Fresh execution: sakana-a 100475: PASS; 84/84 custom modules compiled successfully; receipt Verification/sakana-a100475-fresh-build.json. Compiler pins, source hashes and endpoint axiom outputs are in Verification. Historical receipts and finite checks do not replace fresh replay.

\begin{samepage}
\subsection{OeisA100475.not\_periodic}
\begin{quote}\small\ttfamily theorem not\_periodic \{x : \ensuremath{\mathbb{N}}\} (hx : 0 < x) : \ensuremath{\neg} IsUltimatelyPeriodic (aStartAt x)\end{quote}
{\small Source: proofs/a100475/A100475.lean:17.\par}

{\small Source SHA-256: e22c200d87ed0a15a25d4760e3f37aaf0005b48aa5b734e79fb0776599c04c95\par}

\end{samepage}
\begin{samepage}
\subsection{OeisA100475.conjecture}
\begin{quote}\small\ttfamily theorem conjecture :\newline     False \ensuremath{\leftrightarrow} \ensuremath{\exists} x : \ensuremath{\mathbb{N}}, 0 < x \ensuremath{\wedge} x \ensuremath{\neq} 1 \ensuremath{\wedge} IsUltimatelyPeriodic (aStartAt x)\end{quote}
{\small Source: proofs/a100475/A100475.lean:20.\par}

{\small Source SHA-256: e22c200d87ed0a15a25d4760e3f37aaf0005b48aa5b734e79fb0776599c04c95\par}

{\small\textbf{Sources and attribution:} \href{https://oeis.org/A100475}{1. oeis.org / A100475}; \href{https://github.com/google-deepmind/formal-conjectures/pull/5567}{2. google-deepmind/formal-conjectures}; \href{https://dlmf.nist.gov/27.12\#E2}{3. dlmf.nist.gov / 27.12}; \href{https://mathworld.wolfram.com/Reversal.html}{4. mathworld.wolfram.com / Reversal.html}; \href{https://books.google.com/books/about/Mathematical_Recreations_Essays.html?id=9lJqNJhYc9oC}{5. books.google.com / Mathematical\_Recreations\_Essays.html}; \href{https://github.com/alejandrozu/openmath-2026-judging/blob/d0ed487f487fa7ec814ff705ab55249983fe8707/OpenMath-Judging/review/entries/E02-a100475.txt}{6. alejandrozu/openmath-2026-judging / E02-a 100475.txt}; \href{https://github.com/alejandrozu/openmath-2026-judging}{Central source and adjudication archive}.\par}

\end{samepage}
\clearpage
\section{A three-prime exclusion for the quadratic-residue divisibility criterion}
{\small Provisional mathematical authors: Rachel Teo, Wenyi Wang, Hamdi Abderrahmene\par}\medskip
{\small Rachel Teo  -  Sakana AI.\par}

{\small Wenyi Wang  -  Sakana AI.\par}

{\small Hamdi Abderrahmene  -  Sakana AI.\par}

{\small Affiliations follow the recorded author declarations or public institutional/profile sources. Preferred author names, author order and publication affiliations await author review. This editorial compilation does not establish anthology coauthorship.\par}

\subsection{Abstract}
For n=pqr with three distinct odd primes, the quadratic-residue count a(n), including zero, satisfies a(n)(a(n)\ensuremath{-}1) not dividing n\ensuremath{^2}\ensuremath{-}1. A global arithmetic reduction bounds the first two primes and the quotient; exact discriminant certificates exclude all remaining cases.

\subsection{The residue convention and theorem}
Let a(n) count all square residue classes modulo n, including zero. For distinct odd primes 2<p<q<r, the Chinese remainder theorem gives a(pqr)=(p+1)(q+1)(r+1)/8. The theorem excludes divisibility of n\ensuremath{^2}\ensuremath{-}1 by a(n)(a(n)\ensuremath{-}1) throughout this unbounded squarefree three-prime family. It therefore establishes the corresponding conjectured equivalence on this entire factorization class.

Including zero is essential: a count of only nonzero quadratic residues would change the arithmetic. The theorem says nothing about repeated prime powers or products of four or more distinct primes, and should be described as a uniform infinite subcase of A000224 rather than the full conjecture.

\subsection{Global reduction to a finite certificate}
Set A=(p+1)/2, B=(q+1)/2 and C=(r+1)/2. Under a contrary divisibility assumption, there is a positive integer quotient K satisfying (pqr)\ensuremath{^2}=KABC(ABC\ensuremath{-}1)+1. The proof derives necessary bounds from that exact equation, after explicitly justifying every conversion from natural-number subtraction to ring subtraction.

The key global step bounds p by 381 and gives a p-dependent upper bound on q. With p and q fixed, K lies in an explicit finite interval. The remaining equation is quadratic in C. Writing H=pq and U=AB, its coefficients are 4H\ensuremath{^2}\ensuremath{-}KU\ensuremath{^2}, \ensuremath{-}(4H\ensuremath{^2}\ensuremath{-}KU), and H\ensuremath{^2}\ensuremath{-}1. An integral root forces the corresponding discriminant to be an integer square.

The certificates reject that finite quotient/discriminant domain using strict gaps between consecutive squares, plus a separate congruence exclusion for the exceptional square case. Because the bounds were derived without an upper bound on r, this closes an unbounded theorem. A bounded Pell scan up to a large numerical cutoff would not by itself do so.

\subsection{Formal certificate architecture}
The source splits the argument into square-count transport, necessary bounds, quadratic reduction and generated certificate blocks. Its 171-module record includes each intended block and three final canonical endpoints. A separate Lean 4.33.1 replay is recorded in addition to 4.34.1; tactic-only compatibility changes in Bounds preserve definitions and theorem statements. Several early source comments still say unbuilt draft, so the dated manifest and selected endpoint receipts should govern verification descriptions.

For publication, append a compact definition of the finite domain and the discriminant identity, with machine-readable certificate data in the repository. The long generated leaves belong in supplementary material. Report the fresh common replay status separately from archived receipts and from reviewer sampling. An independent check of 2,024 small prime triples agrees with the formula and exclusion, but is not the universal proof.

\subsection{Prior scope and research contribution}
CRT, quotient/Pell methods and gcd obstructions are background. Earlier records already cover simpler factorization classes and a large bounded search; the claimed delta is the global reduction and exclusion for all ordered triples of distinct odd primes. The previous scan's operational n cutoff must not be confused with an unbounded theorem merely because every local survivor happens to fall below it.

This is a focused computational number-theory paper candidate if the closest pre-event project and specifications do not already entail the theorem. It should state the factorization limitation prominently and credit the baseline methods. The number of generated modules or rejected arithmetic leaves is not a count of separate discoveries.

\subsection{Square-residue counting and the exact quotient equation}
Let a(n) count distinct residues of squares modulo n, including zero. For an odd prime p, the nonzero squares pair the nonzero residues under x\ensuremath{\mapsto}\ensuremath{-}x, so there are (p\ensuremath{-}1)/2 nonzero square residues and (p+1)/2 square residues in total. Chinese remaindering identifies square residues modulo a product of distinct primes with the product of their square-residue sets. Thus, for odd primes p<q<r, set a=(p+1)/2, b=(q+1)/2, c=(r+1)/2, N=pqr and R=abc; then a(N)=R.

Assume, for contradiction, R(R\ensuremath{-}1) divides N\ensuremath{^2}\ensuremath{-}1. Since N>1 and R>1, there is a positive integer K with N\ensuremath{^2}=KR(R\ensuremath{-}1)+1. We will first bound the finite prefix (p,q,K), without bounding the third prime directly, and then exclude every possible integer c. The residue-count bridge is essential: an arithmetic theorem about abc alone would not yet prove the OEIS statement.

\subsection{Lemma: bounding the smallest prime and the quotient}
Since a,b,c\ensuremath{\geq}2, N=(2a\ensuremath{-}1)(2b\ensuremath{-}1)(2c\ensuremath{-}1)\ensuremath{\leq}(2a\ensuremath{-}1)\ensuremath{\cdot}4bc=8R\ensuremath{-}4bc\ensuremath{\leq}8R\ensuremath{-}16, and R\ensuremath{\geq}8. Therefore N\ensuremath{^2}\ensuremath{\leq}(8R\ensuremath{-}16)\ensuremath{^2}<64R(R\ensuremath{-}1)+1. The quotient equation forces K<64, so K\ensuremath{\leq}63. It also gives N\ensuremath{^2}=KR\ensuremath{^2}\ensuremath{-}KR+1<KR\ensuremath{^2}, because KR>1. Hence (N/R)\ensuremath{^2}<K.

The normalized product is N/R=(2\ensuremath{-}1/a)(2\ensuremath{-}1/b)(2\ensuremath{-}1/c). Strict prime order gives a<b<c, so (2\ensuremath{-}1/a)\textasciicircum{}6<(N/R)\ensuremath{^2}<K\ensuremath{\leq}63. If a\ensuremath{\geq}192, Bernoulli's inequality yields (2\ensuremath{-}1/a)\textasciicircum{}6=64(1\ensuremath{-}1/(2a))\textasciicircum{}6\ensuremath{\geq}64(1\ensuremath{-}3/a)\ensuremath{\geq}63, a contradiction. Thus a\ensuremath{\leq}191 and p=2a\ensuremath{-}1\ensuremath{\leq}381. This turns the smallest prime into a genuinely finite variable; no empirical search cutoff is assumed.

\subsection{Lemma: an exact finite bound for the second prime}
Fix p and its half-prime a. Put T=16p\ensuremath{^2}, B=\ensuremath{\lfloor}(T\ensuremath{-}1)/a\ensuremath{^2}\ensuremath{\rfloor}, and d=T\ensuremath{-}Ba\ensuremath{^2}. The preceding-integer floor guarantees 1\ensuremath{\leq}d\ensuremath{\leq}a\ensuremath{^2}, even when T/a\ensuremath{^2} is integral. A strict fixed-prefix estimate gives K<T/a\ensuremath{^2}. To check it, set u=bc, t=(2b\ensuremath{-}1)(2c\ensuremath{-}1), and s=2b+2c\ensuremath{-}1. Then 16u\ensuremath{^2}\ensuremath{-}t\ensuremath{^2}=s(4u+t)>28u>16u/a. Multiplying by p\ensuremath{^2} gives N\ensuremath{^2}<(T/a\ensuremath{^2})R(R\ensuremath{-}1). Comparing with N\ensuremath{^2}=KR(R\ensuremath{-}1)+1 proves K<T/a\ensuremath{^2}. Since K is integral, Ka\ensuremath{^2}<T and K\ensuremath{\leq}B.

For the lower estimate, c\ensuremath{\geq}b and the normalized quotient imply K>(p/a)\ensuremath{^2}(2\ensuremath{-}1/b)\textasciicircum{}4. Bernoulli gives (2\ensuremath{-}1/b)\textasciicircum{}4=16(1\ensuremath{-}1/(2b))\textasciicircum{}4\ensuremath{\geq}16(1\ensuremath{-}2/b). Clearing positive denominators yields T(b\ensuremath{-}2)<Ka\ensuremath{^2}b\ensuremath{\leq}Ba\ensuremath{^2}b=(T\ensuremath{-}d)b. Hence db<2T. Integrality gives b\ensuremath{\leq}\ensuremath{\lfloor}(2T\ensuremath{-}1)/d\ensuremath{\rfloor} and

q\ensuremath{\leq}qBound(p)=2\ensuremath{\lfloor}(32p\ensuremath{^2}\ensuremath{-}1)/d\ensuremath{\rfloor}\ensuremath{-}1, where d=16p\ensuremath{^2}\ensuremath{-}a\ensuremath{^2}\ensuremath{\lfloor}(16p\ensuremath{^2}\ensuremath{-}1)/a\ensuremath{^2}\ensuremath{\rfloor}.

This bound remains finite for every p\ensuremath{\leq}381. Choosing a preceding-integer floor rather than an informal limiting floor prevents a zero-denominator mistake. In the formal development positivity of d is a separate proved lemma.

\subsection{Lemma: a narrow quotient interval and a discriminant certificate}
For a fixed prime pair put H=pq and U=ab. Since c\ensuremath{\geq}b+1, the normalized lower estimate gives K>H\ensuremath{^2}(2b+1)\ensuremath{^2}/[U\ensuremath{^2}(b+1)\ensuremath{^2}]. A sharper limiting upper estimate gives K<4H\ensuremath{^2}/U\ensuremath{^2}. Its cleared difference is

(4H\ensuremath{^2}/U\ensuremath{^2})\ensuremath{\cdot}Uc(Uc\ensuremath{-}1)\ensuremath{-}[H(2c\ensuremath{-}1)]\ensuremath{^2}=(H\ensuremath{^2}/U)[4c(U\ensuremath{-}1)\ensuremath{-}U]>0.

Consequently the exact integer bounds are KLo(p,q)=\ensuremath{\lfloor}H\ensuremath{^2}(2b+1)\ensuremath{^2}/[U\ensuremath{^2}(b+1)\ensuremath{^2}]\ensuremath{\rfloor}+1 and KHi(p,q)=\ensuremath{\lfloor}(4H\ensuremath{^2}\ensuremath{-}1)/U\ensuremath{^2}\ensuremath{\rfloor}. The finite domain now consists of primes 2<p<q with p\ensuremath{\leq}381, q\ensuremath{\leq}qBound(p), and KLo\ensuremath{\leq}K\ensuremath{\leq}KHi. These bounds exclude the unbounded third prime from the enumeration.

Expand the quotient equation as a quadratic in the integer c. Define A=4H\ensuremath{^2}\ensuremath{-}KU\ensuremath{^2}, E=4H\ensuremath{^2}\ensuremath{-}KU and C=H\ensuremath{^2}\ensuremath{-}1. It becomes Ac\ensuremath{^2}\ensuremath{-}Ec+C=0. Completing the square gives (2Ac\ensuremath{-}E)\ensuremath{^2}=D, where D=E\ensuremath{^2}\ensuremath{-}4AC. The finite certificate lists, for each nonexceptional domain triple, an integer z\ensuremath{\geq}0 with z\ensuremath{^2}<D<(z+1)\ensuremath{^2}. Such a strict gap cannot contain an integer square, so that triple has no integer c satisfying the equation. This is stronger than checking a finite collection of third primes.

The archived certificate covers 7,973 domain leaves and has one square-discriminant exception, (p,q,K)=(3,17,31). That case is still impossible: H=51 and U=18, so reducing [H(2c\ensuremath{-}1)]\ensuremath{^2}=K(Uc)(Uc\ensuremath{-}1)+1 modulo 3 gives 0=1. The contradiction does not require c to be prime. The generated Lean certificate proves coverage of the small prime list, the second-prime and quotient intervals, and each strict square gap; an external integer-square-root flag alone would not establish this interface.

\subsection{Assembling the restricted OEIS theorem}
The symbolic inequalities prove that any divisibility counterexample with three distinct odd primes supplies a triple in the finite domain and an integer solution of its quotient equation. The certificate and the exceptional congruence exclude every such solution. Therefore a(pqr)(a(pqr)\ensuremath{-}1) does not divide (pqr)\ensuremath{^2}\ensuremath{-}1 for odd primes p<q<r. The equivalent modular assertion (pqr)\ensuremath{^2}\ensuremath{\equiv}1 modulo a(pqr)(a(pqr)\ensuremath{-}1) is false. Since pqr is composite, the two sides of the restricted prime-characterization equivalence are both false, yielding the submitted restricted conjecture theorem.

This proof does not classify products with repeated prime factors, two-prime composites, products involving 2, or arbitrary integers. Those are outside its hypotheses. A release should attach the generated finite certificate and an executable checker, identify the complete domain formulas above, and state the exact frozen theorem endpoint. The ordinary argument is complete once the explicitly described finite certificate is included; a reader need not trust a search that simply reports no examples.

\subsection{Verification and reproducibility}
Dated partial fresh replay: the 5 October 2026, 08:56:47 UTC qualification verifies 167 of 171 frozen source modules, comprising 166 granular sources and Bounds. The retained Bounds source passed under its original policy before the controller exited 120; its unfinished receipt is preserved. The subsequent source-free initial-gate failure also remains recorded. Four whole-import sources and the three original selected endpoint prints remain unrun, so the full 171-source/3-print scope is unqualified. The mathematical claim remains restricted to three ordered distinct odd primes, using the zero-inclusive quadratic-residue count; it is not the unrestricted A000224 conjecture. Qualification: Verification/A000-actual167-abnormal-controller-exit120-partial-readonly-qualification-20261005T085647463194Z.json. Historical author/internal replay is separate, and these partial checks confer no novelty, score, award or release upgrade. A later bounded dispatch assessment ended without a new Lean child or source/audit attempt. Dated resource-only summary: Verification/root-A000-whole8GiB-bounded60-public-safe-scope-20261005.json.

\begin{samepage}
\subsection{OeisA224.not\_divisible\_three\_primes}
\begin{quote}\small\ttfamily theorem not\_divisible\_three\_primes \{p q r : \ensuremath{\mathbb{N}}\}\newline     (hp : p.Prime) (hq : q.Prime) (hr : r.Prime)\newline     (hp2 : 2 < p) (hpq : p < q) (hqr : q < r) :\newline     \ensuremath{\neg} a (p * q * r) * (a (p * q * r) - 1) \ensuremath{\mid} (p * q * r) \textasciicircum{} 2 - 1\end{quote}
{\small Source: proofs/a000224/OpHack/QuadraticResidue224/Target224.lean:33.\par}

{\small Source SHA-256: 3b67f7494439886af3359a9e684433f49a6dccb7d10258d1d05dcc01de7886d0\par}

\end{samepage}
\begin{samepage}
\subsection{OeisA224.not\_modEq\_three\_primes}
\begin{quote}\small\ttfamily theorem not\_modEq\_three\_primes \{p q r : \ensuremath{\mathbb{N}}\}\newline     (hp : p.Prime) (hq : q.Prime) (hr : r.Prime)\newline     (hp2 : 2 < p) (hpq : p < q) (hqr : q < r) :\newline     \ensuremath{\neg} (p * q * r) \textasciicircum{} 2 \ensuremath{\equiv} 1 [MOD a (p * q * r) * (a (p * q * r) - 1)]\end{quote}
{\small Source: proofs/a000224/OpHack/QuadraticResidue224/Target224.lean:43.\par}

{\small Source SHA-256: 3b67f7494439886af3359a9e684433f49a6dccb7d10258d1d05dcc01de7886d0\par}

\end{samepage}
\begin{samepage}
\subsection{OeisA224.conjecture\_three\_primes}
\begin{quote}\small\ttfamily theorem conjecture\_three\_primes \{p q r : \ensuremath{\mathbb{N}}\}\newline     (hp : p.Prime) (hq : q.Prime) (hr : r.Prime)\newline     (hp2 : 2 < p) (hpq : p < q) (hqr : q < r) :\newline     ((p * q * r).Prime \ensuremath{\wedge} p * q * r \ensuremath{\neq} 2) \ensuremath{\leftrightarrow}\newline       (p * q * r) \textasciicircum{} 2 \ensuremath{\equiv} 1 [MOD a (p * q * r) * (a (p * q * r) - 1)]\end{quote}
{\small Source: proofs/a000224/OpHack/QuadraticResidue224/Target224.lean:50.\par}

{\small Source SHA-256: 3b67f7494439886af3359a9e684433f49a6dccb7d10258d1d05dcc01de7886d0\par}

{\small\textbf{Sources and attribution:} \href{https://oeis.org/A000224}{1. oeis.org / A000224}; \href{https://github.com/umaia1234/agentic-conjectures/blob/96ea0039d9dc3a1940b9d5ec1cbe616ed58d1da2/problems/oeis-a000224/DETAILS.md}{2. umaia1234/agentic-conjectures / DETAILS.md}; \href{https://github.com/umaia1234/agentic-conjectures/blob/96ea0039d9dc3a1940b9d5ec1cbe616ed58d1da2/problems/oeis-a000224/pell_quotient_scan.py}{3. umaia1234/agentic-conjectures / pell\_quotient\_scan.py}; \href{https://github.com/alejandrozu/openmath-2026-judging/blob/d0ed487f487fa7ec814ff705ab55249983fe8707/OpenMath-Judging/review/entries/E02-a000224.txt}{4. alejandrozu/openmath-2026-judging / E02-a 000224.txt}; \href{https://github.com/alejandrozu/openmath-2026-judging}{Central source and adjudication archive}.\par}

\end{samepage}
\clearpage
\section{Factor-preserving deficiency hierarchies in Ferrers graphs}
{\small Provisional mathematical authors: Rachel Teo, Wenyi Wang, Hamdi Abderrahmene\par}\medskip
{\small Rachel Teo  -  Sakana AI.\par}

{\small Wenyi Wang  -  Sakana AI.\par}

{\small Hamdi Abderrahmene  -  Sakana AI.\par}

{\small Affiliations follow the recorded author declarations or public institutional/profile sources. Preferred author names, author order and publication affiliations await author review. This editorial compilation does not establish anthology coauthorship.\par}

\subsection{Abstract}
Under two consecutive matching-union rank equalities in a Ferrers graph, every supplied optimal degree-r factor can be preserved while r\ensuremath{-}e disjoint matchings of the required size are appended. The result determines subsequent ranks and gives balanced and one-defect colour-prefix realizations.

\subsection{Ranks and the preservation theorem}
A Ferrers graph is the bipartite incidence graph of a Young diagram: row lengths are weakly decreasing, and a cell is an allowed row-column edge. Let \ensuremath{\alpha}\_j be the maximum number of edges in a union of j matchings, equivalently the maximum size of an allowed subgraph whose row and column degrees are at most j.

For integers 1\ensuremath{\leq}r\ensuremath{\leq}k and 0\ensuremath{\leq}e<r, assume \ensuremath{\alpha}\_r=rk\ensuremath{-}e and \ensuremath{\alpha}\_(r+1)=(r+1)k\ensuremath{-}r\ensuremath{-}e. Then \ensuremath{\alpha}\_1=k. More strongly, every prescribed optimal degree-r factor F of size rk\ensuremath{-}e is retained, while r\ensuremath{-}e edge-disjoint matching classes, each of size k\ensuremath{-}r and disjoint from F, can be appended. For 0\ensuremath{\leq}h\ensuremath{\leq}r\ensuremath{-}e, the exact later rank is \ensuremath{\alpha}\_(r+h)=rk\ensuremath{-}e+h(k\ensuremath{-}r).

The words every prescribed factor express the useful strengthening. An existence theorem for one specially chosen factor would not imply that an arbitrary already-selected optimum can be extended.

\subsection{Tight corners and matching extensions}
Ferrers shape gives an explicit corner/min-cut formula for the degree-constrained factor rank. Equality in the two adjacent rank bounds forces a tight deficient corner. Its geometry divides the available cells into two arms and constrains how the selected red factor occupies the special corner region.

The proof extracts a budget on the special red edges, applies a capacitated Hall construction to the top arm, and obtains r\ensuremath{-}e matching classes there. It transposes the Ferrers graph and applies the same result to the other arm. The two arms have separated row and column supports, so corresponding matchings can be united without creating a repeated row or column. Each combined class has the required size k\ensuremath{-}r and avoids every edge of the original F.

Appending h classes supplies a lower bound for \ensuremath{\alpha}\_(r+h); the tight corner supplies the matching upper bound. Balanced and one-defect cases additionally organize the initial colour classes with the required prefix sizes. The source bridges the incidence model to the actual YoungDiagram graph and its independence-number notation, preventing a theorem about a surrogate rank definition from being mistaken for the target statement.

\subsection{Balanced and deficient consequences}
At e=0, the result realizes simultaneous first 2r CDS colours: r initial independent sets of size k and r further sets of size k\ensuremath{-}r, with all prefixes optimal. At e=1, it supplies the first 2r\ensuremath{-}1 colours with the stated one-defect profile. For e>1, the submitted theorem does not claim that an arbitrary supplied colouring already has every desired individual early colour size.

The general Latin Tableau problem remains open. The contribution is a restricted structural hierarchy, with a preservation quantifier that is stronger than selecting a convenient optimum. Parameterized sharpness and missing-corner transport discussed elsewhere are outside the selected 44-module formal closure unless separately proved and admitted.

\subsection{Literature and proof records}
Corner formulas, capacitated Hall and bipartite edge colouring are established tools. The greedy boundary e=r\ensuremath{-}1 and several small existential cases are also known. In particular, the recorded September 2026 r=3,e=1 existence comparison should be credited, while distinguishing it from arbitrary supplied-factor preservation.

This publication pass independently replayed all 44 exact frozen authored source modules and the selected audit under Lean 4.33.1 and Mathlib 0df444a360eaa60ab8c11dca51a86af692955474, finishing on 5 October 2026 at 02:41:25 UTC. All source invocations and all four requested endpoint prints passed. The factor-preserving extension, actual Young-diagram later ranks, one-defect independent sets and balanced first-two-r independent sets each use only propext, Classical.choice and Quot.sound. No selected endpoint depends on admitted, native-evaluation or unknown axioms. The archive also retains the entrants' earlier dual-version replays. Some factor-theory full texts remain unchecked, so originality is still a proposed precise extension beyond the credited boundary and existence cases. Formal verification does not settle that comparison or the full Latin Tableau conjecture. Receipt: Verification/sakana-ferrers-fresh-build.json.

\subsection{The factor rank and its corner formula}
Let a Ferrers board have rows 0,...,n\ensuremath{-}1, columns 0,...,W\ensuremath{-}1 and nonincreasing row lengths \ensuremath{\ell}\_u, clipped at W. A cell (u,v) is allowed when v<\ensuremath{\ell}\_u. Let \ensuremath{\alpha}\_s be the maximum number of allowed cells in a set with at most s cells in every row and every column. Bipartite edge colouring identifies such a set with a union of s disjoint matchings. This equality concerns edge sets; it is not a statement that a chosen colouring must already have a prescribed list of colour sizes.

For a corner (i,j), let A(i,j) be the number of allowed cells with u\ensuremath{\geq}i and v\ensuremath{\geq}j. The corner bound is \ensuremath{\alpha}\_s\ensuremath{\leq}s(i+j)+A(i,j). To prove it, partition a feasible set into cells in the first i rows, cells in the first j columns outside those rows, and cells in the remaining corner. The first two parts contribute at most si and sj and the last at most A(i,j). The Ferrers factor/min-cut theorem gives equality at some corner, hence \ensuremath{\alpha}\_s=min\_\{i,j\}\{s(i+j)+A(i,j)\}. The source develops this formula through capacitated matching and actual finite cell sets, including clipped boards and empty row or column sets. We use that standard finite factor theorem below, while proving the additional equality geometry explicitly.

\subsection{Lemma: the two ranks force a nearly square tight corner}
Assume 1\ensuremath{\leq}r\ensuremath{\leq}k, 0\ensuremath{\leq}e<r, \ensuremath{\alpha}\_r=rk\ensuremath{-}e, and \ensuremath{\alpha}\_\{r+1\}=(r+1)k\ensuremath{-}r\ensuremath{-}e. Choose a corner attaining the second minimum and, among all such corners, maximize i+j. Put A=A(i,j). The r-corner lower bound gives rk\ensuremath{-}e\ensuremath{\leq}r(i+j)+A, while equality at r+1 gives (r+1)(i+j)+A=(r+1)k\ensuremath{-}r\ensuremath{-}e. Subtracting yields i+j\ensuremath{\leq}k\ensuremath{-}r, and substituting back yields A\ensuremath{\geq}r\ensuremath{^2}\ensuremath{-}e.

Removing the first remaining row changes the (r+1)-score by r+1 minus the length of that row strip. A strip of more than r+1 would make the new score smaller than its minimum. A strip of exactly r+1 would preserve equality and increase i+j, contradicting maximality. Therefore its length is at most r. The same argument applies to the first remaining column. Ferrers nesting now puts every allowed cell of the corner in the r by r box starting at (i,j), so A\ensuremath{\leq}r\ensuremath{^2}.

If i+j\ensuremath{\leq}k\ensuremath{-}r\ensuremath{-}1, the exact (r+1)-score would give A\ensuremath{\geq}r\ensuremath{^2}+r+1\ensuremath{-}e>r\ensuremath{^2}, a contradiction. Hence i+j+r=k and A=r\ensuremath{^2}\ensuremath{-}e. Since e<r, the corner contains more than r(r\ensuremath{-}1) cells. It must meet the last row and the last column of its r by r box. Ferrers downward closure then makes its first row and first column full. Separate witnesses establish these two strips; no bottom-right cell or completely filled square is claimed when e>0. The remaining corner is therefore a genuine nearly square Ferrers shape with exactly e missing cells.

\subsection{Lemma: every prescribed optimum has the same red budget}
Now fix any feasible r-factor F with |F|=rk\ensuremath{-}e and call its cells red. At the tight corner, the bound |F|\ensuremath{\leq}r(i+j)+A is an equality. Equality in the three-part counting argument forces every first-row capacity and every first-column capacity to be saturated, forces the top-left i by j rectangle to contain no red cells, and forces all allowed cells of the nearly square corner into F. This is a conclusion about an arbitrary prescribed optimum F.

Focus on the top arm, consisting of rows u<i and columns v\ensuremath{\geq}j. Each top row has exactly r red cells. The r special columns j,...,j+r\ensuremath{-}1 are allowed in every top row because the first corner row is full and the upper rows are longer. The corner already places r\ensuremath{^2}\ensuremath{-}e red cells in those columns. Their total r-capacity is r\ensuremath{^2}. Therefore at most e red cells of the top arm lie in the special columns. All its other red cells lie in the ordinary suffix columns v\ensuremath{\geq}j+r. The top-arm ordinary red mass is consequently at least ri\ensuremath{-}e. Transposition gives precisely the same conclusions for the left arm, without changing any red cell.

\subsection{Lemma: a residual Ferrers arm contains enough full matching colours}
We give the capacitated Hall argument responsible for preservation. Consider a general top arm with a rows, r universal special columns, and nested ordinary neighbourhoods of lengths \ensuremath{\lambda}\_u. In each row let R\_u and S\_u be its red ordinary and special neighbours. Assume |R\_u|+|S\_u|\ensuremath{\leq}r, ordinary red column degrees at most r, and \ensuremath{\Sigma}\_u|R\_u|\ensuremath{\geq}ra\ensuremath{-}e. Delete all red cells from the arm. For a row subset T, write N\_j(T) for the number of rows of T allowing ordinary column j, d\_j(T) for its ordinary red degree, and s\_j(T) for its special red degree. A nonred factor with row and column capacities r has the global cut condition

\ensuremath{\Sigma}\_u|R\_u|\ensuremath{\leq}r(a\ensuremath{-}|T|)+\ensuremath{\Sigma}\_\{ordinary j\} min(r,N\_j(T)\ensuremath{-}d\_j(T))+\ensuremath{\Sigma}\_\{special j\} min(r,|T|\ensuremath{-}s\_j(T)).

Here is a direct proof of that condition. Sort T by decreasing ordinary neighbourhood length. Let P be its first min(r,|T|) rows and Q its remaining rows. In an ordinary column with a red edge from Q, nestedness makes every row of P allowed. Since the full ordinary red column has at most r edges, the number of nonred allowed rows in T is at least the number of red Q edges. The latter is also at most r. Thus the ordinary residual capacities cover all ordinary red mass in Q. When a column has no red Q edge, that assertion is immediate.

For each special column j, the nonred rows in P are a subset of its nonred rows in T and number at most r. Hence |P|\ensuremath{\leq}min(r,|T|\ensuremath{-}s\_j(T))+s\_j(P). Summing over the r special columns and using \ensuremath{\Sigma}\_\{u\ensuremath{\in}P\}(|R\_u|+|S\_u|)\ensuremath{\leq}r|P| shows that their residual capacities cover the ordinary red mass in P. Adding the two estimates covers the ordinary red mass inside T. Every row outside T contributes at most r, giving exactly the displayed global cut condition.

Capacitated Hall, equivalently integral max flow in the finite bipartite network, therefore supplies a nonred factor with at least \ensuremath{\Sigma}\_u|R\_u|\ensuremath{\geq}ra\ensuremath{-}e cells and degree at most r. Properly colour its edges with r colours. Each colour is a matching of size at most a. The sum of colour deficits from a is at most e, so at most e colours can be deficient. There are at least r\ensuremath{-}e full matchings, each saturating all a rows, and they are edge disjoint. This proves the residual-arm extension lemma. The argument explains why a total budget of e special red cells loses at most e complete colours, rather than requiring each special column to be free.

\subsection{Completing preservation and the exact rank hierarchy}
Apply the residual-arm lemma to the top arm with a=i and to the transposed left arm with a=j. It supplies r\ensuremath{-}e matchings of size i in the top arm, and r\ensuremath{-}e matchings of size j in the left arm, all avoiding F. Match corresponding colours. Their union is still a matching: the top arm uses rows below index i and columns at least j, while the left arm uses rows at least i and columns below j. These row and column supports are separated. Each union has size i+j=k\ensuremath{-}r. Distinct unions are edge disjoint and retain every edge of the original F.

For 0\ensuremath{\leq}h\ensuremath{\leq}r\ensuremath{-}e, adjoin any h of these new classes to F. The result is a feasible (r+h)-factor of size rk\ensuremath{-}e+h(k\ensuremath{-}r). The same corner bounds \ensuremath{\alpha}\_\{r+h\} above by (r+h)(i+j)+(r\ensuremath{^2}\ensuremath{-}e), which is exactly that size. Therefore \ensuremath{\alpha}\_\{r+h\}=(r+h)k\ensuremath{-}hr\ensuremath{-}e. For the ordinary matching rank, deleting the r corner rows leaves an empty southeast corner, so the (i+r,j) cut gives \ensuremath{\alpha}\_1\ensuremath{\leq}i+r+j=k. Conversely any r-factor decomposes into r matchings, giving rk\ensuremath{-}e\ensuremath{\leq}r\ensuremath{\alpha}\_1. Because e<r and \ensuremath{\alpha}\_1 is integral, \ensuremath{\alpha}\_1\ensuremath{\geq}k. Thus \ensuremath{\alpha}\_1=k.

When e=0, an r-edge-colouring of the original factor has r classes of size k, because their sum is rk and no matching is larger than k. The extension supplies r additional classes of size k\ensuremath{-}r. Every prefix through 2r attains the corresponding rank. When e=1 and r\ensuremath{\geq}2, the original r classes have sizes k except for one class of size k\ensuremath{-}1; order that deficient class last. The r\ensuremath{-}1 new classes have size k\ensuremath{-}r, yielding the advertised one-defect prefixes through 2r\ensuremath{-}1. For e>1 the argument preserves the total original deficit but does not impose a particular distribution among supplied colour classes. Finally, row-column nonincidence of cells is exactly independence in the Young-diagram cell graph, so the factor statements transport to the submitted independence-rank formulation.

\subsection{Verification and reproducibility}
Fresh execution: sakana-ferrers: PASS; 44/44 custom modules compiled successfully; receipt Verification/sakana-ferrers-fresh-build.json. Compiler pins, source hashes and endpoint axiom outputs are in Verification. Historical receipts and finite checks do not replace fresh replay.

\begin{samepage}
\subsection{FerrersArm.TightDefectCorner.defect\_preserves\_factor}
\begin{quote}\small\ttfamily theorem defect\_preserves\_factor \{n W : \ensuremath{\mathbb{N}}\} (\ensuremath{\ell} : Fin n \ensuremath{\to} \ensuremath{\mathbb{N}}) (h\ensuremath{\ell} : Antitone \ensuremath{\ell})\newline     (r k e : \ensuremath{\mathbb{N}}) (hr : 1 \ensuremath{\leq} r) (hk : r \ensuremath{\leq} k) (he : e < r)\newline     (hR : factorRank (W := W) \ensuremath{\ell} r = r*k-e)\newline     (hNext : factorRank (W := W) \ensuremath{\ell} (r+1) = (r+1)*k-r-e)\newline     (F : Finset (Fin n \ensuremath{\times} Fin W)) (hF : Feasible \ensuremath{\ell} r F) (hcard : F.card = r*k-e) :\newline     \ensuremath{\exists} G : Fin (r-e) \ensuremath{\to} Finset (Fin n \ensuremath{\times} Fin W),\newline       (\ensuremath{\forall} c, IsMatching (G c) \ensuremath{\wedge} (G c).card = k-r \ensuremath{\wedge}\newline         \ensuremath{\forall} p \ensuremath{\in} G c, Allowed \ensuremath{\ell} p.1 p.2 \ensuremath{\wedge} p \ensuremath{\notin} F) \ensuremath{\wedge}\newline       (\ensuremath{\forall} c d, c \ensuremath{\neq} d \ensuremath{\to} Disjoint (G c) (G d)) \ensuremath{\wedge}\newline       (\ensuremath{\forall} h \ensuremath{\leq} r-e, factorRank (W := W) \ensuremath{\ell} (r+h) = (r+h)*k-h*r-e) \ensuremath{\wedge}\newline       factorRank (W := W) \ensuremath{\ell} 1 = k\end{quote}
{\small Source: proofs/ferrers/OpHack/Ferrers/DefectExtension.lean:104.\par}

{\small Source SHA-256: c17d8534968c3ac199813212046ba0da17c65fa46e2907460afeca99323b1e20\par}

\end{samepage}
\begin{samepage}
\subsection{FerrersArm.YoungDiagramBridge.defect\_indepNumK\_later\_ranks}
\begin{quote}\small\ttfamily theorem defect\_indepNumK\_later\_ranks (\ensuremath{\mu} : YoungDiagram) (r k e : \ensuremath{\mathbb{N}})\newline     (hr : 1 \ensuremath{\leq} r) (hk : r \ensuremath{\leq} k) (he : e < r)\newline     (hR : FCCompat.indepNumK (FCCompat.YoungDiagram.toSimpleGraph \ensuremath{\mu}) r = r*k-e)\newline     (hNext : FCCompat.indepNumK (FCCompat.YoungDiagram.toSimpleGraph \ensuremath{\mu})\newline       (r+1) = (r+1)*k-r-e) :\newline     \ensuremath{\forall} h \ensuremath{\leq} r-e, FCCompat.indepNumK (FCCompat.YoungDiagram.toSimpleGraph \ensuremath{\mu})\newline       (r+h) = (r+h)*k-h*r-e\end{quote}
{\small Source: proofs/ferrers/OpHack/Ferrers/YoungDiagramDefect.lean:8.\par}

{\small Source SHA-256: c55685043c4b319f834184a218377c836af11200769eeae94fda76c4fc2c2dbc\par}

\end{samepage}
\begin{samepage}
\subsection{FerrersArm.YoungDiagramBridge.one\_defect\_independent\_sets}
\begin{quote}\small\ttfamily theorem one\_defect\_independent\_sets (\ensuremath{\mu} : YoungDiagram) (r k : \ensuremath{\mathbb{N}})\newline     (hr : 2 \ensuremath{\leq} r) (hk : r \ensuremath{\leq} k)\newline     (hR : FCCompat.indepNumK (FCCompat.YoungDiagram.toSimpleGraph \ensuremath{\mu}) r = r*k-1)\newline     (hNext : FCCompat.indepNumK (FCCompat.YoungDiagram.toSimpleGraph \ensuremath{\mu})\newline       (r+1) = (r+1)*k-r-1) :\newline     \ensuremath{\exists} (R : Fin r \ensuremath{\to} Finset (FCCompat.YoungDiagram.Cell \ensuremath{\mu}))\newline       (G : Fin (r-1) \ensuremath{\to} Finset (FCCompat.YoungDiagram.Cell \ensuremath{\mu})),\newline       (\ensuremath{\forall} c, (FCCompat.YoungDiagram.toSimpleGraph \ensuremath{\mu}).IsIndepSet\newline         (R c : Set (FCCompat.YoungDiagram.Cell \ensuremath{\mu})) \ensuremath{\wedge}\newline         (R c).card = (if c.val < r-1 then k else k-1)) \ensuremath{\wedge}\newline       (\ensuremath{\forall} c, (FCCompat.YoungDiagram.toSimpleGraph \ensuremath{\mu}).IsIndepSet\newline         (G c : Set (FCCompat.YoungDiagram.Cell \ensuremath{\mu})) \ensuremath{\wedge} (G c).card = k-r) \ensuremath{\wedge}\newline       (\ensuremath{\forall} c d, c \ensuremath{\neq} d \ensuremath{\to} Disjoint (R c) (R d)) \ensuremath{\wedge}\newline       (\ensuremath{\forall} c d, c \ensuremath{\neq} d \ensuremath{\to} Disjoint (G c) (G d)) \ensuremath{\wedge}\newline       (\ensuremath{\forall} c d, Disjoint (R c) (G d)) \ensuremath{\wedge}\newline       (\ensuremath{\forall} t \ensuremath{\leq} 2*r-1, FCCompat.indepNumK (FCCompat.YoungDiagram.toSimpleGraph \ensuremath{\mu}) t =\newline         if t < r then t*k else t*k-(t-r)*r-1)\end{quote}
{\small Source: proofs/ferrers/OpHack/Ferrers/YoungDiagramDefect.lean:42.\par}

{\small Source SHA-256: c55685043c4b319f834184a218377c836af11200769eeae94fda76c4fc2c2dbc\par}

\end{samepage}
\begin{samepage}
\subsection{FerrersArm.YoungDiagramBridge.first\_two\_r\_independent\_sets}
\begin{quote}\small\ttfamily theorem first\_two\_r\_independent\_sets (\ensuremath{\mu} : YoungDiagram) (r k : \ensuremath{\mathbb{N}})\newline     (hr : 1 \ensuremath{\leq} r) (hk : r \ensuremath{\leq} k)\newline     (hR : FCCompat.indepNumK (FCCompat.YoungDiagram.toSimpleGraph \ensuremath{\mu}) r = r * k)\newline     (hNext : FCCompat.indepNumK (FCCompat.YoungDiagram.toSimpleGraph \ensuremath{\mu})\newline       (r + 1) = (r + 1) * k - r) :\newline     \ensuremath{\exists} R G : Fin r \ensuremath{\to} Finset (FCCompat.YoungDiagram.Cell \ensuremath{\mu}),\newline       (\ensuremath{\forall} c, (FCCompat.YoungDiagram.toSimpleGraph \ensuremath{\mu}).IsIndepSet\newline         (R c : Set (FCCompat.YoungDiagram.Cell \ensuremath{\mu})) \ensuremath{\wedge} (R c).card = k) \ensuremath{\wedge}\newline       (\ensuremath{\forall} c, (FCCompat.YoungDiagram.toSimpleGraph \ensuremath{\mu}).IsIndepSet\newline         (G c : Set (FCCompat.YoungDiagram.Cell \ensuremath{\mu})) \ensuremath{\wedge} (G c).card = k - r) \ensuremath{\wedge}\newline       (\ensuremath{\forall} c d, c \ensuremath{\neq} d \ensuremath{\to} Disjoint (R c) (R d)) \ensuremath{\wedge}\newline       (\ensuremath{\forall} c d, c \ensuremath{\neq} d \ensuremath{\to} Disjoint (G c) (G d)) \ensuremath{\wedge}\newline       (\ensuremath{\forall} c d, Disjoint (R c) (G d)) \ensuremath{\wedge}\newline       (\ensuremath{\forall} t \ensuremath{\leq} 2 * r,\newline         (firstCells R (min t r) \ensuremath{\cup} firstCells G (t - r)).card =\newline           FCCompat.indepNumK (FCCompat.YoungDiagram.toSimpleGraph \ensuremath{\mu}) t)\end{quote}
{\small Source: proofs/ferrers/OpHack/Ferrers/YoungDiagramRestrictedPrefix.lean:158.\par}

{\small Source SHA-256: e5e0e0f73933f360a6fc87e96d2e6172b791acff5ba73aaa0ad0ff836e83654e\par}

{\small\textbf{Sources and attribution:} \href{https://arxiv.org/abs/2408.04086}{1. arXiv source: 2408.04086}; \href{https://www.combinatorics.org/ojs/index.php/eljc/article/download/v32i2p48/pdf/}{2. www.combinatorics.org / pdf}; \href{https://github.com/google-deepmind/formal-conjectures/blob/main/FormalConjectures/Paper/LatinTableau.lean}{3. google-deepmind/formal-conjectures / LatinTableau.lean}; \href{https://github.com/alejandrozu/openmath-2026-judging/blob/d0ed487f487fa7ec814ff705ab55249983fe8707/OpenMath-Judging/review/entries/E02-ferrers-hierarchy.txt}{4. alejandrozu/openmath-2026-judging / E02-ferrers-hierarchy.txt}; \href{https://github.com/alejandrozu/openmath-2026-judging}{Central source and adjudication archive}.\par}

\end{samepage}
\clearpage
\section{A sharp additive-three cover bound for intersecting uniform hypergraphs}
{\small Provisional mathematical authors: Rachel Teo, Wenyi Wang, Hamdi Abderrahmene\par}\medskip
{\small Rachel Teo  -  Sakana AI.\par}

{\small Wenyi Wang  -  Sakana AI.\par}

{\small Hamdi Abderrahmene  -  Sakana AI.\par}

{\small Affiliations follow the recorded author declarations or public institutional/profile sources. Preferred author names, author order and publication affiliations await author review. This editorial compilation does not establish anthology coauthorship.\par}

\subsection{Abstract}
For every finite intersecting r-uniform hypergraph H, 4\ensuremath{\tau}(H)\ensuremath{\leq}|H|+r+3. The degree-at-most-three residual bound closes the explicit additive-three question in Sivashankar's pre-event paper; degree-four peeling yields the unrestricted form and |H|\ensuremath{\geq}3r\ensuremath{-}3 whenever \ensuremath{\tau}(H)=r.

\subsection{Statement and relation to the existing question}
For a finite hypergraph H, \ensuremath{\tau}(H) is the minimum size of a vertex set meeting every edge. Intersecting means that every pair of edges has a common vertex. The new residual theorem asserts 4\ensuremath{\tau}(J)\ensuremath{\leq}|J|+r+3 for an intersecting r-uniform J in which every vertex lies in at most three edges.

The source paper proves the analogous additive-four estimate and explicitly asks whether three is possible. The submitted proof answers that subsidiary question. A standard peeling argument then proves 4\ensuremath{\tau}(H)\ensuremath{\leq}|H|+r+3 without the degree restriction. When \ensuremath{\tau}(H)=r, rearrangement gives |H|\ensuremath{\geq}3r\ensuremath{-}3, and in particular |H|\ensuremath{\geq}21 at r=8. The triangle example supplies sharpness of the additive constant.

\subsection{The residual packing improvement}
The residual proof begins with the source paper's matching and witness machinery. A maximum matching organizes the edges and the uncovered remainder. Its associated link graphs form a three-link packing, encoding the ways covered vertices can meet the uncovered ground set.

The older count leaves a potential equality case consistent with additive four. The new coloured-link budget proves a strict support-sum inequality when the uncovered ground set is odd and large enough relative to the number of packing blocks. The proof separates small configurations and develops structural/equality exclusions for the remaining coloured packings. Combining that strict inequality with the original selected-count identity rules out exactly the obstruction to additive three.

SharpEndpoint initially presents the packing inequality as an explicit argument; the final Main endpoint supplies it from the Packing and Colored modules. This distinction matters: reading the bridge alone could suggest an assumed lemma, while the selected final closure is intended to be unconditional. The independent archived replay reports 18 Main-closure modules and a separate endpoint-definition audit, without importing the Kahn-dependent asymptotic section.

\subsection{Peeling and the general consequence}
If some vertex v has degree at least four, delete every edge containing v and call the remaining hypergraph H\ensuremath{^{\prime}}. It stays intersecting and r-uniform. A minimum cover of H\ensuremath{^{\prime}}, together with v, covers H, so \ensuremath{\tau}(H)\ensuremath{\leq}\ensuremath{\tau}(H\ensuremath{^{\prime}})+1. The deletion removes at least four edges, exactly compensating for the factor four in the cover inequality. Induction on the number of edges lifts the residual additive-three bound to all H.

If no vertex has degree at least four, the residual theorem applies directly. This yields the unrestricted cover inequality and its lower-bound corollaries in the same family. The r=8 specialization and triangle sharpness should be presented as consequences, not independent breakthroughs.

\subsection{Attribution, scope and publication}
Varun Sivashankar's pre-event paper supplies the +4 theorem, matching/witness infrastructure and sharpness examples. Its ancillary Lean prefix is reused with explicit attribution and documented compatibility adjustments. The proposed original delta is the strict packing improvement to +3 and the resulting all-r elementary bound.

The original O(r) Erdős 21 question was already resolved by Kahn, and the source paper's stronger approximately 3.053r asymptotic lower coefficient is not improved here. A title claiming a full solution of Erdős 21 would be misleading. A focused paper answering the explicit residual question, with the supplementary Lean closure and clear source attribution, is the right publication form.

\subsection{The sharp residual inequality and its relation to prior work}
Let H be a finite simple r-uniform intersecting hypergraph and let \ensuremath{\tau}(H) be its minimum transversal size. The result is 4\ensuremath{\tau}(H)\ensuremath{\leq}|H|+r+3. In particular, if \ensuremath{\tau}(H)=r then |H|\ensuremath{\geq}3r\ensuremath{-}3. The proof strengthens the constant 4 in the finite residual inequality of Varun Sivashankar, arXiv:2606.24878v2. Its auxiliary matching and witness setup is credited to that source. The sharper asymptotic coefficient in Sivashankar's work is a separate result and is not improved by this finite additive refinement.

We first prove the residual inequality when every vertex has degree at most 3. For any edge E, intersectingness gives\ensuremath{\Sigma}\_\{F\ensuremath{\in}H\}|E\ensuremath{\cap}F|\ensuremath{\geq}r+|H|\ensuremath{-}1, whereas the degree bound gives\ensuremath{\Sigma}\_\{v\ensuremath{\in}E\}deg(v)\ensuremath{\leq}3r. Thus |H|\ensuremath{\leq}2r+1. This elementary edge-count bound will be used to identify the sole possible equality obstruction.

\subsection{Good triples, witnesses and the selected-incidence inequality}
A good triple is a set of three hyperedges sharing a vertex. In degree at most 3 it is exactly the three-edge block through that vertex. Choose a maximum collection M of pairwise disjoint good triples, where disjointness concerns hyperedge indices, not disjoint hyperedges in the original intersecting hypergraph. Write t=|M|, let U be the unmatched hyperedges, and put u=|U|. Then |H|=3t+u. One vertex covers each matched triple; pairing the remaining intersecting edges covers U in at most\ensuremath{\lceil}u/2\ensuremath{\rceil} vertices. Hence \ensuremath{\tau}(H)\ensuremath{\leq}t+\ensuremath{\lceil}u/2\ensuremath{\rceil}.

Maximum matching leaves no good triple entirely in U. Choose a common vertex \ensuremath{\omega}(\{E,F\}) for each pair of unmatched edges. The chosen vertices are distinct: if two different unmatched pairs had the same witness, their union would contain at least three edges at that vertex, producing a good triple in U. Every selected witness has exactly two incidences with U and at most one with the matched edges.

For each matched hyperedge p, form a link graph G\_p on vertex set U: put the pair \{E,F\} in G\_p exactly when its witness belongs to the three-edge block \{E,F,p\}. No pair belongs to two colour graphs. Within one matched triple \{p,q,r\}, every edge of G\_p meets every edge of G\_q, and similarly for the other two pairs of colours. Otherwise two disjoint unmatched pairs with those colours would yield two disjoint good triples replacing the one original matched triple, increasing the matching. This is an exchange argument in the auxiliary triple system.

Let S count matched/unmatched pairs whose intersection is witnessed by a selected vertex. It is at most the sum of the supports of all G\_p: a selected witness of colour p contributes its two unmatched endpoints to that support. The selected vertices contribute u(u\ensuremath{-}1) incidences with U, twice the number of unmatched pairs. There are 3tu matched/unmatched pairs altogether. An unselected vertex with a incidences in U and b in matched edges satisfies a+b\ensuremath{\leq}3 and a\ensuremath{\leq}2, so ab\ensuremath{\leq}2a. Therefore unselected U incidences are at least(3tu\ensuremath{-}S)/2. Total U incidences equal ur, giving the precise inequality 2u(u\ensuremath{-}1)+3tu\ensuremath{\leq}2ur+S.

\subsection{Lemma: bounds for three cross-intersecting colour graphs}
Consider three edge-disjoint simple graphs on a common u-point ground set. Edges from different colours must intersect. Let s be the sum of their three support sizes and e the total number of edges. The elementary bounds are s\ensuremath{\leq}2e, s\ensuremath{\leq}max(12,u+4), and s\ensuremath{\leq}max(u,12,e+4). We give the structural proof, since the refined bound is what removes the equality case.

If one colour contains two disjoint edges, let W be their four endpoints. Every edge of the other two colours lies in W, since a two-element edge meeting both disjoint pairs has one endpoint in each. If those other edges have no common vertex, every edge of the first colour also lies in W: an edge with an endpoint outside W would need its other endpoint in every edge of the other colours. All three supports then have size at most 4, giving s\ensuremath{\leq}12.

If the union of the other two colours has a common vertex, it has at most two edges: its other endpoints must lie in one of the fixed disjoint pairs not containing the centre. Thus their total support contribution is at most 4, giving s\ensuremath{\leq}u+4. Choose one of those edges as a two-point hitting set for the first colour; every first-colour edge contributes at most one new support point outside it. Hence its support is at most its edge count plus 2. Adding at most twice the other two edge counts gives s\ensuremath{\leq}e+4. If both other colours are empty, the first colour alone gives s\ensuremath{\leq}u.

If every colour is internally intersecting, each is a star or a triangle. When every support has size at most 3, s\ensuremath{\leq}9. Otherwise a colour is a star with at least three distinct leaves. Every edge of another colour must contain the same centre to meet all three star edges. The whole edge union is then a single star, with at mostu\ensuremath{-}1 edges; summing the three coloured supports gives s\ensuremath{\leq}e+3\ensuremath{\leq}u+2. This completes all the claimed bounds, including empty colour classes.

\subsection{The large-ground packing barrier}
Apply the colour lemma separately to each matched triple. More generally, let a three-link packing have t blocks of three colours on a u-point ground set, with global edge disjointness. Write s\_B,e\_B for its block support and edge counts, S=\ensuremath{\Sigma}\_Bs\_B and E=\ensuremath{\Sigma}\_Be\_B. Global edge disjointness gives E\ensuremath{\leq}u(u\ensuremath{-}1)/2.

For u\ensuremath{\geq}10 the three local bounds imply us\_B\ensuremath{\leq}u\ensuremath{^2}+4e\_B. If s\_B\ensuremath{\leq}u+2, use s\_B\ensuremath{\leq}2e\_B and(u\ensuremath{-}2)s\_B\ensuremath{\leq}u\ensuremath{^2}\ensuremath{-}4. If s\_B\ensuremath{\geq}u+3, it is larger than 12 and u, so the refined bound forces s\_B\ensuremath{\leq}e\_B+4. Together with s\_B\ensuremath{\leq}u+4, this gives us\_B\ensuremath{-}4e\_B\ensuremath{\leq}(u\ensuremath{-}4)s\_B+16\ensuremath{\leq}u\ensuremath{^2}. Summing over blocks yields uS\ensuremath{\leq}tu\ensuremath{^2}+4E\ensuremath{\leq}tu\ensuremath{^2}+2u(u\ensuremath{-}1)<u\ensuremath{^2}(t+2). Hence S<u(t+2).

Only odd positive u below 10 can remain. Under t+3\ensuremath{\leq}u and S\ensuremath{\geq}u(t+2), the old support bound S\ensuremath{\leq}t max(12,u+4) reduces the possibilities to(u,t)=(5,2),(7,3),(7,4),(9,5),(9,6). The next section excludes them structurally, rather than relying on an unrecorded brute-force search.

\subsection{The small packing cases}
We need two elementary consequences of support saturation. If a block has s\_B=2e\_B, every colour graph is a matching. If e\_B\ensuremath{\geq}4 and u<2e\_B, at least two colours are nonempty; no colour can have three disjoint edges, because a two-element edge of another colour cannot meet all three. Some colour therefore has exactly two disjoint edges and all other edges lie in their four endpoints. The whole block graph has support at most 4 and e\_B\ensuremath{\leq}6. A saturated six-edge block is exactly K4.

On five ground points, failure of the barrier forces t=2, E=10 and S=20. Both blocks are saturated, have at most 6 edges and hence have at least 4 each. Their union graphs each have support at most 4, but together cover all pairs of K5. That is impossible: for any vertex v, the supports containing v must cover all five vertices through pairs with v. One support of size 4 is insufficient, so v must lie in both supports. Each support would contain all five vertices.

On seven points with four blocks, failure forces E=21,S=42 and saturation of every block. If all four blocks have at least 4 edges, all four supports have size at most 4 and cover all pairs of K7. For each vertex v let m(v) be the set of supports containing it. Every vertex lies in at least two supports, and the total memberships are at most 16, so at least five vertices have exactly two memberships. Distinct such vertices have distinct membership pairs: their two supports already cover all seven vertices, and sharing a second vertex would make their union at most 6. Moreover the membership pair of a twice-covered vertex is not contained in the membership set of any other vertex. There are only six pairs of four supports; at least five distinct pairs occur. Some seventh vertex has at least three memberships, excluding at least three of those pairs. Only three pairs could remain, a contradiction.

If one of the four seven-point blocks has at most 3 edges, the other three must each have 6 to reach 21; they are saturated K4s. Two edge-disjoint K4s on seven vertices share exactly one vertex and cover the ground. No triangle can avoid both: a triangle can meet each support in at most one vertex, yet their union covers the ground. In particular a third disjoint K4 is impossible.

On seven points with three blocks, failure gives total support at least 35, while each block has support at most 12. Thus two blocks have support 12 and the third at least 11. A critical three-colour block with s>max(9,u+4) is confined to four points by the colour case analysis above. With s\ensuremath{\geq}11 it has at least 6 edges, so it is K4. The first two blocks are therefore K4s covering all seven vertices and sharing one vertex z. Any edge avoiding their edges must avoid z. The third block lives on the remaining six points; its support at least 11 exceeds max(9,6+4), so it too is K4, contradicting the same obstruction.

On nine points the local bounds imply 3s\_B\ensuremath{\leq}e\_B+30. For s\_B\ensuremath{\leq}12 this follows from e\_B\ensuremath{\geq}s\_B/2; for s\_B=13 the refined bound gives e\_B\ensuremath{\geq}9. Summing and using E\ensuremath{\leq}36 gives 3S\ensuremath{\leq}36+30t. Failure with t\ensuremath{\leq}6 forces t=6,S=72,E=36 and equality in every local inequality. Equality and s\_B\ensuremath{\leq}2e\_B give s\_B\ensuremath{\geq}12, so all six blocks have s\_B=12,e\_B=6. They are saturated K4s partitioning K9. In any K4 edge partition each vertex degree is divisible by 3, whereas every vertex of K9 has degree 8. This contradiction also excludes the nominal t=5 case. The sharp odd-ground support barrier is now proved for every required size.

\subsection{Removing the equality obstruction and lifting by peeling}
The baseline selected-incidence bound with S\ensuremath{\leq}t max(12,u+4), together with 3t+u\ensuremath{\leq}2r+1, gives u+t\ensuremath{\leq}r+2. Here is the short arithmetic reduction: if u+t\ensuremath{\geq}r+3 then the edge-count bound forces u\ensuremath{\geq}t+5. For u\ensuremath{\geq}8 substitute max(12,u+4)=u+4 into the counting inequality and obtain a contradiction; the cases 5\ensuremath{\leq}u\ensuremath{\leq}7 are immediate integer substitutions. This baseline inequality is part of the credited matching setup.

Suppose 4\ensuremath{\tau}(H)>|H|+r+3. If u=0, the matched cover bound already contradicts this. For positive u, combine \ensuremath{\tau}(H)\ensuremath{\leq}t+\ensuremath{\lceil}u/2\ensuremath{\rceil} and u+t\ensuremath{\leq}r+2. An even u cannot fail the desired inequality. For odd u, failure forces u+t=r+2. The degree-count bound then gives t+3\ensuremath{\leq}u. Substituting r=u+t\ensuremath{-}2 into 2u(u\ensuremath{-}1)+3tu\ensuremath{\leq}2ur+S yields S\ensuremath{\geq}u(t+2). But the actual link graphs form the packing just analyzed, whose support sum is strictly belowu(t+2), and S is at most that support sum. This contradiction proves the sharp residual inequality.

Finally remove the degree-at-most 3 restriction by induction on the number of hyperedges. If a vertex has degree d\ensuremath{\geq}4, delete all edges through it to obtain H\ensuremath{^{\prime}}. A minimum cover of H\ensuremath{^{\prime}} together with that vertex covers H, so \ensuremath{\tau}(H)\ensuremath{\leq}\ensuremath{\tau}(H\ensuremath{^{\prime}})+1. Uniformity and intersectingness pass to H\ensuremath{^{\prime}}, and |H|=|H\ensuremath{^{\prime}}|+d. The inductive bound gives 4\ensuremath{\tau}(H)\ensuremath{\leq}|H\ensuremath{^{\prime}}|+r+3+4\ensuremath{\leq}|H|+r+3. If no such vertex exists, the residual theorem applies. Substituting \ensuremath{\tau}(H)=r gives 3r\ensuremath{\leq}|H|+3.

The additive constant 3 is sharp for the residual inequality: the three edges of a triangle form an intersecting2-uniform hypergraph with \ensuremath{\tau}=2, and 4\ensuremath{\tau}=8=3+2+3. Sharpness of this residual constant should not be confused with equality examples for the bound g(r)\ensuremath{\geq}3r\ensuremath{-}3 at every rank. The source final endpoint instantiates the packing theorem unconditionally; the earlier SharpEndpoint bridge displays it as a hypothesis only to separate modules.

\subsection{Verification and reproducibility}
Fresh execution: sakana-erdos21-sharp-residual: PASS; 18/18 custom modules compiled successfully; receipt Verification/sakana-erdos21-sharp-residual-fresh-build.json. Compiler pins, source hashes and endpoint axiom outputs are in Verification. Historical receipts and finite checks do not replace fresh replay.

\begin{samepage}
\subsection{ErdosLovasz.residual\_cover\_bound\_sharp}
\begin{quote}\small\ttfamily theorem residual\_cover\_bound\_sharp (r : \ensuremath{\mathbb{N}}) (J : Finset (Finset V))\newline     (huni : IsUniform r J) (hint : Intersecting J)\newline     (hdeg : \ensuremath{\forall} v, edgeDeg J v \ensuremath{\leq} 3) :\newline     4 * coverNum J \ensuremath{\leq} J.card + r + 3\end{quote}
{\small Source: proofs/erdos21-sharp-residual/OpHack/Erdos21/Main.lean:25.\par}

{\small Source SHA-256: 9314ae7ba95b6fd9922ccfea882f1d33a6c11c9a40d750109014c8184ab7e414\par}

\end{samepage}
\begin{samepage}
\subsection{ErdosLovasz.cover\_bound\_sharp}
\begin{quote}\small\ttfamily theorem cover\_bound\_sharp (r : \ensuremath{\mathbb{N}}) (H : Finset (Finset V))\newline     (huni : IsUniform r H) (hint : Intersecting H) :\newline     4 * coverNum H \ensuremath{\leq} H.card + r + 3\end{quote}
{\small Source: proofs/erdos21-sharp-residual/OpHack/Erdos21/Main.lean:33.\par}

{\small Source SHA-256: 9314ae7ba95b6fd9922ccfea882f1d33a6c11c9a40d750109014c8184ab7e414\par}

\end{samepage}
\begin{samepage}
\subsection{ErdosLovasz.erdos\_lovasz\_lower\_three}
\begin{quote}\small\ttfamily theorem erdos\_lovasz\_lower\_three (r : \ensuremath{\mathbb{N}}) (H : Finset (Finset V))\newline     (huni : IsUniform r H) (hint : Intersecting H) (hcover : coverNum H = r) :\newline     3 * r \ensuremath{\leq} H.card + 3\end{quote}
{\small Source: proofs/erdos21-sharp-residual/OpHack/Erdos21/Main.lean:41.\par}

{\small Source SHA-256: 9314ae7ba95b6fd9922ccfea882f1d33a6c11c9a40d750109014c8184ab7e414\par}

\end{samepage}
\begin{samepage}
\subsection{ErdosLovasz.erdos\_lovasz\_lower\_three\_sub}
\begin{quote}\small\ttfamily theorem erdos\_lovasz\_lower\_three\_sub (r : \ensuremath{\mathbb{N}}) (H : Finset (Finset V))\newline     (huni : IsUniform r H) (hint : Intersecting H) (hcover : coverNum H = r) :\newline     3 * r - 3 \ensuremath{\leq} H.card\end{quote}
{\small Source: proofs/erdos21-sharp-residual/OpHack/Erdos21/Main.lean:49.\par}

{\small Source SHA-256: 9314ae7ba95b6fd9922ccfea882f1d33a6c11c9a40d750109014c8184ab7e414\par}

\end{samepage}
\begin{samepage}
\subsection{ErdosLovasz.erdos\_lovasz\_eight\_lower}
\begin{quote}\small\ttfamily theorem erdos\_lovasz\_eight\_lower (H : Finset (Finset V))\newline     (huni : IsUniform 8 H) (hint : Intersecting H) (hcover : coverNum H = 8) :\newline     21 \ensuremath{\leq} H.card\end{quote}
{\small Source: proofs/erdos21-sharp-residual/OpHack/Erdos21/Main.lean:56.\par}

{\small Source SHA-256: 9314ae7ba95b6fd9922ccfea882f1d33a6c11c9a40d750109014c8184ab7e414\par}

{\small\textbf{Sources and attribution:} \href{https://arxiv.org/html/2606.24878v2}{1. arXiv source: 2606.24878v2}; \href{https://arxiv.org/abs/2606.24878}{2. arXiv source: 2606.24878}; \href{https://www.erdosproblems.com/21}{3. www.erdosproblems.com / 21}; \href{https://arxiv.org/src/2606.24878v2/anc/lean-proof/RequestProject/Theorem1.lean}{4. arXiv source: src/2606.24878v2/anc/lean-proof/RequestProject/Theorem1.lean}; \href{https://github.com/alejandrozu/openmath-2026-judging/blob/d0ed487f487fa7ec814ff705ab55249983fe8707/OpenMath-Judging/review/entries/E02-erdos21-sharp-residual.txt}{5. alejandrozu/openmath-2026-judging / E02-erdos21-sharp-residual.txt}; \href{https://github.com/alejandrozu/openmath-2026-judging}{Central source and adjudication archive}.\par}

\end{samepage}
\clearpage
\section{A four-term-progression-free set with reciprocal sum exceeding 4.44}
{\small Provisional mathematical authors: Rachel Teo, Wenyi Wang, Hamdi Abderrahmene\par}\medskip
{\small Rachel Teo  -  Sakana AI.\par}

{\small Wenyi Wang  -  Sakana AI.\par}

{\small Hamdi Abderrahmene  -  Sakana AI.\par}

{\small Affiliations follow the recorded author declarations or public institutional/profile sources. Preferred author names, author order and publication affiliations await author review. This editorial compilation does not establish anthology coauthorship.\par}

\subsection{Abstract}
An explicit six-residue augmentation of the base 55 digit construction gives a finite positive four-AP-free set whose exact reciprocal sum exceeds 111/25. This improves the located lower-bound construction while remaining a finite bound, not a reciprocal-divergence theorem.

\subsection{One explicit set for all three assertions}
Let D=\{0,1,2,4,5,9,10,11,14,16,17,18,21,24,30,37,39,41,42,45,47\}. Define T\_j as the nonnegative integers with j base 55 digits drawn from D, allowing leading zeroes, and T\_0=\{0\}. Put M=55\ensuremath{^3} and U=\{26286,26726,26737,120061,120501,120512\}. Set E=T\_22\ensuremath{\cup}\ensuremath{\bigcup}\_(u\ensuremath{\in}U)(u+M T\_21), and A=\{e+1:e\ensuremath{\in}E\}.

The theorem proves that every member of A is positive, A contains no nonconstant four-term arithmetic progression over the integers, and the finite rational sum \ensuremath{\sum}\_(a\ensuremath{\in}A)1/a is strictly greater than 111/25=4.44. These clauses concern the same set of distinct elements. Separate large harmonic sums and AP-free witnesses would not establish the result.

\subsection{Digit descent and augmentation}
The known base 55 alphabet supplies the starting AP-free digital construction. Reducing a hypothetical progression modulo 55 and descending through the digits proves the absence of a nonconstant progression in the digit-restricted part. The six extra residue classes are chosen and checked so that the combined set remains free of forbidden progressions across its pieces.

The formal development includes exact small arithmetic certificates, residue-fibre transport and harmonic estimates. It avoids enumerating an astronomically large set element by element. Grouped levels and rational lower bounds certify the harmonic sum while finset unions preserve distinctness. The finite certificate stages and the original digit alphabet are supporting steps, not additional discoveries.

\subsection{Improvement and limits}
The located Walker construction gives a reciprocal sum above 4.43975, and a later located bound is approximately 4.4397534742. Exceeding 4.44 is a small but explicit improvement. The paper should show where the six-residue augmentation adds mass and why its interactions do not create a four-term progression.

This result does not give unbounded reciprocal sums and does not refute the Erdős-Turán reciprocal-divergence problem. It is a quantitative lower construction for the associated extremal estimation problem. No optimality of the augmentation or global supremum is claimed.

\subsection{Formal endpoint and manuscript form}
The selected endpoint Erdos3Candidate.explicit\_four\_ap\_free\_reciprocal\_gt states positivity, APFree 4 and the rational inequality for one set. On 5 October 2026 at 03:02:16 UTC, an independently checked operational replay passed all 22 authored scientific bodies and the unchanged final audit under Lean 4.34.1 and Mathlib d13f23b723b8a846827a245b89c10fc7d3f11612. Seven broad import headers were replaced by 29 exact pinned tactic and mathematical imports; every statement, proof, option and comment outside those import spans remained byte-identical. The actual final print uses only propext, Classical.choice and Quot.sound, with no admitted, native or unknown axioms. This verifies the selected theorem through the separately labeled import-only route. The original frozen-source replay remains a historical 15-of-22 resource-limited outcome, and the first narrowed route remains a historical 7-of-22 resource-limited outcome; their receipts have not been relabeled or overwritten. The entrants' earlier full Linux archive is another distinct record. Completed route and exact body/output/dependency checks: Verification/sakana-erdos169-fourap-tactic-specific-v2-fresh-build.json and Verification/e169-tactic-specific-v2-completed-qualification.json.

A concise paper should give the explicit alphabet, offsets and block scheme, a human digit-descent proof, and a compact rational-sum certificate. The supplementary repository can carry the full formal derivation. Walker and the intervening augmentation must be credited, with the proposed new improvement restricted to the exact set and 4.44 threshold.

\subsection{An explicit finite set with reciprocal sum greater than 4.44}
Let D=\{0,1,2,4,5,9,10,11,14,16,17,18,21,24,30,37,39,41,42,45,47\}. For j\ensuremath{\geq}0 let T\_j be the nonnegative integers with at most j base-55 digits, all in D; leading zeros are allowed and T\_0=\{0\}. Set M=55\ensuremath{^3}=166375 and U=\{26286,26726,26737,120061,120501,120512\}. Define E=T\_22\ensuremath{\cup}\ensuremath{\bigcup}\_\{u\ensuremath{\in}U\}(u+M T\_21), and A=\{e+1:e\ensuremath{\in}E\}. The result is that A is a finite set of positive integers containing no nonconstant four-term arithmetic progression, and \ensuremath{\sum}\_\{a\ensuremath{\in}A\}1/a>111/25=4.44.

This is a quantitative improvement within a known base-digit construction, credited to Walker in the submitted comparison. It supplies six compatible residue fibers and a certified harmonic increase. A large but finite reciprocal sum is not a divergent reciprocal sum, so this does not settle the original qualitative Erdős problem. The explicit formulas define every element without attempting to print or enumerate the enormous set: |A|=21\ensuremath{^2}\ensuremath{^2}+6\ensuremath{\cdot}21\ensuremath{^2}\ensuremath{^1}=27\ensuremath{\cdot}21\ensuremath{^2}\ensuremath{^1}.

\subsection{Lemma 1: modularly free digits descend to ordinary progressions}
The first exact certificate is that D contains no nonconstant four-term progression modulo 55. It can be checked by testing all pairs a\ensuremath{\in}D and d\ensuremath{\in}\{1,...,54\} and requiring that a,a+d,a+2d,a+3d modulo 55 all belong to D. The resulting list is empty. This tests cyclic progressions, including wraparound, rather than only four-term progressions lying in the interval [0,54].

Suppose a,a+d,a+2d,a+3d all belong to T\_j. Their residues modulo 55 lie in D, so the modular certificate forces 55|d. Their residues are then equal to a single digit r\ensuremath{\in}D. Write a=r+55b and d=55e. Division by 55 gives b,b+e,b+2e,b+3e in T\_\{j\ensuremath{-}1\}. If d\ensuremath{\neq}0 then e\ensuremath{\neq}0. Descending to T\_0=\{0\} is impossible. Hence every T\_j is four-progression-free. The same descent for three digits shows that O=T\_3, regarded as a residue alphabet modulo M, has no nonzero-step modular four-term progression.

\subsection{Lemma 2: the six residues are jointly compatible with the old alphabet}
Each u\ensuremath{\in}U lies strictly between zero and M and is outside O. Its units digit is 7 or 51; its lower two-digit residue is 2086, 2526 or 2537; its third digit is 8 or 39. We prove that O\ensuremath{\cup}U is modularly four-progression-free modulo M. Consider a hypothetical progression a+kt, k=0,1,2,3, with t nonzero modulo M. If it uses no new residue, the preceding old-digit descent already excludes it.

If it uses precisely one new residue u at an endpoint, reduction modulo 55 asks for three consecutive old digits after 7 or after 51. For x=7 and x=51, the set of d\ensuremath{\in}\{0,...,54\} such that x+d,x+2d,x+3d modulo 55 all lie in D is empty. The opposite endpoint follows by reversing the progression. Thus only a single new residue at an interior position remains.

Reverse the progression if necessary so that u is at position one. Then u\ensuremath{-}t,u+t,u+2t are old. Reducing first modulo 55 and then modulo 3025 gives the following complete possible lower-two-digit differences, with the order irrelevant: at residue 2086 they are 1049,1489,1544,1054,1494,1549,1059,1499,1554; at 2526 they are 1489,1494,1499; at 2537 they are 497,1377,502,1382,1502,492,1372. These lists are obtained by testing the 55 possible units differences, lifting each by 55h with 0\ensuremath{\leq}h<55, and testing all three old two-digit memberships.

Write t=d+3025h, using a listed low difference d. Direct integer division for every u and its surviving d gives (u\ensuremath{-}d)//3025=z, (u+d)//3025=z+1 and (u+2d)//3025=z+1, where z=u//3025 is 8 or 39. Consequently the third digits of the three old terms must be z\ensuremath{-}h,z+1+h,z+1+2h modulo 55. For each z\ensuremath{\in}\{8,39\}, the set of h\ensuremath{\in}\{0,...,54\} making all three lie in D is empty. This completes the one-new case, and makes the carry bridge explicit rather than assuming that digitwise progression equations have no carries.

Finally suppose two or more positions i<j use new residues u,v. The gaps j\ensuremath{-}i\ensuremath{\in}\{1,2,3\} are invertible modulo M; their inverses are respectively 1,83188,110917. If u=v, the invertible gap would force t=0 modulo M, a contradiction. Otherwise t is uniquely determined as (v\ensuremath{-}u)/(j\ensuremath{-}i) modulo M, and a=u\ensuremath{-}it modulo M. There are only 6\ensuremath{\cdot}5\ensuremath{\cdot}6=180 active ordered residue-position choices. For each, compute the four residues a+kt modulo M and test membership in the joint enlarged alphabet O\ensuremath{\cup}U. None passes. The joint membership test is essential: testing the old alphabet alone would not exclude a progression with three or four new terms. These exact checks, together with the preceding exhaustive small lists, prove the lemma.

\subsection{Lemma 3: lift the modular certificate through the actual fibers}
Every old element of T\_22 has a unique decomposition r+Mt with r\ensuremath{\in}T\_3 and t\ensuremath{\in}T\_19. Every new element has unique decomposition u+Mt with u\ensuremath{\in}U and t\ensuremath{\in}T\_21. The residue sets T\_3 and U are disjoint, as are the six distinct new residues. Thus the union is disjoint and each fiber parametrization is injective. These facts justify both the stated cardinality and later addition of reciprocal sums.

If E contained a nonconstant four-term progression with difference d, its four residues would lie in O\ensuremath{\cup}U. Modular freeness forces M|d. All four terms therefore lie in one canonical residue fiber. Dividing by M gives a nonconstant progression in T\_19 if the residue is old, or in T\_21 if the residue is new. Both are excluded by Lemma 1. Thus E is four-progression-free. Translating every term by one preserves progression freeness and makes every element positive, proving the corresponding assertions for A.

\subsection{Lemma 4: a harmonic-mean estimate for an entire suffix block}
Uniqueness of base-55 expansions gives |T\_j|=21\textasciicircum{}j. The sum of the digits in D is 433. Averaging the independent digit choices therefore gives the mean of T\_j as \ensuremath{\gamma}(55\textasciicircum{}j\ensuremath{-}1), where \ensuremath{\gamma}=433/(21\ensuremath{\cdot}54)=433/1134. Put c=454/1155. For j\ensuremath{\geq}1 and P=55\textasciicircum{}j\ensuremath{\geq}55, 1+\ensuremath{\gamma}(P\ensuremath{-}1)\ensuremath{\leq}cP; equality occurs at P=55, and the difference increases thereafter because c>\ensuremath{\gamma}.

For a fixed nonnegative prefix p, apply the finite harmonic-mean inequality \ensuremath{\sum}1/z\_i\ensuremath{\geq}N\ensuremath{^2}/\ensuremath{\sum}z\_i to the positive denominators z\_x=p55\textasciicircum{}j+x+1, x\ensuremath{\in}T\_j. Their mean is at most 55\textasciicircum{}j(p+c). Hence \ensuremath{\sum}\_\{x\ensuremath{\in}T\_j\}1/(p55\textasciicircum{}j+x+1)\ensuremath{\geq}(21/55)\textasciicircum{}j/(p+c). This inequality is an ordinary consequence of Cauchy-Schwarz, so it controls a huge suffix block from its cardinality and first moment without enumerating its members.

For a new fiber, u+1\ensuremath{\leq}M. Consequently 1/[M(p55\textasciicircum{}j+x)+u+1]\ensuremath{\geq}1/[M(p55\textasciicircum{}j+x+1)]. The same block estimate divided by M gives \ensuremath{\sum}\_\{x\ensuremath{\in}T\_j\}1/[M(p55\textasciicircum{}j+x)+u+1]\ensuremath{\geq}(21/55)\textasciicircum{}j/[M(p+c)]. Only these positive-length suffix layers use the uniform replacement of u; the singleton prefix contributions below keep the actual six offsets.

\subsection{An exact rational certificate for the final harmonic increase}
Partition T\_22 into its two-digit beginning T\_2 and, for each j=1,...,20, the layers p55\textasciicircum{}j+T\_j, where p=55a+b with a\ensuremath{\in}D\textbackslash{}\{0\} and b\ensuremath{\in}D. This is the partition by highest two-digit prefix. Similarly partition T\_21 into D and, for j=1,...,20, the layers d55\textasciicircum{}j+T\_j with d\ensuremath{\in}D\textbackslash{}\{0\}. The partitions are disjoint by the highest nonzero digit. Let G=\ensuremath{\sum}\_\{j=1\}\textasciicircum{}\{20\}(21/55)\textasciicircum{}j. Applying Lemma 4 and the exact singleton sums gives a lower bound for the reciprocal sum consisting of four small coefficients: old prefix, old tail, new prefix and new tail.

For every positive rational term N/D, floor(QN/D)/Q is a downward bound. At Q=10\textasciicircum{}12 the exact old-prefix floor sum over 55a+b is 3875802807082 and the old-tail floor sum over a\ensuremath{\neq}0 is 913030252293. The latter uses denominators 1155(55a+b)+454 and numerator 1155\ensuremath{\cdot}10\textasciicircum{}12. At Q=10\textasciicircum{}15 the new-prefix floor sum over u\ensuremath{\in}U,d\ensuremath{\in}D, keeping denominator Md+u+1, is 222203111618. The simple uniform new-tail floor sum over u\ensuremath{\in}U,d\ensuremath{\in}D\textbackslash{}\{0\} is 84575370288, with denominator M(1155d+454) and numerator 1155\ensuremath{\cdot}10\textasciicircum{}15.

Therefore \ensuremath{\sum}\_\{a\ensuremath{\in}A\}1/a is at least L=[3875802807082000+222203111618]/10\textasciicircum{}15+[913030252293000+84575370288]G/10\textasciicircum{}15. This exact rational number is greater than 111/25. Its decimal value, provided only for interpretation, is about 4.440007695422332. The proof of the strict comparison uses the finite rational geometric sum, not floating-point rounding. An alternative sharper new-tail certificate appears in the source, but the simpler one already crosses the required threshold and is the one used by the final unconditional endpoint.

The independent checker accompanying this preparation recomputes the base modular exclusions, every candidate list, the carry bridge, all 180 two-new cases, all four floor sums and the final strict rational comparison. Its result is an arithmetic audit, separate from the independent Lean replay. The submitted final endpoint joins positivity, actual finite-set progression freeness and the reciprocal inequality. Early finite-certificate comments must be read together with the later source assembly. The archived Linux and partial original-source Windows records are distinct from the completed, separately qualified scientific-body-identical import-only audit.

\subsection{Current independent verification and the exact finite certificate}
A separate portable exact-integer check completed successfully for the joint 9,267-residue alphabet modulo 166,375: it tested 2,970 base-digit progressions, 1,996,488 one-new-residue cases and all 180 active two-new residue-position choices. Old-only cases follow by the displayed least nonzero base-55 digit descent. The earlier independent rational checker also preserves the strict 111/25 certificate. These arithmetic checks strengthen the ordinary finite-certificate proof while remaining distinct from a completed final Lean axiom audit.

\subsection{Verification and reproducibility}
Fresh execution: sakana-erdos169-fourap: ENVIRONMENT\_BLOCKED\_RESOURCE\_DEPENDENCIES; 15/22 custom modules compiled successfully; 2 unfinished invocations stopped at a recorded resource checkpoint; receipt Verification/sakana-erdos169-fourap-fresh-build.json. Separate qualified import-only route: all 22 unchanged scientific proof bodies and the unchanged final audit pass, with standard kernel axioms only; seven import headers were narrowed under the same source and compiler pins. This does not replace the original frozen-source receipt. Qualification: Verification/e169-tactic-specific-v2-completed-qualification.json. Compiler pins, source hashes and endpoint axiom outputs are in Verification. Historical receipts and finite checks do not replace fresh replay.

\begin{samepage}
\subsection{Erdos3Candidate.explicit\_four\_ap\_free\_reciprocal\_gt}
\begin{quote}\small\ttfamily theorem explicit\_four\_ap\_free\_reciprocal\_gt :\newline     (\ensuremath{\forall} a \ensuremath{\in} sixAFinset, 0 < a) \ensuremath{\wedge}\newline       APFree 4 (sixAFinset : Set \ensuremath{\mathbb{Z}}) \ensuremath{\wedge}\newline       (\ensuremath{\sum} a \ensuremath{\in} sixAFinset, (1 : \ensuremath{\mathbb{Q}}) / (a : \ensuremath{\mathbb{Q}})) > 111 / 25\end{quote}
{\small Source: proofs/erdos169-fourap/OpHack/MainResult.lean:17.\par}

{\small Source SHA-256: cc686e49443c9f23ccbf715be6e9cf71c528736794df230771d33ec84df80ac5\par}

{\small\textbf{Sources and attribution:} \href{https://www.erdosproblems.com/169}{1. www.erdosproblems.com / 169}; \href{https://arxiv.org/abs/2203.06045}{2. arXiv source: 2203.06045}; \href{https://computoergosum.com/en/principia/erdos-169.html}{3. computoergosum.com / erdos-169.html}; \href{https://github.com/alejandrozu/ulam-opdp-difficulty-atlas}{4. alejandrozu/ulam-opdp-difficulty-atlas}; \href{https://github.com/alejandrozu/openmath-2026-judging/blob/d0ed487f487fa7ec814ff705ab55249983fe8707/OpenMath-Judging/review/entries/E02-erdos169-fourap.txt}{5. alejandrozu/openmath-2026-judging / E02-erdos169-fourap.txt}; \href{https://github.com/alejandrozu/openmath-2026-judging}{Central source and adjudication archive}.\par}

\end{samepage}
\clearpage
\section{Exact admissible powers and a sharp logarithmic boundary scale for a Hermitian quartic}
{\small Provisional mathematical authors: Rachel Teo, Wenyi Wang, Hamdi Abderrahmene\par}\medskip
{\small Rachel Teo  -  Sakana AI.\par}

{\small Wenyi Wang  -  Sakana AI.\par}

{\small Hamdi Abderrahmene  -  Sakana AI.\par}

{\small Affiliations follow the recorded author declarations or public institutional/profile sources. Preferred author names, author order and publication affiliations await author review. This editorial compilation does not establish anthology coauthorship.\par}

\subsection{Abstract}
For H\_t(z)=(1+|z|\ensuremath{^2})\textasciicircum{}4\ensuremath{-}t|z|\ensuremath{^4}, the positive exponents producing a holomorphic squared-norm representation are classified exactly on6<t\ensuremath{\leq}7. Near t=8, the least positive admissible exponent has scale log(1/\ensuremath{\delta})/\ensuremath{\delta} with leading constant one, where \ensuremath{\delta}=8\ensuremath{-}t.

\subsection{Squared norms and the finite exponent classification}
A function-level squared norm means a finite sum \ensuremath{\sum}|f\_j(z)|\ensuremath{^2} with holomorphic polynomials f\_j, without requiring the summands to be monomials. For a radial polynomial in |z|\ensuremath{^2}, the formal development proves equivalence with nonnegativity of every coefficient in the one-variable polynomial. The necessity is a Gram/coefficient statement, so the result concerns actual functions rather than only a chosen monomial ansatz.

For 6<t\ensuremath{\leq}7 and a positive integer N, H\_t\textasciicircum{}N has such a representation exactly when N=2, or N\ensuremath{\geq}4, or N=3 and 924\ensuremath{-}210t+18t\ensuremath{^2}\ensuremath{-}t\ensuremath{^3}\ensuremath{\geq}0. The first power fails because its middle coefficient 6\ensuremath{-}t is negative. Exact expansions show that the square and fifth power have nonnegative coefficients throughout this interval. Products of those powers cover every exponent at least four. In the cubic expansion the only possible obstruction is the displayed central coefficient.

\subsection{Lower bounds near the radical boundary}
For 6<t<8, an admissible exponent satisfies 12N(8\ensuremath{-}t)\ensuremath{\geq}t(10\ensuremath{-}t). This inverse-distance bound comes from signed coefficient moments and positivity constraints. It is strengthened near t=8 by a Fourier obstruction: a negative terminal interval in the trigonometric profile must be compensated by the positive majorant allowed by nonnegative coefficients.

The source derives 3t\textasciicircum{}N N\textasciicircum{}(1/4)\ensuremath{\leq}28(16\ensuremath{-}t)\textasciicircum{}N, equivalently 81N t\textasciicircum{}(4N)\ensuremath{\leq}614656(16\ensuremath{-}t)\textasciicircum{}(4N), for 7\ensuremath{\leq}t<8. Taking logarithms and combining with the elementary lower bound gives the quantified leading lower scale. For every \ensuremath{\varepsilon}>0, sufficiently small \ensuremath{\delta}=8\ensuremath{-}t>0 forces \ensuremath{\delta}N\ensuremath{\geq}(1\ensuremath{-}\ensuremath{\varepsilon})log(1/\ensuremath{\delta}) for every positive admissible N.

\subsection{Matching upper bound and least exponent}
The upper bound extracts each coefficient by Fourier integration at a positive radius selected to match that coefficient's mean. The source proves existence of the saddle radius, the local Gaussian expansion of the actual normalized quartic, the modulus gap on outer arcs, and bounds for Gaussian, quartic and edge errors. Small coefficients receive a separate direct finite estimate, and symmetry covers the opposite edge. These steps establish positivity in every degree 0,...,4N uniformly whenever N\ensuremath{\geq}(1+\ensuremath{\varepsilon})log(1/\ensuremath{\delta})/\ensuremath{\delta} and \ensuremath{\delta} is sufficiently small.

Define the least positive admissible index using that nonempty admissible set. The formal squeeze proves both its minimality and 1\ensuremath{-}\ensuremath{\varepsilon}\ensuremath{\leq}\ensuremath{\delta}\ensuremath{\cdot}index/log(1/\ensuremath{\delta})\ensuremath{\leq}1+\ensuremath{\varepsilon} near zero. It excludes exponent zero, whose constant polynomial would otherwise trivialize the minimum. The claimed endpoint is the quantified squeeze, rather than a separately formalized filter-limit statement or eventual-index theorem.

\subsection{Novelty comparison and publication boundary}
The quartic family, threshold t<8, t=7 square and qualitative eventual positivity are known background. The proposed delta is the exact exceptional exponent/cubic transition and the sharp logarithmic scale with leading constant one. The inspected Tan-To appendix studies multiplication by an auxiliary power R\textasciicircum{}m and a different \ensuremath{\Theta}(1/\ensuremath{\alpha}) scale; that does not itself imply this H\_t\textasciicircum{}N result. Their full main theorem and the named Handelman/To-Yeung literature still require comparison.

The archive reports a 40-proof-module standalone closure under Lean 4.34.1. Early uncompiled-draft comments remain in the byte-preserved sources; the current verification register distinguishes archived claims from completed fresh execution. The ordinary-proof expansion now derives the analytic estimates and explicit constants, including variance certificates, the local complex logarithm, the positive Gaussian core, the exact modulus deficit and its sharp boundary coefficient, small-radius control, finite windows and uniform upper-threshold assembly. No additional infinite analytic lemma is left as an unstated premise for the displayed one-variable theorem. The result does not solve the general Hermitian-power problem or provide a formally verified binary homogeneous bridge; historical comparison and author review remain necessary.

\subsection{An exact admissible exponent semigroup}
Put q\_t(x)=(1+x)\textasciicircum{}4\ensuremath{-}tx\ensuremath{^2} and s=t\ensuremath{-}6, so q\_s(x)=1+4x\ensuremath{-}sx\ensuremath{^2}+4x\ensuremath{^3}+x\ensuremath{^4}. A radial polynomial R(|z|\ensuremath{^2}) is a finite sum of squared absolute values of holomorphic polynomials exactly when every coefficient of R is nonnegative. Nonnegative coefficients give the explicit monomial summands. Conversely the coefficient of z\textasciicircum{}k \ensuremath{\bar{z}}\textasciicircum{}k in any holomorphic squared-norm representation is a sum of squared coefficient magnitudes. Polynomial equality on C identifies that coefficient, even if the initial representation contains off-diagonal terms.

For 0<s\ensuremath{\leq}1 the first power fails at coefficient x\ensuremath{^2}. The square has coefficient list [1,8,16\ensuremath{-}2s,8\ensuremath{-}8s,34+s\ensuremath{^2},8\ensuremath{-}8s,16\ensuremath{-}2s,8,1], which is nonnegative. The cube has first-half coefficients [1,12,48\ensuremath{-}3s,76\ensuremath{-}24s,99\ensuremath{-}48s+3s\ensuremath{^2},216\ensuremath{-}24s+12s\ensuremath{^2}], central coefficient 96\ensuremath{-}102s\ensuremath{-}s\ensuremath{^3}, and the reflected first half. Thus the cube is nonnegative exactly when s\ensuremath{^3}+102s\ensuremath{\leq}96.

For the fifth power put h=1\ensuremath{-}s, B=q\_1\ensuremath{^3}+7x\ensuremath{^6}, D=q\_1+x\ensuremath{^2} and E=q\_1\ensuremath{^5}\ensuremath{-}2189x\ensuremath{^1}\ensuremath{^0}. All three have nonnegative coefficient lists. E has first half [1,20,155,580,1135,1784,3970,5000,4205,11580], central coefficient zero, and the reflected half. Direct expansion gives q\_s\ensuremath{^5}=E+5hx\ensuremath{^2}(q\_1\ensuremath{^2})\ensuremath{^2}+10h\ensuremath{^2}x\ensuremath{^4}B+10h\ensuremath{^3}x\ensuremath{^6}q\_1\ensuremath{^2}+5h\ensuremath{^4}x\ensuremath{^8}D+(2189\ensuremath{-}70h\ensuremath{^2}\ensuremath{-}5h\ensuremath{^4}+h\ensuremath{^5})x\ensuremath{^1}\ensuremath{^0}. Because 0\ensuremath{\leq}h\ensuremath{\leq}1 every term is nonnegative. Closure under multiplication now gives every even exponent and every odd exponent at least 5. The positive admissible exponents are exactly 2, all n\ensuremath{\geq}4, and 3 when the central cubic inequality holds. In6<t\ensuremath{\leq}7 the latter is 924\ensuremath{-}210t+18t\ensuremath{^2}\ensuremath{-}t\ensuremath{^3}\ensuremath{\geq}0.

\subsection{A complete moment obstruction near t=8}
Let 6<t<8, \ensuremath{\delta}=8\ensuremath{-}t, A=16\ensuremath{-}t and n>0. Assume q\_t\textasciicircum{}n has nonnegative coefficients. Select coefficients whose index has parity opposite to n and form their centered moments about 2n. The exact zeroth, second and fourth moments are M\ensuremath{_0}=(A\textasciicircum{}n\ensuremath{-}t\textasciicircum{}n)/2, M\ensuremath{_2}=8nA\textasciicircum{}(n\ensuremath{-}1), and M\ensuremath{_4}=[20n/A+384n(n\ensuremath{-}1)/A\ensuremath{^2}]A\textasciicircum{}n+12nt\textasciicircum{}(n\ensuremath{-}1). These identities follow from the centered Euler operator evaluated at 1 and\ensuremath{-}1.

These are moments of nonnegative weights, hence M\ensuremath{_2}\ensuremath{^2}\ensuremath{\leq}M\ensuremath{_0}M\ensuremath{_4}. The power-gap identity gives A\textasciicircum{}n\ensuremath{-}t\textasciicircum{}n\ensuremath{\leq}(A\ensuremath{-}t)nA\textasciicircum{}(n\ensuremath{-}1)=2\ensuremath{\delta}nA\textasciicircum{}(n\ensuremath{-}1), so M\ensuremath{_0}\ensuremath{\leq}\ensuremath{\delta}nA\textasciicircum{}n/A. Also t\textasciicircum{}(n\ensuremath{-}1)\ensuremath{\leq}A\textasciicircum{}(n\ensuremath{-}1), hence M\ensuremath{_4}\ensuremath{\leq}[32n/A+384n(n\ensuremath{-}1)/A\ensuremath{^2}]A\textasciicircum{}n. Substitute and cancel positive quantities to obtain 64A\ensuremath{\leq}\ensuremath{\delta}[32A+384(n\ensuremath{-}1)]. Rearranging gives t(10\ensuremath{-}t)\ensuremath{\leq}12n(8\ensuremath{-}t). The coefficient/squared-norm bridge makes this a necessary condition for the unrestricted holomorphic squared-norm question. Exponent zero is explicitly excluded.

\subsection{A fourth-root Fourier obstruction}
On the unit circle, e\textasciicircum{}(\ensuremath{-}2i\ensuremath{\theta})q\_t(e\textasciicircum{}(i\ensuremath{\theta}))=F\_t(\ensuremath{\theta})=16cos\ensuremath{^4}(\ensuremath{\theta}/2)\ensuremath{-}t. Nonnegative coefficients imply \ensuremath{\int}\ensuremath{_0}\textasciicircum{}\ensuremath{\pi}F\_t(\ensuremath{\theta})\textasciicircum{}n cos(j\ensuremath{\theta})d\ensuremath{\theta}\ensuremath{\geq}0 for every nonnegative integer j by Fourier extraction and symmetry. Choose j=0 for odd n and j=1 for even n; both integrands are negative near\ensuremath{\pi}.

For u=\ensuremath{\pi}\ensuremath{-}\ensuremath{\theta}\ensuremath{\leq}n\textasciicircum{}(\ensuremath{-}1/4)\ensuremath{\leq}1,16sin\ensuremath{^4}(u/2)\ensuremath{\leq}u\ensuremath{^4}\ensuremath{\leq}1/n. Thus |F\_t(\ensuremath{\theta})|\textasciicircum{}n\ensuremath{\geq}(t\ensuremath{-}1/n)\textasciicircum{}n\ensuremath{\geq}(1\ensuremath{-}1/t)t\textasciicircum{}n\ensuremath{\geq}(6/7)t\textasciicircum{}n when t\ensuremath{\geq}7. In the even case cos\ensuremath{\theta}=\ensuremath{-}cosu\ensuremath{\leq}\ensuremath{-}1/2, so both integrands are at most\ensuremath{-}(3/7)t\textasciicircum{}n on an interval of length n\textasciicircum{}(\ensuremath{-}1/4). The source bounds the positive part by the explicit Gaussian majorant A\textasciicircum{}n exp[\ensuremath{-}16n\ensuremath{\theta}\ensuremath{^2}/(9\ensuremath{\pi}\ensuremath{^2})] plus 4\textasciicircum{}n. Its integral plus\ensuremath{\pi} 4\textasciicircum{}n is at most 4A\textasciicircum{}n/\ensuremath{\surd}n; the remainder uses 2\ensuremath{\surd}n\ensuremath{\leq}2\textasciicircum{}n and A\ensuremath{\geq}8.

Comparing the negative interval with that positive mass gives 3t\textasciicircum{}n n\textasciicircum{}(1/4)\ensuremath{\leq}28A\textasciicircum{}n, hence 81nt\textasciicircum{}(4n)\ensuremath{\leq}614656A\textasciicircum{}(4n). Taking logarithms yields log n\ensuremath{\leq}4n log(A/t)+O(1). Since 4log[(8+\ensuremath{\delta})/(8\ensuremath{-}\ensuremath{\delta})]=\ensuremath{\delta}+O(\ensuremath{\delta}\ensuremath{^3}) and the moment bound already gives n\ensuremath{\geq}constant/\ensuremath{\delta}, every admissible positive exponent obeys \ensuremath{\delta}n\ensuremath{\geq}(1\ensuremath{-}\ensuremath{\varepsilon})log(1/\ensuremath{\delta}) for sufficiently small \ensuremath{\delta}, for any \ensuremath{\varepsilon}>0.

\subsection{Quantitative lemmas and assembly for the matching upper bound}
For 0<k\ensuremath{\leq}2N use Cauchy coefficient extraction at a saddle radius 0<r\ensuremath{\leq}1 satisfying N\ensuremath{\mu}(r)=k, with \ensuremath{\mu}(r)=rq\ensuremath{^{\prime}}(r)/q(r). Direct rational calculation gives \ensuremath{\mu}\ensuremath{^{\prime}}(r)=v(r)/r>0 and r/50\ensuremath{\leq}v(r)\ensuremath{\leq}160r on7\ensuremath{\leq}t\ensuremath{\leq}8. The endpoints \ensuremath{\mu}(0)=0, \ensuremath{\mu}(1)=2 give existence; \ensuremath{\mu}(r)\ensuremath{\leq}20r gives Nr\ensuremath{\geq}k/20. Set H(\ensuremath{\theta})=[q(re\textasciicircum{}(i\ensuremath{\theta}))/q(r)]\textasciicircum{}N e\textasciicircum{}(\ensuremath{-}ik\ensuremath{\theta}). A positive real integral of H on[0,\ensuremath{\pi}] proves coefficient positivity.

The proved local complex-log expansion cancels its linear phase at the saddle, leaves Gaussian term\ensuremath{-}Nv(r)\ensuremath{\theta}\ensuremath{^2}/2, and has remainder bounded by 100000Nr|\ensuremath{\theta}|\ensuremath{^3} on |\ensuremath{\theta}|\ensuremath{\leq}1/128. Integrating on the fixed core cutoff 1/20000000 supplies uniform M\ensuremath{_0},c\ensuremath{_0}>0 and core lower bound c\ensuremath{_0}/\ensuremath{\surd}(Nr) for Nr\ensuremath{\geq}M\ensuremath{_0}. Retaining Nr is essential at small-radius saddles.

For r\ensuremath{\geq}1/8, fix \ensuremath{\alpha}<1/4 and \ensuremath{\eta}>0 small enough that \ensuremath{\alpha}+4\ensuremath{\eta}\ensuremath{\leq}1/4. The middle-arc error is at most\ensuremath{\pi}exp(\ensuremath{-}a\_\ensuremath{\eta}N), with a\_\ensuremath{\eta}=coreCutoff\ensuremath{^2}\ensuremath{\eta}\ensuremath{^4}/(40000\ensuremath{\pi}\ensuremath{^6})>0. The near-\ensuremath{\pi} error is at most Cexp(\ensuremath{-}\ensuremath{\alpha}\ensuremath{\delta}N)N\textasciicircum{}(\ensuremath{-}1/4), where C=1+10000\ensuremath{\pi}\ensuremath{^4}. Exact modulus-deficit factorizations for the actual quartic give these estimates, retaining a deficit of \ensuremath{\alpha}\ensuremath{\delta}+(a+b\ensuremath{^2})/10000 near\ensuremath{\pi}; its sharp limiting \ensuremath{\delta} coefficient is 1/4. For r\ensuremath{\leq}1/8, the entire outer arc error is instead at most\ensuremath{\pi}exp(\ensuremath{-}Nr\ensuremath{\cdot}coreCutoff\ensuremath{^2}/\ensuremath{\pi}\ensuremath{^2}).

Thus the complete integral in the compact range is at least c\ensuremath{_0}/\ensuremath{\surd}N minus the middle and quartic errors, and at small radius it is at least c\ensuremath{_0}/\ensuremath{\surd}(Nr) minus the edge error. Choose 1/[4(1+\ensuremath{\varepsilon})]<\ensuremath{\alpha}<1/4. Then N\ensuremath{\geq}(1+\ensuremath{\varepsilon})log(1/\ensuremath{\delta})/\ensuremath{\delta} makes each error a sufficiently small fraction of its positive core for small \ensuremath{\delta}, provided the relevant scale is large.

The bounded remaining coefficient indices have uniform finite-window positivity: for fixed k the explicit sufficient estimate is N\ensuremath{\geq}2k+k\ensuremath{^2}11\textasciicircum{}k, uniformly on7\ensuremath{\leq}t\ensuremath{\leq}8. Palindromicity transfers the lower half to all 0\ensuremath{\leq}k\ensuremath{\leq}4N; later coefficients vanish. Every sufficiently large exponent above the stated logarithmic threshold is therefore admissible. The lower and upper bounds squeeze the least positive admissible index n*(\ensuremath{\delta}) to \ensuremath{\delta}n*(\ensuremath{\delta})/log(1/\ensuremath{\delta})\ensuremath{\to}1.

\subsection{Analytic appendix A: the actual saddle radius and variance}
Write s=t\ensuremath{-}6\ensuremath{\in}[1,2], q(r)=1+4r\ensuremath{-}sr\ensuremath{^2}+4r\ensuremath{^3}+r\ensuremath{^4}, A\ensuremath{_1}(r)=4r\ensuremath{-}2sr\ensuremath{^2}+12r\ensuremath{^3}+4r\ensuremath{^4} and A\ensuremath{_2}(r)=4r\ensuremath{-}4sr\ensuremath{^2}+36r\ensuremath{^3}+16r\ensuremath{^4}. Set \ensuremath{\mu}=A\ensuremath{_1}/q and V=A\ensuremath{_2}/q\ensuremath{-}\ensuremath{\mu}\ensuremath{^2}. Direct expansion gives qA\ensuremath{_2}\ensuremath{-}A\ensuremath{_1}\ensuremath{^2}=4rP, where P=1\ensuremath{-}sr+(9\ensuremath{-}s)r\ensuremath{^2}+20r\ensuremath{^3}+(9\ensuremath{-}s)r\ensuremath{^4}\ensuremath{-}sr\ensuremath{^5}+r\ensuremath{^6}. Thus V=4rP/q\ensuremath{^2} and \ensuremath{\mu}\ensuremath{^{\prime}}(r)=V/r for r>0.

Every lower bound needed for the saddle follows from three exact polynomial certificates: q\ensuremath{-}1=2r(2\ensuremath{-}r)+(2\ensuremath{-}s)r\ensuremath{^2}+4r\ensuremath{^3}+r\ensuremath{^4}; 2P\ensuremath{-}1=(2r\ensuremath{-}1)\ensuremath{^2}+10r\ensuremath{^2}+40r\ensuremath{^3}+2r\ensuremath{^4}((r\ensuremath{-}1)\ensuremath{^2}+6)+2(2\ensuremath{-}s)(r+r\ensuremath{^2}+r\ensuremath{^4}+r\ensuremath{^5}); and A\ensuremath{_1}\ensuremath{-}2r=2r(1\ensuremath{-}r)\ensuremath{^2}+10r\ensuremath{^3}+4r\ensuremath{^4}+2(2\ensuremath{-}s)r\ensuremath{^2}. All terms on their right sides are nonnegative for 0\ensuremath{\leq}r\ensuremath{\leq}1. Dropping negative terms in the definitions also gives q\ensuremath{\leq}10, P\ensuremath{\leq}40 and A\ensuremath{_1}\ensuremath{\leq}20r. Consequently 1\ensuremath{\leq}q\ensuremath{\leq}10, 1/2\ensuremath{\leq}P\ensuremath{\leq}40, r/5\ensuremath{\leq}\ensuremath{\mu}\ensuremath{\leq}20r and r/50\ensuremath{\leq}V\ensuremath{\leq}160r.

The function \ensuremath{\mu} is continuous on [0,1], \ensuremath{\mu}(0)=0, \ensuremath{\mu}(1)=2, and \ensuremath{\mu}\ensuremath{^{\prime}}>0 on (0,1]. Therefore for every integer 0<k\ensuremath{\leq}2N there is a unique 0<r\ensuremath{\leq}1 satisfying N\ensuremath{\mu}(r)=k. The inequality \ensuremath{\mu}\ensuremath{\leq}20r implies Nr\ensuremath{\geq}k/20. This lower bound on the scale is the bridge between coefficient indices and the analytic estimates; N alone is insufficient when r tends to zero.

Put f(\ensuremath{\theta})=q(re\textasciicircum{}\{i\ensuremath{\theta}\})/q(r) and H(\ensuremath{\theta})=f(\ensuremath{\theta})\textasciicircum{}N e\textasciicircum{}\{\ensuremath{-}ik\ensuremath{\theta}\}. Fourier extraction gives [x\textasciicircum{}k]q(x)\textasciicircum{}N=q(r)\textasciicircum{}N/(\ensuremath{\pi}r\textasciicircum{}k)\ensuremath{\cdot}Re\ensuremath{\int}\ensuremath{_0}\textasciicircum{}\ensuremath{\pi} H(\ensuremath{\theta})d\ensuremath{\theta}. The prefactor is positive. It is therefore enough to prove that this real integral is positive at the actual saddle N\ensuremath{\mu}(r)=k.

\subsection{Analytic appendix B: a uniform complex logarithm and positive Gaussian core}
Let u=f(\ensuremath{\theta})\ensuremath{-}1. The elementary bound |e\textasciicircum{}\{ij\ensuremath{\theta}\}\ensuremath{-}1|\ensuremath{\leq}j|\ensuremath{\theta}| and q(r)\ensuremath{\geq}1 give |u|\ensuremath{\leq}(4r+4r\ensuremath{^2}+12r\ensuremath{^3}+4r\ensuremath{^4})|\ensuremath{\theta}|\ensuremath{\leq}24r|\ensuremath{\theta}|. Taylor expansion of each of the four angular exponentials through second order has integral remainder at most j\ensuremath{^3}|\ensuremath{\theta}|\ensuremath{^3}/6, because its third derivative has constant modulus j\ensuremath{^3}. Summing their weighted remainders gives |u\ensuremath{-}i\ensuremath{\mu}\ensuremath{\theta}+(A\ensuremath{_2}/q)\ensuremath{\theta}\ensuremath{^2}/2|\ensuremath{\leq}32r|\ensuremath{\theta}|\ensuremath{^3}. Also |A\ensuremath{_2}/q|\ensuremath{\leq}64r, so |u\ensuremath{-}i\ensuremath{\mu}\ensuremath{\theta}|\ensuremath{\leq}33r\ensuremath{\theta}\ensuremath{^2} on |\ensuremath{\theta}|\ensuremath{\leq}1/128.

On that interval |u|\ensuremath{\leq}3/16<1/4, so the principal logarithm agrees with the convergent series log(1+u)=u\ensuremath{-}u\ensuremath{^2}/2+\ensuremath{\Sigma}\_\{j\ensuremath{\geq}3\}(\ensuremath{-}1)\textasciicircum{}(j+1)u\textasciicircum{}j/j. Its tail has modulus at most |u|\ensuremath{^3}/[3(1\ensuremath{-}|u|)]\ensuremath{\leq}|u|\ensuremath{^3}\ensuremath{\leq}13,824r|\ensuremath{\theta}|\ensuremath{^3}. Further, |u\ensuremath{^2}\ensuremath{-}(i\ensuremath{\mu}\ensuremath{\theta})\ensuremath{^2}|/2\ensuremath{\leq}(33r\ensuremath{\theta}\ensuremath{^2})(44r|\ensuremath{\theta}|)/2\ensuremath{\leq}726r|\ensuremath{\theta}|\ensuremath{^3}. Combining these bounds with the 32r|\ensuremath{\theta}|\ensuremath{^3} angular remainder yields log f=i\ensuremath{\mu}\ensuremath{\theta}\ensuremath{-}V\ensuremath{\theta}\ensuremath{^2}/2+R(\ensuremath{\theta}), with |R(\ensuremath{\theta})|\ensuremath{\leq}14,582r|\ensuremath{\theta}|\ensuremath{^3}\ensuremath{\leq}100,000r|\ensuremath{\theta}|\ensuremath{^3}. This derives the exact generous bound used in the source, rather than assuming a local saddle expansion.

The saddle phase cancels: H(\ensuremath{\theta})=exp(\ensuremath{-}NV\ensuremath{\theta}\ensuremath{^2}/2+NR(\ensuremath{\theta})). Put S=\ensuremath{\surd}(Nr), v\ensuremath{_0}=1/50, v\ensuremath{_1}=160 and C=100,000. Then v\ensuremath{_0}S\ensuremath{^2}\ensuremath{\leq}NV\ensuremath{\leq}v\ensuremath{_1}S\ensuremath{^2} and |NR(\ensuremath{\theta})|\ensuremath{\leq}CS\ensuremath{^2}|\ensuremath{\theta}|\ensuremath{^3}. Fix h=1/20,000,000. The relation 4Ch=v\ensuremath{_0} shows that throughout 0\ensuremath{\leq}\ensuremath{\theta}\ensuremath{\leq}h, |H(\ensuremath{\theta})|\ensuremath{\leq}exp(\ensuremath{-}v\ensuremath{_0}S\ensuremath{^2}\ensuremath{\theta}\ensuremath{^2}/4), by bounding the real error CS\ensuremath{^2}\ensuremath{\theta}\ensuremath{^3}\ensuremath{\leq}v\ensuremath{_0}S\ensuremath{^2}\ensuremath{\theta}\ensuremath{^2}/4.

Here is a fully quantitative positive-core proof. Take T=200 and require S\ensuremath{\geq}max(2CT\ensuremath{^3},T/h)=1.6\ensuremath{\cdot}10\ensuremath{^1}\ensuremath{^2}. On [0,T/S] the logarithmic error has modulus at most CT\ensuremath{^3}/S\ensuremath{\leq}1/2, so its imaginary part has cosine at least 1/2 and Re H\ensuremath{\geq}0. On [0,1/S], NV\ensuremath{\theta}\ensuremath{^2}\ensuremath{\leq}v\ensuremath{_1} and Re(NR)\ensuremath{\geq}\ensuremath{-}1/2, whence Re H\ensuremath{\geq}e\textasciicircum{}(\ensuremath{-}81)/2. This interval contributes at least e\textasciicircum{}(\ensuremath{-}81)/(2S). On [T/S,h], bound a possible negative real part by the Gaussian modulus. The elementary tail inequality \ensuremath{\int}\ensuremath{^{b}}\ensuremath{^{h}}e\textasciicircum{}(\ensuremath{-}c\ensuremath{\theta}\ensuremath{^2})d\ensuremath{\theta}\ensuremath{\leq}e\textasciicircum{}(\ensuremath{-}cb\ensuremath{^2})/(cb) gives a loss at most 4e\textasciicircum{}(\ensuremath{-}v\ensuremath{_0}T\ensuremath{^2}/4)/(v\ensuremath{_0}TS)=e\textasciicircum{}(\ensuremath{-}200)/S. Since e\textasciicircum{}(\ensuremath{-}200)\ensuremath{\leq}e\textasciicircum{}(\ensuremath{-}81)/4, the whole core integral is at least c\ensuremath{_0}/S, where c\ensuremath{_0}=e\textasciicircum{}(\ensuremath{-}81)/4. Thus M\ensuremath{_0}=(1.6\ensuremath{\cdot}10\ensuremath{^1}\ensuremath{^2})\ensuremath{^2} suffices for the condition Nr\ensuremath{\geq}M\ensuremath{_0}. These constants are deliberately conservative and their size affects only the eventual neighborhood, not the leading asymptotic constant.

\subsection{Analytic appendix C: the exact boundary modulus identity}
For r>0 set a=(1\ensuremath{-}r)\ensuremath{^2}/(2r) and b=1+cos\ensuremath{\theta}. Let \ensuremath{\delta}=8\ensuremath{-}t=2\ensuremath{-}s, F=8+16a+4a\ensuremath{^2}+\ensuremath{\delta}, K=2+2a+b+2a(a+2)b, and J=2[a\ensuremath{^2}(4a+12\ensuremath{-}2b)+4a((b\ensuremath{-}1)\ensuremath{^2}+1)+b\ensuremath{^2}(2+b)]. Substituting q(r)=r\ensuremath{^2}F and expanding the real and imaginary parts of q(re\textasciicircum{}\{i\ensuremath{\theta}\}) gives the exact identity 1\ensuremath{-}|f(\ensuremath{\theta})|\ensuremath{^2}=2D, with D=4(2\ensuremath{-}b)(J+\ensuremath{\delta}K)/F\ensuremath{^2}. This finite polynomial identity is the Gap/AngleBridge certificate applied to cos\ensuremath{^2}\ensuremath{\theta}+sin\ensuremath{^2}\ensuremath{\theta}=1; it can also be checked by direct expansion.

For a\ensuremath{\geq}0, 0\ensuremath{\leq}b\ensuremath{\leq}2 and \ensuremath{\delta}\ensuremath{\geq}0, every term of J and K is nonnegative. More precisely J\ensuremath{-}(8a+4b\ensuremath{^2})=2a\ensuremath{^2}(4a+12\ensuremath{-}2b)+8a(b\ensuremath{-}1)\ensuremath{^2}+2b\ensuremath{^3}\ensuremath{\geq}0, and K\ensuremath{\geq}2. For 1/8\ensuremath{\leq}r\ensuremath{\leq}1 we have 0\ensuremath{\leq}a\ensuremath{\leq}4; for 0\ensuremath{\leq}\ensuremath{\delta}\ensuremath{\leq}1, F\ensuremath{\leq}137 and F\ensuremath{^2}\ensuremath{\leq}18,769. Therefore D\ensuremath{\geq}(2\ensuremath{-}b)(a+b\ensuremath{^2}+\ensuremath{\delta})/10,000. Its \ensuremath{\delta}-independent part is at least (2\ensuremath{-}b)(a+b\ensuremath{^2})/10,000. The constants follow by substituting J\ensuremath{\geq}8a+4b\ensuremath{^2} and K\ensuremath{\geq}2 and weakening 8/18,769 to 1/10,000.

To retain the sharp boundary coefficient, choose any 0<\ensuremath{\alpha}<1/4 and let 0<\ensuremath{\eta}\ensuremath{\leq}min(1,(1/4\ensuremath{-}\ensuremath{\alpha})/4). If a,b,\ensuremath{\delta}\ensuremath{\leq}\ensuremath{\eta}, then F\ensuremath{\leq}8+21\ensuremath{\eta} and 4(2\ensuremath{-}b)K\ensuremath{\geq}16\ensuremath{-}8\ensuremath{\eta}. The exact identity (16\ensuremath{-}8\ensuremath{\eta})\ensuremath{-}(1/4\ensuremath{-}4\ensuremath{\eta})(8+21\ensuremath{\eta})\ensuremath{^2}=164\ensuremath{\eta}+(4935/4)\ensuremath{\eta}\ensuremath{^2}+1764\ensuremath{\eta}\ensuremath{^3}\ensuremath{\geq}0 proves 4(2\ensuremath{-}b)K/F\ensuremath{^2}\ensuremath{\geq}\ensuremath{\alpha}. Also 2\ensuremath{-}b\ensuremath{\geq}1. Splitting D into its spatial and \ensuremath{\delta} terms gives D\ensuremath{\geq}\ensuremath{\alpha}\ensuremath{\delta}+(a+b\ensuremath{^2})/10,000. At a=b=\ensuremath{\delta}=0 the \ensuremath{\delta} multiplier is exactly 1/4; this is the origin of the leading asymptotic threshold.

Let u=\ensuremath{\pi}\ensuremath{-}\ensuremath{\theta}. Standard trigonometric bounds give 1+cos\ensuremath{\theta}=1\ensuremath{-}cos u\ensuremath{\geq}2u\ensuremath{^2}/\ensuremath{\pi}\ensuremath{^2} and 1\ensuremath{-}cos\ensuremath{\theta}\ensuremath{\geq}2\ensuremath{\theta}\ensuremath{^2}/\ensuremath{\pi}\ensuremath{^2} for 0\ensuremath{\leq}\ensuremath{\theta}\ensuremath{\leq}\ensuremath{\pi}. In particular b\ensuremath{^2}\ensuremath{\geq}u\ensuremath{^4}/\ensuremath{\pi}\ensuremath{^4}. If r\ensuremath{\geq}1\ensuremath{-}\ensuremath{\eta}/2 and \ensuremath{\theta}\ensuremath{\geq}\ensuremath{\pi}\ensuremath{-}\ensuremath{\eta}, then a,b\ensuremath{\leq}\ensuremath{\eta}, so |f(\ensuremath{\theta})|\ensuremath{\leq}exp(\ensuremath{-}\ensuremath{\alpha}\ensuremath{\delta}\ensuremath{-}u\ensuremath{^4}/(10,000\ensuremath{\pi}\ensuremath{^4})). The passage from D to this modulus bound uses |f|\ensuremath{^2}=1\ensuremath{-}2D\ensuremath{\leq}e\textasciicircum{}(\ensuremath{-}2D), followed by a nonnegative square root.

On the remaining compact outer region r\ensuremath{\geq}1/8, \ensuremath{\theta}\ensuremath{\geq}h, if either r\ensuremath{\leq}1\ensuremath{-}\ensuremath{\eta}/2 or u\ensuremath{\geq}\ensuremath{\eta}, then a+b\ensuremath{^2}\ensuremath{\geq}\ensuremath{\eta}\ensuremath{^4}/(8\ensuremath{\pi}\ensuremath{^4}). The radial case follows from a\ensuremath{\geq}\ensuremath{\eta}\ensuremath{^2}/8; the angular case follows from the fourth-power bound. Combining this with 2\ensuremath{-}b\ensuremath{\geq}2h\ensuremath{^2}/\ensuremath{\pi}\ensuremath{^2} gives D\ensuremath{\geq}h\ensuremath{^2}\ensuremath{\eta}\ensuremath{^4}/(40,000\ensuremath{\pi}\ensuremath{^6})=a\_\ensuremath{\eta}>0, and hence |f|\ensuremath{\leq}e\textasciicircum{}(\ensuremath{-}a\_\ensuremath{\eta}). The near-\ensuremath{\pi} rectangle and its complement now give uniform bounds for every compact-radius saddle, including r=1.

\subsection{Analytic appendix D: small radii and the outer integral errors}
For 0<r\ensuremath{\leq}1/8 a separate algebraic certificate retains the scale r. Write c=cos\ensuremath{\theta} and E=8r\ensuremath{-}16r\ensuremath{^2}\ensuremath{-}16r\ensuremath{^3}\ensuremath{-}16r\ensuremath{^5}\ensuremath{-}16r\ensuremath{^6}. The exact squared-modulus factorization is q(r)\ensuremath{^2}\ensuremath{-}|q(re\textasciicircum{}\{i\ensuremath{\theta}\})|\ensuremath{^2}=(1\ensuremath{-}c)T, where T\ensuremath{-}E is the following sum of nonnegative terms: 4r\ensuremath{^2}(4\ensuremath{-}s(1+c))+8r\ensuremath{^3}(2\ensuremath{-}s)+8r\ensuremath{^3}(2c+1)\ensuremath{^2}+16r\ensuremath{^4}c\ensuremath{^2}(1+c)+64r\ensuremath{^4}(1+c)+8r\ensuremath{^5}(2\ensuremath{-}s)+8r\ensuremath{^5}(2c+1)\ensuremath{^2}+4r\ensuremath{^6}(4\ensuremath{-}s(1+c))+8r\ensuremath{^7}. Here s\ensuremath{\leq}2 and \ensuremath{-}1\ensuremath{\leq}c\ensuremath{\leq}1 make every coefficient factor nonnegative.

Moreover E\ensuremath{-}4r=4r(1\ensuremath{-}4r\ensuremath{-}4r\ensuremath{^2}\ensuremath{-}4r\ensuremath{^4}\ensuremath{-}4r\ensuremath{^5})\ensuremath{\geq}0 for r\ensuremath{\leq}1/8, and q(r)\ensuremath{\leq}2 by dropping its negative quadratic term. Thus D=(1\ensuremath{-}c)T/(2q(r)\ensuremath{^2})\ensuremath{\geq}r(1\ensuremath{-}c)/2\ensuremath{\geq}r\ensuremath{\theta}\ensuremath{^2}/\ensuremath{\pi}\ensuremath{^2}, so |f(\ensuremath{\theta})|\ensuremath{\leq}exp(\ensuremath{-}r\ensuremath{\theta}\ensuremath{^2}/\ensuremath{\pi}\ensuremath{^2}). On \ensuremath{\theta}\ensuremath{\geq}h this gives |H(\ensuremath{\theta})|\ensuremath{\leq}exp(\ensuremath{-}Nrh\ensuremath{^2}/\ensuremath{\pi}\ensuremath{^2}) and a total outer loss at most \ensuremath{\pi}exp(\ensuremath{-}Nrh\ensuremath{^2}/\ensuremath{\pi}\ensuremath{^2}). The proof has retained Nr in both the Gaussian core and the error.

At compact radii, the middle region contributes a loss at most \ensuremath{\pi}e\textasciicircum{}(\ensuremath{-}a\_\ensuremath{\eta}N). Near \ensuremath{\pi} the loss is at most e\textasciicircum{}(\ensuremath{-}\ensuremath{\alpha}\ensuremath{\delta}N)\ensuremath{\int}\ensuremath{_0}\textasciicircum{}\ensuremath{\eta}e\textasciicircum{}(\ensuremath{-}cNu\ensuremath{^4})du with c=1/(10,000\ensuremath{\pi}\ensuremath{^4}). Put S\ensuremath{_4}=N\textasciicircum{}(1/4) and split at 1/S\ensuremath{_4}. The first part has length at most 1/S\ensuremath{_4}. In the second part, (S\ensuremath{_4}u)\ensuremath{^4}\ensuremath{\geq}S\ensuremath{_4}u, so the integrand is bounded by e\textasciicircum{}(\ensuremath{-}cS\ensuremath{_4}u), whose integral is at most 1/(cS\ensuremath{_4}). Therefore the total loss is at most Qe\textasciicircum{}(\ensuremath{-}\ensuremath{\alpha}\ensuremath{\delta}N)N\textasciicircum{}(\ensuremath{-}1/4), where Q=1+10,000\ensuremath{\pi}\ensuremath{^4}. This elementary bound requires no Gamma-function integral.

The entire normalized coefficient integral is consequently at least c\ensuremath{_0}N\textasciicircum{}(\ensuremath{-}1/2)\ensuremath{-}\ensuremath{\pi}e\textasciicircum{}(\ensuremath{-}a\_\ensuremath{\eta}N)\ensuremath{-}Qe\textasciicircum{}(\ensuremath{-}\ensuremath{\alpha}\ensuremath{\delta}N)N\textasciicircum{}(\ensuremath{-}1/4) for r\ensuremath{\geq}1/8 and Nr\ensuremath{\geq}M\ensuremath{_0}. For r\ensuremath{\leq}1/8 and Nr\ensuremath{\geq}M\ensuremath{_0} it is at least c\ensuremath{_0}(Nr)\textasciicircum{}(\ensuremath{-}1/2)\ensuremath{-}\ensuremath{\pi}e\textasciicircum{}(\ensuremath{-}Nrh\ensuremath{^2}/\ensuremath{\pi}\ensuremath{^2}). These are the actual integrals of the polynomial, rather than a formal model with an unproved application premise.

\subsection{Analytic appendix E: finite coefficient windows and uniform threshold assembly}
For a fixed positive coefficient index k, write q=1+g with g=4x\ensuremath{-}sx\ensuremath{^2}+4x\ensuremath{^3}+x\ensuremath{^4}. The absolute sum of the coefficients of g is at most 11. Hence every coefficient of g\textasciicircum{}j has absolute value at most 11\textasciicircum{}j; g\textasciicircum{}j has no coefficient below degree j and its degree-j coefficient is 4\textasciicircum{}j. The binomial expansion gives [x\textasciicircum{}k]q\textasciicircum{}N=binom(N,k)4\textasciicircum{}k+\ensuremath{\Sigma}\_\{j<k\}binom(N,j)[x\textasciicircum{}k]g\textasciicircum{}j.

If N\ensuremath{\geq}2k, all binomial coefficients in the last sum are at most binom(N,k\ensuremath{-}1). Its absolute value is therefore at most k\ensuremath{\cdot}11\textasciicircum{}k\ensuremath{\cdot}binom(N,k\ensuremath{-}1). The adjacent-binomial ratio is binom(N,k)/binom(N,k\ensuremath{-}1)=(N\ensuremath{-}k+1)/k. Thus the explicit sufficient condition N\ensuremath{\geq}2k+k\ensuremath{^2}11\textasciicircum{}k makes that ratio strictly larger than k11\textasciicircum{}k and proves coefficient positivity, uniformly in 1\ensuremath{\leq}s\ensuremath{\leq}2. The zeroth coefficient is always 1. For any fixed K, taking the maximum of these finite thresholds proves positivity for all k\ensuremath{\leq}K.

Fix \ensuremath{\varepsilon}>0 and choose 1/[4(1+\ensuremath{\varepsilon})]<\ensuremath{\alpha}<1/4; then choose \ensuremath{\eta} as in appendix C. At N\ensuremath{\geq}N\_\ensuremath{\delta}=(1+\ensuremath{\varepsilon})log(1/\ensuremath{\delta})/\ensuremath{\delta}, the compact quartic loss divided by its core is bounded by (Q/c\ensuremath{_0})N\textasciicircum{}(1/4)e\textasciicircum{}(\ensuremath{-}\ensuremath{\alpha}\ensuremath{\delta}N). For sufficiently small \ensuremath{\delta} this function of N is decreasing throughout N\ensuremath{\geq}N\_\ensuremath{\delta}, since its logarithmic derivative is 1/(4N)\ensuremath{-}\ensuremath{\alpha}\ensuremath{\delta}<0. At N\_\ensuremath{\delta} its value is a constant times log(1/\ensuremath{\delta})\textasciicircum{}(1/4)\ensuremath{\delta}\textasciicircum{}(\ensuremath{\alpha}(1+\ensuremath{\varepsilon})\ensuremath{-}1/4), which tends to zero. This proves the required error bound uniformly for every larger exponent, not just at one chosen exponent.

The middle loss/core ratio is (\ensuremath{\pi}/c\ensuremath{_0})\ensuremath{\surd}N e\textasciicircum{}(\ensuremath{-}a\_\ensuremath{\eta}N), which tends to zero. At small radii, the corresponding ratio is (\ensuremath{\pi}/c\ensuremath{_0})\ensuremath{\surd}(Nr)e\textasciicircum{}(\ensuremath{-}Nrh\ensuremath{^2}/\ensuremath{\pi}\ensuremath{^2}), which tends to zero as Nr increases. Choose a finite scale M large enough for the positive core and this last error, and choose K>20M. For k>K the saddle relation gives Nr\ensuremath{\geq}k/20>M. At compact radii N\ensuremath{\geq}8M\ensuremath{_0} suffices for the core, and the two compact errors are small by the preceding uniform estimates. For k\ensuremath{\leq}K use the finite window bound. Shrinking \ensuremath{\delta} ensures N\_\ensuremath{\delta} exceeds all fixed exponent bounds just selected.

This proves strict positivity for all lower-half coefficient indices 0\ensuremath{\leq}k\ensuremath{\leq}2N. The quartic is palindromic, so its Nth power is palindromic of degree 4N; replacing k by 4N\ensuremath{-}k proves every remaining coefficient positive. Coefficients above degree 4N vanish. The coefficient/squared-norm criterion now supplies the function-level holomorphic squared-norm representation. Together with the previously proved Fourier lower bound this gives the normalized least-index limit with constant one. When choosing the least integer exponent, apply the upper argument with a smaller internal \ensuremath{\varepsilon} and absorb the ceiling error \ensuremath{\delta}/log(1/\ensuremath{\delta}), which tends to zero.

These appendices transcribe the algebraic certificates and analytic estimates behind the selected frozen Lean assembly into an ordinary proof. The enormous conservative local constants are sufficient uniform constants and are not presented as optimized numerical bounds. Historical comparison with earlier Hermitian-power positivity theorems, fresh formal replay and contributor approval remain separate from completion of this ordinary argument.

\subsection{Verification and reproducibility}
Fresh execution: sakana-opdp 13: PASS; 40/40 custom modules compiled successfully; receipt Verification/sakana-opdp13-fresh-build.json. Compiler pins, source hashes and endpoint axiom outputs are in Verification. Historical receipts and finite checks do not replace fresh replay.

\begin{samepage}
\subsection{PowerSemigroup.function\_exact\_exponent\_semigroup}
\begin{quote}\small\ttfamily theorem function\_exact\_exponent\_semigroup \{t : \ensuremath{\mathbb{R}}\} (ht6 : 6 < t) (ht7 : t \ensuremath{\leq} 7)\newline     (n : \ensuremath{\mathbb{N}}) (hn : 0 < n) :\newline     FunctionSquaredNorm (fun z : \ensuremath{\mathbb{C}} =>\newline       (((1+Complex.normSq z)\textasciicircum{}4-t*(Complex.normSq z)\textasciicircum{}2)\textasciicircum{}n : \ensuremath{\mathbb{R}})) \ensuremath{\leftrightarrow}\newline       n=2 \ensuremath{\vee} 4\ensuremath{\leq}n \ensuremath{\vee} (n=3 \ensuremath{\wedge} 0\ensuremath{\leq}924-210*t+18*t\textasciicircum{}2-t\textasciicircum{}3)\end{quote}
{\small Source: proofs/opdp13/OpHack/PowerSemigroup/Evaluation.lean:116.\par}

{\small Source SHA-256: 9924251f0abcd263e38621ad0b96cadbd0a33f6c353d41ecf3e16e30a860ebfc\par}

\end{samepage}
\begin{samepage}
\subsection{PowerMoment.function\_exponent\_lower\_bound}
\begin{quote}\small\ttfamily theorem function\_exponent\_lower\_bound (t : \ensuremath{\mathbb{R}}) (ht : 6 < t) (ht8 : t < 8)\newline     (n : \ensuremath{\mathbb{N}}) (hn : 0 < n)\newline     (h : PowerSemigroup.FunctionSquaredNorm (fun z : \ensuremath{\mathbb{C}} =>\newline       (((1+Complex.normSq z)\textasciicircum{}4-t*(Complex.normSq z)\textasciicircum{}2)\textasciicircum{}n : \ensuremath{\mathbb{R}}))) :\newline     t*(10-t) \ensuremath{\leq} 12*(n : \ensuremath{\mathbb{R}})*(8-t)\end{quote}
{\small Source: proofs/opdp13/OpHack/PowerMoment/Bound.lean:86.\par}

{\small Source SHA-256: 556b24cda4326f34d2b8e723b629b76a7930ff6a4c115f831942550fdd5fa830\par}

\end{samepage}
\begin{samepage}
\subsection{PowerFourier.function\_logarithmic\_growth}
\begin{quote}\small\ttfamily theorem function\_logarithmic\_growth : \ensuremath{\forall} \ensuremath{\varepsilon} : \ensuremath{\mathbb{R}}, 0 < \ensuremath{\varepsilon} \ensuremath{\to} \ensuremath{\exists} \ensuremath{\eta} : \ensuremath{\mathbb{R}}, 0 < \ensuremath{\eta} \ensuremath{\wedge}\newline     \ensuremath{\forall} (t : \ensuremath{\mathbb{R}}) (n : \ensuremath{\mathbb{N}}), 7 \ensuremath{\leq} t \ensuremath{\to} t < 8 \ensuremath{\to} 0 < n \ensuremath{\to}\newline       PowerSemigroup.FunctionSquaredNorm (fun z : \ensuremath{\mathbb{C}} =>\newline         (((1+Complex.normSq z)\textasciicircum{}4-t*(Complex.normSq z)\textasciicircum{}2)\textasciicircum{}n : \ensuremath{\mathbb{R}})) \ensuremath{\to} 8-t < \ensuremath{\eta} \ensuremath{\to}\newline       (1-\ensuremath{\varepsilon})*Real.log (1/(8-t)) \ensuremath{\leq} (8-t)*(n : \ensuremath{\mathbb{R}})\end{quote}
{\small Source: proofs/opdp13/OpHack/PowerFourier/Growth.lean:59.\par}

{\small Source SHA-256: dab12267b81d6895e9cfc646fedf3b3180a6ded6987210111863cd98c118e1be\par}

\end{samepage}
\begin{samepage}
\subsection{PowerUpper.function\_near\_boundary\_upper}
\begin{quote}\small\ttfamily theorem function\_near\_boundary\_upper :\newline     \ensuremath{\forall} \ensuremath{\varepsilon} : \ensuremath{\mathbb{R}}, 0 < \ensuremath{\varepsilon} \ensuremath{\to} \ensuremath{\exists} \ensuremath{\delta}0 : \ensuremath{\mathbb{R}}, 0 < \ensuremath{\delta}0 \ensuremath{\wedge} \ensuremath{\delta}0 \ensuremath{\leq} 1 \ensuremath{\wedge}\newline       \ensuremath{\forall} (\ensuremath{\delta} : \ensuremath{\mathbb{R}}) (N : \ensuremath{\mathbb{N}}), 0 < \ensuremath{\delta} \ensuremath{\to} \ensuremath{\delta} < \ensuremath{\delta}0 \ensuremath{\to}\newline         (1+\ensuremath{\varepsilon})*Real.log (1/\ensuremath{\delta})/\ensuremath{\delta} \ensuremath{\leq} (N : \ensuremath{\mathbb{R}}) \ensuremath{\to}\newline         PowerSemigroup.FunctionSquaredNorm (fun z : \ensuremath{\mathbb{C}} =>\newline           (((1+Complex.normSq z)\textasciicircum{}4-(8-\ensuremath{\delta})*(Complex.normSq z)\textasciicircum{}2)\textasciicircum{}N : \ensuremath{\mathbb{R}}))\end{quote}
{\small Source: proofs/opdp13/OpHack/PowerUpper/Function.lean:9.\par}

{\small Source SHA-256: c0b278ad2bc4d2d0297a99649e208205dcbc4136c8a3bf89c75604995a7a8e90\par}

\end{samepage}
\begin{samepage}
\subsection{PowerUpper.LeastIndex.leastPositiveIndex\_squeeze}
\begin{quote}\small\ttfamily theorem leastPositiveIndex\_squeeze :\newline     \ensuremath{\forall} \ensuremath{\varepsilon} : \ensuremath{\mathbb{R}}, 0 < \ensuremath{\varepsilon} \ensuremath{\to} \ensuremath{\exists} \ensuremath{\delta}0 : \ensuremath{\mathbb{R}}, 0 < \ensuremath{\delta}0 \ensuremath{\wedge} \ensuremath{\delta}0 \ensuremath{\leq} 1 \ensuremath{\wedge}\newline       \ensuremath{\forall} \ensuremath{\delta} : \ensuremath{\mathbb{R}}, 0 < \ensuremath{\delta} \ensuremath{\to} \ensuremath{\delta} < \ensuremath{\delta}0 \ensuremath{\to}\newline         IsLeastAdmissible \ensuremath{\delta} (leastPositiveIndex \ensuremath{\delta}) \ensuremath{\wedge}\newline         1-\ensuremath{\varepsilon} \ensuremath{\leq} \ensuremath{\delta}*(leastPositiveIndex \ensuremath{\delta} : \ensuremath{\mathbb{R}})/Real.log (1/\ensuremath{\delta}) \ensuremath{\wedge}\newline         \ensuremath{\delta}*(leastPositiveIndex \ensuremath{\delta} : \ensuremath{\mathbb{R}})/Real.log (1/\ensuremath{\delta}) \ensuremath{\leq} 1+\ensuremath{\varepsilon}\end{quote}
{\small Source: proofs/opdp13/OpHack/PowerUpper/LeastIndex.lean:40.\par}

{\small Source SHA-256: 41aa021b2f3b58b80de296aa22bf5ca917ffb9853bdaa740af8c9dcb36b442c0\par}

{\small\textbf{Sources and attribution:} \href{https://aimath.org/WWN/complexitycrmap/complexitycrmap.pdf}{1. aimath.org / complexitycrmap.pdf}; \href{https://arxiv.org/abs/1012.2479}{2. arXiv source: 1012.2479}; \href{https://doi.org/10.1007/s12220-019-00334-9}{3. DOI: 10.1007/s12220-019-00334-9}; \href{https://doi.org/10.1007/s12220-025-02228-5}{4. DOI: 10.1007/s12220-025-02228-5}; \href{https://github.com/alejandrozu/openmath-2026-judging/blob/d0ed487f487fa7ec814ff705ab55249983fe8707/OpenMath-Judging/review/entries/E02-opdp13-quartic.txt}{5. alejandrozu/openmath-2026-judging / E02-opdp13-quartic.txt}; \href{https://github.com/alejandrozu/openmath-2026-judging}{Central source and adjudication archive}.\par}

\end{samepage}
\clearpage
\section{The sharp half-turn defect gap in full no-three-in-line grid configurations}
{\small Provisional mathematical authors: Rachel Teo, Wenyi Wang, Hamdi Abderrahmene\par}\medskip
{\small Rachel Teo  -  Sakana AI.\par}

{\small Wenyi Wang  -  Sakana AI.\par}

{\small Hamdi Abderrahmene  -  Sakana AI.\par}

{\small Affiliations follow the recorded author declarations or public institutional/profile sources. Preferred author names, author order and publication affiliations await author review. This editorial compilation does not establish anthology coauthorship.\par}

\subsection{Abstract}
A no-three-collinear set of 2n points in an n\ensuremath{\times}n grid is either invariant under the half-turn or has at least four points whose partners are absent. An exact ten-point 5\ensuremath{\times}5 example attains four, giving the global minimum positive one-sided defect.

\subsection{Statement and counting convention}
Let A\ensuremath{\subseteq}\{0,...,n\ensuremath{-}1\}\ensuremath{^2} contain exactly 2n points, with no three distinct points collinear. Let R(x,y)=(n\ensuremath{-}1\ensuremath{-}x,n\ensuremath{-}1\ensuremath{-}y), and define the one-sided defect D=\{p\ensuremath{\in}A:R(p)\ensuremath{\notin}A\}. The theorem is |D|=0 or |D|\ensuremath{\geq}4. It is important to specify one-sided defect: |A\ensuremath{\triangle}R(A)|=2|D|, so the sharp symmetric-difference value is eight.

The theorem addresses the historically posed near-half-turn question for full grid configurations. It does not prove that a 2n-point configuration exists for every n, nor that defect four is attainable at every order.

\subsection{Geometric proof of the gap}
No row or column can contain more than two points. Since there are n rows and 2n points, each row contains exactly two, and the same holds for every column. Consequently the sum of the centred position vectors of all points of A is zero: its centroid is the grid centre.

Cancel all opposite pairs under R, together with the central fixed point if present. The remaining unmatched vectors, namely D, still sum to zero. A single unmatched vector would have to be zero and would therefore be fixed, a contradiction. Two unmatched vectors summing to zero would be opposite and therefore paired, another contradiction.

Suppose there were three unmatched points. The remaining centrally invariant part has 2n\ensuremath{-}3 points, an odd number, so it contains the grid centre. For n\ensuremath{\geq}3 this part has at least three points and therefore also contains a nonzero opposite pair. That pair and the centre are three collinear points, contradicting the hypothesis. The cases n=0,1,2 are handled separately. This centroid/parity geometry is necessary; row/column balance alone permits a two-miss alternating rectangle and does not prove the gap.

\subsection{Exact sharp witness}
The 5\ensuremath{\times}5 witness is \{(1,0),(2,0),(0,1),(4,1),(3,2),(4,2),(0,3),(1,3),(2,4),(3,4)\}. It has ten points. Exact determinants show that every triple of distinct points is noncollinear, and direct half-turn pairing gives four unmatched points. The formal proof checks those integer conditions and transports them to the grid geometry.

Independent ordinary checking confirmed all 120 distinct triples. Exhaustive enumeration through several small grids is consistent with the theorem but is supplementary, because the centroid argument proves all n. The universal theorem and witness are in the same contribution family.

\subsection{Priority and formal status}
The Bielefeld source records four as a minimum known example and asks about the nearest nonsymmetric cases. The no-positive-defect-below-four obstruction appears to answer that subsidiary question. Because its proof is elementary, historical and folklore comparison remains especially important; the absence of a located proof is not certification of first discovery.

The archive records a clean four-module standalone release build with standard axioms. This is a natural short note with a complete ordinary proof, a simple diagram and a formal supplement. The original source site and known defect-four examples should be credited, and the paper title should avoid a claim to solve the broad no-three-in-line problem.

\subsection{Verification and reproducibility}
Fresh execution: sakana-opdp 43: PASS; 4/4 custom modules compiled successfully; receipt Verification/sakana-opdp43-fresh-build.json. Compiler pins, source hashes and endpoint axiom outputs are in Verification. Historical receipts and finite checks do not replace fresh replay.

\begin{samepage}
\subsection{NoThreeSymmetry.grid\_half\_turn\_gap}
\begin{quote}\small\ttfamily theorem grid\_half\_turn\_gap \{n : \ensuremath{\mathbb{N}}\} (A : Finset (GridPoint n))\newline     (h : GridNoThree A) (hcard : A.card = 2*n) :\newline     (gridDefect A).card = 0 \ensuremath{\vee} 4 \ensuremath{\leq} (gridDefect A).card\end{quote}
{\small Source: proofs/opdp43/OpHack/NoThree/Main.lean:8.\par}

{\small Source SHA-256: f8262d41b0933681b1db8dafb3d4302471dc28650b8ca950e3e1fc01f203c3d2\par}

\end{samepage}
\begin{samepage}
\subsection{NoThreeSymmetry.sharp\_half\_turn\_gap}
\begin{quote}\small\ttfamily theorem sharp\_half\_turn\_gap :\newline     (\ensuremath{\forall} (n : \ensuremath{\mathbb{N}}) (A : Finset (GridPoint n)), GridNoThree A \ensuremath{\to} A.card = 2*n \ensuremath{\to}\newline       (gridDefect A).card = 0 \ensuremath{\vee} 4 \ensuremath{\leq} (gridDefect A).card) \ensuremath{\wedge}\newline     (\ensuremath{\exists} A : Finset (GridPoint 5), A.card = 10 \ensuremath{\wedge} GridNoThree A \ensuremath{\wedge} (gridDefect A).card = 4)\end{quote}
{\small Source: proofs/opdp43/OpHack/NoThree/Main.lean:56.\par}

{\small Source SHA-256: f8262d41b0933681b1db8dafb3d4302471dc28650b8ca950e3e1fc01f203c3d2\par}

{\small\textbf{Sources and attribution:} \href{https://wwwhomes.uni-bielefeld.de/achim/no3in/symmetry_remarks.html}{1. wwwhomes.uni-bielefeld.de / symmetry\_remarks.html}; \href{https://www.unsolvedmath.com/problems/GREEN-045}{2. www.unsolvedmath.com / GREEN-045}; \href{https://github.com/alejandrozu/openmath-2026-judging/blob/d0ed487f487fa7ec814ff705ab55249983fe8707/OpenMath-Judging/review/entries/E02-opdp43-halfturn.txt}{3. alejandrozu/openmath-2026-judging / E02-opdp43-halfturn.txt}; \href{https://github.com/alejandrozu/openmath-2026-judging}{Central source and adjudication archive}.\par}

\end{samepage}
\clearpage
\section{A uniform entropy bound for fibres of k\ensuremath{\sigma}(k)}
{\small Provisional mathematical authors: Rachel Teo, Wenyi Wang, Hamdi Abderrahmene\par}\medskip
{\small Rachel Teo  -  Sakana AI.\par}

{\small Wenyi Wang  -  Sakana AI.\par}

{\small Hamdi Abderrahmene  -  Sakana AI.\par}

{\small Affiliations follow the recorded author declarations or public institutional/profile sources. Preferred author names, author order and publication affiliations await author review. This editorial compilation does not establish anthology coauthorship.\par}

\subsection{Abstract}
The number f(n) of positive integers k satisfying k\ensuremath{\sigma}(k)=n has an all-target upper bound exp((C+\ensuremath{\varepsilon})log n/log log n), where C=log(2539/1000)/3+log(10/9)/2\ensuremath{\approx}0.36327036. This refines the located log 3/3 coefficient while leaving the requested zero-constant asymptotic open.

\subsection{The exact arithmetic count}
Let \ensuremath{\sigma}(k) be the sum of positive divisors of k, and let f(n)=\#\{k\ensuremath{\geq}1:k\ensuremath{\sigma}(k)=n\}. Since k\ensuremath{\sigma}(k)\ensuremath{\geq}k for positive k, the solutions form an explicit finite subset of\{0,...,n\}; the source proves identity with the canonical formal-conjecture count, including its positivity convention.

For every \ensuremath{\varepsilon}>0, the result gives a threshold N such that all n\ensuremath{\geq}N satisfy f(n)\ensuremath{\leq}exp((C+\ensuremath{\varepsilon})log n/log log n), with C=log(2539/1000)/3+log(10/9)/2. The numerical comparison is all-target: C is smaller than log 3/3\ensuremath{\approx}0.366204. A separately located log 2/2 bound for odd targets does not subsume a theorem for all n.

\subsection{Powerful cores and an exact pattern bound}
Split k into its powerful core, containing prime exponents at least two, and a coprime squarefree residual part. The squarefree map s\ensuremath{\mapsto}s\ensuremath{\sigma}(s) is injective; the formal argument removes the largest prime and treats the small primes separately. Thus a solution in a fixed target fibre can be encoded by its core exponent pattern without multiplying by an uncontrolled number of squarefree choices.

The equation k\ensuremath{\sigma}(k)=n constrains both divisibility and size of that core. The development establishes exact weighted exponent-pattern bounds, with local weights and a cubic estimate at a fixed rational parameter. These finite inequalities yield the 2539/1000 constant, rather than introducing an unproved entropy heuristic.

A logarithmic cutoff divides target primes into small and large ranges. Small-prime choices are bounded crudely but contribute a lower-order term; large-prime exponents obey a log-size budget. The proven weighted finite estimate is then converted to the exponential form, and a real asymptotic bridge absorbs lower-order terms into \ensuremath{\varepsilon}. Every premise of the intermediate conditional bridge is discharged by the final SigmaFiberAsymptotic theorem.

\subsection{Mathematical and formal boundaries}
The requested Erdős 1060 conclusion has a little-o exponent, corresponding to an arbitrarily small leading constant. A fixed positive C does not establish that conclusion. The result is a uniform quantitative partial improvement and should be labelled that way.

Fresh independent replay completed 2026-10-05T01:08:02+00:00: all ten frozen authored Lean modules compiled under Lean 4.33.1 and Mathlib 0df444a360eaa60ab8c11dca51a86af692955474. The actual printed axiom lists for canonical\_count\_eventually\_upper, canonical\_count\_upper and entropyConstant\_lt\_old contain only propext, Classical.choice and Quot.sound; all three requested prints are present, with no admitted, native-evaluation or unrecognized axioms. Source hashes, commands, individual outputs and guarded resource measurements are preserved in Verification/sakana-erdos1060-uniform-partial-fresh-build.json. This verifies the displayed uniform partial theorem, while literature priority, author approval and the full zero-constant target remain separate.

\subsection{Publication strategy and attribution}
Powerful-core decomposition, squarefree injectivity strategy and cutoff counting are prior machinery. The candidate original contribution is the exact weighted pattern improvement and its all-target entropy coefficient. The closest Kominers and related sources should be compared on their domains, constants and quantitative effectiveness before announcing best-known status.

A standalone analytic/computational number-theory note should supply the local rational inequality, the core injection and a transparent cutoff argument. The related cube-sum paper reuses some of this formal infrastructure, which should be acknowledged rather than counted as an independent discovery of the same lemmas.

\subsection{The exact arithmetic fiber and the theorem proved}
Write \ensuremath{\sigma}(k) for the sum of the positive divisors of k, and let F(n)=\#\{k\ensuremath{\geq}1:k\ensuremath{\sigma}(k)=n\}. For n\ensuremath{\geq}1 every member of this fiber satisfies k\ensuremath{\leq}n, so this is also the canonical bounded count \#\{k\ensuremath{\leq}n:k\ensuremath{\sigma}(k)=n\}. Put C=log(2539/1000)/3+log(10/9)/2, with natural logarithms. The result proved here is: for every \ensuremath{\varepsilon}>0 there is N(\ensuremath{\varepsilon}) such that, for every integer n\ensuremath{\geq}N(\ensuremath{\varepsilon}), F(n)\ensuremath{\leq}exp((C+\ensuremath{\varepsilon})log n/log log n). The quantifier is uniform in all n. The exact inequality (2539/1000)\ensuremath{^2}(10/9)\ensuremath{^3}<9 implies C<log(3)/3. In particular this improves the positive constant in the familiar powerful-core upper bound. It does not give F(n)=exp(o(log n/log log n)), which would require that every positive constant can be used.

The proof has three independent bridges. First the arithmetic fiber injects into exponent patterns of powerful cores. Second a weighted finite count controls those patterns with an exact rational inequality. Third a moving prime cutoff removes the small-prime and rounding errors. We give the bridges explicitly so that the final asymptotic statement is not resting on an unproved finite-counting hypothesis.

\subsection{Lemma 1: squarefree injectivity and the powerful-core injection}
For a squarefree integer s, multiplicativity gives s\ensuremath{\sigma}(s)=\ensuremath{\prod}\_\{p|s\}p(p+1). These products determine their prime supports. Indeed, let p>3 be the largest prime in the union of two supports with equal products. Divisibility by p forces a factor q(q+1) on each side with q\ensuremath{\leq}p. If p divides q then q=p. If p divides q+1 and q<p, then q+1\ensuremath{\leq}p, hence q=p\ensuremath{-}1. This is impossible for primes p>3: p is odd, so p\ensuremath{-}1 is an even integer greater than two. Thus p occurs on both sides. Cancel its positive factor p(p+1), and continue inductively. When only 2 and 3 remain, the four possible products are 1, 6, 12 and 72; they are distinct. Consequently s\ensuremath{\mapsto}s\ensuremath{\sigma}(s) is injective on squarefree positive integers.

Factor k uniquely as k=P(k)s(k), where P(k)=\ensuremath{\prod}\_\{v\_p(k)\ensuremath{\geq}2\}p\textasciicircum{}\{v\_p(k)\} and s(k) is the product of the primes occurring to exponent exactly one. The factors P(k) and s(k) are coprime; s(k) is squarefree. The divisor-sum function is multiplicative on coprime arguments, so k\ensuremath{\sigma}(k)=P(k)\ensuremath{\sigma}(P(k))\ensuremath{\cdot}s(k)\ensuremath{\sigma}(s(k)). If two members of the same fiber have the same core P, cancellation of P\ensuremath{\sigma}(P) leaves equal squarefree products, and the preceding lemma gives equal residuals. Thus k\ensuremath{\mapsto}P(k) is injective on a fixed fiber.

If n=\ensuremath{\prod}p\textasciicircum{}\{a\_p\}, then k divides n and each exponent of its powerful core belongs to \{0\}\ensuremath{\cup}\{2,...,a\_p\}. Also \ensuremath{\sigma}(k)\ensuremath{\geq}k, so k\ensuremath{^2}\ensuremath{\leq}n and therefore P(k)\ensuremath{^2}\ensuremath{\leq}n. In logarithmic form every core pattern e=(e\_p) satisfies \ensuremath{\sum}e\_p log p\ensuremath{\leq}(log n)/2. The image of the fiber need not be all patterns satisfying these conditions. Counting all admissible patterns gives an upper bound, and the formal source also proves equality between the fiber count and the actual image of the core map. No surjectivity onto the larger majorant is required.

\subsection{Lemma 2: an exact weighted local inequality}
For a\ensuremath{\geq}1 set L\_a(x)=1+\ensuremath{\sum}\_\{j=2\}\textasciicircum{}a x\textasciicircum{}j, where an empty sum is zero. Put x=9/10 and B=2539/1000. The crucial exact estimate is L\_a(x)\ensuremath{^3}\ensuremath{\leq}B\textasciicircum{}a for every a\ensuremath{\geq}1. The cases a=1 and a=2 are direct rational calculations. For a=3, L\_3(x)=1+x\ensuremath{^2}+x\ensuremath{^3}=B, so equality holds. The remaining initial calculation is L\_4(x)\ensuremath{^3}\ensuremath{\leq}B\ensuremath{^4}, again an inequality between rational numbers.

To prove every larger exponent at once, observe that for a\ensuremath{\geq}3, L\_a(x)\ensuremath{\geq}B and x\textasciicircum{}\{a+1\}\ensuremath{\leq}x\ensuremath{^4}. Since L\_\{a+1\}(x)=L\_a(x)+x\textasciicircum{}\{a+1\}, these two inequalities give B L\_\{a+1\}(x)\ensuremath{\leq}L\_a(x)(B+x\ensuremath{^4})=L\_a(x)L\_4(x). Cubing, applying the induction hypothesis and the a=4 calculation, and cancelling B\ensuremath{^3} gives L\_\{a+1\}(x)\ensuremath{^3}\ensuremath{\leq}B\textasciicircum{}\{a+1\}. This is a uniform argument over all exponents; it is not a numerical search over a selected exponent range.

More generally, let S be a finite set of exponent patterns satisfying e\_p\ensuremath{\in}\{0\}\ensuremath{\cup}\{2,...,a\_p\} and \ensuremath{\sum}e\_p\ensuremath{\leq}m. Since 0<x\ensuremath{\leq}1, each term x\textasciicircum{}\{\ensuremath{\sum}e\_p\} is at least x\textasciicircum{}m. Summing and then enlarging S to the full product of coordinate choices gives |S|x\textasciicircum{}m\ensuremath{\leq}\ensuremath{\prod}\_p L\_\{a\_p\}(x). Cubing and using the local estimate yields |S|\ensuremath{^3}x\textasciicircum{}\{3m\}\ensuremath{\leq}B\textasciicircum{}\{\ensuremath{\sum}a\_p\}. This elementary weighted count is the finite source of the improved constant.

\subsection{Lemma 3: the finite prime-cutoff bound}
Fix n\ensuremath{\geq}2 and a real y>1. Separate target primes p\ensuremath{\leq}y from p>y. The small-prime coordinates are counted without a weight: the set \{0\}\ensuremath{\cup}\{2,...,a\_p\} has exactly a\_p members, including when a\_p=1. Write Q=\ensuremath{\prod}\_\{p\ensuremath{\leq}y, p|n\}a\_p and S=\ensuremath{\sum}\_\{p>y, p|n\}a\_p. On the large-prime coordinates the core budget gives \ensuremath{\sum}\_\{p>y\}e\_p\ensuremath{\leq}log n/(2log y). Set m=ceil(log n/(2log y)). Apply the weighted count to those coordinates and the unweighted product to the others. The resulting finite inequality is F(n)\ensuremath{^3}x\textasciicircum{}\{3m\}\ensuremath{\leq}Q\ensuremath{^3}B\textasciicircum{}S.

The target factorization gives S log y\ensuremath{\leq}\ensuremath{\sum}a\_p log p=log n. Hence S\ensuremath{\leq}R=log n/log y. Moreover m\ensuremath{\leq}R/2+1. Taking logarithms of the positive upper bound, with the case F(n)=0 immediate, gives log F(n)\ensuremath{\leq}log Q+(log B/3\ensuremath{-}log x/2)R\ensuremath{-}log x. The coefficient is precisely C, and the last term is D=log(10/9).

No prime number theorem is needed to estimate Q. There are at most y distinct positive integers p\ensuremath{\leq}y, and every target exponent satisfies a\_p\ensuremath{\leq}log n/log 2. Thus log Q\ensuremath{\leq}y log(1+log n/log 2). We obtain the unconditional finite estimate F(n)\ensuremath{\leq}exp(y log(1+log n/log 2)+C log n/log y+D). It is valid for every n\ensuremath{\geq}2 and every real y>1. In the formal chain, SigmaFiberEntropy proves this estimate on the actual arithmetic fiber before SigmaFiberAsymptotic calls the analytic bridge.

\subsection{The cutoff asymptotics and the strict improvement}
Put L=log n and choose y=L/(log L)\ensuremath{^3} once L is sufficiently large. Then y>1 and log y=log L\ensuremath{-}3log log L. Since log log L/log L tends to zero, (log L)/(log y) tends to one. Normalize the exponent in the finite estimate by L/log L. Its large-prime term tends to C. Its small-prime term becomes log(1+L/log 2)/(log L)\ensuremath{^2} and tends to zero: the numerator differs from log L by a bounded quantity, while the denominator is quadratic in log L. Its rounding term becomes D log L/L and also tends to zero. The normalized total therefore tends to C.

For every \ensuremath{\varepsilon}>0 this convergence makes the finite exponent at most (C+\ensuremath{\varepsilon})L/log L beyond a threshold, which proves the theorem for all sufficiently large integers n. Finally, 6C=log(B\ensuremath{^2}(10/9)\ensuremath{^3})<log 9=2log 3. Division by six gives C<log 3/3. This comparison uses exact rational arithmetic inside the logarithm, rather than rounded decimal constants. The structural injection and the cutoff method should be distinguished from the particular weighted improvement when crediting earlier literature.

\subsection{Fresh independent formal verification of the exact theorem}
Fresh independent replay completed 2026-10-05T01:08:02+00:00: all ten frozen authored Lean modules compiled under Lean 4.33.1 and Mathlib 0df444a360eaa60ab8c11dca51a86af692955474. The actual printed axiom lists for canonical\_count\_eventually\_upper, canonical\_count\_upper and entropyConstant\_lt\_old contain only propext, Classical.choice and Quot.sound; all three requested prints are present, with no admitted, native-evaluation or unrecognized axioms. Source hashes, commands, individual outputs and guarded resource measurements are preserved in Verification/sakana-erdos1060-uniform-partial-fresh-build.json. This verifies the displayed uniform partial theorem, while literature priority, author approval and the full zero-constant target remain separate.

\subsection{Verification and reproducibility}
Fresh execution: sakana-erdos1060-uniform-partial: PASS; 10/10 custom modules compiled successfully; receipt Verification/sakana-erdos1060-uniform-partial-fresh-build.json. Compiler pins, source hashes and endpoint axiom outputs are in Verification. Historical receipts and finite checks do not replace fresh replay.

\begin{samepage}
\subsection{Erdos1060.canonical\_count\_eventually\_upper}
\begin{quote}\small\ttfamily theorem canonical\_count\_eventually\_upper \{\ensuremath{\varepsilon} : \ensuremath{\mathbb{R}}\} (h\ensuremath{\varepsilon} : 0 < \ensuremath{\varepsilon}) :\newline     \ensuremath{\forall}\ensuremath{^{f}} n : \ensuremath{\mathbb{N}} in atTop,\newline       (\#\{k \ensuremath{\leq} n | k * ArithmeticFunction.sigma 1 k = n\} : \ensuremath{\mathbb{R}}) \ensuremath{\leq}\newline         Real.exp (((Real.log (2539 / 1000) / 3 + Real.log (10 / 9) / 2) + \ensuremath{\varepsilon}) *\newline           Real.log (n : \ensuremath{\mathbb{R}}) / Real.log (Real.log (n : \ensuremath{\mathbb{R}})))\end{quote}
{\small Source: proofs/erdos1060-uniform-partial/src/SigmaFiberAsymptotic.lean:38.\par}

{\small Source SHA-256: 440f2f6aaa1d5dd04611fe855077352a3b75171d6e1ec823939ede8609e2d45a\par}

\end{samepage}
\begin{samepage}
\subsection{Erdos1060.canonical\_count\_upper}
\begin{quote}\small\ttfamily theorem canonical\_count\_upper \{\ensuremath{\varepsilon} : \ensuremath{\mathbb{R}}\} (h\ensuremath{\varepsilon} : 0 < \ensuremath{\varepsilon}) :\newline     \ensuremath{\exists} N : \ensuremath{\mathbb{N}}, \ensuremath{\forall} n : \ensuremath{\mathbb{N}}, N \ensuremath{\leq} n \ensuremath{\to}\newline       (\#\{k \ensuremath{\leq} n | k * ArithmeticFunction.sigma 1 k = n\} : \ensuremath{\mathbb{R}}) \ensuremath{\leq}\newline         Real.exp (((Real.log (2539 / 1000) / 3 + Real.log (10 / 9) / 2) + \ensuremath{\varepsilon}) *\newline           Real.log (n : \ensuremath{\mathbb{R}}) / Real.log (Real.log (n : \ensuremath{\mathbb{R}})))\end{quote}
{\small Source: proofs/erdos1060-uniform-partial/src/SigmaFiberAsymptotic.lean:49.\par}

{\small Source SHA-256: 440f2f6aaa1d5dd04611fe855077352a3b75171d6e1ec823939ede8609e2d45a\par}

\end{samepage}
\begin{samepage}
\subsection{Erdos1060Partial.entropyConstant\_lt\_old}
\begin{quote}\small\ttfamily theorem entropyConstant\_lt\_old : entropyConstant < Real.log 3 / 3\end{quote}
{\small Source: proofs/erdos1060-uniform-partial/src/EntropyConversion.lean:45.\par}

{\small Source SHA-256: a9c629fe04c5f4252e92ebdef87ba5942f689fb64a0ca7e1fbc46f8f3957dc14\par}

{\small\textbf{Sources and attribution:} \href{https://www.erdosproblems.com/1060}{1. www.erdosproblems.com / 1060}; \href{https://scottkom.com/assets/articles/Kominers_ksigmak.pdf}{2. scottkom.com / Kominers\_ksigmak.pdf}; \href{https://github.com/google-deepmind/formal-conjectures/blob/main/FormalConjectures/ErdosProblems/1060.lean}{3. google-deepmind/formal-conjectures / 1060.lean}; \href{https://github.com/alejandrozu/ulam-opdp-difficulty-atlas}{4. alejandrozu/ulam-opdp-difficulty-atlas}; \href{https://github.com/alejandrozu/openmath-2026-judging/blob/d0ed487f487fa7ec814ff705ab55249983fe8707/OpenMath-Judging/review/entries/E02-erdos1060-uniform-partial.txt}{5. alejandrozu/openmath-2026-judging / E02-erdos1060-uniform-partial.txt}; \href{https://github.com/alejandrozu/openmath-2026-judging}{Central source and adjudication archive}.\par}

\end{samepage}
\clearpage
\section{A log(5/3) uniform bound for representations as two nonnegative cubes}
{\small Provisional mathematical authors: Rachel Teo, Wenyi Wang, Hamdi Abderrahmene\par}\medskip
{\small Rachel Teo  -  Sakana AI.\par}

{\small Wenyi Wang  -  Sakana AI.\par}

{\small Hamdi Abderrahmene  -  Sakana AI.\par}

{\small Affiliations follow the recorded author declarations or public institutional/profile sources. Preferred author names, author order and publication affiliations await author review. This editorial compilation does not establish anthology coauthorship.\par}

\subsection{Abstract}
For ordered nonnegative representations x\ensuremath{^3}+y\ensuremath{^3}=n, the formal endpoint gives the eventual bound exp((log(5/3)+\ensuremath{\varepsilon})log n/log log n). The entrant-supplied ordinary proof optimizes the same weighted ternary coding method to coefficient C=log 3\ensuremath{-}H((2\ensuremath{-}\ensuremath{\surd}2)/3)\ensuremath{-}((2\ensuremath{-}\ensuremath{\surd}2)/3)log 2\ensuremath{\approx}0.4695030882. A gcd transfer and prime cutoff remove primitive and cube-free restrictions. The stronger ordinary endpoint is not asserted as Lean-checked.

\subsection{Count and uniform theorem}
Let R(n)=\#\{(x,y)\ensuremath{\in}\ensuremath{\mathbb{N}}\ensuremath{^2}:x\ensuremath{^3}+y\ensuremath{^3}=n\}, with ordered pairs and zero coordinates included. For every \ensuremath{\varepsilon}>0 there is N such that R(n)\ensuremath{\leq}exp((log(5/3)+\ensuremath{\varepsilon})log n/log log n) for every n\ensuremath{\geq}N. No cubefree-target restriction or gcd-one restriction appears in the final theorem.

The coefficient log(5/3)\ensuremath{\approx}0.510826 improves the earlier same-family log(7/4)\ensuremath{\approx}0.559616 endpoint. The located approximately0.556477 comparison applies only to cubefree targets. It therefore cannot be substituted as a stronger all-target theorem. The result still falls far short of a fixed polylogarithmic bound, which remains the original target.

\subsection{Primitive ratios as a separated ternary code}
For a primitive positive solution, the ratio of its coordinates is a unit modulo the relevant prime powers, with cube equal to minus one. The unit cube-root lemma gives at most three choices at each prime-power coordinate, including the ramified prime three. Distinct solutions are encoded by distinct ternary words.

If two encoded words agree at selected prime-power coordinates, those moduli divide a cross determinant of the two solutions. An upper bound for the determinant and the target-size product force a weighted separation property. This converts an arithmetic representation problem into a weighted code problem without assuming the desired representation bound.

Elias/Plotkin localization bounds the size of a separated code in a weighted ball. An exact binomial block certificate uses blocks of 500 coordinates and support 90, producing the rational base 5/3 with an explicit finite multiplicative constant. The constant is harmless for the eventual leading exponent, but its proof is retained rather than replaced by a decimal numerical guess.

\subsection{From primitive pairs to the canonical count}
A general positive pair has gcd d and a primitive pair representing n/d\ensuremath{^3}. The source proves the exact gcd-fibre identity and sums over positive cube divisors d. Coordinatewise multiplicative weights bound this sum uniformly across targets, including those with repeated prime factors.

The prime cutoff and elementary log budgets control the resulting product. They reuse background modules from the k\ensuremath{\sigma}(k) development, while the cubic encoding and local-factor estimates are specific to this result. Finally at most two boundary pairs have a zero coordinate, and their contribution is absorbed. A semantic bridge identifies the resulting count with the original ordered cube-value representation definition; it does not import the open conjecture as a hypothesis.

\subsection{Verification, prior work and manuscript obligations}
The source contains 33 proof modules and two semantic/axiom audit modules in the recorded replay. Some arithmetic files were reconstructed after a workspace replacement and are clearly labelled variants; the dated fresh archived receipt supersedes their stale pending comments, without pretending the missing historical bytes were recovered. The final transitive axioms in those records are standard only.

Merikoski's ratio-root/determinant methods and classical coding/entropy bounds must be credited. The proposed new delta is the exact 5/3 formal weighted-code constant, the optimized ordinary coefficient 0.4695030882..., and their all-target transfer. This is a quantitative partial-bound paper, not a solution of Erdős 829 or a signed-integer representation theorem. A human exposition should display the weighted-distance condition, the finite block inequality and the gcd factorization so that the formal chain remains intelligible.

\subsection{Fresh independent verification of the displayed finite/formal scope}
Fresh independent replay completed 2026-10-05T01:22:59Z: all 31 frozen authored sources compiled, including the archived entrant-delivered reconstructed variants, under exact Lean 4.33.1 and Mathlib 0df444a360eaa60ab8c11dca51a86af692955474. Both actual requested axiom prints, canonicalCubeCount\_eventually\_exp and canonicalCubeCount\_eventual\_threshold, are present and contain only propext, Classical.choice and Quot.sound, with no admitted, native-evaluation or unknown axioms. This verifies the displayed all-target log(5/3) endpoint; the sharper ordinary coefficient .46950308823856934 remains a separately explained source result. The replay proves the delivered reconstructed proof, without claiming recovery of the first lost source bytes. Receipt: Verification/sakana-erdos829-log-five-thirds-fresh-build.json.

\subsection{The representation count and its precise scope}
Let R(n)=\#\{(x,y)\ensuremath{\in}\ensuremath{\mathbb{N}}\ensuremath{^2}:x\ensuremath{^3}+y\ensuremath{^3}=n\}, where the pairs are ordered and zero coordinates are allowed. The claimed improvement is that, for every \ensuremath{\varepsilon}>0, R(n)\ensuremath{\leq}exp((log(5/3)+\ensuremath{\varepsilon})log n/log log n) for every sufficiently large integer n. This is a pointwise bound over all targets, including targets with repeated prime factors. It is a positive-constant upper bound, rather than the n\textasciicircum{}\{o(1)\} or boundedness conclusions appearing in stronger versions of the representation problem. The proof first treats positive primitive pairs and then restores the common divisor and zero-coordinate cases.

The arithmetic encoding into modular cubic roots follows an established coding approach. The contribution described here is the explicit weighted ternary Elias estimate, its finite 500-coordinate certificate, and the complete transfer back to the canonical all-target count. The proposed coefficient log(5/3) is about 0.510826. A bound proved only for cube-free targets has different quantifiers and does not by itself settle this all-target statement. Any claim of priority should be checked against that distinction and the full literature, rather than only compared by decimal coefficients.

\subsection{Lemma 1: primitive solutions form a separated weighted ternary code}
Suppose x,y>0, gcd(x,y)=1 and x\ensuremath{^3}+y\ensuremath{^3}=n. Then x and y are each coprime to n: a prime dividing x and n would divide y\ensuremath{^3}, hence y, contradicting primitivity. For every full prime power p\textasciicircum{}a dividing n, the unit r\_p=\ensuremath{-}x/y modulo p\textasciicircum{}a satisfies r\_p\ensuremath{^3}=1. There are at most three such units. For odd p the unit group modulo p\textasciicircum{}a is cyclic, so the kernel of the cubing map has at most three elements, including when p=3. For p=2 the unit group has order a power of two; cubing is an automorphism because three is coprime to that order. Thus its only cube root of one is one. Label the available roots independently by symbols in the three-element alphabet \ensuremath{\mathbb{Z}}/3\ensuremath{\mathbb{Z}}.

Give the coordinate p weight w\_p=a log p. Its total weight is W=log n. Consider two distinct positive primitive solutions (x,y) and (u,v). Their integer determinant D=xv\ensuremath{-}uy is nonzero. Indeed, D=0 would give equal reduced positive fractions x/y=u/v, which forces x=u and y=v. Each coordinate product xv and uy has cube strictly less than n\ensuremath{^2} because every positive coordinate has cube less than n. Positivity therefore gives |D|\ensuremath{^3}<n\ensuremath{^2}: the absolute difference of two nonnegative products is at most their larger product.

If two codewords agree at p, their modular ratios agree, and multiplication by the invertible denominators gives p\textasciicircum{}a|D. The agreeing full prime powers are pairwise coprime, so their product A divides |D|. Consequently A\ensuremath{^3}<n\ensuremath{^2} and log A<(2/3)log n. The sum of the weights of the disagreeing coordinates is W\ensuremath{-}log A>W/3. This also proves that the encoding is injective. Thus actual primitive solutions, rather than hypothetical modular data, form a weighted ternary code with strict distance greater than one third of the total weight.

\subsection{Lemma 2: a weighted radius-19/100 ball contains at most 45 separated words}
For a ternary word z, write d\_w(z,0)=\ensuremath{\sum}\_\{i:z\_i\ensuremath{\neq}0\}w\_i. Let a code C of size M lie in the ball d\_w(z,0)\ensuremath{\leq}\ensuremath{\alpha}W, where 0\ensuremath{\leq}\ensuremath{\alpha}\ensuremath{\leq}2/3. At coordinate i let n\_0,n\_1,n\_2 be the symbol counts and t=n\_1+n\_2. The number of ordered mismatches at that coordinate is M\ensuremath{^2}\ensuremath{-}n\_0\ensuremath{^2}\ensuremath{-}n\_1\ensuremath{^2}\ensuremath{-}n\_2\ensuremath{^2}. Since n\_0=M\ensuremath{-}t and n\_1\ensuremath{^2}+n\_2\ensuremath{^2}\ensuremath{\geq}t\ensuremath{^2}/2, this is at most 2Mt\ensuremath{-}(3/2)t\ensuremath{^2}. Completing the square around t=\ensuremath{\alpha}M gives the affine tangent bound (2\ensuremath{-}3\ensuremath{\alpha})Mt+(3/2)\ensuremath{\alpha}\ensuremath{^2}M\ensuremath{^2}.

Sum these bounds with the nonnegative coordinate weights. The ball condition gives \ensuremath{\sum}\_i w\_i t\_i=\ensuremath{\sum}\_\{z\ensuremath{\in}C\}d\_w(z,0)\ensuremath{\leq}M\ensuremath{\alpha}W. Since 2\ensuremath{-}3\ensuremath{\alpha}\ensuremath{\geq}0, the total ordered pair distance is at most (2\ensuremath{\alpha}\ensuremath{-}(3/2)\ensuremath{\alpha}\ensuremath{^2})M\ensuremath{^2}W. Separation, including the diagonal correction, gives the lower bound M(M\ensuremath{-}1)W/3. At \ensuremath{\alpha}=19/100 the upper coefficient is 6517/20000. Therefore (M\ensuremath{^2}\ensuremath{-}M)/3\ensuremath{\leq}(6517/20000)M\ensuremath{^2}, equivalently (449/20000)M\ensuremath{^2}\ensuremath{\leq}M. If M>0, then M\ensuremath{\leq}20000/449<45; the stated safe integer bound M\ensuremath{\leq}45 follows. Keeping the factor 1/3 on the diagonal correction is essential to this computation.

Let B be the full weighted ball of radius 19W/100 in the ambient group (\ensuremath{\mathbb{Z}}/3\ensuremath{\mathbb{Z}})\textasciicircum{}s. For each z, consider pairs (c,b)\ensuremath{\in}C\ensuremath{\times}B with c+b=z. Projection to b is injective on this fiber, and translation preserves weighted coordinate mismatch. Hence the b-projections are a separated code in B and the fiber has at most 45 elements. Double-counting over all 3\textasciicircum{}s values of z yields |C||B|\ensuremath{\leq}45\ensuremath{\cdot}3\textasciicircum{}s. This argument accommodates arbitrary nonnegative weights; equal weights are not assumed.

\subsection{Lemma 3: a uniform ball volume and the exact 500/90 bridge}
Choose uniformly a t-element subset I of the s coordinates. The sum of its weights has average tW/s, because each coordinate occurs in exactly binom(s\ensuremath{-}1,t\ensuremath{-}1) subsets. Markov counting shows that at least (1\ensuremath{-}100t/(19s))binom(s,t) of those supports have total weight at most 19W/100, provided 100t<19s. Every such support gives 2\textasciicircum{}t words in the ball, by assigning either nonzero symbol at each support coordinate. Distinct supports give disjoint sets of words. Thus (19s\ensuremath{-}100t)binom(s,t)2\textasciicircum{}t\ensuremath{\leq}19s|B|.

Put m=floor(s/500) and t=90m. Then 500m\ensuremath{\leq}s<500m+500 and 19s\ensuremath{-}100t\ensuremath{\geq}s. A blockwise choice of 90 positions in each of m disjoint 500-position blocks proves binom(s,90m)\ensuremath{\geq}binom(500,90)\textasciicircum{}m. The ball-volume inequality therefore gives [binom(500,90)2\textasciicircum{}90]\textasciicircum{}m\ensuremath{\leq}19|B|. Combine it with packing to obtain |C|[binom(500,90)2\textasciicircum{}90]\textasciicircum{}m\ensuremath{\leq}855\ensuremath{\cdot}3\textasciicircum{}s.

The exact finite certificate is 3\textasciicircum{}500\ensuremath{\leq}(5/3)\textasciicircum{}500 binom(500,90)2\textasciicircum{}90. After clearing denominators it is simply the integer inequality 9\textasciicircum{}500\ensuremath{\leq}5\textasciicircum{}500 binom(500,90)2\textasciicircum{}90. This can be checked with integer arithmetic, and the source proves it by exact rational normalization. Raising it to m and using s\ensuremath{\leq}500m+500 gives 3\textasciicircum{}s\ensuremath{\leq}3\textasciicircum{}500(5/3)\textasciicircum{}\{500m\}[binom(500,90)2\textasciicircum{}90]\textasciicircum{}m. Cancelling the positive block factor proves |C|\ensuremath{\leq}855\ensuremath{\cdot}3\textasciicircum{}500(5/3)\textasciicircum{}s. The zero-coordinate ambient cube has one word and is covered separately by the same constant. Applied to the preceding arithmetic code, this bounds the primitive positive count by a constant times (5/3)\textasciicircum{}\{\ensuremath{\omega}(n)\}.

\subsection{Lemma 4: restore all common divisors through an exact Euler product}
Every positive representation has a unique common divisor d=gcd(x,y). Writing x=du and y=dv produces a primitive positive representation of n/d\ensuremath{^3}, and necessarily d\ensuremath{^3}|n. Conversely every such primitive representation scales back uniquely. Therefore the positive ordered count is exactly the sum over d\ensuremath{^3}|n of the primitive counts of n/d\ensuremath{^3}. This bridge is needed because the theorem is not restricted to cube-free targets.

Write B\_0=5/3 and n=\ensuremath{\prod}p\textasciicircum{}\{a\_p\}. The weight B\_0\textasciicircum{}\{\ensuremath{\omega}(n/d\ensuremath{^3})\} factors coordinatewise. If d has exponent j at p, with 0\ensuremath{\leq}j\ensuremath{\leq}floor(a\_p/3), the local contribution is one when 3j=a\_p and B\_0 otherwise. The sum of all local contributions is H(a)=floor(a/3)B\_0+(1 if a\ensuremath{\equiv}0 mod 3, otherwise B\_0). Hence the weighted divisor sum is at most \ensuremath{\prod}\_\{p|n\}H(a\_p). In fact the ordinary exponent parametrization gives equality; the formal upper-bound proof only needs injection of divisor exponent coordinates.

For every a\ensuremath{\geq}0, H(a)\ensuremath{\leq}B\_0\textasciicircum{}a and H(a)\ensuremath{\leq}2(a+1). The first estimate has direct cases a=0,1,2; for larger a use H(a+3)=H(a)+B\_0, H(a)\ensuremath{\geq}1 and 1+B\_0\ensuremath{\leq}B\_0\ensuremath{^3}. Induction in steps of three then gives H(a+3)\ensuremath{\leq}H(a)B\_0\ensuremath{^3}\ensuremath{\leq}B\_0\textasciicircum{}\{a+3\}. The linear bound follows directly from the formula. The source uses the harmless primitive prefactor max(1,\ensuremath{\omega}(n))3\textasciicircum{}507, since 855\ensuremath{\cdot}3\textasciicircum{}500\ensuremath{\leq}3\textasciicircum{}507. It obtains R\_+(n)\ensuremath{\leq}max(1,\ensuremath{\omega}(n))3\textasciicircum{}507\ensuremath{\prod}H(a\_p), where R\_+ counts only positive coordinates.

\subsection{The all-target cutoff and zero-coordinate completion}
Set U=1+log n/log 2 and fix a real cutoff y>1. There are at most y target primes p\ensuremath{\leq}y, and each a\_p\ensuremath{\leq}log n/log 2. For these primes the linear estimate gives log H(a\_p)\ensuremath{\leq}log(2U)\ensuremath{\leq}2log U, because n\ensuremath{\geq}2 implies U\ensuremath{\geq}2. For primes p>y, the exponential estimate and \ensuremath{\sum}\_\{p>y\}a\_p\ensuremath{\leq}log n/log y give \ensuremath{\sum}log H(a\_p)\ensuremath{\leq}log(5/3)log n/log y. The prefactor max(1,\ensuremath{\omega}(n)) is at most U, and log U\ensuremath{\leq}y log U. Consequently R\_+(n)\ensuremath{\leq}exp(3y log U+log(5/3)log n/log y+507log 3).

Choose y=L/(log L)\ensuremath{^3} with L=log n. Divide the exponent by L/log L. The small-prime contribution tends to zero; the cutoff ratio log L/log y tends to one; and the constant term tends to zero. Thus, for every \ensuremath{\varepsilon}>0, R\_+(n)\ensuremath{\leq}exp((log(5/3)+\ensuremath{\varepsilon})log n/log log n) for all sufficiently large n. There are at most two representations with a zero coordinate, namely (0,a) and (a,0) when n=a\ensuremath{^3}. Apply the positive estimate with \ensuremath{\varepsilon}/2 and absorb this bounded addition into the remaining \ensuremath{\varepsilon}/2 margin, whose exponential grows without bound. This proves the stated bound on R(n).

The final canonical-count bridge observes that the map (x,y)\ensuremath{\mapsto}(x\ensuremath{^3},y\ensuremath{^3}) is injective, so counting pairs of cube values and counting pairs of their nonnegative roots give exactly the same number. The manuscript therefore reaches the originally defined representation function, not just an intermediate code cardinality. Some frozen source modules describe reconstructed bytes after a workspace replacement. That provenance must be retained in the reproducibility record, while the current independent replay manifest determines whether their formal endpoints have actually been rebuilt.

\subsection{The optimized ordinary theorem: coefficient 0.4695030882...}
The entrant-supplied ordinary exposition proves a stronger coefficient than the rational formal endpoint. Put a=(2\ensuremath{-}\ensuremath{\surd}2)/3, H(t)=\ensuremath{-}t log t\ensuremath{-}(1\ensuremath{-}t)log(1\ensuremath{-}t), and C\_*=log 3\ensuremath{-}H(a)\ensuremath{-}a log 2=0.46950308823856934.... The ordinary theorem is that for every \ensuremath{\varepsilon}>0, R(n)\ensuremath{\leq}exp((C\_*+\ensuremath{\varepsilon})log n/log log n) for all sufficiently large n. The submitted Lean theorem remains the weaker coefficient log(5/3). The following optimization explains the stronger ordinary result without attributing its verification to Lean.

For any fixed 0<\ensuremath{\alpha}<a, the preceding localized average-distance calculation yields M(M\ensuremath{-}1)/3\ensuremath{\leq}(2\ensuremath{\alpha}\ensuremath{-}(3/2)\ensuremath{\alpha}\ensuremath{^2})M\ensuremath{^2}. Its denominator 1\ensuremath{-}6\ensuremath{\alpha}+(9/2)\ensuremath{\alpha}\ensuremath{^2} is positive, so M\ensuremath{\leq}K\_\ensuremath{\alpha}=[1\ensuremath{-}6\ensuremath{\alpha}+(9/2)\ensuremath{\alpha}\ensuremath{^2}]\textasciicircum{}\{-1\}; increasing K\_\ensuremath{\alpha} to an integer gives a harmless safe fiber bound. Double-counting translated ball incidences therefore gives |C|V(\ensuremath{\alpha}W)\ensuremath{\leq}K\_\ensuremath{\alpha} 3\textasciicircum{}s, uniformly in every choice of weights.

Fix 0<\ensuremath{\beta}<\ensuremath{\alpha} and take t=floor(\ensuremath{\beta}s). The same first-moment support count gives V(\ensuremath{\alpha}W)\ensuremath{\geq}(1\ensuremath{-}t/(\ensuremath{\alpha}s))binom(s,t)2\textasciicircum{}t. Its prefactor is bounded below by the positive constant 1\ensuremath{-}\ensuremath{\beta}/\ensuremath{\alpha}. Stirling's formula gives log binom(s,floor(\ensuremath{\beta}s))=H(\ensuremath{\beta})s+O(log(s+1)), with constants depending only on \ensuremath{\beta}. Thus for every \ensuremath{\eta}>0 there is K\_\{\ensuremath{\alpha},\ensuremath{\beta},\ensuremath{\eta}\}, independent of s and of the coordinate weights, such that |C|\ensuremath{\leq}K\_\{\ensuremath{\alpha},\ensuremath{\beta},\ensuremath{\eta}\}exp((log 3\ensuremath{-}H(\ensuremath{\beta})\ensuremath{-}\ensuremath{\beta}log 2+\ensuremath{\eta})s). The zero-coordinate and finitely many small-coordinate cases can be absorbed in that one constant.

The function H(t)+t log 2 is increasing for 0<t<2/3, since its derivative is log(2(1\ensuremath{-}t)/t). First choose \ensuremath{\beta}<\ensuremath{\alpha}<a close enough to a, and then choose \ensuremath{\eta}>0 small enough. For every B>exp(C\_*) the preceding estimate consequently gives P(n)\ensuremath{\leq}K\_B B\textasciicircum{}\{\ensuremath{\omega}(n)\}, with one constant K\_B independent of the target. This order of choices is what makes the arithmetic transfer uniform. The coefficient is approached from above; no estimate is asserted at a radius whose localized denominator vanishes.

Choose exp(C\_*)<B<2. We may also ensure B>3/2. To see the latter threshold lies below exp(C\_*), note a<1/5 and the monotonicity just established. Hence C\_*>log 3\ensuremath{-}H(1/5)\ensuremath{-}(1/5)log 2. Exponentiating five times reduces comparison with log(3/2) to the exact inequality 3\textasciicircum{}5(1/5)(4/5)\textasciicircum{}4/2>(3/2)\textasciicircum{}5. Thus B\ensuremath{^3}\ensuremath{\geq}B+1 and B\ensuremath{^3}\ensuremath{\geq}2, which are the only local-factor inequalities required.

The gcd decomposition now gives R\_+(n)\ensuremath{\leq}K\_B\ensuremath{\sum}\_\{d\ensuremath{^3}|n\}B\textasciicircum{}\{\ensuremath{\omega}(n/d\ensuremath{^3})\}. The local factor for exponent e=3q+r is qB+1 if r=0, and (q+1)B otherwise. It is at most B\textasciicircum{}e and at most 2(e+1). For the exponential bound check e=0,1,2,3 directly; when e\ensuremath{\geq}1 the factor is at least B and T\_\{e+3\}=T\_e+B\ensuremath{\leq}2T\_e\ensuremath{\leq}B\ensuremath{^3}T\_e. The e=0 to e=3 step uses B+1\ensuremath{\leq}B\ensuremath{^3}. These inequalities prove the all-exponent claim by induction.

The same elementary prime cutoff then bounds the logarithm of this Euler product by (log B+o(1))log n/log log n; log K\_B is negligible. The small-prime part uses only the number of integers up to y and the linear local bound, with y=log n/(log log n)\ensuremath{^3}. Let B decrease to exp(C\_*), and absorb the at most two zero-coordinate representations. This proves the stronger ordinary theorem. It continues to leave the original polylogarithmic target open.

The archived verification receipt states that the five reconstructed modules were part of the entrants' final active source set and passed their Oct 2 clean replay alongside 28 historical modules and two audits. The entrant declaration attributes all supplied results to their team. Their old pending comments do not override that later receipt. The independent current replay should be reported for the actual frozen reconstructed bytes, without suggesting that the lost historical versions were recovered or inventing organizer authorship.

\subsection{Fresh independent verification of the displayed finite/formal scope}
Fresh independent replay completed 2026-10-05T01:22:59Z: all 31 frozen authored sources compiled, including the archived entrant-delivered reconstructed variants, under exact Lean 4.33.1 and Mathlib 0df444a360eaa60ab8c11dca51a86af692955474. Both actual requested axiom prints, canonicalCubeCount\_eventually\_exp and canonicalCubeCount\_eventual\_threshold, are present and contain only propext, Classical.choice and Quot.sound, with no admitted, native-evaluation or unknown axioms. This verifies the displayed all-target log(5/3) endpoint; the sharper ordinary coefficient .46950308823856934 remains a separately explained source result. The replay proves the delivered reconstructed proof, without claiming recovery of the first lost source bytes. Receipt: Verification/sakana-erdos829-log-five-thirds-fresh-build.json.

\subsection{Verification and reproducibility}
Fresh execution: sakana-erdos829-log-five-thirds: PASS; 31/31 custom modules compiled successfully; receipt Verification/sakana-erdos829-log-five-thirds-fresh-build.json. Compiler pins, source hashes and endpoint axiom outputs are in Verification. Historical receipts and finite checks do not replace fresh replay.

\begin{samepage}
\subsection{Erdos829FiveThirds.canonicalCubeCount\_eventually\_exp}
\begin{quote}\small\ttfamily theorem canonicalCubeCount\_eventually\_exp \{\ensuremath{\varepsilon} : \ensuremath{\mathbb{R}}\} (h\ensuremath{\varepsilon} : 0<\ensuremath{\varepsilon}) :\newline     \ensuremath{\forall}\ensuremath{^{f}} n : \ensuremath{\mathbb{N}} in atTop, (canonicalCubeCount n : \ensuremath{\mathbb{R}}) \ensuremath{\leq}\newline       Real.exp ((Real.log (5/3 : \ensuremath{\mathbb{R}})+\ensuremath{\varepsilon})*Real.log (n : \ensuremath{\mathbb{R}})/Real.log (Real.log (n : \ensuremath{\mathbb{R}})))\end{quote}
{\small Source: proofs/erdos829-log-five-thirds/FiveThirdsCanonical.lean:8.\par}

{\small Source SHA-256: 51541747f8104588463d8b7b051d69692876cee605c68a2f50df9da132e6a006\par}

\end{samepage}
\begin{samepage}
\subsection{Erdos829FiveThirds.canonicalCubeCount\_eventual\_threshold}
\begin{quote}\small\ttfamily theorem canonicalCubeCount\_eventual\_threshold \{\ensuremath{\varepsilon} : \ensuremath{\mathbb{R}}\} (h\ensuremath{\varepsilon} : 0<\ensuremath{\varepsilon}) :\newline     \ensuremath{\exists} N : \ensuremath{\mathbb{N}}, \ensuremath{\forall} n : \ensuremath{\mathbb{N}}, N\ensuremath{\leq}n \ensuremath{\to} (canonicalCubeCount n : \ensuremath{\mathbb{R}}) \ensuremath{\leq}\newline       Real.exp ((Real.log (5/3 : \ensuremath{\mathbb{R}})+\ensuremath{\varepsilon})*Real.log (n : \ensuremath{\mathbb{R}})/Real.log (Real.log (n : \ensuremath{\mathbb{R}})))\end{quote}
{\small Source: proofs/erdos829-log-five-thirds/FiveThirdsCanonical.lean:14.\par}

{\small Source SHA-256: 51541747f8104588463d8b7b051d69692876cee605c68a2f50df9da132e6a006\par}

{\small\textbf{Sources and attribution:} \href{https://www.erdosproblems.com/829}{1. www.erdosproblems.com / 829}; \href{https://github.com/google-deepmind/formal-conjectures/blob/main/FormalConjectures/ErdosProblems/829.lean}{2. google-deepmind/formal-conjectures / 829.lean}; \href{https://github.com/google-deepmind/formal-conjectures/blob/main/FormalConjecturesForMathlib/Combinatorics/Additive/Convolution.lean}{3. google-deepmind/formal-conjectures / Convolution.lean}; \href{https://arxiv.org/html/2112.03617}{4. arXiv source: 2112.03617}; \href{https://arxiv.org/html/2207.12487v2}{5. arXiv source: 2207.12487v2}; \href{https://github.com/alejandrozu/ulam-opdp-difficulty-atlas}{6. alejandrozu/ulam-opdp-difficulty-atlas}; \href{https://github.com/alejandrozu/openmath-2026-judging/blob/d0ed487f487fa7ec814ff705ab55249983fe8707/OpenMath-Judging/review/entries/E02-erdos829-log-five-thirds.txt}{7. alejandrozu/openmath-2026-judging / E02-erdos829-log-five-thirds.txt}; \href{https://github.com/alejandrozu/openmath-2026-judging}{Central source and adjudication archive}.\par}

\end{samepage}
\clearpage
\section{Finite order obstructions for edge-resilient four-critical graphs}
{\small Provisional mathematical authors: Rachel Teo, Wenyi Wang, Hamdi Abderrahmene\par}\medskip
{\small Rachel Teo  -  Sakana AI.\par}

{\small Wenyi Wang  -  Sakana AI.\par}

{\small Hamdi Abderrahmene  -  Sakana AI.\par}

{\small Affiliations follow the recorded author declarations or public institutional/profile sources. Preferred author names, author order and publication affiliations await author review. This editorial compilation does not establish anthology coauthorship.\par}

\subsection{Abstract}
A finite vertex-critical four-chromatic graph that stays four-chromatic after deletion of any r or fewer edges satisfies 27r+19\ensuremath{\leq}4n, together with stronger quadratic and integral independent-set filters. These refine explicit finite lower bounds, without improving the prior eventual asymptotic restriction.

\subsection{Graph hypotheses and inequalities}
Let G be a finite simple graph on n vertices, with chromatic number four. Assume that deletion of any vertex lowers that number, and that deletion of any set of at most r edges does not lower it, where r\ensuremath{\geq}1. The submitted theorem gives 27r+19\ensuremath{\leq}4n.

Two further necessary conditions are 27(r+1)(n\ensuremath{-}1)\ensuremath{\leq}(2n+1)\ensuremath{^2} and 9b(r+1)\ensuremath{\leq}(2b+1)(n\ensuremath{-}b), where b=floor((n+1)/3). The integral filter can exclude an order allowed by the linear inequality: at r=4 it requires n\ensuremath{\geq}33, whereas the linear inequality alone gives n\ensuremath{\geq}32. These are one family of obstructions, not three separate original results.

\subsection{Criticality supplies a finite code}
Deleting a vertex of a four-critical graph produces a three-colouring of the remainder. The proof extracts a large independent colour class and encodes neighbourhood information on that class. Edge resilience forces separation between the relevant codewords: too few differences would produce a critical edge set of size at most r.

The finite-code theorem compares a lower bound for total pairwise distance with an upper bound obtained by coordinate counts. This is a Plotkin/Cauchy-Schwarz argument. Optimizing the independent-set parameter and retaining its integer size gives the quadratic and exact filters; a final numerical derivation yields the clean linear inequality.

The formal graph module supplies every hypothesis of the numeric code lemma from the actual canonical criticality definitions. The final theorem is not a numeric inequality assuming an unproved graph-to-code model. The target audit checks deletion and finite-cardinality conventions against the registered problem.

\subsection{An ordinary derivation of the code and order bounds}
Put s=r+1\ensuremath{\geq}2. For any vertex v, choose a proper three-colouring of G\ensuremath{-}v. Each of its three colour classes contains at least s neighbours of v: otherwise delete the fewer than s edges from v into that class and give v its colour, producing a three-colouring after too few deletions. Thus every vertex has degree at least 3s.

Choose a maximum independent set A of size a and put H=V(G)\ensuremath{\setminus}A, |H|=h=n\ensuremath{-}a. A three-colouring after deleting one vertex shows n\ensuremath{-}1\ensuremath{\leq}3a. For each v\ensuremath{\in}A let C\_v\ensuremath{\subseteq}H be its nonneighbours in H. Independence of A and the degree bound imply |C\_v|\ensuremath{\leq}h\ensuremath{-}3s. If u,v\ensuremath{\in}A are distinct, take a three-colouring of G\ensuremath{-}u and look at neighbours x of u having the colour of v. There are at least s of them. Such an x is outside A and cannot be adjacent to v, so it belongs to C\_v\ensuremath{\setminus}C\_u. Therefore |C\_v\ensuremath{\setminus}C\_u|\ensuremath{\geq}s in both directed orders.

Here is the complete numerical lemma. Let a subsets C\_v of an h-element ground set satisfy those weight and directed separation bounds. Write t\_x=\#\{v:x\ensuremath{\in}C\_v\}, T=\ensuremath{\sum}t\_x and Q=\ensuremath{\sum}t\_x\ensuremath{^2}. Double counting directed differences gives a(a\ensuremath{-}1)s\ensuremath{\leq}\ensuremath{\sum}\_(u,v)|C\_v\ensuremath{\setminus}C\_u|=aT\ensuremath{-}Q. Also T\ensuremath{\leq}a(h\ensuremath{-}3s) and T\ensuremath{^2}\ensuremath{\leq}hQ by Cauchy-Schwarz.

If h\ensuremath{\leq}6s, then T\ensuremath{\leq}a(h\ensuremath{-}3s)\ensuremath{\leq}ah/2. The function aT\ensuremath{-}T\ensuremath{^2}/h is increasing for 0\ensuremath{\leq}T\ensuremath{\leq}ah/2, so a(a\ensuremath{-}1)s\ensuremath{\leq}a\ensuremath{^2}(h\ensuremath{-}3s)\ensuremath{-}a\ensuremath{^2}(h\ensuremath{-}3s)\ensuremath{^2}/h=3a\ensuremath{^2}s(h\ensuremath{-}3s)/h. Rearrangement gives 9as\ensuremath{\leq}(2a+1)h. If h>6s, that inequality is immediate from(2a+1)h>6s(2a+1)>9as. This proves the code bound in all cases. The separation of two codewords also gives h\ensuremath{\geq}4s, because one C\_v has at least s members while its weight is at most h\ensuremath{-}3s.

The degree bound gives n\ensuremath{\geq}3s+1\ensuremath{\geq}7; hence a\ensuremath{\geq}ceil((n\ensuremath{-}1)/3)\ensuremath{\geq}2, so two codewords do exist. Combining h\ensuremath{\geq}4s with n\ensuremath{-}1\ensuremath{\leq}3a and n=a+h yields 2n\ensuremath{\geq}12s\ensuremath{-}1 and therefore n\ensuremath{\geq}6s. Now divide the code bound by a>0. The function(2a+1)(n\ensuremath{-}a)/a=2n\ensuremath{-}2a+n/a\ensuremath{-}1 is decreasing in positive a. Substituting the smaller integer b=floor((n+1)/3) gives 9bs\ensuremath{\leq}(2b+1)(n\ensuremath{-}b). Substituting the real lower value(n\ensuremath{-}1)/3 instead gives 27s(n\ensuremath{-}1)\ensuremath{\leq}(2n+1)\ensuremath{^2}.

Finally n\ensuremath{\geq}6s\ensuremath{\geq}12. If27s\ensuremath{\geq}4n+9, the quadratic inequality would force(4n+9)(n\ensuremath{-}1)\ensuremath{\leq}(2n+1)\ensuremath{^2}, or n\ensuremath{\leq}10, a contradiction. Therefore 27s\ensuremath{\leq}4n+8. Replacing s by r+1 yields 27r+19\ensuremath{\leq}4n, as claimed. This derives all three displayed inequalities from the graph hypotheses, without an assumed coding model.

\subsection{What the finite improvement accomplishes}
The located prior explicit bound is n\ensuremath{\geq}5r+6. After integer rounding, the new linear condition ties at r=1 and improves the finite lower requirement for larger r. However, Skottova and Steiner also prove an eventual restriction r=O(n/(log n)\textasciicircum{}c), which is stronger asymptotically. The new theorem does not supersede that result or construct any graph meeting its obstruction.

In particular it does not solve existence for r\ensuremath{\geq}2, cover k>4, or address infinite graphs. Known r=1 constructions are prior evidence, not a new contribution of this packet.

\subsection{Formal evidence and paper form}
The recorded independent isolated Lean 4.34.1 replay compiled the three proof modules and a separate target audit from unchanged frozen sources, with an initially empty own-output directory. Third-party caches were reused and all final endpoints report standard axioms. Fresh replay outcomes belong in the common register.

A short extremal-graph paper should state the three inequalities, explain the criticality-to-code separation, and give small-r comparisons. The precise theorem is more useful than an inflated conjecture-solution headline. Credit the source definitions, prior critical-edge-set paper and classical finite coding inequalities.

\subsection{Verification and reproducibility}
Fresh execution: sakana-erdos944-order-bound: PASS; 3/3 custom modules compiled successfully; receipt Verification/sakana-erdos944-order-bound-fresh-build.json. Compiler pins, source hashes and endpoint axiom outputs are in Verification. Historical receipts and finite checks do not replace fresh replay.

\begin{samepage}
\subsection{Erdos944Bounds.erdos944\_order\_linear}
\begin{quote}\small\ttfamily theorem erdos944\_order\_linear \{r : \ensuremath{\mathbb{N}}\} (hr : 1 \ensuremath{\leq} r)\newline     (hG : G.IsCritical 4)\newline     (hE : \ensuremath{\forall} edges : Set (Sym2 V), G.IsCriticalEdges edges \ensuremath{\to} r < edges.ncard) :\newline     27*r + 19 \ensuremath{\leq} 4 * Fintype.card V\end{quote}
{\small Source: proofs/erdos944-order-bound/OpHack/Erdos944Bounds.lean:42.\par}

{\small Source SHA-256: 8f9e544d0e9f465f8d92d0c0b878ab2bd9b542c0c969eaa4e2d7c908acec27f4\par}

\end{samepage}
\begin{samepage}
\subsection{Erdos944Bounds.erdos944\_order\_quadratic}
\begin{quote}\small\ttfamily theorem erdos944\_order\_quadratic \{r : \ensuremath{\mathbb{N}}\} (hr : 1 \ensuremath{\leq} r)\newline     (hG : G.IsCritical 4)\newline     (hE : \ensuremath{\forall} edges : Set (Sym2 V), G.IsCriticalEdges edges \ensuremath{\to} r < edges.ncard) :\newline     27 * (r+1) * (Fintype.card V - 1) \ensuremath{\leq} (2 * Fintype.card V + 1)\textasciicircum{}2\end{quote}
{\small Source: proofs/erdos944-order-bound/OpHack/Erdos944Bounds.lean:28.\par}

{\small Source SHA-256: 8f9e544d0e9f465f8d92d0c0b878ab2bd9b542c0c969eaa4e2d7c908acec27f4\par}

\end{samepage}
\begin{samepage}
\subsection{Erdos944Bounds.erdos944\_order\_exact}
\begin{quote}\small\ttfamily theorem erdos944\_order\_exact \{r : \ensuremath{\mathbb{N}}\} (hr : 1 \ensuremath{\leq} r)\newline     (hG : G.IsCritical 4)\newline     (hE : \ensuremath{\forall} edges : Set (Sym2 V), G.IsCriticalEdges edges \ensuremath{\to} r < edges.ncard) :\newline     9 * ((Fintype.card V + 1) / 3) * (r+1) \ensuremath{\leq}\newline       (2 * ((Fintype.card V + 1) / 3) + 1) *\newline         (Fintype.card V - (Fintype.card V + 1) / 3)\end{quote}
{\small Source: proofs/erdos944-order-bound/OpHack/Erdos944Bounds.lean:54.\par}

{\small Source SHA-256: 8f9e544d0e9f465f8d92d0c0b878ab2bd9b542c0c969eaa4e2d7c908acec27f4\par}

{\small\textbf{Sources and attribution:} \href{https://www.erdosproblems.com/944}{1. www.erdosproblems.com / 944}; \href{https://arxiv.org/html/2508.08703}{2. arXiv source: 2508.08703}; \href{https://github.com/google-deepmind/formal-conjectures/blob/main/FormalConjectures/ErdosProblems/944.lean}{3. google-deepmind/formal-conjectures / 944.lean}; \href{https://github.com/KitaKen1/erdos-944-dirac-k4-lean}{4. KitaKen1/erdos-944-dirac-k4-lean}; \href{https://arxiv.org/abs/2606.18462}{5. arXiv source: 2606.18462}; \href{https://github.com/alejandrozu/ulam-opdp-difficulty-atlas}{6. alejandrozu/ulam-opdp-difficulty-atlas}; \href{https://github.com/alejandrozu/openmath-2026-judging/blob/d0ed487f487fa7ec814ff705ab55249983fe8707/OpenMath-Judging/review/entries/E02-erdos944-order-bound.txt}{7. alejandrozu/openmath-2026-judging / E02-erdos944-order-bound.txt}; \href{https://github.com/alejandrozu/openmath-2026-judging}{Central source and adjudication archive}.\par}

\end{samepage}
\clearpage
\section{A sharp collapsed-fibre likelihood classification for a four-visible restricted Boltzmann machine}
{\small Provisional mathematical authors: Rachel Teo, Wenyi Wang, Hamdi Abderrahmene\par}\medskip
{\small Rachel Teo  -  Sakana AI.\par}

{\small Wenyi Wang  -  Sakana AI.\par}

{\small Hamdi Abderrahmene  -  Sakana AI.\par}

{\small Affiliations follow the recorded author declarations or public institutional/profile sources. Preferred author names, author order and publication affiliations await author review. This editorial compilation does not establish anthology coauthorship.\par}

\subsection{Abstract}
For pure four-parity data and a four-visible, one-hidden spin RBM, collapsed parameters are local likelihood maxima precisely below an explicit tanh\ensuremath{^2} threshold. At equality and beyond they are saddles in the full parameter space; every collapsed point remains globally suboptimal and non-strict.

\subsection{The model and full-space theorem}
For s\ensuremath{\in}\{\ensuremath{-}1,1\}\ensuremath{^4} and 0<|\ensuremath{\rho}|\ensuremath{\leq}1, take data distribution Q\_\ensuremath{\rho}(s)=(1+\ensuremath{\rho}s\ensuremath{_1}s\ensuremath{_2}s\ensuremath{_3}s\ensuremath{_4})/16. The one-hidden spin RBM has normalized probabilities proportional to exp(b\ensuremath{\cdot}s)cosh(c+w\ensuremath{\cdot}s), with b,w\ensuremath{\in}\ensuremath{\mathbb{R}}\ensuremath{^4} and c\ensuremath{\in}\ensuremath{\mathbb{R}}. Its expected log likelihood is \ensuremath{\ell}\_\ensuremath{\rho}. The collapsed fibre consists of \ensuremath{\theta}\ensuremath{_0}=(0,c\ensuremath{_0},0), which all represent the uniform visible distribution.

Every affine directional derivative at \ensuremath{\theta}\ensuremath{_0} is zero. The sharp theorem says \ensuremath{\theta}\ensuremath{_0} is a local maximum in the full nine-dimensional parameter space if and only if tanh\ensuremath{^2}(c\ensuremath{_0})<(3+2|\ensuremath{\rho}|)/(3+6|\ensuremath{\rho}|). At equality or above, every full parameter neighbourhood contains both a strictly higher and a strictly lower likelihood point. Every \ensuremath{\theta}\ensuremath{_0} also has some finite RBM with strictly better likelihood and distinct equal-likelihood points arbitrarily nearby.

\subsection{Why higher-order analysis is needed}
The collapsed parametrization is singular. Ordinary stationarity and a quadratic Hessian test miss the parity interaction, whose first nonzero signal appears at higher order. The proof expands log cosh and decomposes the visible functions into Walsh/parity components. It controls nuisance visible biases and normalizing terms by quantitative Jensen and likelihood-domination estimates.

The quartic geometry produces the threshold. An independent ordinary fourth-order derivation agrees with the displayed expression, but that calculation alone does not prove a full neighbourhood maximum. The source's StableRegion and LocalMaximum lemmas supply uniform control over arbitrary small parameter perturbations, rather than just along one symmetric ray.

Outside the threshold, a quartic escape direction yields a higher point, while another direction gives a lower point. At the equality case, the proof moves along the likelihood-flat collapsed fibre to nearby parameters on the unstable side and transports saddle behaviour back to every neighbourhood. This is the subtle boundary step; it should be explained explicitly rather than inferred from a strictly signed quartic coefficient that vanishes there.

\subsection{Non-strictness and global suboptimality}
Changing the hidden bias while keeping b=w=0 leaves the visible distribution uniform, which creates a flat direction and makes collapsed local maxima non-strict. A suitable finite coupled RBM captures part of the nonzero parity signal and improves the likelihood. Thus a collapsed local maximum can be nonglobal even though it has nearby points of equal value.

The theorem concerns parameter-space likelihood only. It does not assert a distribution-space maximum, an attracting training basin, strict trapping, EM convergence or a lower bound on optimization time. It fixes exactly four visible spins, one hidden unit and pure parity data. Ordinary explorations for larger dimensions or general four-wise-uniform data are outside the selected formal endpoint.

\subsection{Prior setting, formal closure and publication}
Collapsed stationarity, singular mixture critical fibres and familiar RBM interaction formulas are prior mathematics. The Fukumizu-Akaho-Amari comparison identifies a degenerate situation where the second-derivative splitting matrix vanishes. The proposed delta is the exact higher-order threshold, including full-neighbourhood and boundary statements.

This publication pass independently replayed the exact 15-source frozen closure under Lean 4.33.1 and Mathlib 0df444a360eaa60ab8c11dca51a86af692955474, finishing on 5 October 2026 at 01:36:52 UTC. All sources and all three selected axiom prints passed. sharp\_parity\_classification, collapsed\_local\_max\_iff and nonstable\_parameter\_saddle each use only propext, Classical.choice and Quot.sound, with no admitted, native-evaluation or unknown axioms. This verifies the displayed full-neighbourhood classification for four visible spins, one hidden unit and pure four-parity data, including its equality boundary; broader models and training claims remain outside its scope. The dated archive's separate internal 17-module post-recovery run is preserved as historical evidence rather than being relabeled as this independent run. Receipt: Verification/sakana-opdp89-parity-sharp-fresh-build.json.

\subsection{The finite statistical model and the collapsed fibre}
Use visible spins s=(s\ensuremath{_1},s\ensuremath{_2},s\ensuremath{_3},s\ensuremath{_4})\ensuremath{\in}\{\ensuremath{-}1,1\}\ensuremath{^4} and one hidden spin. The data distribution is Q\_\ensuremath{\rho}(s)=[1+\ensuremath{\rho}s\ensuremath{_1}s\ensuremath{_2}s\ensuremath{_3}s\ensuremath{_4}]/16, where 0<|\ensuremath{\rho}|\ensuremath{\leq}1. Let U denote uniform averaging over the sixteen states. For visible bias b\ensuremath{\in}R\ensuremath{^4}, hidden bias c\ensuremath{\in}R and weights w\ensuremath{\in}R\ensuremath{^4}, the visible model has potential f(s)=b\ensuremath{\cdot}s+log cosh(c+w\ensuremath{\cdot}s) and probability P(s)=exp(f(s))/\ensuremath{\Sigma}\_x exp(f(x)). Its population log likelihood is L=\ensuremath{\Sigma}\_s Q\_\ensuremath{\rho}(s)log P(s). Additive constants in f do not change P.

At b=w=0 the model is uniform for every finite c, so the whole collapsed fibre has likelihood \ensuremath{-}log 16. The data have zero moments for all nonempty spin products of degree less than four. Differentiating the finite sixteen-state likelihood at a collapsed point therefore gives zero in every parameter direction: visible-bias and weight derivatives are linear spin features, and the hidden-bias derivative is constant. Thus the fibre is stationary. This elementary stationarity is not the distinctive result. The theorem classifies precisely when it is locally maximal in all nine parameters, including highly unequal and zero weights.

\subsection{Lemma: a quantitative likelihood penalty in Walsh coordinates}
For a finite potential f write g=f\ensuremath{-}U[f], let \ensuremath{\tau}=U[s\ensuremath{_1}s\ensuremath{_2}s\ensuremath{_3}s\ensuremath{_4}f] be its top Walsh coefficient, and let \ensuremath{\beta}\_ij=U[s\_i s\_j f] be its six pair coefficients. Orthogonality of spin products gives U[g\ensuremath{^2}]\ensuremath{\geq}\ensuremath{\Sigma}\_\{i<j\}\ensuremath{\beta}\_ij\ensuremath{^2}. Because Q\_\ensuremath{\rho} differs from U only in its four-spin moment,

L(f)+log 16=\ensuremath{\rho}\ensuremath{\tau}\ensuremath{-}log U[exp(g)].

If |g(s)|\ensuremath{\leq}\ensuremath{\omega}, then log U[exp(g)]\ensuremath{\geq}e\textasciicircum{}(\ensuremath{-}2\ensuremath{\omega})U[g\ensuremath{^2}]/2. One way to prove this uniform bound is to set \ensuremath{\psi}(t)=log U[exp(tg)] for 0\ensuremath{\leq}t\ensuremath{\leq}1. We have \ensuremath{\psi}(0)=\ensuremath{\psi}\ensuremath{^{\prime}}(0)=0 and \ensuremath{\psi}\ensuremath{^{\prime\prime}}(t)=Var\_t(g), the variance under the exponential tilt. Every tilted state mass is at least e\textasciicircum{}(\ensuremath{-}2\ensuremath{\omega}) times its uniform mass. Using Var\_t(g)=inf\_a E\_t[(g\ensuremath{-}a)\ensuremath{^2}] therefore gives \ensuremath{\psi}\ensuremath{^{\prime\prime}}(t)\ensuremath{\geq}e\textasciicircum{}(\ensuremath{-}2\ensuremath{\omega})Var\_U(g). Integrating twice yields the claimed factor one half. This finite argument requires no asymptotic Taylor remainder.

Consequently L(f)+log 16\ensuremath{\leq}\ensuremath{\rho}\ensuremath{\tau}\ensuremath{-}e\textasciicircum{}(\ensuremath{-}2\ensuremath{\omega})\ensuremath{\Sigma}\_\{i<j\}\ensuremath{\beta}\_ij\ensuremath{^2}/2. The visible linear term b\ensuremath{\cdot}s contributes to neither \ensuremath{\tau} nor \ensuremath{\beta}\_ij; it remains covered through the nonnegative total normalization penalty. This observation is what permits a bound in the complete parameter neighbourhood.

\subsection{Lemma: finite differences dominate parity gain}
Put h(x)=log cosh x. Suppose on [c\ensuremath{-}\ensuremath{\Sigma}\_i|w\_i|,c+\ensuremath{\Sigma}\_i|w\_i|] we have h\ensuremath{^{\prime\prime}}\ensuremath{\geq}m>0 and |h\ensuremath{^{(4)}}|\ensuremath{\leq}M. Repeated symmetric finite differences give |\ensuremath{\tau}|\ensuremath{\leq}M|w\ensuremath{_1}w\ensuremath{_2}w\ensuremath{_3}w\ensuremath{_4}|. More explicitly, the four-spin Walsh coefficient of h(c+w\ensuremath{\cdot}s) is (w\ensuremath{_1}w\ensuremath{_2}w\ensuremath{_3}w\ensuremath{_4}/16) times the integral of h\ensuremath{^{(4)}}(c+\ensuremath{\Sigma}\_iw\_i t\_i) over [\ensuremath{-}1,1]\ensuremath{^4}; that cube has volume 16.

The same two-variable difference identity, averaged over the other two spins, shows that \ensuremath{\beta}\_ij has the sign of w\_iw\_j and |\ensuremath{\beta}\_ij|\ensuremath{\geq}m|w\_iw\_j|. Hence \ensuremath{\Sigma}\ensuremath{\beta}\_ij\ensuremath{^2}\ensuremath{\geq}m\ensuremath{^2}E(w), where E(w)=\ensuremath{\Sigma}\_\{i<j\}w\_i\ensuremath{^2}w\_j\ensuremath{^2}. Finally 6|w\ensuremath{_1}w\ensuremath{_2}w\ensuremath{_3}w\ensuremath{_4}|\ensuremath{\leq}E(w). For example, adding the three squares (w\ensuremath{_1}w\ensuremath{_2}\ensuremath{\pm}w\ensuremath{_3}w\ensuremath{_4})\ensuremath{^2}, (w\ensuremath{_1}w\ensuremath{_3}\ensuremath{\pm}w\ensuremath{_2}w\ensuremath{_4})\ensuremath{^2} and (w\ensuremath{_1}w\ensuremath{_4}\ensuremath{\pm}w\ensuremath{_2}w\ensuremath{_3})\ensuremath{^2} proves both signs of this inequality.

Combining these facts with the likelihood penalty gives

L(b,c,w)+log 16\ensuremath{\leq}[|\ensuremath{\rho}|M/6\ensuremath{-}e\textasciicircum{}(\ensuremath{-}2\ensuremath{\omega})m\ensuremath{^2}/2]E(w).

This estimate holds for signed, unequal and vanishing weights. If only one weight is nonzero then E(w)=0, and the inequality still gives nonincrease. A proof based solely on an equal-weight ray or a radial fourth-order expansion would leave these thin directions uncontrolled.

\subsection{Proof of local maximality below the exact threshold}
Write u=tanh\ensuremath{^2}c and q=1\ensuremath{-}u>0. The derivatives are h\ensuremath{^{\prime\prime}}(c)=q and h\ensuremath{^{(4)}}(c)=\ensuremath{-}2q(1\ensuremath{-}3u). Strict negativity of the coefficient in the preceding bound at the collapsed point is equivalent to |\ensuremath{\rho}||h\ensuremath{^{(4)}}(c)|<3[h\ensuremath{^{\prime\prime}}(c)]\ensuremath{^2}, or 2|\ensuremath{\rho}||1\ensuremath{-}3u|<3(1\ensuremath{-}u). For 0<|\ensuremath{\rho}|\ensuremath{\leq}1 and 0\ensuremath{\leq}u<1, this is equivalent to u<(3+2|\ensuremath{\rho}|)/(3+6|\ensuremath{\rho}|). On u\ensuremath{\leq}1/3 the inequality is automatically strict; on u>1/3 ordinary rearrangement gives the displayed threshold.

Assume the strict threshold at c\ensuremath{_0}. By continuity choose a small \ensuremath{\eta}>0 so that m=h\ensuremath{^{\prime\prime}}(c\ensuremath{_0})\ensuremath{-}\ensuremath{\eta}>0, M=|h\ensuremath{^{(4)}}(c\ensuremath{_0})|+\ensuremath{\eta}, and |\ensuremath{\rho}|M/6<e\textasciicircum{}(\ensuremath{-}2\ensuremath{\eta})m\ensuremath{^2}/2. Continuity of the derivatives gives an interval around c\ensuremath{_0} on which h\ensuremath{^{\prime\prime}}\ensuremath{\geq}m and |h\ensuremath{^{(4)}}|\ensuremath{\leq}M. For (b,c,w) sufficiently close to (0,c\ensuremath{_0},0), its entire hidden input interval stays inside that interval. Since the state space is finite and every centered potential is zero at the collapsed point, the same neighbourhood also has |g(s)|\ensuremath{\leq}\ensuremath{\eta} for all states. The quantitative bound then proves L(b,c,w)\ensuremath{\leq}\ensuremath{-}log 16 throughout that neighbourhood. This is a local maximum in the full nine-dimensional parameter space. No unproved uniform Taylor estimate remains in this argument.

\subsection{The converse, equality boundary and nonglobality}
To see the sharp converse, fix c and write t=tanh c, q=1\ensuremath{-}t\ensuremath{^2}. Choose the ray w=\ensuremath{\varepsilon}a and b=\ensuremath{-}t\ensuremath{\varepsilon}a. It cancels the first-order visible potential. Expanding log cosh and the finite normalization gives zero likelihood derivatives through degree three, and the fourth-order coefficient

\ensuremath{-}q\ensuremath{^2}E(a)/2+2\ensuremath{\rho}q(3t\ensuremath{^2}\ensuremath{-}1)a\ensuremath{_1}a\ensuremath{_2}a\ensuremath{_3}a\ensuremath{_4}.

Take a\ensuremath{_1}=a\ensuremath{_2}=a\ensuremath{_3}=1 and a\ensuremath{_4} with sign \ensuremath{\rho}. Then E(a)=6 and the coefficient is q[(3+6|\ensuremath{\rho}|)t\ensuremath{^2}\ensuremath{-}(3+2|\ensuremath{\rho}|)]. Above the threshold it is positive. Smoothness and Taylor's theorem, or the source's four repeated mean-value arguments on the explicit jets, give strict ascent for sufficiently small positive \ensuremath{\varepsilon}. Thus an exterior collapsed point is not a local maximum.

At equality the ray's fourth coefficient vanishes, so that calculation alone does not decide the boundary. Instead use the flat fibre: threshold equality implies c\ensuremath{_0}\ensuremath{\neq}0, and there are c values arbitrarily close to c\ensuremath{_0} with larger |c| and hence tanh\ensuremath{^2}c above the threshold. Each such collapsed point has exactly the same likelihood as the boundary point, and each has arbitrarily nearby strict-ascent points. Every neighbourhood of the boundary therefore contains a point of higher likelihood. This proves that equality belongs to the saddle region.

Every neighbourhood of any collapsed point also contains a lower-likelihood point: set w=0 and perturb one visible bias to \ensuremath{\varepsilon}\ensuremath{\neq}0. The data's one-spin mean is zero, giving L=\ensuremath{-}log 16\ensuremath{-}log cosh \ensuremath{\varepsilon}<\ensuremath{-}log 16. At and beyond the threshold, higher and lower points therefore occur in every neighbourhood, which is the stated parameter-saddle property. Below threshold the maxima are non-strict, because changing c along the collapsed fibre produces distinct points of exactly equal likelihood. Finally the threshold is less than one for |\ensuremath{\rho}|>0, so some finite hidden bias lies outside it and admits ascent above the common collapsed likelihood. Thus every collapsed point, including the stable local maxima, is globally suboptimal.

These conclusions give the complete finite-fibre classification: stationarity everywhere; local maximality if and only if the strict threshold holds; saddle behaviour at equality and beyond; non-strictness along the fibre; and nonglobality. The statement concerns the specified one-hidden-spin, four-visible-spin parity model and does not assert a classification of arbitrary RBMs or all stationary points.

\subsection{Fresh independent verification of the sharp classification}
This publication pass independently replayed the exact 15-source frozen closure under Lean 4.33.1 and Mathlib 0df444a360eaa60ab8c11dca51a86af692955474, finishing on 5 October 2026 at 01:36:52 UTC. All sources and all three selected axiom prints passed. sharp\_parity\_classification, collapsed\_local\_max\_iff and nonstable\_parameter\_saddle each use only propext, Classical.choice and Quot.sound, with no admitted, native-evaluation or unknown axioms. This verifies the displayed full-neighbourhood classification for four visible spins, one hidden unit and pure four-parity data, including its equality boundary; broader models and training claims remain outside its scope. The dated archive's separate internal 17-module post-recovery run is preserved as historical evidence rather than being relabeled as this independent run. Receipt: Verification/sakana-opdp89-parity-sharp-fresh-build.json.

\subsection{Verification and reproducibility}
Fresh execution: sakana-opdp89-parity-sharp: PASS; 15/15 custom modules compiled successfully; receipt Verification/sakana-opdp89-parity-sharp-fresh-build.json. Compiler pins, source hashes and endpoint axiom outputs are in Verification. Historical receipts and finite checks do not replace fresh replay.

\begin{samepage}
\subsection{ParityFour.sharp\_parity\_classification}
\begin{quote}\small\ttfamily theorem sharp\_parity\_classification (\ensuremath{\rho} c\ensuremath{_0} : \ensuremath{\mathbb{R}}) (h\ensuremath{\rho} : 0 < |\ensuremath{\rho}|) (h\ensuremath{\rho}1 : |\ensuremath{\rho}| \ensuremath{\leq} 1) :\newline     (\ensuremath{\forall} (b : Fin 4 \ensuremath{\to} \ensuremath{\mathbb{R}}) (d : \ensuremath{\mathbb{R}}) (w : Fin 4 \ensuremath{\to} \ensuremath{\mathbb{R}}),\newline       HasDerivAt (fun t => objective \ensuremath{\rho} (affineParameterRay c\ensuremath{_0} b d w t)) 0 0) \ensuremath{\wedge}\newline     (IsLocalMax (objective \ensuremath{\rho}) (nullParameter c\ensuremath{_0}) \ensuremath{\leftrightarrow}\newline       Real.tanh c\ensuremath{_0} \textasciicircum{} 2 < (3+2*|\ensuremath{\rho}|)/(3+6*|\ensuremath{\rho}|)) \ensuremath{\wedge}\newline     ((3+2*|\ensuremath{\rho}|)/(3+6*|\ensuremath{\rho}|) \ensuremath{\leq} Real.tanh c\ensuremath{_0} \textasciicircum{} 2 \ensuremath{\to} ParameterSaddle \ensuremath{\rho} (nullParameter c\ensuremath{_0})) \ensuremath{\wedge}\newline     (\ensuremath{\exists} \ensuremath{\theta} : Parameters, objective \ensuremath{\rho} (nullParameter c\ensuremath{_0}) < objective \ensuremath{\rho} \ensuremath{\theta}) \ensuremath{\wedge}\newline     (\ensuremath{\forall} U : Set Parameters, U \ensuremath{\in} nhds (nullParameter c\ensuremath{_0}) \ensuremath{\to}\newline       \ensuremath{\exists} \ensuremath{\theta} : Parameters, \ensuremath{\theta} \ensuremath{\in} U \ensuremath{\wedge} \ensuremath{\theta} \ensuremath{\neq} nullParameter c\ensuremath{_0} \ensuremath{\wedge}\newline         objective \ensuremath{\rho} \ensuremath{\theta} = objective \ensuremath{\rho} (nullParameter c\ensuremath{_0}))\end{quote}
{\small Source: proofs/opdp89-parity-sharp/SharpClassification.lean:40.\par}

{\small Source SHA-256: 15ef1ad86c585c87e7ac4fabdc3a58702c915e69dbf6aa80d701258bc9efc602\par}

\end{samepage}
\begin{samepage}
\subsection{ParityFour.collapsed\_local\_max\_iff}
\begin{quote}\small\ttfamily theorem collapsed\_local\_max\_iff (\ensuremath{\rho} c\ensuremath{_0} : \ensuremath{\mathbb{R}}) (h\ensuremath{\rho} : 0 < |\ensuremath{\rho}|) (h\ensuremath{\rho}1 : |\ensuremath{\rho}| \ensuremath{\leq} 1) :\newline     IsLocalMax (objective \ensuremath{\rho}) (nullParameter c\ensuremath{_0}) \ensuremath{\leftrightarrow}\newline       Real.tanh c\ensuremath{_0} \textasciicircum{} 2 < (3+2*|\ensuremath{\rho}|)/(3+6*|\ensuremath{\rho}|)\end{quote}
{\small Source: proofs/opdp89-parity-sharp/SharpClassification.lean:22.\par}

{\small Source SHA-256: 15ef1ad86c585c87e7ac4fabdc3a58702c915e69dbf6aa80d701258bc9efc602\par}

\end{samepage}
\begin{samepage}
\subsection{ParityFour.nonstable\_parameter\_saddle}
\begin{quote}\small\ttfamily theorem nonstable\_parameter\_saddle (\ensuremath{\rho} c\ensuremath{_0} : \ensuremath{\mathbb{R}}) (h\ensuremath{\rho} : 0 < |\ensuremath{\rho}|) (h\ensuremath{\rho}1 : |\ensuremath{\rho}| \ensuremath{\leq} 1)\newline     (houtside : (3+2*|\ensuremath{\rho}|)/(3+6*|\ensuremath{\rho}|) \ensuremath{\leq} Real.tanh c\ensuremath{_0} \textasciicircum{} 2) :\newline     ParameterSaddle \ensuremath{\rho} (nullParameter c\ensuremath{_0})\end{quote}
{\small Source: proofs/opdp89-parity-sharp/SharpClassification.lean:33.\par}

{\small Source SHA-256: 15ef1ad86c585c87e7ac4fabdc3a58702c915e69dbf6aa80d701258bc9efc602\par}

{\small\textbf{Sources and attribution:} \href{http://aimpl.org/boltzmann/2/}{1. aimpl.org / 2}; \href{https://raw.githubusercontent.com/alejandrozu/opdp-week-math-100/main/data/problems.json}{2. raw.githubusercontent.com / problems.json}; \href{https://proceedings.neurips.cc/paper_files/paper/2002/file/2d3acd3e240c61820625fff66a19938f-Paper.pdf}{3. proceedings.neurips.cc / 2d3acd3e240c61820625fff66a19938f-Paper.pdf}; \href{https://arxiv.org/html/2605.19178}{4. arXiv source: 2605.19178}; \href{https://github.com/alejandrozu/openmath-2026-judging/blob/d0ed487f487fa7ec814ff705ab55249983fe8707/OpenMath-Judging/review/entries/E02-opdp89-parity-local-maxima.txt}{5. alejandrozu/openmath-2026-judging / E02-opdp89-parity-local-maxima.txt}; \href{https://github.com/alejandrozu/openmath-2026-judging}{Central source and adjudication archive}.\par}

\end{samepage}
\clearpage
\section{A modulus 13 obstruction to diagonal-pullback optimality for a four-point L-pattern}
{\small Provisional mathematical authors: Rachel Teo, Wenyi Wang, Hamdi Abderrahmene\par}\medskip
{\small Rachel Teo  -  Sakana AI.\par}

{\small Wenyi Wang  -  Sakana AI.\par}

{\small Hamdi Abderrahmene  -  Sakana AI.\par}

{\small Affiliations follow the recorded author declarations or public institutional/profile sources. Preferred author names, author order and publication affiliations await author review. This editorial compilation does not establish anthology coauthorship.\par}

\subsection{Abstract}
An explicit 79-point subset of (\ensuremath{\mathbb{Z}}/13\ensuremath{\mathbb{Z}})\ensuremath{^2} avoids a prescribed four-point L-pattern and exceeds every diagonal pullback of a cyclic four-AP-free set. Exact cyclic sharpness gives pullback capacity 78. This refutes a finite auxiliary ansatz, without addressing reciprocal divergence.

There exists an L-free B\ensuremath{\subseteq}(\ensuremath{\mathbb{Z}}/13\ensuremath{\mathbb{Z}})\ensuremath{^2} with |B|=79, while every diagonal pullback S\_A from a cyclic four-AP-free A has |S\_A|\ensuremath{\leq}78.

\subsection{The finite model and ansatz}
A strong L-pattern in (\ensuremath{\mathbb{Z}}/13\ensuremath{\mathbb{Z}})\ensuremath{^2} consists of(x,y),(x,y+d),(x,y+2d),(x+d,y) with d\ensuremath{\neq}0. The retained explicit set B has 79 points and contains no such pattern. For a cyclic four-AP-free A\ensuremath{\subseteq}\ensuremath{\mathbb{Z}}/13\ensuremath{\mathbb{Z}}, the diagonal pullback S\_A=\{(x,y):x\ensuremath{-}y\ensuremath{\in}A\} is L-free and has 13|A| points.

An exhaustive exact certificate proves that every cyclic four-AP-free subset has at most six points, and\{0,1,2,4,5,7\} attains six. Consequently every pullback S\_A has at most 78 points, whereas B has 79. The diagonal-pullback construction therefore fails to be optimal in this finite auxiliary extremal problem.

\subsection{Witness and exact upper certificate}
The 79-point witness is stored literally as a finite set, and its cardinality and all forbidden-pattern tests are kernel-decided. The cyclic upper bound is a separate Boolean-tree certificate covering every subset of the 13-element cyclic group, with cardinality transport to the ordinary set statement.

The useful final step is source-faithful transport. Fugu3PullbackCore uses the original diagonal convention x\ensuremath{-}y rather than an interchangeable-looking surrogate, and PullbackObstruction proves equivalence of the original curried predicates with the certificate predicates. It then combines the 79-point witness, the upper bound six and the identity |S\_A|=13|A| in one endpoint.

This proves a strict separation from every member of the proposed ansatz, not simply from one tested seed. It does not prove that 79 is the optimum among all L-free subsets.

\subsection{Interpretation and limits}
The construction is a useful warning against a plausible reduction: one-dimensional progression avoidance need not yield extremal two-dimensional L-pattern avoidance. The gap is only one point, but the universal comparison over all diagonal pullbacks makes its meaning precise.

There is no infinite reciprocal-divergent progression-free set in this theorem, and no statement forcing long arithmetic progressions. It must not be announced as a solution of Erdős 3. The research value is a finite obstruction to a specific auxiliary strategy and a reproducible exact extremal certificate.

\subsection{Formal status and short-note structure}
The archive records the cyclic sharpness checks and a seven-module fresh source-only bridge replay under Lean 4.33.1 with standard axioms. The selected source states retained\_witness\_beats\_every\_diagonal\_pullback. The literal 79-point set and the six-point example should be included as a table or supplementary JSON so that independent reproduction is easy.

A short paper can present the model, the ansatz, the witnesses, the complete finite upper-bound method and the faithful bridge. Historical priority and the original proposed ansatz should be attributed. The full million-byte Boolean decision tree belongs in the repository rather than the printed note.

\subsection{Fresh independent verification of the displayed finite/formal scope}
Fresh independent replay verifies both finite scopes separately under exact Lean 4.33.1 and Mathlib 0df444a360eaa60ab8c11dca51a86af692955474. The six-source base completed 2026-10-05T01:26:46Z with all five selected theorem prints; the seven-source pullback extension completed 2026-10-05T01:30:52Z with all three selected prints. Every actual list uses only standard kernel axioms; B\_card needs only propext and Quot.sound, and the other endpoints additionally use Classical.choice. No admitted, native-evaluation or unrecognized axioms occur. These are the explicit finite witness, cyclic extremum and diagonal-pullback comparison theorems, not a solution of the full Erdős 3 problem. The independent finite enumerations and these fresh formal checks support the same scoped conclusions. Receipts: Verification/sakana-FOCUS-E3-fresh-build.json and Verification/sakana-FOCUS-E3-pullback-fresh-build.json.

\subsection{Exact finite theorem and the diagonal construction}
Work in the prime cyclic group C=\ensuremath{\mathbb{Z}}/13\ensuremath{\mathbb{Z}}. A subset A is cyclic four-AP-free if it contains no four points a,a+d,a+2d,a+3d with d\ensuremath{\neq}0. A subset S of C\ensuremath{^2} is strong-L-free if it contains no four points (x,y),(x,y+d),(x,y+2d),(x+d,y) with d\ensuremath{\neq}0. The prescribed orientation and equal step are part of the model. This result is not a statement about every rotated or rescaled Euclidean L-pattern.

Theorem. There is an explicit strong-L-free B\ensuremath{\subseteq}C\ensuremath{^2} with |B|=79. For every cyclic four-AP-free A\ensuremath{\subseteq}C, the diagonal pullback S\_A=\{(x,y)\ensuremath{\in}C\ensuremath{^2}:x\ensuremath{-}y\ensuremath{\in}A\} is strong-L-free and has |S\_A|=13|A|\ensuremath{\leq}78. In particular |B|>|S\_A| for every member of this ansatz. The theorem is a universal finite separation from the specified construction, rather than a comparison with a single chosen seed.

\subsection{Two elementary pullback lemmas}
Lemma 1 (avoidance transfer). Suppose a strong L were present in S\_A. Put s=x\ensuremath{-}y. The four differences x\ensuremath{-}y, x\ensuremath{-}(y+d), x\ensuremath{-}(y+2d), (x+d)\ensuremath{-}y are s,s\ensuremath{-}d,s\ensuremath{-}2d,s+d. Reordered as s\ensuremath{-}2d,s\ensuremath{-}d,s,s+d, they form a cyclic four-term arithmetic progression with nonzero step d in A. This contradicts the hypothesis. Therefore S\_A is strong-L-free.

Lemma 2 (exact fiber count). The map A\ensuremath{\times}C\ensuremath{\to}S\_A given by (a,y)\ensuremath{\mapsto}(a+y,y) is a bijection. Its inverse sends (x,y) to (x\ensuremath{-}y,y), whose first coordinate belongs to A by the definition of S\_A. Consequently |S\_A|=|A||C|=13|A|. The x\ensuremath{-}y convention is deliberately retained, since the final formal bridge uses that exact original definition.

\subsection{The explicit 79-point construction}
Specify B by the following thirteen rows, where the numbers after x are exactly the y-coordinates belonging to that row, all read modulo 13: x=0: 0,2,3,5,6,7,9; x=1: 0,1,3,4,7,12; x=2: 1,2,6,7,9; x=3: 1,4,6,8,10,11; x=4: 0,3,4,7,8,11; x=5: 0,1,3,5,7,10,11; x=6: 1,2,4,5,10,12; x=7: 4,5,7,8,10,12; x=8: 2,6,8,9,11,12; x=9: 1,2,5,6,7,9; x=10: 1,2,3,8,11,12; x=11: 0,2,3,8,10,12; x=12: 2,4,5,6,9,10. The row sizes are 7,6,5,6,6,7,6,6,6,6,6,6,6 and sum to 79. All listed coordinates are distinct within a row, so this also proves the cardinality directly.

Lemma 3 (finite L avoidance). For each x\ensuremath{\in}\{0,...,12\}, y\ensuremath{\in}\{0,...,12\}, d\ensuremath{\in}\{1,...,12\}, test membership of (x,y),(x,y+d),(x,y+2d),(x+d,y), reducing coordinates modulo 13. There are exactly 13\ensuremath{\cdot}13\ensuremath{\cdot}12=2,028 tests. No test has all four memberships true for the displayed rows. These tests are a complete proof obligation because they enumerate precisely every parameter triple in the definition of strong-L-free. The supplied R17SplitReplay partitions them by the thirteen x values, verifies each row in the kernel, then joins the universal statement.

The independent ordinary checker reads the literal witness from that source, verifies 79 distinct points, and executes the same fully exhaustive predicate from its mathematical definition. Its result is recorded with the exact source hash in final\_four\_independent\_certificate\_audit.json. The construction may have been found by search, but its validity rests on the finite exhaustive certificate; no probabilistic or numerical search assumption enters the theorem.

\subsection{Cyclic sharpness at modulus 13}
Lemma 4 (cyclic upper bound). Every cyclic four-AP-free subset of C has at most six elements. Represent a subset A by the 13-bit integer m=\ensuremath{\Sigma}\_\{a\ensuremath{\in}A\}2\textasciicircum{}a. For every start a\ensuremath{\in}\{0,...,12\} and step d\ensuremath{\in}\{1,...,12\}, define q(a,d)=\ensuremath{\Sigma}\_\{k=0\}\ensuremath{^3}2\textasciicircum{}((a+kd) mod 13). Because 13 is prime and d\ensuremath{\neq}0, its four coordinates are distinct. A contains that progression exactly when the bitwise intersection m AND q(a,d) equals q(a,d). Thus the following finite algorithm proves the lemma: enumerate all integers 0\ensuremath{\leq}m<8192; for each m with more than six set bits, verify that at least one of the 156 parameter masks is contained in m.

Reversal of an AP does not change its four-element set, and the generated parameter masks give 78 distinct four-element sets. Deduplication is optional; checking all 156 parameters has the same meaning. The independent enumeration checks all 8,192 subsets, finds 1,886 AP-free subsets, and finds maximum cardinality six. The source BooleanTree is a kernel-checkable exhaustive decision tree for the same upper assertion, with CardBridge converting Boolean membership indicators back to the ordinary finite-set cardinality. The printed algorithm supplies a small independent way to verify its finite conclusion.

The upper bound is sharp: A\ensuremath{_0}=\{0,1,2,4,5,7\} passes every cyclic AP test and has six elements. Its pullback has 78 points by Lemma 2 and is strong-L-free by Lemma 1. Combining Lemmas 1-4 with the explicit B yields the strict inequality 79>13|A| for every eligible A. This is the entire ordinary argument.

\subsection{The exact boundary of the obstruction}
The separation disproves optimality of diagonal pullbacks in this one finite strong-L model. It leaves open the maximum size of arbitrary strong-L-free subsets of C\ensuremath{^2}; a 79-point witness is a lower bound, not an upper bound. It also supplies no infinite progression-free set with divergent reciprocal sum and proves no Erdős-Turán implication. Its mathematical value is to identify a concrete failure of an auxiliary extremal ansatz that could otherwise guide a search in the wrong direction.

The archived reconstructed row partition and cyclic certificate are preserved as delivered source artifacts. The supplemental PullbackObstruction explicitly identifies the original curried avoidance predicates with the conjunction predicates used in the finite certificates and then proves retained\_witness\_beats\_every\_diagonal\_pullback. The source chronology and the original witness chronology belong in the evidence record; the latest complete finite bridge should not be silently described as the earliest submission version.

\subsection{Fresh independent verification of the displayed finite/formal scope}
Fresh independent replay verifies both finite scopes separately under exact Lean 4.33.1 and Mathlib 0df444a360eaa60ab8c11dca51a86af692955474. The six-source base completed 2026-10-05T01:26:46Z with all five selected theorem prints; the seven-source pullback extension completed 2026-10-05T01:30:52Z with all three selected prints. Every actual list uses only standard kernel axioms; B\_card needs only propext and Quot.sound, and the other endpoints additionally use Classical.choice. No admitted, native-evaluation or unrecognized axioms occur. These are the explicit finite witness, cyclic extremum and diagonal-pullback comparison theorems, not a solution of the full Erdős 3 problem. The independent finite enumerations and these fresh formal checks support the same scoped conclusions. Receipts: Verification/sakana-FOCUS-E3-fresh-build.json and Verification/sakana-FOCUS-E3-pullback-fresh-build.json.

\subsection{Verification and reproducibility}
Fresh execution: sakana-FOCUS-E3: PASS; 6/6 custom modules compiled successfully; receipt Verification/sakana-FOCUS-E3-fresh-build.json. sakana-FOCUS-E3-pullback: PASS; 7/7 custom modules compiled successfully; receipt Verification/sakana-FOCUS-E3-pullback-fresh-build.json. Compiler pins, source hashes and endpoint axiom outputs are in Verification. Historical receipts and finite checks do not replace fresh replay.

\begin{samepage}
\subsection{Erdos3PullbackBridge.card\_diagonal\_pullback\_le\_78}
\begin{quote}\small\ttfamily theorem card\_diagonal\_pullback\_le\_78 (A : Finset (ZMod 13))\newline     (hA : Fugu3Lshape.APFree4 A) : (Fugu3Lshape.SA A).card \ensuremath{\leq} 78\end{quote}
{\small Source: Focus\_Evidence/evidence/FOCUS-E3/pullback-bridge/source/PullbackObstruction.lean:31.\par}

{\small Source SHA-256: 042f95ddc1c1630d43710d987dbca399cc4bc1a736004290a2c5425b41f029c8\par}

\end{samepage}
\begin{samepage}
\subsection{Erdos3PullbackBridge.retained\_witness\_beats\_every\_diagonal\_pullback}
\begin{quote}\small\ttfamily theorem retained\_witness\_beats\_every\_diagonal\_pullback :\newline     Erdos3R17.B.card = 79 \ensuremath{\wedge} Fugu3Lshape.LFree Erdos3R17.B \ensuremath{\wedge}\newline     \ensuremath{\forall} A : Finset (ZMod 13), Fugu3Lshape.APFree4 A \ensuremath{\to}\newline       Fugu3Lshape.LFree (Fugu3Lshape.SA A) \ensuremath{\wedge}\newline       (Fugu3Lshape.SA A).card = 13 * A.card \ensuremath{\wedge}\newline       (Fugu3Lshape.SA A).card \ensuremath{\leq} 78 \ensuremath{\wedge}\newline       (Fugu3Lshape.SA A).card < Erdos3R17.B.card\end{quote}
{\small Source: Focus\_Evidence/evidence/FOCUS-E3/pullback-bridge/source/PullbackObstruction.lean:38.\par}

{\small Source SHA-256: 042f95ddc1c1630d43710d987dbca399cc4bc1a736004290a2c5425b41f029c8\par}

{\small\textbf{Sources and attribution:} \href{https://github.com/alejandrozu/ulam-opdp-difficulty-atlas}{1. alejandrozu/ulam-opdp-difficulty-atlas}; \href{https://github.com/alejandrozu/openmath-2026-judging/blob/d0ed487f487fa7ec814ff705ab55249983fe8707/OpenMath-Judging/review/entries/E02-FOCUS-E3.txt}{2. alejandrozu/openmath-2026-judging / E02-FOCUS-E3.txt}; \href{https://github.com/alejandrozu/openmath-2026-judging}{Central source and adjudication archive}.\par}

\end{samepage}
\clearpage
\section{One-factor rigidity at three sparse 3\ensuremath{\times}3 matrix-multiplication seeds}
{\small Provisional mathematical authors: Rachel Teo, Wenyi Wang, Hamdi Abderrahmene\par}\medskip
{\small Rachel Teo  -  Sakana AI.\par}

{\small Wenyi Wang  -  Sakana AI.\par}

{\small Hamdi Abderrahmene  -  Sakana AI.\par}

{\small Affiliations follow the recorded author declarations or public institutional/profile sources. Preferred author names, author order and publication affiliations await author review. This editorial compilation does not establish anthology coauthorship.\par}

\subsection{Abstract}
At each of three stored 23-product, 139-support HKS seeds, all three paired-factor coefficient matrices have column rank 23. Consequently fixing any two factor lists uniquely determines the third, excluding a support reduction through a change to only one list at those seeds.

For each archived HKS 139 seed, rank\_Q(M\_UV)=rank\_Q(M\_UW)=rank\_Q(M\_VW)=23; fixing two factor lists uniquely determines the third.

\subsection{Tensor factors and the precise obstruction}
A 23-product bilinear scheme writes the 3\ensuremath{\times}3 multiplication tensor as a sum of 23 simple tensors u\_i\ensuremath{\otimes}v\_i\ensuremath{\otimes}w\_i. Fixing u and v turns the coefficient equations into a linear system for w; analogous statements hold for the other pairs. The paired-factor matrix has 81 rows and 23 columns, with column i formed by the pairwise products of the chosen factor coefficients.

At each of three literal known HKS 139 seeds, the UV, UW and VW matrices have column rank 23 over the rationals. Each corresponding linear system is injective. Since the stored third list is one solution, it is the unique solution. A change to only that list cannot lower support while preserving the tensor identity.

\subsection{Certificates and quantifiers}
The certificate supplies exact rank/injectivity evidence for all nine seed/pair combinations, followed by the uniqueness consequence. Independent rational elimination reproduced rank 23 for each of those matrices. The seeds themselves have total support 139 and are background constructions.

The theorem fixes the selected two factor lists exactly. It does not prove local rigidity when multiple lists vary, global inequivalence of all schemes, absence of a continuous family, a minimum support 139, or impossibility of support 138 or 137. Chandragupt's 138 scheme is compatible with this result because it changes the construction jointly.

\subsection{Research value and formal status}
The underlying linear algebra is classical. The proposed new contribution is a specific exact method obstruction at useful sparse seeds, showing why a one-leg search cannot improve them. Its scientific weight is consequently narrower than a new multiplication scheme; this distinction should be reflected in both placement and paper wording.

On 5 October 2026, all three frozen seed projects completed fresh Lean 4.33.1 replay at the exact Mathlib revision 0df444a360eaa60ab8c11dca51a86af692955474. Each project compiled all 71 source modules and printed all six selected declarations: the 18 requested prints used only propext, Classical.choice and Quot.sound, with no admitted, native-evaluation or unknown axioms. The audited statements establish rational injectivity of each literal 81-by-23 paired-factor matrix and equality of arbitrary remaining-factor matrices X and Y whenever B X = B Y. Full column rank 23 and the one-factor obstruction follow by ordinary linear algebra and the tensor coefficient correspondence explained here. The scope fixes the other two factor lists exactly; it does not establish multi-factor local rigidity, a minimum support of 139, or the impossibility of support 138 or 137. A standalone note should display the seed representation, the exact left-inverse certificate and the remaining-factor bridge. The three dated replay receipts and exact source statements are preserved in the common verification register.

\subsection{Attribution and publication scope}
Credit HKS for the seed decompositions and identify their literal versions. The note should explain why the tensor coefficient equations give the paired linear system, then prove the uniqueness corollary. Repository certificates can replace lengthy printed rational elimination tables while keeping the finite theorem reproducible.

If subject review finds the same rank check already published or too routine as a mathematical novelty claim, it belongs in the nonnovel/useful-certificate paper. That possibility does not affect the correctness of the exact rank evidence.

\subsection{Paired-factor matrices and the exact uniqueness claim}
Let U,V,W each be a list of 23 vectors in \ensuremath{\mathbb{Q}}\ensuremath{^9}. A 23-product 3\ensuremath{\times}3 matrix multiplication construction is represented by the simple-tensor sum \ensuremath{\Sigma}\_\{t=1\}\textasciicircum{}\{23\}u\_t\ensuremath{\otimes}v\_t\ensuremath{\otimes}w\_t. The sparse HKS 139 seeds studied here are previously known constructions; each has a total of 139 nonzero coordinates across its three lists. This chapter proves an obstruction to changing a single list at those fixed seeds. It does not propose those known seeds as new multiplication schemes.

For the pair U,V define B\_UV\ensuremath{\in}\ensuremath{\mathbb{Q}}\textasciicircum{}(81\ensuremath{\times}23) by (B\_UV)\_\{(i,j),t\}=(u\_t)\_i(v\_t)\_j, for 0\ensuremath{\leq}i,j<9. If X is the 23\ensuremath{\times}9 matrix whose tth row is w\_t, then flattening the first two tensor slots identifies the tensor with B\_UV X. Thus replacing W by another list Y, while retaining U and V exactly, preserves the tensor if and only if B\_UV X=B\_UV Y. The other two paired-factor matrices B\_UW and B\_VW arise by permuting the slots.

Theorem. At each of the three stored HKS 139 seeds, B\_UV, B\_UW and B\_VW all have column rank 23 over \ensuremath{\mathbb{Q}}. For each of those nine cases, B\_pair X=B\_pair Y implies X=Y for matrices X,Y with 23 rows and any common finite column set. In particular, a valid third factor list is uniquely determined by the two fixed stored factor lists. It cannot be changed to reduce support while preserving those exact two lists.

\subsection{Exact rank certificates and a direct proof of injectivity}
Lemma 1 (selected-row left inverse). Let B be an m\ensuremath{\times}r matrix over a field. Choose r row indices and let M be the resulting r\ensuremath{\times}r matrix. If an r\ensuremath{\times}r matrix L satisfies LM=I\_r, then B has trivial kernel. Indeed Bz=0 implies Mz=0 by restriction to the selected rows, and multiplication by L gives z=LMz=0. Since B has r columns, its column rank is r. This proof requires no floating-point rank tolerance and no assumption about the condition number.

The submitted certificates are even simpler than a generic rational inverse. Every selected minor M and inverse L has integer entries and satisfies the literal integer identity LM=I\ensuremath{_2}\ensuremath{_3}. For each of the three seeds, there are certificates for UV, UW and VW. Their inverses are verified row by row in 23 separate kernel modules and assembled into left\_inverse\_uv, left\_inverse\_uw and left\_inverse\_vw. Casting the identity \ensuremath{\mathbb{Z}}\ensuremath{\to}\ensuremath{\mathbb{Q}} gives the rational left inverse used in Lemma 1.

For seed zero, with zero-based flattened row index 9i+j, the UV pivot rows are 0,1,2,8,12,13,14,15,19,27,28,29,32,39,40,49,54,55,56,63,67,72,73. The UW pivot rows are 0,2,3,4,6,9,10,12,15,18,21,24,28,34,35,37,40,48,59,64,65,72,75. The VW pivot rows are 0,1,2,8,12,13,14,16,17,18,27,28,29,37,39,40,46,51,54,55,56,57,75. The two additional seeds have their own literal pivot lists and inverse matrices; they are separate certificates, rather than an extrapolation from seed zero.

The independent checker reads the literal U,V,W arrays and L matrices in each KernelData file, reconstructs all 81\ensuremath{\times}23 paired matrices from the product formula, selects the certified 23 rows, and multiplies L by M using exact integers. All nine products equal I\ensuremath{_2}\ensuremath{_3}. It also checks that every pivot list contains 23 distinct indices and that every input seed has total support 139. This provides an independent arithmetic audit of precisely the rank premises in the formal endpoints.

\subsection{Uniqueness for the remaining factor and its consequences}
Lemma 2 (columnwise uniqueness). Suppose B has trivial kernel and BX=BY. For every column j of X and Y, B(X\_j\ensuremath{-}Y\_j)=0. Lemma 1 gives X\_j\ensuremath{-}Y\_j=0. Applying this to every column proves X=Y. The original source states this lemma with an arbitrary column type p, so the assertion concerns the full third factor array and is not confined to one selected output coordinate.

Apply Lemma 2 to B\_UV at a stored seed. The old W is already a solution of the tensor equations. Any replacement W\ensuremath{^{\prime}} producing the same tensor with those U,V must equal W coordinate by coordinate. Consequently its support is exactly the old support. The same reasoning applies to U when V,W are fixed, and to V when U,W are fixed. This proves the announced one-factor search obstruction in all three positions at all three seeds.

Termwise rescalings or permutations that also change a fixed list fall outside these quantifiers. Likewise an improvement that moves two or three lists together is not forbidden. A jointly different support-138 scheme is consistent with all nine certificates. The result proves neither a global lower bound 139 nor a neighborhood theorem for arbitrary simultaneous perturbations. Full column rank of a paired matrix is also different from injectivity of the derivative of the complete three-factor polynomial map; no such stronger Jacobian assertion is established here.

\subsection{A small reproducible note about an exact search obstruction}
The ordinary proof consists of the paired-factor flattening, Lemma 1, nine finite integer identities, and Lemma 2. It is complete at that scope. The appropriate supplement contains the three input schemes, pivot lists and inverse matrices, so a reader can reproduce all nine products without installing Lean. The corresponding kernel certificates then give an additional independently specified proof path under the pinned environment.

The scientific contribution should be described as an exact obstruction to a one-factor optimization strategy at named useful seeds. The linear algebra is standard, and the constructions it analyzes are known. The publication question is whether the exact seed obstruction and its documented search consequence were previously recorded. A concise note can be worthwhile even if its contest contribution is small, but its title and theorem must keep the fixed-seed/fixed-two-list hypotheses visible.

\subsection{Verification and reproducibility}
Fresh execution: sakana-FOCUS-MATRIX-seed0: PASS; 71/71 custom modules compiled successfully; receipt Verification/sakana-FOCUS-MATRIX-seed0-fresh-build.json. sakana-FOCUS-MATRIX-seed1: PASS; 71/71 custom modules compiled successfully; receipt Verification/sakana-FOCUS-MATRIX-seed1-fresh-build.json. sakana-FOCUS-MATRIX-seed2: PASS; 71/71 custom modules compiled successfully; receipt Verification/sakana-FOCUS-MATRIX-seed2-fresh-build.json. Compiler pins, source hashes and endpoint axiom outputs are in Verification. Historical receipts and finite checks do not replace fresh replay.

\begin{samepage}
\subsection{RecoveredMMTKernel.remaining\_factor\_unique\_uv}
\begin{quote}\small\ttfamily theorem remaining\_factor\_unique\_uv \{p : Type*\}\newline     (X Y : Matrix (Fin 23) p \ensuremath{\mathbb{Q}})\newline     (h : (B\_uv.map (Int.castRingHom \ensuremath{\mathbb{Q}})) * X = (B\_uv.map (Int.castRingHom \ensuremath{\mathbb{Q}})) * Y) : X = Y\end{quote}
{\small Source: Focus\_Evidence/evidence/FOCUS-MATRIX/kernel/KernelFinal.lean:192.\par}

{\small Source SHA-256: a0e3432e3fa233f5d79d56cf8c90138310752bc3ed52a5f7f1eaab44bfab3fec\par}

\end{samepage}
\begin{samepage}
\subsection{RecoveredMMTKernel.remaining\_factor\_unique\_uw}
\begin{quote}\small\ttfamily theorem remaining\_factor\_unique\_uw \{p : Type*\}\newline     (X Y : Matrix (Fin 23) p \ensuremath{\mathbb{Q}})\newline     (h : (B\_uw.map (Int.castRingHom \ensuremath{\mathbb{Q}})) * X = (B\_uw.map (Int.castRingHom \ensuremath{\mathbb{Q}})) * Y) : X = Y\end{quote}
{\small Source: Focus\_Evidence/evidence/FOCUS-MATRIX/kernel/KernelFinal.lean:204.\par}

{\small Source SHA-256: a0e3432e3fa233f5d79d56cf8c90138310752bc3ed52a5f7f1eaab44bfab3fec\par}

\end{samepage}
\begin{samepage}
\subsection{RecoveredMMTKernel.remaining\_factor\_unique\_vw}
\begin{quote}\small\ttfamily theorem remaining\_factor\_unique\_vw \{p : Type*\}\newline     (X Y : Matrix (Fin 23) p \ensuremath{\mathbb{Q}})\newline     (h : (B\_vw.map (Int.castRingHom \ensuremath{\mathbb{Q}})) * X = (B\_vw.map (Int.castRingHom \ensuremath{\mathbb{Q}})) * Y) : X = Y\end{quote}
{\small Source: Focus\_Evidence/evidence/FOCUS-MATRIX/kernel/KernelFinal.lean:216.\par}

{\small Source SHA-256: a0e3432e3fa233f5d79d56cf8c90138310752bc3ed52a5f7f1eaab44bfab3fec\par}

\end{samepage}
\begin{samepage}
\subsection{RecoveredMMTKernelSeed1.remaining\_factor\_unique\_uv}
\begin{quote}\small\ttfamily theorem remaining\_factor\_unique\_uv \{p : Type*\}\newline     (X Y : Matrix (Fin 23) p \ensuremath{\mathbb{Q}})\newline     (h : (B\_uv.map (Int.castRingHom \ensuremath{\mathbb{Q}})) * X = (B\_uv.map (Int.castRingHom \ensuremath{\mathbb{Q}})) * Y) : X = Y\end{quote}
{\small Source: Focus\_Evidence/evidence/FOCUS-MATRIX/additional-seeds/seed\_1/proof/Seed1KernelFinal.lean:192.\par}

{\small Source SHA-256: c3e50cc5b7cfae9880b802e9ba2af987f2504cf6d1e9c39ece761d14b80af3e1\par}

\end{samepage}
\begin{samepage}
\subsection{RecoveredMMTKernelSeed1.remaining\_factor\_unique\_uw}
\begin{quote}\small\ttfamily theorem remaining\_factor\_unique\_uw \{p : Type*\}\newline     (X Y : Matrix (Fin 23) p \ensuremath{\mathbb{Q}})\newline     (h : (B\_uw.map (Int.castRingHom \ensuremath{\mathbb{Q}})) * X = (B\_uw.map (Int.castRingHom \ensuremath{\mathbb{Q}})) * Y) : X = Y\end{quote}
{\small Source: Focus\_Evidence/evidence/FOCUS-MATRIX/additional-seeds/seed\_1/proof/Seed1KernelFinal.lean:204.\par}

{\small Source SHA-256: c3e50cc5b7cfae9880b802e9ba2af987f2504cf6d1e9c39ece761d14b80af3e1\par}

\end{samepage}
\begin{samepage}
\subsection{RecoveredMMTKernelSeed1.remaining\_factor\_unique\_vw}
\begin{quote}\small\ttfamily theorem remaining\_factor\_unique\_vw \{p : Type*\}\newline     (X Y : Matrix (Fin 23) p \ensuremath{\mathbb{Q}})\newline     (h : (B\_vw.map (Int.castRingHom \ensuremath{\mathbb{Q}})) * X = (B\_vw.map (Int.castRingHom \ensuremath{\mathbb{Q}})) * Y) : X = Y\end{quote}
{\small Source: Focus\_Evidence/evidence/FOCUS-MATRIX/additional-seeds/seed\_1/proof/Seed1KernelFinal.lean:216.\par}

{\small Source SHA-256: c3e50cc5b7cfae9880b802e9ba2af987f2504cf6d1e9c39ece761d14b80af3e1\par}

\end{samepage}
\begin{samepage}
\subsection{RecoveredMMTKernelSeed2.remaining\_factor\_unique\_uv}
\begin{quote}\small\ttfamily theorem remaining\_factor\_unique\_uv \{p : Type*\}\newline     (X Y : Matrix (Fin 23) p \ensuremath{\mathbb{Q}})\newline     (h : (B\_uv.map (Int.castRingHom \ensuremath{\mathbb{Q}})) * X = (B\_uv.map (Int.castRingHom \ensuremath{\mathbb{Q}})) * Y) : X = Y\end{quote}
{\small Source: Focus\_Evidence/evidence/FOCUS-MATRIX/additional-seeds/seed\_2/proof/Seed2KernelFinal.lean:192.\par}

{\small Source SHA-256: edf0277c1191d282ca7259e8000c26029e89795902d6a42fb78de4e34c66fa26\par}

\end{samepage}
\begin{samepage}
\subsection{RecoveredMMTKernelSeed2.remaining\_factor\_unique\_uw}
\begin{quote}\small\ttfamily theorem remaining\_factor\_unique\_uw \{p : Type*\}\newline     (X Y : Matrix (Fin 23) p \ensuremath{\mathbb{Q}})\newline     (h : (B\_uw.map (Int.castRingHom \ensuremath{\mathbb{Q}})) * X = (B\_uw.map (Int.castRingHom \ensuremath{\mathbb{Q}})) * Y) : X = Y\end{quote}
{\small Source: Focus\_Evidence/evidence/FOCUS-MATRIX/additional-seeds/seed\_2/proof/Seed2KernelFinal.lean:204.\par}

{\small Source SHA-256: edf0277c1191d282ca7259e8000c26029e89795902d6a42fb78de4e34c66fa26\par}

\end{samepage}
\begin{samepage}
\subsection{RecoveredMMTKernelSeed2.remaining\_factor\_unique\_vw}
\begin{quote}\small\ttfamily theorem remaining\_factor\_unique\_vw \{p : Type*\}\newline     (X Y : Matrix (Fin 23) p \ensuremath{\mathbb{Q}})\newline     (h : (B\_vw.map (Int.castRingHom \ensuremath{\mathbb{Q}})) * X = (B\_vw.map (Int.castRingHom \ensuremath{\mathbb{Q}})) * Y) : X = Y\end{quote}
{\small Source: Focus\_Evidence/evidence/FOCUS-MATRIX/additional-seeds/seed\_2/proof/Seed2KernelFinal.lean:216.\par}

{\small Source SHA-256: edf0277c1191d282ca7259e8000c26029e89795902d6a42fb78de4e34c66fa26\par}

{\small\textbf{Sources and attribution:} \href{https://github.com/alejandrozu/openmath-2026-judging/blob/d0ed487f487fa7ec814ff705ab55249983fe8707/OpenMath-Judging/review/entries/E02-FOCUS-MATRIX.txt}{1. alejandrozu/openmath-2026-judging / E02-FOCUS-MATRIX.txt}; \href{https://github.com/alejandrozu/openmath-2026-judging}{Central source and adjudication archive}.\par}

\end{samepage}
\clearpage
\section{Star edge colourings of generalized Petersen graphs and certified reductions towards the six-colour conjecture}
{\small Provisional mathematical authors: Woohyuk Kang\par}\medskip
{\small Woohyuk Kang  -  HTPeo; Kyung Hee University, CS\&E.\par}

{\small Affiliations follow the recorded author declarations or public institutional/profile sources. Preferred author names, author order and publication affiliations await author review. This editorial compilation does not establish anthology coauthorship.\par}

\subsection{Abstract}
The latest construction gives five-star edge colourings of GP(n,k) for 1\ensuremath{\leq}k\ensuremath{\leq}15 and n\ensuremath{\geq}2k+1, except GP(3,1), alongside other infinite families, small cubic multigraph classifications and conditional reduction theorems. The general six-colour conjecture remains open.

\ensuremath{\chi}\ensuremath{^{\prime}}\_s(GP(n,k))\ensuremath{\leq}5 for 1\ensuremath{\leq}k\ensuremath{\leq}15, n\ensuremath{\geq}2k+1, (n,k)\ensuremath{\neq}(3,1); additionally selected infinite-family colourings, finite six-colour cases and explicitly conditional reduction theorems.

\subsection{Definitions and principal unconditional scope}
A star edge colouring is a proper edge colouring with no two-coloured path or cycle of four edges. Parallel edges in a loopless multigraph are distinct and adjacent. The Dvořák-Mohar-Šámal conjecture asks whether every subcubic multigraph has star chromatic index at most six; the located general published bound is seven.

The strongest current generalized-Petersen endpoint states that GP(n,k) has a five-star edge colouring for 1\ensuremath{\leq}k\ensuremath{\leq}15 and n\ensuremath{\geq}2k+1, except(n,k)=(3,1). Here GP(n,k) has an outer n-cycle, n spokes and inner edges of step k. The condition n>2k makes the construction an ordinary simple cubic graph. This theorem includes parameter regions not covered by the earlier even-order or gcd-based results.

The packet also contains colourings of flower and Goldberg snarks, Möbius and inflation families, and a matching-colour-class six-colour construction for GP(n,2). These should be itemized as selected endpoints, with known subfamilies subtracted from the proposed novelty. The whole set remains one competition contribution family, even though a paper may naturally have several theorems.

\subsection{Periodic blocks, seams and local-window lifting}
The GP(n,k) construction colours repeated blocks and inserts a finite seam determined by the residue of n modulo a selected period. Each k has an explicit ordinary graph colouring recipe, small exceptional tables and representative windows. For example, the k=4 development repeats a period-five block, handles 9\ensuremath{\leq}n<19 directly and compares later windows with representatives 25+(n mod 5).

A forbidden four-edge path or four-cycle touches only a bounded range of blocks. The BlockStar.star\_of\_windows theorem embeds such a candidate into a window of length 3D+1, where D bounds the relevant edge jump. A finite representative-window check therefore rules out every forbidden pattern in all large n. This is an infinite proof supported by finite local certificates, rather than an inference that a finite experimental table probably continues.

The independent ordinary audit regenerated 1,366 GP instances across k=1,...,15, including all small orders through 100 and large examples near 1000, and checked proper colouring and every relevant four-edge walk. Fifty-one flower-snark examples also passed. All GP JSON tables matched the actual Lean literals. These 1,417 tests support semantic fidelity and the finite data, while the formal local-window theorem supplies the intended unbounded conclusion.

\subsection{Small multigraphs and the reduction programme}
The unconditional finite endpoints include six-star edge colourability of all bridgeless cubic loopless multigraphs with at most 14 vertices, and of all subcubic multigraphs with at most seven vertices. They are not interchangeable: lifting a cubic census result to arbitrary subcubic graphs can increase the number of vertices, so the 14-vertex threshold cannot simply be copied.

The reduction library organizes a general loopless-subcubic theorem into cubic and leaf statements. Its basic equivalence says the full conjecture follows from six-colourability of bridgeless cubic graphs together with the six-colourability of their appropriate leaf attachments. Components, bridges and a degree-two core are handled by restriction and gluing lemmas. Two-edge-cut decomposition and digon insertion then concentrate the unresolved cases on more structured cubic hosts.

The library also introduces perfect-matchings-as-colour-classes, far-exchange sets and pole-dominance data. An MC colouring has one prescribed perfect matching as its sixth colour class. EX1-goodness asks for compatible choices preserving an edge's membership status; pole dominance encodes the boundary colour compatibility needed to glue cubic pieces. Theorems prove implications from precisely stated remaining hypotheses. They do not prove those hypotheses globally.

\subsection{Prior work and the actual new regions}
The 2024 Omoomi-Vahid Dastjerdi paper proves all gcd(n,k)\ensuremath{\geq}3 and selected gcd-two cases, and discusses earlier Zhu-Shao even/gcd-one classes. It does not establish all odd n\ensuremath{\geq}7 for k=2. That odd k=2 region is a particularly clean candidate new infinite scope in the current GPAll theorem. The latest October 2 source replaces the disputed partial k=2 and k=3 arguments described in an older novelty memo.

Known covers of Petersen and K4, selected even odd-k classes and previously proved periodic families should be acknowledged as background. Newness of a definition alone does not make every theorem about it a substantive obstacle removal. The contribution needs a theorem-by-theorem comparison demonstrating which parameter regions and gluing/reduction capabilities are new.

An optional strong-five-colour hypothesis may be false in light of the Deng-Liu-Feng 2025 infinite cubic family requiring six colours. Its exact graph assumptions require full-text comparison. An implication from that hypothesis can remain logically valid, even if the hypothesis is false; it would then not be a route proving the conjecture. This concern does not invalidate the explicit GP and flower endpoints.

\subsection{Formal record and manuscript structure}
An independently qualified composite import-only route checks the exact scientific scope of all 228 authored modules, using 99 preserved original-owned outputs through read-only V1 copies, 122 unaffected V3-owned passes and seven new V4 source checks. Four import headers differ from the frozen original source; every scientific statement, proof body, option and comment remains identical. The original frozen-source replay remains a separate 119-module checkpoint. The original organizer audit exited 1 after printing 120 of its 123 requested endpoints because it omitted the import needed to expose three Star6Equiv declarations; that actual failure and its raw results are preserved. A separate organizer supplement adds only import Star6Equiv, repeats all 123 original requests and completes their fresh axiom audit. 120 use standard kernel axioms, 3 are axiom-free, and 0 depend on recognized native evaluation trust; none contain selected admissions or unrecognized axioms. No selected endpoint uses native evaluation trust. A separate imported finite certificate uses native\_decide and lies outside this selected endpoint set. These checks verify the submitted finite constructions, infinite-family constructions and conditional implications as typed. They do not prove the five remaining infinite structural premises or the general DMS conjecture. Competition p=.15 is retained.

The ordinary-proof expansion supplies the unconditional generalized-Petersen construction's exact period,seam and exception tables for 1\ensuremath{\leq}k\ensuremath{\leq}15, including 519 finite instances and 20,612 windows checked independently. It also explains the elementary I/II reduction, every middle-seam obstruction in two-/three-cut gluing, the prescribed matching-bit induction and the actual universal host/marking census bridge. The delivered base12,simple14,b14d assembly supersedes the older partial census snapshot. The five infinite structural premises remain open and the current full fresh replay is separately tracked. The user retains p=.15 for contest assessment; it is not a fraction of the general conjecture proved.

This is a substantial formal combinatorics preprint candidate. A readable main text plus repository tables and formal closure is preferable to a manuscript whose evidence is only a module count. The general DMS conjecture, all k\ensuremath{\geq}16 and the unrestricted cubic/subcubic cases remain open in this submission.

\subsection{The coloring target and the matching-preserving refinement}
All graphs in this chapter are finite loopless multigraphs, and parallel edges are distinct edges. A star edge coloring is a proper edge coloring with no bicolored path or cycle of four edges. The Dvořák-Mohar-Šámal target asks for six colors on every subcubic multigraph. The submitted work does not prove that assertion. Its main contribution is a collection of reductions and finite/class results that isolates explicit remaining structural statements. The organizer has retained event credit p=.15 for this family; that administrative quantity is not a mathematical fraction of a theorem proved.

Call a six-color star coloring MC if its sixth color class is a perfect matching. A cubic graph X is EX1-good if, for every edge g and each membership bit t\ensuremath{\in}\{0,1\} realized by some perfect matching, there is an MC coloring whose sixth matching has [g\ensuremath{\in}P]=t. This prescribed-bit requirement makes a recursive gluing argument stronger than a statement that X merely has some six-color star coloring.

Let G be the connected bridgeless loopless cubic multigraphs. Let D\_class consist of those X\ensuremath{\in}G for which every two-edge-cut side has size two on at least one side; these are called two-cut-reduced. A simple cubic graph is cyclically four-edge-connected when deletion of at most three edges leaves at most one component containing a cycle. These are exact structural scopes: the two far-exchange assumptions below concern cyclically four-edge-connected simple hosts, while pole and leaf assumptions concern the larger two-cut-reduced multigraph class.

\subsection{Digon insertion, leaf attachment and far exchanges}
For a simple cubic host H and a subset D\_edges of its edges, form H\textasciicircum{}\{D\_edges\} by replacing each marked edge ab by a path a-u-v-b with two distinct parallel u-v edges. This adds two vertices per marked edge and preserves cubic degree. For a cubic G\_0 and an edge g=st, the leaf graph T(G\_0,g) deletes g, inserts a new vertex x joined to s and t, and attaches a new leaf \ensuremath{\ell} to x. An eligible g is the unique edge between its ends and lies in no two-edge cut.

Let P be a perfect matching of X. An edge set C disjoint from P is a far-exchange set if every P edge meets at most one C edge; no edge outside P\ensuremath{\cup}C that is parallel to no C edge has both endpoints incident with C edges; and any two P edges meeting no C edge have edge distance at least three. Here edge distance is the minimum length of a path between endpoints of the two edges. These conditions are designed to make C the fifth color class while P is the sixth.

FE-EXT-D10 states that for every cyclically four-edge-connected simple cubic H with at least ten vertices, every D\_edges with |V(H)|+2|D\_edges|\ensuremath{\geq}16, every perfect matching P of X=H\textasciicircum{}\{D\_edges\}, and every far-exchange C for P, there is a star six-coloring with sixth class exactly P and fifth class exactly C. FE-EXIST-D10 states, on the same H and D\_edges scope, that for every edge g and each bit t\ensuremath{\in}\{0,1\} there are P,C with [g\ensuremath{\in}P]=t and C far-exchange for P. These are separate existence and completion requirements; neither is asserted proved universally.

Their conjunction immediately implies EX1-goodness on that host class: given g,t, apply existence to choose P,C and then completion to color them. The chosen sixth class has the required membership bit. This is a genuine conditional implication, and the original universal six-color theorem cannot be obtained by simply removing either premise. The analogous far-exchange existence statement on every two-cut-reduced host is false in the source's reported counterexample, so extending its scope would invalidate the reduction route.

\subsection{The pole-dominance interface and the triangle exception}
At a cubic vertex v with three distinct neighbors x\_1,x\_2,x\_3, delete v and attach a separate port edge p\_i at each x\_i. An MC pole coloring is a star six-coloring in which each original remaining vertex meets exactly one sixth-color edge. For each port define Col\_i as the two colors on the non-port edges incident with x\_i. Define Blk\_i as those colors on an edge x\_i r for which r meets another edge of the same color as p\_i. These are exactly the local colors capable of participating in a length-four alternating crossing of a seam.

An admissible outside datum chooses pairwise distinct port colors a\_1,a\_2,a\_3, with a\_i=6 at a specified matching port i, and sets B\_t\ensuremath{\subseteq}O\_t of two outside colors. At the matching port, O\_i avoids the other two port colors and the sixth color, and B\_i=O\_i. At a nonmatching port t, sixth color belongs to O\_t, a\_t does not, and B\_t avoids sixth color and the other nonmatching port colors. The datum also excludes the simultaneous cross-incidences a\_k\ensuremath{\in}O\_j and a\_j\ensuremath{\in}O\_k. A compatible pole coloring has p\_t color a\_t and satisfies Col\_t\ensuremath{\cap}O\_t\ensuremath{\cap}(Blk\_t\ensuremath{\cup}B\_t)=\ensuremath{\varnothing} at each port.

A pole is dominant if every admissible outside datum at every matching port has a compatible coloring; and, for each internal edge g and each bit s, there is some matching-port index i for which every admissible datum admits a compatible coloring with [color(g)=6]=s. The quantifier order in the second clause is significant: one port is chosen before the outside data vary. POLE asserts this dominance for every eligible three-neighbor vertex of every two-cut-reduced host with at least ten vertices.

The source's three-cut gluing lemma uses an EX1-good side and a dominant pole side, prescribing the three seam colors and ruling out the alternating crossings through the displayed compatibility condition. One minimum triangle configuration falls outside that ordinary pole interface. TD-TRI handles a triangle y\_1,y\_2,y\_3 forming a cyclic three-cut side in a non-cyclically-four-connected simple host, when exactly the edge y\_2y\_3 carries a digon insertion. Its required coloring has the three specific edges e\_1,y\_2u,vy\_3 in color six and preserves the prescribed matching bit for every outside edge g whose bit is realizable by a perfect matching constrained to contain those three edges. The realizability condition is retained; it is not an assertion about arbitrary assignments of bits.

\subsection{The reduction theorem and its remaining finite premises}
IID states that for every two-cut-reduced cubic G\_0 with at least six vertices and every eligible edge g, the leaf graph T(G\_0,g) has a star six-coloring. The ordinary main reduction is: FE-EXT-D10, FE-EXIST-D10, POLE, TD-TRI and IID imply star six-colorability of every finite loopless subcubic multigraph, subject to the source's established finite census inputs and gluing lemmas. The proof first composes the two far-exchange statements, uses separately certified small cyclically-four-connected hosts, and then inducts across cyclic three-cuts using dominance and the triangle exception. This yields EX1-goodness for all sufficiently large two-cut-reduced hosts.

The next reduction removes two-edge cuts while preserving the realizable matching bit, yielding the bridgeless cubic statement. A separate leaf reduction transfers IID to arbitrary bridgeless cubic hosts. Finally a general subcubic graph is reduced through components and bridge gluing. Once every remaining bridge is pendant, deleting leaves leaves a bridgeless core with degrees two and three. With at least two degree-two vertices one completes it to a cubic bridgeless graph by joining the attached leaves cyclically, using a digon in the two-vertex case. With exactly one degree-two vertex, suppress it to obtain a cubic host and use the leaf graph. Restriction of a star coloring to a subgraph preserves both properness and the forbidden-path condition.

The historical ROOT-CS4 statement kept BASE12,SIMPLE14,B14D as explicit finite premises. The current delivered source discharges them by universal host classification, kernel-checked marking coverage, sound coloring certificates and isomorphism transport, as explained below. Therefore the current implication may be stated with the five structural premises alone. Those five premises remain open, so this is still not an unconditional proof of DMS. A historical incomplete certificate count in the older author manuscript is not the status of the delivered current source assembly; present fresh replay is tracked separately.

Other useful conditional endpoints include the reduction from a sharp five-color statement for connected bridgeless simple cubic graphs except K\_\{3,3\} and the prism, and a closed-neighborhood missing-color condition. Those stronger hypotheses also remain open. A logical reduction theorem can be a publishable result in its own right, but only when its hypotheses, its finite trusted inputs and its gluing proof are stated completely and compared to prior reductions.

\subsection{Finite achievements, useful class results and publication boundary}
The delivered finite results include all two-cut-reduced cubic cases of order ten,twelve and fourteen, with universal census completeness supplied in the same source chain. The small-host digon insertion reduction uses 451 representative markings of K\_4,K\_\{3,3\}, Q\_3 and the Wagner graph and checked marking-orbit coverage. This replaces the older manuscript snapshot's partial fourteen-vertex assembly. The source also proves the eligible leaf condition through fourteen vertices; larger finite pole experiments and the universal structural hypotheses must retain their respective audit boundaries.

A further formal theorem gives a six-color star edge coloring of GP(k,2), k\ensuremath{\geq}5, with the spokes exactly the sixth color class. This prescribed matching feature is stronger than merely saying the graph has six colors. The source also contains five-color constructions for flower and Goldberg snark families and certain covers and generalized Petersen families. Some of those overlap published results, and two universal generalized-Petersen five-color claims are audit-disputed and excluded by the original author manuscript. They should not re-enter the accepted-result anthology through a broad family summary.

The surviving bottlenecks are concrete. The far-exchange completion needs compatible local data even when exchange neighborhoods crowd. Pole dominance must accommodate all admissible seam data. The triangle route must handle matching configurations not transferred by the available one-crossing gluing. IID must be established above the finite range. Evidence from large searches is valuable for choosing lemmas, but a search with no counterexample does not establish any of these universal statements.

This exposition is adapted with attribution from the preserved star 6\_partial.tex and its accompanying statement records. The paper itself says its mathematics was generated with a multi-agent proof-search system and separates model audits, kernel checks, and computational evidence. Woohyuk Kang is credited for the DMS result in the pinned team contribution table; provisional human authorship still requires author approval. The original detailed manuscript should be retained beside the new synthesis so that definitions, fact identifiers, side-branch disputes and proof-source boundaries remain traceable.

\subsection{A complete elementary proof of the two cubic sufficient conditions}
Let I mean that every finite bridgeless loopless cubic multigraph has a star six-coloring. Let II mean that T(G,g) has a star six-coloring for every such cubic G and every edge g, including parallel edges. The source states I\ensuremath{\wedge}II if and only if every loopless subcubic multigraph has a star six-coloring. The reverse implication follows by applying the subcubic assertion to the two displayed graph types. The following elementary argument expands the source's forward proof outline and makes its bridge step explicit.

We first record a gluing observation for pendant edges. Suppose A and B are star six-colored graphs with designated pendant edges ua and vb, where a and b are leaves. Delete a and b and replace their two pendant edges by one edge uv. After a permutation of colors in B, the two former pendant edges can have the same color \ensuremath{\lambda}. Let S\_A be the colors of the other edges at u and S\_B the colors of the other edges at v. Each set has size at most two and avoids \ensuremath{\lambda}. There are five remaining colors, so a color permutation fixing \ensuremath{\lambda} can send S\_B into the complement of S\_A: that complement has at least three elements. Extend this injection on S\_B to a permutation of all five colors. Thus S\_A and the relabeled S\_B are disjoint.

Give uv color \ensuremath{\lambda}. The result is proper at u and v because \ensuremath{\lambda} belongs to neither adjacent color set. Any four-edge path lying on one side, or consisting of uv and three edges on one side, has the same colored pattern as a path in one of the original pendant-edge graphs. It is therefore not bicolored. A four-edge path that uses edges on both sides of uv has one or two edges before the bridge and the complementary two or one after it. If it were bicolored, the edges immediately adjacent to uv on its two sides would have the same color, contrary to S\_A\ensuremath{\cap}S\_B=\ensuremath{\varnothing}. A cycle cannot contain the new bridge. This exhausts all forbidden patterns and proves that the glued coloring is a star coloring. The argument requires the original pendant edges to have been colored; arbitrary colorings of A\ensuremath{-}a and B\ensuremath{-}b alone need not extend.

Now induct strongly on the number of edges of a connected loopless subcubic H; different components can be colored separately. If a bridge uv has at least one edge on each side, split it and replace it by a new pendant edge in each component. Each resulting graph has fewer edges than H. Apply the induction hypothesis, then the preceding gluing observation. Thus it remains to treat graphs in which every bridge is pendant. If deleting the leaves leaves an edgeless core, H is a star of degree at most three and has an immediate coloring with three colors. Otherwise its core C is connected and bridgeless, and every core vertex has degree two or three. Indeed a bridge remaining in C would be a non-pendant bridge of H. A core vertex of degree one would make its incident core edge such a bridge. All leaves of H are attached to degree-two core vertices, at most one leaf at each.

If C has no degree-two vertex, H=C and I applies. If it has k\ensuremath{\geq}2 degree-two vertices, attach one leaf to each, using an existing leaf of H when present. Join these k leaves in a cycle when k\ensuremath{\geq}3, or by two parallel edges when k=2. The resulting graph G is loopless and cubic. It is bridgeless: old core edges remain on their old cycles, new leaf-cycle edges lie on the new cycle or digon, and every attachment edge lies on a cycle formed from a path in C between two attachment vertices, the corresponding attachment edges and a path in the new leaf cycle. H embeds as a subgraph of G. Color G by I and restrict the coloring to H.

Finally suppose C has exactly one degree-two vertex y. Its two neighbors are distinct. If both edges went to the same neighbor x, then x would have degree three, and its remaining edge would be a bridge leaving the pair \{x,y\}, contradicting bridgelessness. Suppress y, replacing its two-edge path by a new edge g. The resulting G\_0 is loopless, cubic and bridgeless; cycles are preserved by suppression, and parallel edges cause no difficulty. If H already has a leaf at y, then H is exactly the corresponding leaf graph T(G\_0,g); otherwise it is a subgraph of it. Apply II and restrict. These cases finish the induction. Restricting a star coloring preserves properness and cannot introduce a new bicolored four-edge path or cycle.

This proves the I/II sufficient-condition theorem as an ordinary graph argument. It does not prove I or II. The further EX1, two-cut and three-cut reductions are what seek to derive those conditions from the five stated open structural hypotheses. Their separate matching-preservation and seam obligations are not eliminated by this elementary bridge lemma.

\subsection{The finite certificate argument and its exact completeness boundary}
For a fixed finite edge list E, a candidate certificate supplies a color function c:E\ensuremath{\to}\{1,...,6\}. A direct checker verifies that distinct edges sharing a vertex have different colors. It then enumerates four-edge paths and four-edge cycles and rejects a certificate when the four colors use only two values. To verify MC, it counts the color-six edges incident with every vertex and requires that count to be exactly one. To verify a prescribed edge bit, it also checks c(g)=6 exactly when t=1. These conditions are finite and are exactly the mathematical conditions for an MC star coloring with that bit. A soundness theorem proves that a successful Boolean check implies those properties; kernel evaluation of the literal data supplies the successful checks in the indicated Lean modules.

There are two logically separate universal quantifiers. The certificate theorem covers every graph literally listed, every edge of each listed graph and the specified membership bits. A graph-class theorem also requires that every graph in the claimed class be isomorphic to one of those lists. Isomorphism transports the coloring, perfect matching and prescribed edge bit, so the latter completeness claim, once supplied, bridges the two statements. Kernel checking of every listed coloring alone does not establish that no graph was omitted. Nor does byte equality with a census file establish completeness of the external census generator.

BASE12 concerns ten- and twelve-vertex two-cut-reduced bridgeless cubic multigraphs and the realizable matching bit at each edge. SIMPLE14 covers the 341 simple three-edge-connected cubic hosts on fourteen vertices. B14D covers 984 non-simple two-cut-reduced marking representatives on fourteen vertices. The older author paper described a partial R\_0 through R\_791 certificate snapshot. In the current delivered source all groups are assembled in d14, including the universal host and marking completeness bridge. The historical partial count is retained as provenance, not a current missing-proof claim.

The small-host digon result has a stronger integrated reduction. It uses 451 representative edge markings: one for K\_4, thirteen for K\_\{3,3\}, 129 for the cube Q\_3 and 308 for the Wagner graph. The total is 1+13+129+308=451. Here the selected source includes the reduction to those markings as well as their color certificates. The scope remains cyclically four-edge-connected simple hosts on at most eight vertices with enough inserted digons to give at least sixteen vertices. It is not a claim that every arbitrary marking or every larger host is covered by those 451 literal instances.

The historical layer 28 implication exposes BASE12,SIMPLE14,B14D as premises. Current base12,simple14,b14d prove each; d1012\_of/d14\_of contain the actual classification-to-certificate linkage. The final current chain retains the open structural assumptions but not an external census premise. Successful fresh replay of that assembly remains distinct from reading or reusing archived compilation receipts.

\subsection{An unconditional delivered theorem: Petersen parameters k from one to fifteen}
The current pinned pack 4 source contains a separate unconditional construction theorem: for 1\ensuremath{\leq}k\ensuremath{\leq}15 and n\ensuremath{\geq}2k+1, excluding the pair (n,k)=(3,1), GP(n,k) admits a star edge coloring with five colors. In the k=1 row this means n\ensuremath{\geq}4. This is the theorem gp\_star5 in GPAll.lean, obtained by the fifteen exact GPnK.star5 constructors. It is distinct from the open all-subcubic DMS assertion and from the older disputed claims about unrestricted generalized Petersen parameters. Historical originality of each individual k case must still be compared with earlier work; formal correctness does not mean that all fifteen cases are new.

The graph has vertices u\_j,v\_j for j modulo n and three edge types at each block j: outer u\_ju\_\{j+1\}, spoke u\_jv\_j and inner v\_jv\_\{j+k\}. In the source these are vertices 2j,2j+1 and edge IDs 3j,3j+1,3j+2. A color pattern is a triple (outer,spoke,inner) with entries in \{0,1,2,3,4\}. Thus there is a literal explicit coloring of the actual Petersen graph, rather than a theorem about an abstract structure that is assumed to be colorable. The five-color result does not prescribe the spokes as one color class; the separate six-color GP(n,2) matching-preserving theorem has that additional property.

For each k the data consist of a period p, a seam base length s, a transition threshold T, literal small-graph patterns for n\_min\ensuremath{\leq}n<T, and p residue-specific rows A\_r. Each A\_r has length p+s+r. Its first p triples are a common period word P, and the last s+r triples are the seam word for r. For n\ensuremath{\geq}T put r=n mod p and L=n\ensuremath{-}r\ensuremath{-}s. At block j<L use P\_\{j mod p\}; at j\ensuremath{\geq}L use the seam triple at offset j\ensuremath{-}L. In source notation this is idx(n,j)=j mod p in the periodic part and p+j\ensuremath{-}L in the seam, followed by tabs[r][idx(n,j)]. For n<T the entire n-triple coloring is the corresponding smallTab row.

The complete literal period, seam and exception data for all fifteen cases have been extracted without alteration into dms\_gp\_periodic\_certificate\_data.json. It records the SHA-256 of every GPnKT.lean and GPnK.lean against the frozen build plan, their exact source paths, and line locators for tabs, smallTab, idx, chi, Wsmall, rep and star5. Together with the ordinary reduction below and the independent checker, this is a reproducible finite certificate for the construction. No SAT solver or search heuristic is part of the correctness argument.

\subsection{The period/seam case table and a fully printed example}
Read each tuple in the following table as (n\_min,p,s,T,B,N). The checked finite range is n\_min\ensuremath{\leq}n<N; an arbitrary larger graph uses the representative n\ensuremath{^{\prime}}=B+(n mod p). The window spans block offsets 0 through 3k, so it contains 3k+1 consecutive cyclic blocks. The tuple rows are: k=1: (4,4,4,12,12,16); k=2: (5,5,5,15,20,25); k=3: (7,4,4,15,20,24); k=4: (9,5,5,19,25,30); k=5: (11,6,6,23,30,36); k=6: (13,7,7,27,35,42); k=7: (15,8,8,31,40,48); k=8: (17,9,9,35,45,54); k=9: (19,10,10,39,50,60); k=10: (21,6,12,39,48,54); k=11: (23,8,16,47,64,72); k=12: (25,9,18,52,63,72); k=13: (27,8,16,51,64,72); k=14: (29,10,20,59,80,90); k=15: (31,8,16,55,72,80).

Every row satisfies N=B+p, B divisible by p, B\ensuremath{\geq}T and B\ensuremath{\geq}s+3k+p. These simple arithmetic checks are part of the extracted certificate. They ensure that every representative is in the finite checked range, uses the period/seam branch, and has enough untouched periodic space to represent any local window. The longest row, k=15, has n\_min=31, p=8, s=16, T=55, B=72, N=80 and a 46-block window; the finite source checks exactly the graphs 31\ensuremath{\leq}n<80.

For k=1 the common period is P=[(2,0,1),(1,3,0),(0,2,3),(4,1,2)]. The seams for residues r=0,1,2,3 are respectively [(2,0,1),(3,4,0),(1,4,2),(4,3,4)]; [(3,1,0),(2,4,1),(0,4,3),(1,4,2),(4,3,4)]; [(3,1,0),(0,4,2),(2,3,1),(0,4,3),(1,4,2),(4,3,4)]; and [(3,1,0),(1,4,2),(0,4,3),(2,4,1),(0,4,3),(1,4,2),(4,3,4)]. For n\ensuremath{\geq}12 repeat the period through block n\ensuremath{-}r\ensuremath{-}5 and append the appropriate seam, of length 4+r. The finite exception list covers 4\ensuremath{\leq}n<12. This prints one whole infinite construction in ordinary notation; the other fourteen have the identical algorithm with the exact data in the supplement.

The locator pattern is uniform: GPnKT.lean contains tabs, smallTab, idx and chi; GPnK.lean contains Wsmall, rep and star5; its imported GPnKW\_i.lean files partition the finite checked interval. GPk.gp in GPPeriodic.lean supplies the graph wiring and BlockStar.star\_of\_windows supplies the local-to-global theorem. For example GPn15.Wc1 in GPn15W1.lean checks 31\ensuremath{\leq}n<40 and the remaining eight W files extend the checks through n=79. All these are theorem/code locators in the frozen pack 4 source, not references to a later guessed construction.

\subsection{Ordinary finite-certificate proof of the infinite Petersen construction}
The local checker can be described directly. Regard block positions temporarily as integers on a line. The outer, spoke and inner edges have jumps 1,0,k and their prescribed endpoint vertex types. At each of the two vertices in the central block position 2k, list the three incident half-edges. Every incident pair e\_2\ensuremath{\neq}e\_3 must have different colors. For every edge e\_1 incident at the other end of e\_2 and every edge e\_4 incident at the other end of e\_3, reject the pattern when c(e\_1)=c(e\_3) and c(e\_4)=c(e\_2). Such a rejected pattern is exactly the alternating four-edge obstruction centered at the chosen vertex. The source check is slightly conservative because it examines all such half-edge choices; this strengthens its sufficiency.

All edge owners in those tests lie between positions 0 and 3k. Indeed the two central incident edges have owner positions between k and 2k; their other endpoints lie between k and 3k; and any edge incident at either other endpoint has owner between 0 and 3k. Therefore the list of 3k+1 color triples determines every properness and four-edge test centered at that vertex. Every actual four-edge path or cycle has a central vertex between its second and third edges, so it is excluded by one such test. Reading these windows cyclically modulo n proves a star coloring whenever every cyclic starting window passes. This is the ordinary content of BlockStar.star\_of\_windows, whose formal soundness proof also permits a short graph's window to wrap more than once.

For n\_min\ensuremath{\leq}n<N, check every starting position j=0,...,n\ensuremath{-}1 directly against the literal data. The source proves those finite Boolean checks with kernel decide. A separate standard-library Python checker has independently checked both all local windows and all four-edge walks with distinct edges on the actual finite Petersen graphs. All 519 parameter instances and all 20,612 local windows pass. It also checks every color is in the stated five-color set, every exception row has exactly n triples, all residue rows have exactly p+s+r triples, and their first p triples agree. These are concrete finite verifications, not a search without a counterexample.

For n\ensuremath{\geq}N set r=n mod p, n\ensuremath{^{\prime}}=B+r and L=n\ensuremath{-}r\ensuremath{-}s. Fix a cyclic starting position j<n. If j+3k<L, the entire window lies in the periodic region and does not wrap. Set j\ensuremath{^{\prime}}=j mod p. Since j\ensuremath{^{\prime}}+3k\ensuremath{\leq}p\ensuremath{-}1+3k<B\ensuremath{-}s, its window in the representative also lies in the periodic region. The two windows use exactly the same period indices at every offset.

If j+3k\ensuremath{\geq}L, set j\ensuremath{^{\prime}}=j+n\ensuremath{^{\prime}}\ensuremath{-}n. This is nonnegative because j\ensuremath{\geq}n\ensuremath{-}r\ensuremath{-}s\ensuremath{-}3k and hence j\ensuremath{^{\prime}}\ensuremath{\geq}B\ensuremath{-}s\ensuremath{-}3k\ensuremath{\geq}p; also j\ensuremath{^{\prime}}<n\ensuremath{^{\prime}}. The shift n\ensuremath{^{\prime}}\ensuremath{-}n is a multiple of p. Before either window wraps, the shift preserves a periodic index modulo p and preserves the seam offset relative to L, whose representative counterpart is B\ensuremath{-}s. The wrap thresholds agree because j\ensuremath{^{\prime}}+h\ensuremath{-}n\ensuremath{^{\prime}}=j+h\ensuremath{-}n. After wrapping, both windows have the same small nonnegative block position, at most 3k, which is in the periodic region since B\ensuremath{-}s>3k. Thus their complete index lists agree for every offset 0\ensuremath{\leq}h\ensuremath{\leq}3k.

In either case the arbitrary n-window equals a checked window of the finite representative B+r, which lies between B and N\ensuremath{-}1. The local checker therefore passes at every position in the arbitrary graph. Its soundness proves the five-color star property for every n\ensuremath{\geq}N; the finite range proves all smaller admitted n. This completes the infinite construction by a finite certificate and a proved period/seam reduction. It has no external graph-census completeness premise, because the graph family is defined explicitly for every integer n,k. The proof neither solves any of the five open DMS structural statements nor establishes any of the open premises in ROOT-CS4.

\subsection{Two-cut seam lemma: an exhaustive four-edge obstruction argument}
Use \ensuremath{\mu} for the matching color, the sixth color in ordinary notation and Lean color 5. Let \{e\_1,e\_2\} be a two-edge cut, with e\_i=a\_i b\_i, distinct a\_1,a\_2 on side A and distinct b\_1,b\_2 on side B. A pole consists of a side's internal edges and both cut edges, whose remote ends are leaves in that pole. Suppose the two pole colorings are star colorings and agree on the cut edges. Their union is proper: every pair of adjacent edges is visible in a pole containing its common vertex. If neither middle edge of a four-edge walk is a cut edge, the central three vertices are on one side and all four edges lie in its pole. Thus only a cut in position two or three needs a new argument; reversing the walk reduces position three to position two. This includes four-cycles as well as paths.

First suppose both cut edges belong to the global perfect matching M, which is precisely the \ensuremath{\mu} color class. Require the colors of nonmatching edges at a\_i to be disjoint from those at b\_i, for each i. A potentially alternating walk with e\_i in position two has nonmatching edges in positions one and three because matching edges cannot be adjacent. Alternation would give these two edges the same color, contradicting that disjointness. This proves glue\_M. The matching is obtained by joining the two internal covers and the two cut edges; each endpoint then has its unique matching edge and all other vertices retain theirs.

Now suppose neither cut edge belongs to M. Require (F1) different colors on e\_1,e\_2 whenever a matching edge joins a\_1,a\_2 or b\_1,b\_2, and (F2) different colors on the nonmatching internal edges at opposite ends of each cut edge. Consider a potentially alternating walk f\_1,e\_i,f\_3,f\_4 oriented from A to B. If f\_1 is nonmatching, so is f\_3 because their colors agree, and F2 contradicts that agreement. Otherwise f\_1 and f\_3 are matching. If f\_4 is a cut edge, it must be the other cut edge, and f\_3 is a matching rung joining b\_1,b\_2; F1 contradicts equality of the colors of e\_i and f\_4.

In the remaining case f\_4 is internal to B. Append the unique matching edge m at its far endpoint v\_4. Then e\_i,f\_3,f\_4,m is an alternating four-edge walk entirely in the B pole: positions one and three have the assumed equal nonmatching color, while positions two and four have color \ensuremath{\mu}. In the original four-edge path or four-cycle, v\_4 differs from both endpoints v\_2,v\_3 of f\_3, so m\ensuremath{\neq}f\_3. The far endpoint of m cannot be v\_2 or v\_3, since matching uniqueness there would force m=f\_3, and cannot be the remote A endpoint of e\_i because m is internal to B. Looplessness excludes its being v\_4. Thus the appended walk is a genuine forbidden path or cycle, contradicting the B pole's star property. This proves glue\_F and exhausts all seam positions. Source locators are MhFact\_c64b6bd62b78d841.lean:glue\_prop,glue\_mid,glue\_M,glue\_F. No general assertion that arbitrary pole colorings always glue is being made.

\subsection{Prescribing the edge bit in two-cut gluing}
For this prescribed-bit gluing theorem, additionally assume at most one internal edge joins a\_1,a\_2 and at most one joins b\_1,b\_2. These are essential hypotheses, later derived by the minimal-two-cut selection; distinct cut endpoints alone do not imply them. The edge closure G\_A adds an edge g\_A=a\_1a\_2 to the internal A edges; define G\_B similarly. The digon closure G\_A\textasciicircum{}D instead retains the two cut edges and adds two parallel edges b\_1b\_2 at their remote ends. All these closures are connected bridgeless loopless cubic graphs. A perfect matching of a closure containing g\_A induces a pole cover in which both cut edges are matching. A matching avoiding g\_A induces a cover in which neither cut edge is matching. The two types must agree across the seam. Every edge of a connected bridgeless cubic multigraph can be included in and avoided by a perfect matching, the Petersen-Schönberger property PStat proved in the source. Consequently EX1-goodness permits either bit at any requested closure edge.

To preserve a request on an internal B edge, color G\_B by EX1 with that edge bit. For a request on a cut edge, impose the bit on g\_B instead: the two cut edges have the same matching status after gluing. If g\_B is matching, choose an EX1 coloring of G\_A with g\_A matching. The nonmatching color sets at each closure endpoint have two elements. In either closure their two endpoint sets cannot be identical: identical sets would force two common colors across a matching edge, and the star property forbids equal colors on distinct nonmatching edges at its opposite ends. At most one rung can be a single edge visible at both ends. A permutation fixing \ensuremath{\mu} can make the endpoint sets disjoint across both seams. After normalizing the first pairs to \{0,1\}, the remaining statement is the finite g2mChk test on 5\textasciicircum{}4=625 tuples and the 5!=120 permutations of the nonmatching colors. Duplicate entries and identical endpoint sets are explicitly excluded; every admitted tuple passes. The seam lemma then gives an MC coloring with the requested bit.

If g\_B is not matching and there is no matching rung b\_1b\_2, select a nonrung internal edge incident with a\_1 and require it to be matching in G\_A. This forces g\_A to be nonmatching and excludes a matching rung a\_1a\_2. At each seam end there is exactly one nonmatching internal edge. Align the two closure-edge colors \ensuremath{\gamma} by a permutation fixing \ensuremath{\mu}, then permute the other four colors so that the two B internal colors avoid the at most two A internal colors. Such a permutation exists because the complement in four colors has at least two elements. The source also checks all 4\textasciicircum{}4=256 normalized cases in g2fChk. The two seam edges may share \ensuremath{\gamma} because no matching rung is present. Apply glue\_F.

The remaining case is a matching rung on side B. Here use EX1 on G\_A\textasciicircum{}D with one remote digon edge matching. Its two pendant seam colors are distinct and avoid \ensuremath{\mu}: if they were equal, the other, nonmatching parallel digon edge and the pendant edges would form a forbidden alternating four-edge walk using the matching parallel edge. On the B pole, initially both pendant edges have the old closure color \ensuremath{\gamma}. Recolor only e\_2 to \ensuremath{\gamma}\ensuremath{^{\prime}}. Let f\_1 and f\_2 be the nonmatching internal edges at b\_1 and b\_2; let h be the other nonmatching edge at the far endpoint of f\_2. Choose \ensuremath{\gamma}\ensuremath{^{\prime}} outside \{\ensuremath{\gamma},c(f\_1),c(f\_2),c(h),\ensuremath{\mu}\}. Six colors suffice because at most five are forbidden. A walk containing the recolored pendant edge can only contain it first or last. If its next edge is the rung, its third edge has color \ensuremath{\gamma} or c(f\_1); if its next edge is f\_2, its third edge has color \ensuremath{\mu} or c(h). None equals \ensuremath{\gamma}\ensuremath{^{\prime}}. Thus the recoloring preserves the star property and leaves every internal B edge, its matching bit and the \ensuremath{\mu} class unchanged.

Map the two A pendant colors to \ensuremath{\gamma},\ensuremath{\gamma}\ensuremath{^{\prime}} while fixing \ensuremath{\mu}. The B nonmatching internal color at either seam avoids \ensuremath{\gamma},\ensuremath{\gamma}\ensuremath{^{\prime}},\ensuremath{\mu}. A permutation fixing these three colors can move those two B colors off the corresponding A colors: among the three remaining colors one avoids two exclusions when the B colors agree; when they differ, each has at least two allowed images and distinct images can be chosen. This is perm\_4c. Both F1 and F2 now hold, so glue\_F finishes the rung case. Exchanging A,B covers a requested internal A edge. The resulting exact theorem ex1\_2c\ensuremath{^{\prime\prime}} assumes at most one internal rung on each side, EX1 for G\_A,G\_B and additionally for G\_A\textasciicircum{}D when B has a rung, and G\_B\textasciicircum{}D when A has a rung. The argument does not extend one fixed arbitrary perfect matching; it produces a matching with the one requested edge bit. The finite permutation checks have also been independently re-enumerated in dms\_gluing\_census\_exposition\_audit.json.

\subsection{Small two-cut interfaces and the well-founded two-cut induction}
Small closures need stronger interface certificates when EX1-goodness fails. The source provides two exact recipes. In an M certificate the cut edges have color \ensuremath{\mu} and the internal matching covers A except a\_1,a\_2. Let S\_i be the two nonmatching colors at a\_i. TSmall means S\_1\ensuremath{\neq}S\_2, equivalently |S\_1\ensuremath{\cap}S\_2|\ensuremath{\leq}1, and TMeet means S\_1\ensuremath{\cap}S\_2 is nonempty. In an F certificate the internal matching covers all A and both cut edges avoid \ensuremath{\mu}. Feq requires equal seam colors and no matching rung; Fneq requires unequal seam colors; Fnc additionally requires that each seam endpoint's nonmatching internal color avoid the other seam edge's color.

Recipe D supplies an M/TMeet certificate and an Fnc certificate, and for every internal edge and each requested bit a certificate of one of these two types preserving that bit. Recipe E supplies M/TSmall, Feq and Fneq certificates, and bit coverage by M/TSmall or Feq. This is a finite collection of color functions with a universal finite check on every internal edge. From a digon closure, restricting a matching either gives M/TMeet or Fneq; swapping the two parallel remote edges handles the second digon orientation. Hence D on A plus EX1 for G\_B\textasciicircum{}D suffices. From an edge closure with at most one rung, a matching through the closure edge gives M/TSmall; avoiding it gives Feq without a matching rung or, after the preceding recoloring, Fnc with a matching rung. Hence E on A plus EX1 for G\_B suffices. In each case the bit is taken from the side containing the requested edge and all other choices supply only a compatible interface. These are Cut2.sr\_a and Cut2.sr\_b in MhFact\_657e581f1145250f.lean.

Assume H: every two-cut-reduced connected bridgeless cubic graph of order at least ten is EX1-good. Strongly induct on the even vertex count n to prove EX1 for every connected bridgeless cubic graph of order at least ten. If the graph is two-cut-reduced, H applies. Otherwise choose a two-cut whose two sides have at least four vertices, minimizing the smaller side A. Then 4\ensuremath{\leq}a\ensuremath{\leq}b, both a,b are even, the A edge closure is two-cut-reduced, and there is at most one rung on either side. The four closure orders are a,b,a+2,b+2, all strictly below n because the opposite side has at least four vertices. Thus no induction invocation is circular.

If a\ensuremath{\geq}10, use H on the reduced A closure, the induction hypothesis on B and both required digon closures, and ex1\_2c\ensuremath{^{\prime\prime}}. If a\ensuremath{\in}\{4,6,8\} and b\ensuremath{\geq}10, the larger closures are supplied by induction; the proved SmallFacts give an EX1 closure or a D/E interface on A, and the preceding recipes handle the exceptional four-vertex digon obstruction. If both sides have order at most eight, the same recipes handle all cases except the stated residual a=4,b=6 interface, which is a separately certified finite table. The source SmallFacts has exactly five clauses s2,s3,s4,s5,tabc and is proved by sfS2,sfS3,sfS4,sfS5,sfTabc; it is not a sixth open infinite hypothesis. This establishes ex1red. Together with the separately certified cubic cases of order at most eight, it gives the cubic coloring conclusion conditional on H.

\subsection{Three-cut compatibility, dominance and bit-preserving gluing}
For a three-cut e\_i=y\_iw\_i with all y\_i distinct and all w\_i distinct, define Col\_i to be colors on internal edges at that seam endpoint. Define Blk\_i to contain \ensuremath{\kappa} when an internal edge of color \ensuremath{\kappa} from the endpoint reaches a vertex incident with another pole edge of color c(e\_i). The exact gluing condition is Col\_i(A)\ensuremath{\cap}Col\_i(B)\ensuremath{\cap}(Blk\_i(A)\ensuremath{\cup}Blk\_i(B))=\ensuremath{\varnothing} for every i. It is sufficient without further unmentioned seam conditions. The two poles agree on e\_i, so properness is inherited. Distinct cut edges cannot be consecutive in a walk because their endpoints are distinct. If a middle edge is e\_i and the walk crosses A\ensuremath{\to}B, its first and third edges are internal on opposite sides. In an alternating walk their common color \ensuremath{\kappa} lies in both Col sets, and its fourth edge certifies \ensuremath{\kappa}\ensuremath{\in}Blk\_i(B). The reverse orientation puts \ensuremath{\kappa} in Blk\_i(A). A cut in the other middle position is handled by reversing the walk; when no middle edge is a cut, all four edges are in one pole. These cases prove glue 3.

Contract A to a new cubic hub joined to w\_1,w\_2,w\_3; write X/A for the resulting B closure. It is connected, bridgeless and cubic. A cut of size two in this contraction lifts to a cut of size two in X, so two-cut reduction is preserved. Its order is b+1. An MC coloring has exactly one matching-color hub edge, determining a port i. Its two outside colors O\_t and blocking colors B\_t form the admissible datum stated earlier: properness and the absence of alternating walks through the hub give each avoidance and cross-incidence clause. A compatible A pole therefore glues by the displayed seam criterion and retains a unique matching-color edge at every original vertex.

Dominance has two different jobs. D1 says that for each matching port every admissible outside datum has a compatible coloring. D2 says that for each internal edge g and bit t there exists a port i before the outside datum varies, such that every admissible datum at i has a compatible coloring with [c(g)=\ensuremath{\mu}]=t. If the request lies in A, choose that D2 port, use EX1 on X/A to put its hub edge in \ensuremath{\mu}, and then use D2 on the actual outside datum. If the request lies in B or on a cut edge, use EX1 on X/A to preserve that request, take its actual matching port and use D1. This proves the exact brick theorem: dominant A pole plus EX1-good B contraction implies EX1-good X. Reversing the sides gives the symmetric theorem. Quantifiers cannot be interchanged: merely finding a convenient coloring for one outside datum would not prove dominance.

The source also gives a useful conditional induction without yet solving its core. A three-cut is reducible when a\ensuremath{\geq}9 and either b\ensuremath{\geq}9, the B contraction is already EX1-good, or b\ensuremath{\leq}7 and the B pole is dominant. Under POLE, the A pole is dominant by applying POLE to its opposite contraction, of order a+1\ensuremath{\geq}10. In the first case apply strong induction to the B contraction of smaller order b+1; in the second use its supplied EX1 theorem; in the third apply induction to the A contraction of smaller order a+1 and use the flipped brick. Every contraction preserves the graph class, and a,b\ensuremath{\geq}3 ensure strict descent. Consequently, proving EX1 for the graphs with no reducible three-cut would prove H under POLE. This is layer\_3cut, a conditional reduction, rather than a proof that all cores satisfy EX1.

\subsection{Why the small cubic census is exhaustive in the delivered source}
The current source contains a universal census bridge; it does not rely only on checking an externally generated list. Let S be the simple, three-edge-connected cubic graphs, with vertex count on incident vertices. The removable-edge theorem says that every member of S of even order n\ensuremath{\geq}6 has an edge uv whose deletion and suppression of u,v yields another member of S of order n\ensuremath{-}2. The reduction removes uv and replaces the other u-edges uu\_1,uu\_2 by u\_1u\_2, and the other v-edges vv\_1,vv\_2 by v\_1v\_2. Conversely, subdividing two distinct edges of the reduced graph and joining the new subdivision vertices reconstructs the original graph. Vertex and edge labels need not be canonical; explicit isomorphisms perform the transport.

Here is the cut argument behind removability. For any proper nonempty side T in a three-edge-connected graph, |\ensuremath{\delta}(T)|\ensuremath{\geq}3. Crossing counts are submodular: |\ensuremath{\delta}(A\ensuremath{\cap}T)|+|\ensuremath{\delta}(A\ensuremath{\cup}T)|\ensuremath{\leq}|\ensuremath{\delta}(A)|+|\ensuremath{\delta}(T)|. Choose A with a minimum number of vertices among nontrivial three-cut sides. If a nontrivial three-cut T separates u,v\ensuremath{\in}A, uncross T with A, or with its complement when the union is the whole vertex set. Both resulting nonempty cuts have at least three edges, so equality forces a three-cut side properly contained in A. Minimality says that side is a singleton. Applying the complementary uncrossing shows A\ensuremath{\cap}T=\{u\} and A minus T=\{v\}. This would give |A|\ensuremath{\leq}2, whereas the cubic handshake identity 3|A|=2|E(A)|+3 forces |A| odd, and a nontrivial side has at least three vertices. Thus no nontrivial three-cut separates two vertices of A.

If the graph contains a triangle, its three vertices give a minimum nontrivial three-cut side, and choose a triangle edge uv. If it is triangle-free and has a nontrivial three-cut, choose any internal edge of a minimum side; such an edge exists by the same handshake count. If it has no nontrivial three-cut, choose any edge. The preceding uncrossing excludes a bad three-cut separating u,v in the first two cases; the last case excludes it directly. The exact crossing-count identity for the suppressed graph then shows that a cut with fewer than three edges would lift to just such a bad original three-cut, so none exists. In the triangle-free cases the new edges cannot be parallel to old edges, because that would make a triangle. Nor can the two new edges have the same unordered endpoint pair: then u and v would share a neighbor, making a triangle through uv. In the triangle case, an old edge parallel to either new edge, or equality of the two new unordered endpoint pairs, would put at least five edges on at most four vertices, leaving a cut of size at most two to the nonempty remainder when n\ensuremath{\geq}6. This also contradicts three-edge connectivity. Looplessness and cubic degree follow directly from the four distinct incident edges. These cases prove removability without assuming an exhaustive external graph generator.

Start at order four: a simple cubic graph on four vertices is K\_4. The sound insertion-table theorem scl\_step says that a complete list at order n plus a certificate for every ordered pair of distinct edges of every listed host gives a complete list at n+2. Each table entry supplies an isomorphism from the inserted graph to a claimed representative. The packed permutation is checked to be bijective; equality of adjacency bitsets, simplicity, and equal edge counts make the induced edge map a bijection as well. The source iterates this step through orders 4,6,8,10,12,14, whose list sizes are 1,2,4,14,57,341. The literal tables are kernel-checked certificates; the searches that discovered them are unnecessary to the proof. This proves universal classification of S in those six orders, not just membership of the listed graphs. Locators are removable and scl\_step/scl\_stepA in MhFact\_54ed0ff9a1ed26b6.lean, scl 4/6/8/10 in MhFact\_4f913d5550e57c85.lean, scl 12 in MhFact\_fbae65436c8cc8bf.lean and scl 14 in MhFact\_e1599bf3f3c26bdb.lean.

\subsection{From universal host classification to BASE12, SIMPLE14 and B14D}
Every two-cut-reduced connected bridgeless cubic multigraph X of order at least ten is isomorphic to Q\textasciicircum{}D for a simple three-edge-connected cubic host Q. Contract each digon gadget, replacing its two internal vertices and four edges by its marked host edge. Two-cut reduction excludes a longer chain of consecutive gadgets: a side containing two consecutive digons would have four vertices and a two-edge boundary with a larger opposite side. The surviving host retains cubic degree, connectivity and bridgelessness. Any two-cut in the host would lift to a forbidden nontrivial two-cut in X; and a parallel pair in a cubic host with at least four vertices gives such a two-cut. The possible two-vertex triple-edge host has at most three marked edges and reconstructs a graph of at most eight vertices, excluded here. Thus Q belongs to S, has even order q\ensuremath{\geq}4, and |V(X)|=q+2|D|. The formal theorem hasm\_converse constructs the representation and to\_dig turns it into isomorphism transport for the finite certificates.

For a fixed labeled host with m edges, encode D by an m-bit integer. The source orbit coverage covAllB considers every bit vector of the required weight and gives a checked host automorphism taking it to a listed representative code. This is an exhaustive finite quantifier, not an appeal to a graph-isomorphism search tool. The sound theorem digClassB extends that host isomorphism to a digon-insertion isomorphism. digClassZ handles the zero marking. At total order ten or twelve, the representative counts grouped by host order q=4,6,8,10,12 are 5,24,56,78,57, totaling 220. At total order fourteen, the counts for q=4,6,8,10,12,14 are 1,21,124,356,482,341, totaling 1,325. The 341 zero-marked fourteen-vertex hosts are the simple graphs; the other 984 representatives are the non-simple cases. These totals concern marking orbits of all admitted host lists, not merely a count of colorings.

Each representative supplies GoodH: an MC star six-coloring realizing both bits at every edge, together with six-color leaf-graph certificates at every edge whose ends are digon-free. The representative checker verifies the neighbor list, the star condition on all four-edge walks, the unique matching-color edge at each vertex and the edge-bit coverage. Its soundness is repChk\_sound. An eligible edge of X has no parallel partner and lies in no two-edge cut. Both ends must be digon-free, since a digon at one end would isolate its two vertices with a two-edge cut containing that edge. Hence the leaf certificate applies after isomorphism. All relevant properties transport under the explicit vertex/edge bijections. This is concl\_of\_iso, not a new assumption about unlisted graphs.

The source assembly d1012\_of calls to\_dig, the complete host classifiers, digClassB/Z and concl\_of\_iso; d1012 supplies all its actual GoodH arguments and proves EX1 and the eligible leaf assertion for every X of order ten or twelve. Thus base12:BASE12 is unconditional at this delivered source scope. Similarly d14 supplies all six host-order certificate groups and proves both properties at order fourteen. Its simple14 and b14d corollaries discharge SIMPLE14 and B14D, with no remaining external enumeration premise. The older author manuscript's snapshot covered only R\_0 through R\_791 of the 984 non-simple representatives. That historical limitation should be retained as provenance of that old draft; it is superseded by the delivered current assembly gB4,gB6,gB8,gB10,gB12,gB14 and d14. Fresh replay of every required module is still a distinct verification step, and no replay status is inferred from reading a theorem statement.

Combining these census theorems with the two-cut induction yields all bridgeless cubic loopless multigraphs of order at most fourteen, including disconnected ones by components. The leaf reduction yields T(X,g) for every connected bridgeless cubic host of order at most fourteen, even when g is in a digon or a two-cut. The delivered bounded subcubic reduction may double the original vertex count, so the unconditional arbitrary-subcubic conclusion is order at most seven. These precise thresholds are Star6.star6\_cubic\_bridgeless\_le14, matching\_colour\_class\_10\_14, star 6\_leaf\_le14 and star 6\_subcubic\_le7; one must not replace the seven by fourteen.

\subsection{The five-premise conditional induction and its sole triangle exception}
Assume exactly the five structural statements FE-EXT-D10, FE-EXIST-D10, POLE, TD-TRI and IID stated above. Their truth is not established here. The two far-exchange statements compose to give EX1 on every digon insertion of a cyclically four-edge-connected simple host of order at least ten and total order at least sixteen. For host order at most eight the separately proved SMALLHOST-D theorem supplies the same result from 451 kernel-covered marking representatives of K\_4,K\_\{3,3\},Q\_3 and the Wagner graph. These are finite completed inputs, not additional open universal assumptions.

Strongly induct on n to prove H for all two-cut-reduced cubic graphs of order at least ten. The even orders ten,twelve,fourteen are the census bases just proved. For n\ensuremath{\geq}16 pass by hasm\_converse to Q\textasciicircum{}D. If Q is cyclically four-edge-connected, use the preceding large-host or small-host theorem. Otherwise choose a cyclic three-cut in Q. Its cut ends are distinct: two at the same cubic vertex would leave a smaller cut on a cyclic side, contrary to three-edge connectivity. Lift the cut to Q\textasciicircum{}D. Its two side orders a,b are odd, by 3a=2|E(A)|+3. Their hub contractions remain two-cut-reduced cubic graphs and have orders a+1,b+1.

If a,b\ensuremath{\geq}9, POLE gives dominance on one contraction and induction gives EX1 on the other contraction, whose order is smaller than n; apply brick. If one side has at most seven vertices, classify its contraction among the certified cubic graphs of order at most eight. Every case either has EX1, giving a brick using POLE on the larger opposite contraction, or has a certified dominant pole, giving a brick using induction on the larger side. The only residual labeled apex configuration is Rep\_6 at its two exceptional apex vertices. The source proves that this configuration corresponds precisely to a triangle in Q with one inserted digon on its edge opposite the apex. If a cut edge is marked, shift the cut through that digon, increasing the small side from five to seven vertices, where the ordinary small-side cases apply. If no cut edge is marked, the residual five-vertex triangle-with-digon piece is exactly the TD-TRI scope.

Write that piece's nine pole edges in the order (e\_i,e\_j,e\_k,y\_iy\_j,y\_iy\_k,y\_ju,d\_1,d\_2,vy\_k), where d\_1,d\_2 are parallel uv edges. Contracting the far side gives the six-vertex cubic graph with edge list [(0,1),(0,2),(0,3),(1,2),(1,3),(2,4),(4,5),(4,5),(5,3)]. Four checked MC color rows, in Lean colors 0 through 5, are [0,5,2,1,5,3,5,0,4], [0,5,2,1,5,3,0,5,4], [0,2,5,5,1,4,5,0,3] and [0,2,5,5,1,4,0,5,3]. Each has one color-5 edge at each vertex and no alternating four-edge walk. These give the matching-port j and k cases under far-side dominance. The exceptional matching-port i requires the constrained matching \{e\_i,y\_ju,vy\_k\}, supplied by TD rather than assumed to follow from ordinary EX1.

For an outside edge request, D2 on the dominant far-side pole chooses a port. Ports j,k use the displayed finite rows and then compatible gluing. At port i the source's admissible datum construction produces a matching of the far-side contraction through the i hub edge with the requested bit; lifting it by adding y\_ju and vy\_k gives a realizable matching containing the three constrained piece edges. TD-TRI then gives the required MC coloring preserving that bit. For a piece edge or cut edge, the displayed rows realize every bit except those forcing the constrained triple, which are again supplied by TD. This is a finite nine-edge case split in tri\_brick\_td, not a proof of TD-TRI. Matching lift and restriction explicitly preserve every outside edge bit. It completes the residual case and hence the induction H, conditional on the five assumed premises.

Now apply ex1red to obtain EX1 for all sufficiently large bridgeless cubic graphs. Small graphs are covered separately, so condition I follows. The source leaf induction gives condition II from H and IID: remove a two-cut side of at least four vertices not wholly containing g, use a two-sided interface supplied by H and the finite small cases, and apply induction to the smaller closure containing g or its closure-edge image. If no such side remains, parity forces every two-cut side not containing g to have exactly two vertices. The graph is then two-cut-reduced; a parallel partner of g would produce a reducible complementary side, and a two-cut containing g would force both sides to have two vertices, contrary to order at least six. Thus g is eligible and IID applies. Orders at most four are finite base cases. Finally the already proved ordinary I/II reduction gives DMS. In symbols the delivered theorem is (FEEXTD10\ensuremath{\wedge}FEEXISTD10\ensuremath{\wedge}POLE\ensuremath{\wedge}TDTRI\ensuremath{\wedge}IID)\ensuremath{\Rightarrow}DMS, with its finite premises discharged in the current source. The implication is proved; the conjunction is not.

\subsection{Current scope, audit status and remaining mathematical work}
The earlier draft's census-to-graph gap is absent from the current delivered Lean chain: universal host classification, marking coverage, color-certificate soundness, isomorphism transport and finite-hypothesis discharge are all supplied by named theorems. This ordinary exposition explains that chain and the exhaustive two-/three-cut seam arguments. It should be read with the preserved literal certificates and proof source, whose file hashes and exact theorem locators are recorded in dms\_gluing\_census\_exposition\_audit.json. The short historical author paper and current delivered archive are different snapshots. Source inspection alone is not a completed replay; the separately qualified composite route and the preserved partial original-source route are distinguished in the formal record above.

What remains mathematically open is the truth of the five infinite structural statements, or a different route that avoids any false one. Their conditional conclusion must never be described as an unconditional proof of DMS. In the latest source FEEXISTD10 can be restricted to the remaining larger host range FEEXISTD18, and IID to IID16, because additional finite theorems cover the intervening cases. Those sharpened reductions do not establish the remaining universal instances. The unconditional GP(n,k) theorem for k from one to fifteen, the bounded cubic/subcubic results and the logical reduction theorems are the publication scope. Their historical novelty and human author approval remain separate from formal correctness. The conditional proof is source-faithful editorial exposition with finite supplements, not a claim that an organizer supplied new proofs of the open premises.

\subsection{Verification and reproducibility}
Fresh execution: htpeo-dms-current: SCHEDULED\_RESOURCE\_CHECKPOINT; 119/228 custom modules compiled successfully; receipt Verification/htpeo-dms-current-fresh-build.json. htpeo-dms-core: PASS; 1/1 custom modules compiled successfully; receipt Verification/htpeo-dms-core-fresh-build.json. Separate qualified body-identical composite route: all 228 scientific bodies are checked through 99 previous original-owned outputs via read-only V1 copies, 122 unaffected V3-owned passes and seven new V4 source checks. The original frozen-source 119-module checkpoint is not replaced. The original organizer audit exited 1 with 120/123 prints; a separate supplement adds only import Star6Equiv and audits all 123 unchanged requests, with 120 standard-kernel and three axiom-free endpoints. A separate imported finite certificate, MGraph.k4subdiv\_star6 in StarCert, uses native\_decide outside that selected set; the whole portfolio is not asserted standard-only. The five infinite structural premises and general DMS conjecture remain open, and p=.15 is unchanged. Qualification: Verification/dms-v4-vector-composite-supplement-final-qualification-20261005.json. Compiler pins, source hashes and endpoint axiom outputs are in Verification. Historical receipts and finite checks do not replace fresh replay.

\begin{samepage}
\subsection{gp\_star5}
\begin{quote}\small\ttfamily theorem gp\_star5 (k n : Nat) (hk : 1 \ensuremath{\leq} k) (hK : k \ensuremath{\leq} 15) (hn : 2 * k + 1 \ensuremath{\leq} n) (hex : \ensuremath{\neg} (n = 3 \ensuremath{\wedge} k = 1)) :\newline     \ensuremath{\exists} c : Fin (gp n k).m \ensuremath{\to} Fin 5, (gp n k).Star 5 c\end{quote}
{\small Source: entries/dms-star6/artifact/lean/pack4/src/GPAll.lean:25.\par}

{\small Source SHA-256: e831e584837dc1e0a2e4ef26bf7875bebd3a18ad54deff326b3c68d1af970893\par}

\end{samepage}
\begin{samepage}
\subsection{gp\_family}
\begin{quote}\small\ttfamily theorem gp\_family (k n : Nat) (hk : 1 \ensuremath{\leq} k) (hK : k \ensuremath{\leq} 15) (hn : 2 * k + 1 \ensuremath{\leq} n) (hex : \ensuremath{\neg} (n = 3 \ensuremath{\wedge} k = 1)) :\newline     StarFam (gp n k) 5\end{quote}
{\small Source: entries/dms-star6/artifact/lean/pack4/src/GPAll.lean:45.\par}

{\small Source SHA-256: e831e584837dc1e0a2e4ef26bf7875bebd3a18ad54deff326b3c68d1af970893\par}

{\small\textbf{Sources and attribution:} \href{https://arxiv.org/html/2410.15024v1}{1. arXiv source: 2410.15024v1}; \href{https://arxiv.org/abs/1011.3376}{2. arXiv source: 1011.3376}; \href{https://doi.org/10.7151/dmgt.2195}{3. DOI: 10.7151/dmgt.2195}; \href{https://link.springer.com/article/10.1007/s40840-025-01875-9}{4. link.springer.com / s40840-025-01875-9}; \href{https://github.com/alejandrozu/ulam-opdp-difficulty-atlas}{5. alejandrozu/ulam-opdp-difficulty-atlas}; \href{https://github.com/alejandrozu/openmath-2026-judging/blob/d0ed487f487fa7ec814ff705ab55249983fe8707/OpenMath-Judging/review/entries/E08-DMS.txt}{6. alejandrozu/openmath-2026-judging / E08-DMS.txt}; \href{https://github.com/alejandrozu/openmath-2026-judging}{Central source and adjudication archive}.\par}

\end{samepage}
\clearpage
\section{An exact 39-line arrangement with 471 triangular faces}
{\small Provisional mathematical authors: Woohyuk Kang\par}\medskip
{\small Woohyuk Kang  -  HTPeo; Kyung Hee University, CS\&E.\par}

{\small Affiliations follow the recorded author declarations or public institutional/profile sources. Preferred author names, author order and publication affiliations await author review. This editorial compilation does not establish anthology coauthorship.\par}

\subsection{Abstract}
Thirty-nine explicit rational straight lines form a simple arrangement with 471 bounded triangular faces. Exact geometry independently confirms the data and improves the recorded pre-event 470 construction by one; no optimality is claimed.

There exists a simple arrangement of 39 ordinary straight lines with 471 bounded triangular faces; hence K(39)\ensuremath{\geq}471.

\subsection{Kobon triangles and the exact theorem}
In the ordinary straight-line Kobon problem, count bounded triangular faces formed by the arrangement, with disjoint interiors. A triple of lines can determine a geometric triangle whose interior is crossed by other lines; such a triangle is not an arrangement face and must not be counted. The submitted 39-line certificate has 471 actual bounded triangular faces.

The arrangement is simple: no parallel pairs, no duplicate lines and no triple concurrence. All 741 pair intersections are distinct. The freshly replayed formal endpoints carry the rational construction into actual real triangle faces and pairwise interior-disjointness.

\subsection{Independent exact geometry}
The ordinary audit normalizes line equations, computes every rational intersection and orders the intersection points on each line. A candidate triangle is counted only when each of its three sides joins consecutive arrangement vertices. This adjacency condition excludes a side cut by another intersection and ensures the bounded triangular-face interpretation.

The independent recount gives 471, using reviewer-owned exact arithmetic and without executing the entrant's search program. Its simple-arrangement checks agree with the source. The full line list and the deterministic face list belong in supplementary JSON, accompanied by a small visualization for intuition. A plot alone cannot certify the count because close intersections may be visually indistinguishable.

\subsection{Baseline comparison and originality}
The pre-event organizer source already contains a simple39-line470-triangle arrangement and a matching SimpleLowerBound39470 theorem. An independent recount gives 470 for those same 741 distinct-intersection conventions. The proposed improvement is therefore one triangle, not an improvement from 468 obtained by overlooking the stronger pre-event certificate.

The historical straight-line/pseudoline distinction also matters. Generic doubling formulas in the located literature can concern pseudolines; they do not automatically prove stretchability of the doubled result. The cited straight-line recursive sequences do not already establish 39\ensuremath{\geq}471. This clears those particular comparisons, while a specialist historical check remains necessary before a numerical-record headline.

\subsection{Publication and formal boundary}
This is a concise exact-construction note, with a natural theorem K(39)\ensuremath{\geq}471. It does not assert K(39)=471, attainment of the general upper bound, or a new construction at every order. Keep the 39-line lower certificate separate from the known 18-line 93 reconstruction and unrelated upper-bound conjectures.

This publication pass independently replayed all 55 exact frozen authored source modules and the selected audit under Lean 4.33.1 and Mathlib 0df444a360eaa60ab8c11dca51a86af692955474, finishing on 5 October 2026 at 02:18:01 UTC. All invocations and all seven selected axiom prints passed. The concrete certificate uses only propext, sol\_simple is axiom-free, and the actual real-face lower bound, interior-disjointness and face-checker equivalence use only standard kernel axioms. No selected theorem depends on admitted, native-evaluation or unknown axioms. This closes the requested formal geometric verification of the 39-line/471-face construction; its historical originality and competition placement remain separate. Receipt: Verification/htpeo-kobon471-fresh-build.json.

\subsection{The exact line construction}
Let K(n) be the largest number of nonoverlapping triangular bounded faces obtainable from n distinct straight lines in the real plane. Here a triangular face has three supporting arrangement lines, three distinct pairwise intersection vertices, and an open interior meeting no arrangement line. The submitted theorem is the lower bound K(39)\ensuremath{\geq}471. The construction uses exact integer line coefficients, with no reliance on a plotted image or on floating-point tolerances.

In their stored order, the 39 lines have equations ax+by+c=0 with the following triples (a,b,c): (17161722537,999852726795,-127435541096); (22983693025,999735840037,-316196247610); (45670989237,998956535962,-471645285334); (94115157790,995561317586,-619750011076); (187926156025,982183160048,-766457490638); (331697758175,943385709676,-848196419199); (475600349348,879661473352,-843271769478); (-617409467007,-786641945265,778587187470); (-723884368047,-689921315584,647370952204); (-791940903741,-610597744003,482778979846); (-839112733346,-543957554168,266701231093); (-859633577426,-510911061302,112222169536); (-868955541211,-494890156902,-64658080768); (-867784656583,-496940428824,-188748162882); (-876491741189,-481416895868,-338731588296); (-896172287080,-443706245021,-513049342444); (-919742395709,-392522515959,-668703708621); (-958055287236,-286583437412,-820458209792); (-982602471473,-185721250951,-898096419089); (-999886196670,15086209235,-969689450013); (985200663602,-171404937032,951334830743); (968093210130,-250590375912,895580815198); (940679865061,-339295433905,755205631360); (911608781025,-411058913489,589025620706); (897872663908,-440255243474,437919208839); (891692807329,-452641069013,264544679621); (893476899830,-449109150954,140986222138); (886814812192,-462124971060,-33319910766); (869339155735,-494215977389,-240677165983); (836979314362,-547234526807,-412611362015); (773380286739,-633942372841,-590183916150); (630904029917,-775860879948,-753795449928); (465045543929,-885286756973,-808570424607); (-319225892954,947678652956,800152717693); (-180806126256,983518756663,687691427594); (-53896644347,998546519561,499919062709); (-11533130317,999933491241,389309310392); (15081678171,999886265024,150910488147); (18508631601,999828700606,-12114696146). Each line is proper because (a,b)\ensuremath{\neq}(0,0). All coefficients have absolute value below the hill bound 10\ensuremath{^3}\ensuremath{^0}. The large integral scale reflects rational coordinates encoded without rounding; it does not affect the geometric meaning.

\subsection{A determinant criterion for an empty triangular interior}
For coefficient triples p=(a,b,c), q=(d,e,f), write w(p,q)=ae\ensuremath{-}bd and P(p,q)=(bf\ensuremath{-}ce, cd\ensuremath{-}af, ae\ensuremath{-}bd). If w\ensuremath{\neq}0, the pairwise intersection is (P\ensuremath{_1}/w,P\ensuremath{_2}/w). For another line l, write D(l,p,q)=l\ensuremath{\cdot}P(p,q), the 3\ensuremath{\times}3 determinant of coefficient rows. Evaluating the affine form of l at p\ensuremath{\cap}q gives D(l,p,q)/w(p,q). Consequently s(l,p,q)=D(l,p,q)w(p,q) has the same sign as that affine value: their ratio is w(p,q)\ensuremath{^2}>0.

Lemma 1 (nondegenerate triangle). If p,q,r are pairwise nonparallel and D(p,q,r)\ensuremath{\neq}0, their pairwise intersections A=p\ensuremath{\cap}q, B=p\ensuremath{\cap}r, C=q\ensuremath{\cap}r are distinct and noncollinear. A coincidence would be a point on all three lines and make their coefficient determinant zero. If A,B,C were collinear, the line p already contains A,B; the third point C would therefore also lie on p, giving the same forbidden concurrency. The three open side halfplanes oriented toward the opposite vertices have intersection exactly the open triangle.

Lemma 2 (affine sign criterion). Let A,B,C be noncollinear and l a proper line. The line l misses the open triangle exactly when its affine values at A,B,C are all nonnegative or all nonpositive. For the forward implication, if some vertex value is positive and another is negative, take a small positive weight on the third vertex and adjust positive weights on the first two so their weighted affine sum is zero; the resulting strictly positive barycentric combination lies in the open triangle on l, a contradiction. For the reverse implication, the value at any interior point is a weighted sum with three strictly positive weights. If the vertex values have one weak sign, this sum has that strict sign unless all three values vanish. A proper affine form cannot vanish at three noncollinear points, so that exception is impossible.

It follows that a triple p,q,r bounds a triangular arrangement face if and only if its three pair determinants are nonzero, its triple determinant is nonzero, and for every arrangement line l the three integers s(l,p,q),s(l,p,r),s(l,q,r) are all \ensuremath{\geq}0 or all \ensuremath{\leq}0. This criterion uses only integer additions, products and comparisons. The three supporting lines themselves pass it, since each vanishes at two of the vertices and has a single nonzero value at the opposite vertex.

\subsection{The complete face census and independent evaluation}
Lemma 3 (finite certificate). For the literal 39-line list, every one of the 741 pair determinants is nonzero and every one of the 9,139 triple determinants is nonzero. Enumerate all increasing index triples i<j<k<39, apply the criterion of Lemma 2 to each, and retain precisely the passing triples. The resulting list has 471 entries. Because the increasing triples enumerate all possible supporting triples once, and the criterion is an equivalence, this is an exact count of the triangular faces of this arrangement.

A minimal checker computes P(p,q) once for every pair, then for each increasing triple and each of the 39 lines evaluates the three signed determinants. A failing mixed-sign line rejects the triple immediately. No numerical intersection coordinates, angular sort, or geometric tolerance is required. The independent ordinary checker for this editorial draft implements this determinant criterion from its mathematical definition and obtains exactly 471 passing triples. It preserves all 39 coefficient triples and the complete 471-face list in final\_four\_independent\_certificate\_audit.json.

The submitted Lean packet has a second route through the actual hill convention. On each line, arrangement intersections are ordered by x, or by y for a vertical line, equal points are merged, and a proposed side is accepted when its endpoint ranks differ by one. In a simple arrangement, a different line crossing the triangle interior must cross two of its open sides; side adjacency forbids those crossings. Conversely a triangular face has no intervening arrangement intersection on its sides. The packet proves both the determinant checker count and the transcription of the evaluator count equal 471, and that both produce the identical reported face list.

\subsection{Why the counted triangles are nonoverlapping}
Lemma 4 (cell characterization). In the complement of finitely many lines, each fixed vector of strict affine signs determines a convex open set: an intersection of open halfplanes. Any nonempty such set is a connected component of the complement. The triangle cut out by a passing triple lies wholly in one sign cell, because every line misses its interior. The three supporting halfplanes already confine that cell to the triangle, and the remaining halfplanes contain every interior point by Lemma 2. Therefore the triangle is exactly that cell.

Two distinct sign cells have disjoint interiors. If two passing triples had intersecting open triangles, their cells would coincide. A nondegenerate triangular cell has precisely three supporting boundary lines, so the two triples would consist of the same three lines. In the increasing-index convention, they would be the same triple, contrary to distinctness. Thus all 471 interiors are pairwise disjoint. This proves the usual nonoverlapping-triangle interpretation of K(39)\ensuremath{\geq}471, rather than merely counting arbitrary line triples that form triangles.

The source Faces.lean gives a direct affine proof of the same disjointness implication and assembles an explicit family of 471 distinct index triples and corresponding real vertices. Main.lean states both kobon\_lower\_bound and the stronger displayed family theorem kobon\_nonoverlapping. The ordinary determinant and cell argument here explains the geometry behind those endpoints in full.

\subsection{What the one-triangle improvement does and does not establish}
The frozen pre-event organizer material contains a straight-line 39-line construction with 470 triangles. The concrete comparison is therefore a gain of one triangle, 470\ensuremath{\to}471. It would be misleading to compare only against a weaker 468 entry when the 470 certificate is already available. This result is a new-construction candidate at that order; historical literature and repositories must still be checked before claiming it is the first such lower bound.

No upper certificate K(39)\ensuremath{\leq}471 is supplied. The general upper bound 481 remains distinct from this lower construction. Pseudoline arrangements and doubling formulas that have not been shown stretchable do not automatically give straight-line constructions of the same size. The present proof concerns this exact straight-line rational arrangement, at order 39, and its exact triangular bounded faces. Search history, the unrelated 18-line known construction, and other order-specific bounds do not enlarge that theorem.

\subsection{Fresh independent kernel verification of the geometric construction}
This publication pass independently replayed all 55 exact frozen authored source modules and the selected audit under Lean 4.33.1 and Mathlib 0df444a360eaa60ab8c11dca51a86af692955474, finishing on 5 October 2026 at 02:18:01 UTC. All invocations and all seven selected axiom prints passed. The concrete certificate uses only propext, sol\_simple is axiom-free, and the actual real-face lower bound, interior-disjointness and face-checker equivalence use only standard kernel axioms. No selected theorem depends on admitted, native-evaluation or unknown axioms. This closes the requested formal geometric verification of the 39-line/471-face construction; its historical originality and competition placement remain separate. Receipt: Verification/htpeo-kobon471-fresh-build.json.

\subsection{Verification and reproducibility}
Fresh execution: htpeo-kobon 471: PASS; 55/55 custom modules compiled successfully; receipt Verification/htpeo-kobon471-fresh-build.json. Compiler pins, source hashes and endpoint axiom outputs are in Verification. Historical receipts and finite checks do not replace fresh replay.

\begin{samepage}
\subsection{Kobon.kobon\_lower\_bound}
\begin{quote}\small\ttfamily theorem kobon\_lower\_bound :\newline     \ensuremath{\exists} L : List Line, L.length = 39 \ensuremath{\wedge} (\ensuremath{\forall} l \ensuremath{\in} L, l.a \ensuremath{\neq} 0 \ensuremath{\vee} l.b \ensuremath{\neq} 0) \ensuremath{\wedge}\newline       L.Pairwise (fun p q => \ensuremath{\neg} \ensuremath{\forall} P : \ensuremath{\mathbb{R}} \ensuremath{\times} \ensuremath{\mathbb{R}}, OnLine p P \ensuremath{\leftrightarrow} OnLine q P) \ensuremath{\wedge}\newline       \{t : \ensuremath{\mathbb{N}} \ensuremath{\times} \ensuremath{\mathbb{N}} \ensuremath{\times} \ensuremath{\mathbb{N}} | t.1 < t.2.1 \ensuremath{\wedge} t.2.1 < t.2.2 \ensuremath{\wedge} t.2.2 < 39 \ensuremath{\wedge}\newline         IsTriFace L (nth L t.1) (nth L t.2.1) (nth L t.2.2)\}.ncard = 471\end{quote}
{\small Source: entries/hills/kobon-n39-lavaskiller/lean/KobonCert/Main.lean:29.\par}

{\small Source SHA-256: eaa912ef13c79c074821a1a674f51de8dcf0c2b4a628e8a670d4633a9a579311\par}

\end{samepage}
\begin{samepage}
\subsection{Kobon.kobon\_nonoverlapping}
\begin{quote}\small\ttfamily theorem kobon\_nonoverlapping :\newline     \ensuremath{\exists} (L : List Line) (T : Finset (\ensuremath{\mathbb{N}} \ensuremath{\times} \ensuremath{\mathbb{N}} \ensuremath{\times} \ensuremath{\mathbb{N}})) (tri : \ensuremath{\mathbb{N}} \ensuremath{\times} \ensuremath{\mathbb{N}} \ensuremath{\times} \ensuremath{\mathbb{N}} \ensuremath{\to} (\ensuremath{\mathbb{R}} \ensuremath{\times} \ensuremath{\mathbb{R}}) \ensuremath{\times} (\ensuremath{\mathbb{R}} \ensuremath{\times} \ensuremath{\mathbb{R}}) \ensuremath{\times} (\ensuremath{\mathbb{R}} \ensuremath{\times} \ensuremath{\mathbb{R}})),\newline       L.length = 39 \ensuremath{\wedge} (\ensuremath{\forall} l \ensuremath{\in} L, l.a \ensuremath{\neq} 0 \ensuremath{\vee} l.b \ensuremath{\neq} 0) \ensuremath{\wedge}\newline       L.Pairwise (fun p q => \ensuremath{\neg} \ensuremath{\forall} P : \ensuremath{\mathbb{R}} \ensuremath{\times} \ensuremath{\mathbb{R}}, OnLine p P \ensuremath{\leftrightarrow} OnLine q P) \ensuremath{\wedge}\newline       T.card = 471 \ensuremath{\wedge}\newline       (\ensuremath{\forall} t \ensuremath{\in} T, t.1 < t.2.1 \ensuremath{\wedge} t.2.1 < t.2.2 \ensuremath{\wedge} t.2.2 < 39 \ensuremath{\wedge}\newline         TriWitness L (nth L t.1) (nth L t.2.1) (nth L t.2.2) (tri t).1 (tri t).2.1 (tri t).2.2) \ensuremath{\wedge}\newline       (\ensuremath{\forall} t \ensuremath{\in} T, \ensuremath{\forall} t' \ensuremath{\in} T, t \ensuremath{\neq} t' \ensuremath{\to}\newline         Disjoint (openTri (tri t).1 (tri t).2.1 (tri t).2.2)\newline           (openTri (tri t').1 (tri t').2.1 (tri t').2.2))\end{quote}
{\small Source: entries/hills/kobon-n39-lavaskiller/lean/KobonCert/Main.lean:40.\par}

{\small Source SHA-256: eaa912ef13c79c074821a1a674f51de8dcf0c2b4a628e8a670d4633a9a579311\par}

{\small\textbf{Sources and attribution:} \href{https://github.com/alejandrozu/kobon-proof/commit/99fdc8ec1ef8b1fb22c3da32b011b7361762e958}{1. alejandrozu/kobon-proof}; \href{https://arxiv.org/html/0706.0723v1}{2. arXiv source: 0706.0723v1}; \href{https://zegalur.github.io/line-order/gallery/kobon.html}{3. zegalur.github.io / kobon.html}; \href{https://github.com/alejandrozu/ulam-opdp-difficulty-atlas}{4. alejandrozu/ulam-opdp-difficulty-atlas}; \href{https://github.com/alejandrozu/openmath-2026-judging/blob/d0ed487f487fa7ec814ff705ab55249983fe8707/OpenMath-Judging/review/entries/E08-KOBON39.txt}{5. alejandrozu/openmath-2026-judging / E08-KOBON 39.txt}; \href{https://github.com/alejandrozu/openmath-2026-judging}{Central source and adjudication archive}.\par}

\end{samepage}
\clearpage
\section{A rational 23-product 3\ensuremath{\times}3 matrix-multiplication scheme of support 138}
{\small Provisional mathematical authors: Chandragupt Sharma\par}\medskip
{\small Chandragupt Sharma  -  Cluster Innovation Centre, University of Delhi.\par}

{\small Affiliations follow the recorded author declarations or public institutional/profile sources. Preferred author names, author order and publication affiliations await author review. This editorial compilation does not establish anthology coauthorship.\par}

\subsection{Abstract}
An exact rational decomposition of the 3\ensuremath{\times}3 multiplication tensor uses 23 products and 138 nonzero factor coefficients. A diagonal gauge leaves 126 unit coefficients and 12 dyadic coefficients. Exact tensor and polynomial checks confirm the finite scheme and a parameter family; rank 22, support 137 impossibility and practical speedup are not established.

The 3\ensuremath{\times}3 multiplication tensor admits a rational 23-term decomposition with factor supports 47,47,44, totalling 138.

\subsection{Tensor statement and support measure}
For 3\ensuremath{\times}3 matrices A and B, a 23-product bilinear scheme forms 23 linear expressions in A,23 in B, multiplies the corresponding pairs, and recombines the products into AB. The three factor lists U,V,W each have 23 vectors with nine coefficients. The total number of nonzero coefficients in those lists is the support measure used here.

The submitted rational scheme has support 47+47+44=138. All 729 tensor coefficient equations hold exactly. It therefore gives a valid 23-product multiplication algorithm and tensor rank at most 23. It does not construct 22 products or prove a matching tensor-rank lower bound. Compared with the located 139-support baseline, the proposed numerical advance is one nonzero coefficient.

\subsection{Exact identities and rational lifting}
The finite certificate expands each rank-one term and checks every triple of input/output coefficient indices against the multiplication tensor. The archived convention places A and B in row order and the C/trace leg in column order, which must be respected when turning the tensor into output formulas.

Reviewer-owned Fraction arithmetic confirms the 138 scheme, its gauged version and a separate 143 scheme. Every rational coefficient matches the scaled Lean literal. The core finite endpoint uses ordinary decide and reports no added mathematical axioms for validity and support. A fresh independent build under the installed Lean 4.31.0 checks those same archived core statements; this earlier alternate-version check is distinct from the subsequently completed pinned 4.34.1 replay recorded in the current verification register.

The search uses support patterns and exact rational lifting. The paper should state the literal factor lists or a compact structured representation so a reader can reproduce the algorithm directly, then describe how the search reached that support rather than presenting the search log as a proof.

\subsection{Gauges, coefficients and the parameter family}
A diagonal gauge preserves each simple tensor while redistributing scalar factors among its three lists. The gauged 138 scheme has 126 coefficients equal to\ensuremath{\pm}1 and 12 equal to\ensuremath{\pm}2 or\ensuremath{\pm}1/2. Those are coefficient shifts, not 12 additional bilinear products. Its naive addition count is 83 versus a recorded 84 baseline, while the reported common-subexpression count 77 is worse than a cited optimized 70; no overall practical speedup is demonstrated.

An independent symbolic computation clears denominators in the reported family and checks all 729 polynomial identities in Q[s]. Away from s=0 and s=\ensuremath{-}2, every specialization is valid and the support remains 138; s=\ensuremath{-}1 recovers the stored scheme. Polynomial identity also establishes validity for real parameters outside the poles, although the universal-family statement is not part of the supplied finite Lean certificate.

Reducing the support pattern with every nonzero replaced by one over GF(2) gives four failing tensor equations. This excludes a sign-only rational scheme with precisely that support pattern. It does not give a literal mod-two reduction of dyadic coefficients and does not rule out other patterns or fields.

\subsection{Structural claims that need separate proof}
The numerical tangent dimension seven and six diagonal gauge directions suggest local family structure. They do not prove a global one-dimensional quotient component, classify every parameter point, establish all projective boundary limits, minimize the number of nonunit coefficients or prove support 137 impossible. Failed searches at 137 are evidence about the tested search, not an impossibility theorem.

Rank triples of factor matrices can distinguish certain constructions, but a complete inequivalence claim needs the precise action and invariant proof. The 143 scheme is a valid auxiliary construction and may have methodological interest; it does not improve the 138/139 frontier and should not be counted as a second numerical award.

\subsection{Prior comparison and paper readiness}
Credit HKS, Smirnov and other 23-product constructions used as baselines, and compare support under the same factor and gauge conventions. The proposed new claim is the certified 138 scheme and its carefully scoped exact family, not minimum-rank matrix multiplication.

A focused algebra/computational mathematics paper can lead with the finite algorithm, include exact proof and rational parameter formulas, and separate empirical local structure from theorems. Source-pin reproducibility and prior-equivalence review remain the final gates. The current evidence makes this one of the clearest original construction candidates in the portfolio.

\subsection{From the trace tensor to an actual 23-product algorithm}
Use the trace convention T(A,B,C)=tr(ABC)=\ensuremath{\sum}\_\{i,j,k=0\}\ensuremath{^2} A\_\{ij\}B\_\{jk\}C\_\{ki\}. Each factor table u,v,w has 23 rows of nine rational coefficients; the coordinate order in all three tables is row-major in its own matrix. The claimed identity is T=\ensuremath{\sum}\_\{\ensuremath{\ell}=0\}\textasciicircum{}\{22\}U\_\ensuremath{\ell}(A)V\_\ensuremath{\ell}(B)W\_\ensuremath{\ell}(C), where U\_\ensuremath{\ell}(A)=\ensuremath{\sum}\_\{i,j\}u\_\{\ensuremath{\ell},3i+j\}A\_\{ij\}, and V,W are defined the same way on B,C.

For a,b,c\ensuremath{\in}\{0,...,8\}, write a=3i+j, b=3j\ensuremath{^{\prime}}+k, c=3k\ensuremath{^{\prime}}+i\ensuremath{^{\prime}}. The identity is equivalent to the 729 exact equations \ensuremath{\sum}\_\ensuremath{\ell} u\_\{\ensuremath{\ell},a\}v\_\{\ensuremath{\ell},b\}w\_\{\ensuremath{\ell},c\}=1 if j=j\ensuremath{^{\prime}}, k=k\ensuremath{^{\prime}} and i=i\ensuremath{^{\prime}}, and zero otherwise. Necessity follows by coefficient extraction in independent variables; sufficiency follows by expanding the 23 products and collecting every coefficient. This is the ordinary proof bridge between a stored factor array and matrix multiplication. It does not depend on a numerical residual tolerance.

To implement the multiplication algorithm, compute m\_\ensuremath{\ell}=U\_\ensuremath{\ell}(A)V\_\ensuremath{\ell}(B) using 23 nonscalar input products. Then (AB)\_\{ik\}=\ensuremath{\sum}\_\ensuremath{\ell} w\_\{\ensuremath{\ell},3k+i\}m\_\ensuremath{\ell}. The transposed output indexing is forced by the C\_\{ki\} trace leg. Reading that index as 3i+k without changing conventions would produce the wrong output map. The archived exact factor lists are therefore part of the result, rather than optional search data.

\subsection{The support-138 certificate and its arithmetic meaning}
The stored scheme has 47 nonzero u coefficients, 47 nonzero v coefficients and 44 nonzero w coefficients, for total support 138. The exact rational certificate evaluates each of the 729 sums above and finds its required value. An independent Fraction computation and the integer-scaled Lean certificate corroborate this finite identity; the current verification register records the later pinned-toolchain replay. Thus the scheme proves rank at most 23 for this tensor and the stated support count. It does not give a rank-22 algorithm or prove that 138 is the minimum support of all rank-23 algorithms.

The discovery route first found a support-capped valid pattern over GF(3), with seed 3158 in the preserved walk, and then lifted it numerically before snapping and solving the remaining linear tensor equations exactly over the rationals. The proof is the exact final identity; floating-point convergence and the search seed are provenance. The final coefficients are signed powers of two. Termwise rescaling u\_\ensuremath{\ell}\ensuremath{\to}a\_\ensuremath{\ell}u\_\ensuremath{\ell}, v\_\ensuremath{\ell}\ensuremath{\to}b\_\ensuremath{\ell}v\_\ensuremath{\ell}, w\_\ensuremath{\ell}\ensuremath{\to}w\_\ensuremath{\ell}/(a\_\ensuremath{\ell}b\_\ensuremath{\ell}), with nonzero scalars, leaves each rank-one tensor unchanged and preserves its support.

The recorded gauged representative has 126 coefficients equal to \ensuremath{\pm}1 and twelve equal to \ensuremath{\pm}2 or \ensuremath{\pm}1/2. They are coefficients in the same 23 products. Computing each linear form independently costs support\ensuremath{-}(23+23+9)=83 additions, because all 23 input forms in each leg and all nine output forms are nonempty. Common-subexpression sharing can change this count and the order of operations, while coefficient scaling can change floating-point growth. Hence a one-entry support improvement over 139 does not on its own prove a practical runtime or stability improvement.

\subsection{An exact family, rather than a numerical tangent conjecture}
The source reports a rational one-parameter family whose parameter is s=u[12][6]. All coefficients not listed as moving remain fixed at the stored support-138 values. Among the moving u entries are u[3][6]=\ensuremath{-}8(s+2), u[7][1]=u[10][4]=u[20][7]=s/[8(s+2)], u[7][2]=u[10][5]=u[20][8]=1/[4(s+2)], u[9][8]=u[11][6]=u[18][7]=\ensuremath{-}8s, and u[12][6]=s. These formulas show directly where the exceptional parameter values arise.

The remaining v and w entries in the supplied family file are rational functions with denominators powers of s and s+2. The exact proof substitutes all literal functions into each of the 729 tensor equations, clears denominators, and checks that every numerator polynomial is identically zero over Q[s]. This proves the tensor identity over Q(s), and hence at every real or rational specialization s\ensuremath{\neq}0,\ensuremath{-}2. Each prescribed nonzero entry remains nonzero at those specializations, so support stays 138. The specialization s=\ensuremath{-}1 recovers the stored rational scheme. A finite set of sampled numerical points would not alone establish this universal statement; the cleared polynomial identity is the necessary bridge.

The Jacobian calculation at the stored scheme gives a seven-dimensional tangent space, six reported diagonal gauge directions and one remaining direction. The explicit curve is a valid family irrespective of whether this local numerical dimension describes the entire algebraic component. It does not prove that all solutions with that support lie on the curve, that a quotient is globally one-dimensional, or that all projective boundary degenerations have been classified. Searches forcing one entry to zero and residuals tending to zero with growing coefficients are likewise evidence of the tested search, not exact support-137 impossibility proofs.

\subsection{A useful inequivalence invariant and the characteristic-two obstruction}
View each of the nine-coefficient factors as a 3\ensuremath{\times}3 matrix. Under a sandwich symmetry of the matrix multiplication tensor, a factor is multiplied on the left and right by invertible 3\ensuremath{\times}3 matrices, which preserves its rank. Nonzero term rescaling also preserves rank, and cyclic slot permutation or transpose preserves the unordered rank triple of each term. Consequently the multiset of unordered factor-rank triples is an invariant of this standard equivalence action.

For the 138 scheme the source's exact rank calculation gives fifteen terms of type (1,1,1), four of type (1,1,2), and four of type (2,2,2). For the compared Smirnov 139 representative it gives eleven, eight and four respectively. These multisets differ, so the two schemes are inequivalent under the stated sandwich, slot and transpose action. This is a precise structural distinction. A broader group or global classification claim would need its own theorem identifying the allowed transformations.

Replacing every nonzero of the 138 pattern by one over GF(2) gives four failed tensor equations. Any rational scheme with exactly that pattern and all nonzero coefficients \ensuremath{\pm}1 would reduce to this all-one pattern over GF(2), while the integer matrix multiplication identities would reduce correctly. Thus the failed pattern excludes that precise sign-only class. It does not reduce the actual dyadic scheme modulo two: denominators divisible by two are not invertible there. It also does not exclude other support patterns or coefficients over another field.

The main paper should lead with the exact scheme and its finite identity, then the rational family and scoped invariant, with empirical geometry and search observations explicitly separated. The independent symbolic and rational certificates make these algebraic results concrete. Priority still requires comparison of equivalent schemes under matching conventions and contributor approval for the new manuscript; the preserved author paper and factor data should remain available in the central source index.

\subsection{Verification and reproducibility}
Fresh pinned replay PASS: all seven archived custom Lean modules compile under Lean 4.34.1 (commit 5045d005), and all 12 selected endpoint axiom audits pass. Validity and support endpoints are axiom-free; tensor identities use propext and Quot.sound. No submitted compiled artifact was used. Full commands, source hashes, timing and output: Verification/chandragupt-pinned-build.json.

\begin{samepage}
\subsection{MM3.scheme138\_identities}
\begin{quote}\small\ttfamily theorem scheme138\_identities :\newline     \ensuremath{\forall} a b c : Nat, a < 9 \ensuremath{\to} b < 9 \ensuremath{\to} c < 9 \ensuremath{\to}\newline       coeff (terms Scheme138.u Scheme138.v Scheme138.w) a b c =\newline         Scheme138.su * Scheme138.sv * Scheme138.sw * target a b c\end{quote}
{\small Source: lean/MM3/Certificates.lean:22.\par}

{\small Source SHA-256: 78cd77c0514b515d10e28a7bed4505a20ad4c2124cc35346fee04c2b589c9bfa\par}

\end{samepage}
\begin{samepage}
\subsection{MM3.scheme138\_support}
\begin{quote}\small\ttfamily theorem scheme138\_support : support Scheme138.u Scheme138.v Scheme138.w = 138\end{quote}
{\small Source: lean/MM3/Certificates.lean:28.\par}

{\small Source SHA-256: 78cd77c0514b515d10e28a7bed4505a20ad4c2124cc35346fee04c2b589c9bfa\par}

\end{samepage}
\begin{samepage}
\subsection{MM3.scheme138Gauged\_identities}
\begin{quote}\small\ttfamily theorem scheme138Gauged\_identities :\newline     \ensuremath{\forall} a b c : Nat, a < 9 \ensuremath{\to} b < 9 \ensuremath{\to} c < 9 \ensuremath{\to}\newline       coeff (terms Scheme138Gauged.u Scheme138Gauged.v Scheme138Gauged.w) a b c =\newline         Scheme138Gauged.su * Scheme138Gauged.sv * Scheme138Gauged.sw * target a b c\end{quote}
{\small Source: lean/MM3/Certificates.lean:38.\par}

{\small Source SHA-256: 78cd77c0514b515d10e28a7bed4505a20ad4c2124cc35346fee04c2b589c9bfa\par}

\end{samepage}
\begin{samepage}
\subsection{MM3.scheme138Gauged\_support}
\begin{quote}\small\ttfamily theorem scheme138Gauged\_support :\newline     support Scheme138Gauged.u Scheme138Gauged.v Scheme138Gauged.w = 138\end{quote}
{\small Source: lean/MM3/Certificates.lean:44.\par}

{\small Source SHA-256: 78cd77c0514b515d10e28a7bed4505a20ad4c2124cc35346fee04c2b589c9bfa\par}

{\small\textbf{Sources and attribution:} \href{https://github.com/ChandraguptSharma07/matrix-multiplication-tensor-3x3/tree/479c2416b6261ee841052eef4881df0e40054f3f}{1. ChandraguptSharma07/matrix-multiplication-tensor-3x3}; \href{https://doi.org/10.1134/S0965542513120129}{2. DOI: 10.1134/S0965542513120129}; \href{https://arxiv.org/abs/1905.10192}{3. arXiv source: 1905.10192}; \href{https://arxiv.org/abs/2212.01175}{4. arXiv source: 2212.01175}; \href{https://github.com/alejandrozu/openmath-2026-judging/blob/d0ed487f487fa7ec814ff705ab55249983fe8707/OpenMath-Judging/review/entries/E04-SUPPORT138.txt}{5. alejandrozu/openmath-2026-judging / E04-SUPPORT 138.txt}; \href{https://github.com/alejandrozu/openmath-2026-judging}{Central source and adjudication archive}.\par}

\end{samepage}
\clearpage
\section{Clebsch structure and exact stochastic refinements for monochromatic K\_4}
{\small Provisional mathematical authors: Luke Van Seters, Madhan Jothimani\par}\medskip
{\small Luke Van Seters  -  Klaviyo (provisional; see affiliation register).\par}

{\small Madhan Jothimani  -  Northeastern University (provisional; see affiliation register).\par}

{\small Affiliations follow the recorded author declarations or public institutional/profile sources. Preferred author names, author order and publication affiliations await author review. This editorial compilation does not establish anthology coauthorship.\par}

\subsection{Abstract}
A structural decomposition of the published 768-vertex PPSS construction leads to a 192-class stochastic family and a 3840-class refinement with exact upper bound 0.030138887566497221.... The latest packet supplies finite-realization and counting proofs, with native evaluation explicitly in the arithmetic trust boundary. Event admission and exact comparison with the earlier Technion construction remain separate questions.

For every n\ensuremath{\geq}4 there exists a red/blue colouring of K\_n with monochromatic K\_4 proportion at most 8450462766487926638466333426306607129/280384030360880691940646801777885184000; hence the corresponding upper-limit bound holds.

\subsection{Two exact construction theorems}
The compact 192-class construction at p=32/41 and h=22/41 has exact density 1013294255057839/33620705806123008\ensuremath{\approx}0.030138994133588287. The refined 3840-class construction has exact density 8450462766487926638466333426306607129/280384030360880691940646801777885184000\ensuremath{\approx}0.030138887566497220. Both are below the reported McKay numerical value and the rounded 0.030139 threshold in the Technion announcement.

The scientific proposal is broader than a single decimal: the public PPSS graph is decomposed into twelve 64-vertex families, each built from sixteen four-vertex fibres over the Clebsch graph, with inter-family rules on \ensuremath{\mathbb{Z}}3\ensuremath{\times}\ensuremath{\mathbb{Z}}2\ensuremath{\times}\ensuremath{\mathbb{Z}}2. The stochastic 192-class family makes two formerly fractional patterns adjustable. A pentagon refinement and two balanced sign splits give the 3840-class table. These are latest-research results at the frozen October 3 revision; the earlier competition leaderboard packet must not silently acquire later results.

\subsection{A direct finite-realization proof}
For any finite symmetric probability table W, choose a class independently and uniformly for every vertex of K\_n. Conditioned on the classes, colour each edge independently using the corresponding table entry as its red probability. For four distinct vertices, the class labels are independent, allowing repetitions, and the six colour trials are conditionally independent.

The expected normalized ordered injective-four-tuple density is the ordered class average of the six red probabilities' product plus the blue complements' product. Every monochromatic unordered four-set has exactly 24 orderings, and the total ordered injective tuple count is 24 times binomial(n,4); therefore this is also the usual unordered-four-set proportion. That equivalence is an ordinary 24-ordering observation, not a separately selected Lean bijection theorem. Some outcome cannot exceed the expectation, proving the actual finite coloring bound at every n\ensuremath{\geq}4. Repeated class labels are legitimate because graph vertices remain distinct. The same uniform finite bound controls the upper limit of extremal minima, and any ordinary limit if supplied separately.

\subsection{Centered refinement and exact counting}
The embedded 3840 witness stores 1248 fractional 20\ensuremath{\times}20 blocks and derives the remaining blocks by the Clebsch rule. Zero row means in the centered blocks allow a substantial algebraic reduction of the six-edge expansion. A graph term with a leaf vanishes after summing that vertex; only the empty/base term and the surviving subgraphs without degree-one vertices remain. For K4 those include its four triangles, three four-cycles, six diamonds and the full perturbation K4.

Eleven general realization/counting prototype theorems have standard-kernel axiom scope. The fixed 3840-class sparse-count correctness, raw\_total and density instantiations additionally inherit native finite-data facts, including hard\_blocks. Native evaluation checks the literal integer expression; kernel rational arithmetic then converts 519196432373018212667371525712277942005760 over 65536\textasciicircum{}6\ensuremath{\cdot}3840\textasciicircum{}4 to the stated fraction. The complete instantiated endpoints therefore trust the native compiler/evaluator as well as the kernel; calling their last rational step kernel-checked does not remove that inherited trust.

\subsection{Priority, theorem scope and release}
A value below an announced rounded threshold does not establish that it improves the announcement's unknown exact construction. Both statements can be below 0.030139, with either one smaller. Obtain the earlier exact value and conventions before a chronological record claim. The PPSS structural decomposition and a transparent new family may have independent mathematical interest even if the numerical result is surpassed.

The official Random Structures \& Algorithms 2025 program already advertised Noam Feinstein's Improved Constructions for Ramsey Multiplicities, jointly with Chaim Even-Zohar. Its abstract describes random constructions improving c\_4 and c\_5 and emphasizes their symmetry and structure, but supplies no exact density, manuscript or construction certificate. Therefore the comparison needed here predates the January 21, 2026 Technion talk, and the unresolved question concerns both the exact numerical bound and possible overlap of construction ideas. Program: \url{https://www.dmg.tuwien.ac.at/rsa2025/}.

The proposed p=0.05 recommendation for Luke and Madhan's structural contribution remains conditional on the chairs' assessment of independent contribution after comparing the earlier construction. Neither the rounded threshold nor this source search establishes independent priority.

The paper's ordinary algebraic-minimum analysis, structural relabeling checks and exact rational certificates have different verification boundaries. Global minimum and degree-seventeen claims should not be labelled native-certified merely because a rational point is certified. The source does not prove convergence of the finite minimum sequence in its selected asymptotic module; its upper-limit and every-n endpoints do not require convergence.

For this collection, preserve the authors' own detailed manuscript and figures as an attributed source supplement and use this chapter as the cross-portfolio explanation. Confirm author approval, affiliations, the event-production chronology and final selected closure before release. Contest points for a frozen earlier result and scientific significance of this later packet are separate decisions.

\subsection{Independent exact integer recount}
An independent reviewer-owned integer calculation now matches the complete 3840-class raw numerator 519196432373018212667371525712277942005760 and the reduced density stated above. It reads the frozen literal blocks without executing the entrant's program. Zero row/column sums, the base transitivity and sparse multiplicities are checked. Permuting the four coarse vertices reduces 6912 ordered supported tetrahedra to 288 unordered ones. Bounded signed-integer limb contractions and ordinary arbitrary-size integer recombination avoid floating-point arithmetic and overflow.

The source-data run completed in 14.437 seconds with peak resident memory 122916864 bytes. The portable checker, pickle-free inputs, per-source hashes and full contribution receipt are supplied in Verification/Independent checkers. This confirms the arithmetic independently; it does not itself close the separate unchanged Lean replay, authenticate the discovery date or settle comparison with the earlier Technion construction.

\subsection{Actual fresh verification of the novel construction scope}
The complete source-identical fresh replay finished on 5 October 2026 at 02:18:36 UTC, and an independent root review qualified its exact scope. All 51 frozen source modules passed, retaining 33 receipt-matched outputs from this reviewer's earlier cold-source run during the native continuation. All 35 requested unique axiom prints were captured: eleven general statements use standard kernel axioms, twenty-three concrete statements inherit the expected native-evaluation axioms, and the remaining concrete ValidTemplate.countSubgraphs\_numerator theorem also uses standard kernel axioms. No selected print contains an admitted or unknown axiom. The exact 192/3840 density, literal-count correspondence, every-n coloring and unconditional upper-limit endpoints are verified at their displayed scopes; ordinary-limit statements retain their convergence hypotheses. The earlier 31-output milestone is preserved as dated evidence. Receipt: Verification/luke-k4-ramsey-current-native-fresh-build.json.

Concrete finite computation retains native compiler/evaluator trust. The source-identical native image and initialization review were accepted with a Windows image-observer limitation (WinError 299): continuous observation of every native image is not claimed, and the error timing remains unresolved in the receipt. The independent exact-integer 3840 recount gives separate arithmetic evidence. Full formal completion does not establish a global numerical record, independent priority over Feinstein/Even-Zohar's earlier construction, discovery chronology or publication author approval.

\subsection{The exact bound and its finite-coloring consequence}
For a symmetric N\ensuremath{\times}N probability table W=(w\_\{ij\}) with entries in [0,1], define F(W)=N\textasciicircum{}\{-4\}\ensuremath{\sum}\_\{i,j,k,l\}[w\_\{ij\}w\_\{ik\}w\_\{il\}w\_\{jk\}w\_\{jl\}w\_\{kl\}+(1\ensuremath{-}w\_\{ij\})(1\ensuremath{-}w\_\{ik\})(1\ensuremath{-}w\_\{il\})(1\ensuremath{-}w\_\{jk\})(1\ensuremath{-}w\_\{jl\})(1\ensuremath{-}w\_\{kl\})]. The sum includes repeated class indices. The diagonal entry describes an edge between distinct vertices that happen to receive the same class; it is therefore part of the probability table, not an omitted diagonal in the density calculation.

Give each vertex of K\_n an independent uniform class and color each edge independently, conditional on the classes, with its table probability. Any four distinct vertices have independent class labels and conditionally independent six edges, so their monochromatic probability is exactly F(W). Linearity of expectation bounds the normalized ordered injective-four-tuple density for every n\ensuremath{\geq}4. Each unordered four-set has 24 orderings, so dividing both its monochromatic count and total count by 24 gives the usual four-set proportion. This equivalence is an ordinary observation, not a separately selected formal bijection theorem. Averaging an optimal coloring of K\_\{n+1\} over its n-vertex subsets shows the minimum proportion m\_4(n) is nondecreasing; boundedness therefore gives its limit c\_4\ensuremath{\leq}F(W). This monotonicity and limit-existence argument is ordinary exposition. The selected Lean Asymptotic module proves the unconditional upper-limit bound; its ordinary-limit statement assumes convergence and does not itself prove this averaging or existence lemma.

The final explicit table has N=3840 and denominator Q=65536. Its exact density is U=8450462766487926638466333426306607129/280384030360880691940646801777885184000\ensuremath{\approx}0.030138887566497220. Thus m\_4(n)\ensuremath{\leq}U for every n\ensuremath{\geq}4 and c\_4\ensuremath{\leq}U. This is below the publicly announced threshold 0.030139, but an announcement without its exact value does not determine whether U improves that earlier construction. The structural contribution and this unresolved numerical comparison must remain separate.

\subsection{Twelve families and the 192-class Clebsch rule}
Let V=F\_2\textasciicircum{}4 and S=\{e\_1,e\_2,e\_3,e\_4,e\_1+e\_2+e\_3+e\_4\}; put C=\{0\}\ensuremath{\cup}S. The Clebsch graph connects x,y when x+y\ensuremath{\in}S. Let H=Z\_3\ensuremath{\times}Z\_2\ensuremath{\times}Z\_2 index twelve families u=(a,e,s). For d=u\ensuremath{-}v define the family type as Z for identical families; H\_type when d\_s=d\_a=0,d\_e=1; X when d\_s=d\_e=0,d\_a\ensuremath{\neq}0; P when d\_s=1,d\_a=0; and Z in all remaining cases. This rule is symmetric because the defining conditions are invariant under changing the sign of d.

On the 192 classes H\ensuremath{\times}V define B\_\{p,h\}((u,x),(v,y)) by z=x+y. Type Z has red probability zero on C and one off C. Type X is its complement. Type P has probability p on C and zero off C. Type H\_type has probability h at z=0, one on S and zero off C. This completely specifies a symmetric probability table for p,h\ensuremath{\in}[0,1]. Every row has 87 entries equal to one, twelve equal to p, one equal to h and 92 equal to zero.

The submitted labeling identifies the 768 vertices of the PPSS graph with H\ensuremath{\times}V\ensuremath{\times}\{1,2,3,4\}. Within a family the blue graph is the Clebsch graph with each vertex replaced by a blue four-clique. Between different families the red proportions across the sixteen pairs of each four-by-four fiber block follow these same four types, except that the P same-position block may be one or three quarters, and the H same-position block is one half. The structure theorem is a finite entry-by-entry comparison with the public original matrix and an explicit bijective labeling. Averaging the fiber blocks is a new table; it is not merely a relabeling and need not preserve K\_4 density.

\subsection{A polynomial density and a compact rational theorem}
Translation invariance on H\ensuremath{\times}V lets one fix the first of the four class indices, reducing the density sum from 192\textasciicircum{}4 to 192\textasciicircum{}3 terms. Each term is a polynomial of degree at most six in p,h. The source gives F(B\_\{p,h\})=P(p,h)/7077888, where P=251655\ensuremath{-}97344p+66888p\ensuremath{^2}\ensuremath{-}7008p\ensuremath{^3}+1224p\ensuremath{^4}\ensuremath{-}6552h+9216ph\ensuremath{-}4176p\ensuremath{^2}h+2304p\ensuremath{^3}h\ensuremath{-}144p\ensuremath{^4}h+990h\ensuremath{^2}\ensuremath{-}576ph\ensuremath{^2}+648p\ensuremath{^2}h\ensuremath{^2}\ensuremath{-}288p\ensuremath{^3}h\ensuremath{^2}+72p\ensuremath{^4}h\ensuremath{^2}\ensuremath{-}16h\ensuremath{^3}+3h\ensuremath{^4}.

The polynomial certificate can be obtained by literal symbolic expansion or exact interpolation on \{0,1,...,6\}\ensuremath{^2}: degree at most six in each variable guarantees uniqueness on that grid, and the arithmetic evaluation remains meaningful even when a grid point is outside the probability range. Substitution at p=32/41,h=22/41 gives F=1013294255057839/33620705806123008\ensuremath{\approx}0.030138994133588287. Together with the probabilistic realization lemma this proves a compact explicit upper bound. The Lean development separately evaluates the literal table density, so its finite rational endpoint does not depend on trusting the interpolation routine.

The author manuscript also isolates the unique interior minimum of P by a rational Gröbner basis, an irreducible degree-seventeen elimination polynomial, a Sturm count, isolating intervals and comparison with all four boundary edges. Its approximate minimum is 0.030138977289665338 at p\ensuremath{\approx}0.7791808330,h\ensuremath{\approx}0.5342651880. These algebraic optimization claims have their own computer-algebra certificates. The rational construction theorem is already sufficient for the compact upper bound, regardless of whether one uses the stronger minimum classification.

\subsection{Pentagon and balanced-sign refinements}
Split each base class into five subclasses indexed by Z\_5. Assign each fractional pair a phase \ensuremath{\gamma}\_\{ij\}, with \ensuremath{\gamma}\_\{ji\}=\ensuremath{-}\ensuremath{\gamma}\_\{ij\}; declare (a,b) on-pattern when a\ensuremath{-}b\ensuremath{-}\ensuremath{\gamma}\_\{ij\}=\ensuremath{\pm}1 mod 5. A P pair is either inactive with constant p\_0/Q, or active with x/Q on-pattern and one otherwise. An H pair has zero on-pattern and y/Q otherwise. Constant zero/one blocks stay constant. The preserved phase data contain 272 inactive and 880 active P pairs and 96 H pairs. This is an explicit finite table, rather than an unproved closed-form phase rule.

For Q=9,p\_0=7,x=4,y=8, every active P block averages 7/9 and every H block averages 8/15. The exact density certificate gives 4198776398959/139314069504000\ensuremath{\approx}0.030138925766. Further rational tuning yields the 960-class parent W\_0 with density about 0.030138903565. Keeping block means alone does not fix the K\_4 density, because its six probabilities are multiplied before averaging.

For any symmetric amplitude A with |A\_\{ij\}|\ensuremath{\leq}min(w\_\{ij\},1\ensuremath{-}w\_\{ij\}), split each class into two signs and set W\ensuremath{^{\prime}}((i,s),(j,t))=w\_\{ij\}+st A\_\{ij\}. The range condition keeps this in [0,1] and each block retains mean w\_\{ij\}. Expand the six edge factors for four independently chosen signs. A term survives sign averaging exactly when every vertex has even degree in the selected edge subset of K\_4. The nonempty possibilities are its four triangles and three four-cycles. Blue replaces A by \ensuremath{-}A, so its triangle terms change sign and its cycle terms do not.

Consequently F(W\ensuremath{^{\prime}})\ensuremath{-}F(W)=T\_3(A)+T\_4(A), with T\_3=4N\textasciicircum{}\{-4\}\ensuremath{\sum}\_\{i,j,k,l\}A\_\{ij\}A\_\{jk\}A\_\{ki\}(w\_\{il\}w\_\{jl\}w\_\{kl\}\ensuremath{-}q\_\{il\}q\_\{jl\}q\_\{kl\}) and T\_4=3N\textasciicircum{}\{-4\}\ensuremath{\sum}\_\{i,j,k,l\}A\_\{ij\}A\_\{jk\}A\_\{kl\}A\_\{li\}(w\_\{ik\}w\_\{jl\}+q\_\{ik\}q\_\{jl\}), where q=1\ensuremath{-}w. This identity includes repeated coarse indices. Scaling A by \ensuremath{\varepsilon} produces cubic and quartic changes, with no linear or quadratic change. Thus a refinement improvement can be invisible to the Hessian in the unrefined parameter family.

Two successive sign refinements give 1920 and then 3840 classes. Equivalently each of the 192 coarse classes has twenty subclasses: five pentagon labels and two independent binary sign labels. The source records all 1248 fractional 20\ensuremath{\times}20 blocks, with integer entries in [0,65536], symmetry by transposition, and row/column sums equal to twenty times the coarse numerator. The coarse means are those of B\_\{51064/65536,35139/65536\}. These checks establish a valid exact witness before any density value is evaluated.

\subsection{The exact character-count proof and reproducibility boundary}
For an integer table r=QW put D(r)=\ensuremath{\sum}\_\{i,j,k,l\}\ensuremath{\prod}\_\{s<t\}r\_\{i\_s i\_t\}. Then F(W)=[D(r)+D(Q\ensuremath{-}r)]/(N\textasciicircum{}4Q\textasciicircum{}6). The balanced signs give a free action of G=F\_2\ensuremath{^2} with 960 orbits. More generally, if |G|=K and r((u,x),(v,y))=f\_\{uv\}(x+y), define Fourier blocks \ensuremath{\hat{f}}\_c(u,v)=\ensuremath{\sum}\_\{z\ensuremath{\in}G\}f\_\{uv\}(z)(\ensuremath{-}1)\textasciicircum{}\{c\ensuremath{\cdot}z\}. Fourier inversion gives r=K\textasciicircum{}\{-1\}\ensuremath{\sum}\_c \ensuremath{\hat{f}}\_c \ensuremath{\chi}\_c(x)\ensuremath{\chi}\_c(y).

Substitute this inversion into all six factors of D. It contributes K\textasciicircum{}\{-6\}. Summing the four sign variables gives K\textasciicircum{}4 precisely when the labels incident to each vertex sum to zero, and zero otherwise. Thus D=K\textasciicircum{}\{-2\}\ensuremath{\sum}\_\{c\ensuremath{\in}C\_bal\}\ensuremath{\sum}\_\{u\_1,...,u\_4\}\ensuremath{\prod}\_\{s<t\}\ensuremath{\hat{f}}\_\{c\_\{st\}\}(u\_s,u\_t). The balanced labels are exactly (a+b,a+d,b+d,a,b,d), with a,b,d\ensuremath{\in}G. For K=4 there are 64 such labels, reduced by vertex-permutation symmetry to eleven representatives. This derivation is the ordinary proof behind the accelerated independent counter.

Each contraction is computed modulo several primes and reconstructed uniquely by the Chinese remainder theorem after the modulus product exceeds twice an explicit absolute bound. The code's floating-point matrix contractions are exact integer dot products in the stated range, since it checks n(p\ensuremath{-}1)\ensuremath{^2}<2\textasciicircum{}53. Shape, range, symmetry, sign invariance and Fourier reconstruction are checked on every input entry before that reduction is used.

The final certificate gives D(r)=259147196226119060009588075500947710392320 and D(Q\ensuremath{-}r)=260049236146899152657783450211330231613440, with denominator N\textasciicircum{}4Q\textasciicircum{}6=17226794825372509712833339501233265704960000. Their sum reduces to U. The eleven general realization/counting prototypes have standard-kernel axiom scope. The fixed 3840-class CountCorrect, raw\_total and density instantiations inherit native finite-data facts such as hard\_blocks, so their complete proof trust includes compiler/runtime evaluation. Their final kernel rational arithmetic does not remove that inherited native trust. The preserved character-counter certificate uses a different reduction and records the same exact totals; the current replay register separately identifies independently rerun computations. The current replay register determines what has been freshly rebuilt; an author receipt or source scan alone is not substituted for it.

This account is adapted with attribution from Luke Van Seters and Madhan Jothimani's preserved paper at revision 2cf216ba885fbb1b2d1b2d4920efeaa6ba438842. The original paper, witness files and certificate should accompany the new short synthesis. The source's production dates, Oct 3 numerical Lean additions and the organizer's late-intake rules belong in a separately dated event provenance record. None changes the mathematical content of the exact witness or establishes comparison with an unpublished earlier value.

\subsection{Verification and reproducibility}
Fresh execution: luke-k4-ramsey-current: PASS\_WITH\_NATIVE\_EVALUATION; 51/51 custom modules compiled successfully; receipt Verification/luke-k4-ramsey-current-native-fresh-build.json; isolated source-identical native adapter; compiler/native-evaluation trust is explicit; the earlier direct attempt is preserved separately; positive native identity/initialization evidence has a separate root review, with one WinError 299 retained, unknown error timing and incomplete continuous memory observation. Compiler pins, source hashes and endpoint axiom outputs are in Verification. Historical receipts and finite checks do not replace fresh replay.

{\small Independent exact integer recounts confirm both the 192-class density 1013294255057839/33620705806123008 and the complete 3840-class raw numerator and fraction above. Receipts: Verification/Independent checkers/luke\_clebsch192\_independent\_recount.json and luke3840\_independent\_exact\_recount.json. The counts do not themselves certify Lean execution or historical priority.\par}

\begin{samepage}
\subsection{K4Ramsey.Clebsch192.exists\_coloring}
\begin{quote}\small\ttfamily For every n\ensuremath{\geq}4 an actual colouring has density\ensuremath{\leq}1013294255057839/33620705806123008.\end{quote}
{\small Source: lean/K4Ramsey/Constructions/Clebsch192/Bound.lean.\par}

\end{samepage}
\begin{samepage}
\subsection{K4Ramsey.Final3840.exists\_coloring\_exact}
\begin{quote}\small\ttfamily For every n\ensuremath{\geq}4 an actual colouring has density\ensuremath{\leq}8450462766487926638466333426306607129/280384030360880691940646801777885184000.\end{quote}
{\small Source: lean/K4Ramsey/Constructions/Final3840/NumericalBound.lean.\par}

{\small\textbf{Sources and attribution:} \href{https://github.com/lukevs/k4-ramsey/tree/2cf216ba885fbb1b2d1b2d4920efeaa6ba438842}{1. lukevs/k4-ramsey}; \href{https://github.com/lukevs/k4-ramsey/blob/2cf216ba885fbb1b2d1b2d4920efeaa6ba438842/paper/main.tex}{2. lukevs/k4-ramsey / main.tex}; \href{https://github.com/lukevs/k4-ramsey/blob/2cf216ba885fbb1b2d1b2d4920efeaa6ba438842/docs/verification.md}{3. lukevs/k4-ramsey / verification.md}; \href{https://math.technion.ac.il/events/noam-feinstein/}{4. math.technion.ac.il / noam-feinstein}; \href{https://www.dmg.tuwien.ac.at/rsa2025/}{5. www.dmg.tuwien.ac.at / rsa 2025}; \href{https://github.com/alejandrozu/openmath-2026-judging}{Central source and adjudication archive}.\par}

\end{samepage}
\clearpage
\section{A generic matrix pivot with a vanishing eigenphase Walsh angle}
{\small Provisional mathematical authors: Matt Bowring\par}\medskip
{\small Matt Bowring  -  Purdue University, Department of Mechanical Engineering (documented January 2026; current preferred publication affiliation to confirm).\par}

{\small Affiliations follow the recorded author declarations or public institutional/profile sources. Preferred author names, author order and publication affiliations await author review. This editorial compilation does not establish anthology coauthorship.\par}

\subsection{Abstract}
For arbitrary finite complex matrix dimension D=4(q+1), invertible top blocks X and Y with pairwise distinct eigenvalue crossing angles admit a one-qubit pivot whose polar-factor multiplexor has an eigenphase labeling with an exactly vanishing balanced Walsh angle. The construction combines a discrete half-plane rotation with continuous normalized eigenphase lifting and an intermediate value theorem. It is a genuine restricted matrix-existence theorem and a proposed partial advance toward improved unitary synthesis. The complete circuit bridge, nongeneric nodes and universal improved CNOT count remain unproved in the inspected formal chain; fresh pinned replay verifies the selected generic theorem using only standard kernel axioms; novelty and author review remain pending.

For D=4(q+1), invertible complex D-by-D blocks X,Y, a full eigenvalue list of X inverse times Y and pairwise crossing-angle separation modulo pi, there exist an unrestricted real angle alpha and an angle beta in [0,pi], an ordinary unitary one-qubit gate, pivoted blocks satisfying the lifted block multiplication identity, and a real phase list of zxzC(Xprime,Yprime) whose finalAngle(q+1) equals zero. No unitarity assumption on the entire block matrix U is needed. The theorem concludes no circuit or CNOT-count bound.

\subsection{Attribution, source and verification status}
Contributor: Matt Bowring, sole human author declared in the supplied correspondence. Provisional publication affiliation: Purdue University, Department of Mechanical Engineering (documented Jan 2026; current preferred publication affiliation to confirm). AI assistance generated the formal code under his direction; human mathematical responsibility and the distinction between a proof and its generated presentation remain explicit.

Source snapshot: bowrango/lean-unitary, commit 7014b63ee2cfd01c641a7f539c73b664a2aa3d66, recorded 4 October 2026 at 19:22 UTC. A fresh exact replay finished 5 October 2026 at 01:38:41 UTC under Lean 4.30.0-rc2 and Mathlib c1e30e172c8fda21e6776bf1f10351e882ee31b9. All 26 authored sources compiled. The six selected generic/count/conditional-bridge theorem prints use only propext, Classical.choice and Quot.sound; count\_closed uses only propext and Quot.sound. No selected print contains sorryAx, native-evaluation or unknown axioms. A separate baseline audit prints sorryAx for all eight Claims endpoints, including Claims.synth; the baseline source has 13 admitted placeholders. Compilation success therefore verifies the selected scoped theorem and does not verify the universal synthesis headline. Receipt: Verification/matt-unitary-current-fresh-build.json.

\subsection{The actual matrix theorem}
The actual formal result is a finite-dimensional matrix theorem. Let D=4(q+1). Write a block matrix U=[X Y;Z W] with complex D-by-D blocks, and suppose X and Y are invertible. Let u\_1,...,u\_D list the eigenvalues of X inverse times Y, counted with their algebraic multiplicities. Assume that no two different entries in this list have arguments differing by an integer multiple of pi: geometrically, no two lie on the same line through the origin. This is the explicitly restricted generic regime. The theorem is stated for an arbitrary finite index type of cardinality D, so it is an infinite family of matrix dimensions rather than a test of one numerical instance.

The pivot is the ordinary one-qubit matrix g=Rz(alpha)Ry(beta). Multiplying U by the lifted g produces top blocks Xg=e\textasciicircum{}\{-i alpha/2\}cos(beta/2)X+e\textasciicircum{}\{i alpha/2\}sin(beta/2)Y and Yg=-e\textasciicircum{}\{-i alpha/2\}sin(beta/2)X+e\textasciicircum{}\{i alpha/2\}cos(beta/2)Y. The formal development proves this block multiplication identity and the unitarity of g with ordinary complex matrices. It does not replace the matrix by an abstract object satisfying the desired conclusion.

The associated multiplexor is defined by its actual polar-factor formula
C(beta)=-i Xg* (Xg Xg*)\textasciicircum{}\{-1/2\}(Yg Yg*)\textasciicircum{}\{-1/2\}Yg,
where * denotes conjugate transpose and inverse square roots use the continuous functional calculus. For invertible Xg and Yg the library proves C(beta) is unitary. It proves the endpoint symmetry C(pi)=C(0)*, hence C(pi)=C(0)\textasciicircum{}\{-1\}, and continuity of C along the path under the no-real-crossing condition. Those identities explain why eigenphases can be used to interpolate between opposite endpoint imbalances.

\subsection{A half-plane rotation avoids singularities}
The first angle alpha is chosen by a discrete intermediate-value argument. As a rotation sweeps the nonzero eigenvalues of X inverse times Y, their upper-half-plane count changes one at a time because the crossing angles are distinct. Opposite ends of a half-turn have complementary counts. Since D is even, some intervening interval has exactly D/2 values in each half-plane and no eigenvalue on the real axis. This simultaneously prevents singularities in Xg,Yg and balances the endpoint determinant phase. The exact hypotheses matter: coincident crossings are outside this theorem and are not removed by calling the input generic without explanation.

\subsection{Continuous normalized eigenphases and a vanishing Walsh angle}
The second angle beta is found by continuous eigenphase lifting. The development constructs normalized ordered lists of real phases representing a unitary matrix's eigenvalues. A normalized list is monotone and spans at most 2pi; its total sum is tied to a continuous argument sigma(beta) of det C(beta). Existence follows by choosing eigenvalue arguments and shifting branches by integer multiples of 2pi. Uniqueness at a fixed sum, stability under small changes, continuity of characteristic-polynomial root lists up to permutation, and continuity of the determinant phase establish continuity of the selected normalized list. This addresses the genuine eigenvalue-collision/branch problem that an informal appeal to continuity of individual principal arguments would miss.

At beta=pi the normalized list equals the negative of its initial list in reversed order. A middle-half imbalance therefore changes sign. The real intermediate value theorem supplies beta* in [0,pi] with zero imbalance. A combinatorial permutation places the equal-size selected subset in the positive half of a Walsh row, and integer branch shifts preserve the complex eigenvalues. Consequently the final angle phi=(1/D)sum\_j h\_j theta\_j, with h\_j=+1 on the first D/2 indices and -1 on the second D/2, is exactly zero. The endpoint LeanUnitary.Pivot.pivot\_merge concludes existence of an unrestricted real alpha and a beta in [0,pi], the pivot, the block identity, the phase list of zxzC(Xg,Yg), and that vanishing Walsh angle.

\subsection{What the theorem accomplishes}
This is a meaningful structural partial result in the proposed synthesis programme. Its theorem is stronger than a numerical fit, a local Jacobian calculation or parameter counting: within the stated regime it proves an actual pivot exists in every admitted dimension. Following the successful exact replay, and subject to comparison with previous matrix/pivot methods, it can support a short research note devoted to generic phase balancing. Useful supporting results include continuous normalized eigenphase lifts and the explicit non-singularity criterion for the polar path. The submitted endpoint's name includes the word merge, but the formal conclusion itself does not contain a circuit, a CNOT gate, or a gate-count inequality.

\subsection{Relation to established Block-ZXZ and Walsh methods}
The polar factor C=\ensuremath{-}iP(X)*P(Y) is inherited from Block-ZXZ synthesis (De Vos-De Baerdemacker, arXiv:1512.07240v3, section 5, equation 5; Krol-Al-Ars, arXiv:2403.13692v2, section 4.1), and the Walsh/Gray-code angle transform is inherited from Möttönen et al. (quant-ph/0404089v3, equations 3-5). Conditional originality credit concerns only the combined generic pivot-existence theorem proved here, subject to comparison with prior phase-balancing methods.

The closest inspected primary papers do not give the combined balanced generic pivot, normalized continuous eigenphase lift and zero-Walsh existence argument. This limited comparison supports considering that lemma as an original candidate; it does not certify global historical priority or complete the missing circuit construction.

\subsection{Dated affiliation and pre-event research history}
Provisional publication affiliation: Purdue University, Department of Mechanical Engineering (documented Jan 2026; current preferred publication affiliation to confirm). The author's own website identifies him as a quantum software engineer and graduate student. His coauthored arXiv paper An Electronic Ising Machine, version 2 dated 9 January 2026, explicitly gives the Department of Mechanical Engineering at Purdue University for Matt Bowring. This is dated documentary evidence; author approval of the current preferred affiliation remains pending.

Bowring's personal post Thinking about optimal unitary synthesis, dated 12 September 2024, describes an earlier programme on sparse uniformly controlled rotations and approximate error bounds within Quantum Shannon Decomposition, motivated by his unitaryGate work. He explicitly describes omitted work as incomplete. The inspected post does not state or prove the current balanced polar-pivot existence theorem or its normalized continuous eigenphase and zero-Walsh argument. It establishes pre-event research history, rather than priority for this exact completed result. The conditional P2=.05 recommendation remains confined to the combined generic theorem; the date of an earlier research programme alone does not change that recommendation.

\subsection{The proposed global circuit consequence remains conditional}
The universal synthesis claim remains separate. The proposed recurrence is c\_2=3 and c\_\{n+1\}=4c\_n+3*2\textasciicircum{}n-6, giving c\_n=(21*4\textasciicircum{}n-72*2\textasciicircum{}n+96)/48. Compared with the established 22/48 Block-ZXZ recurrence, this would save (4\textasciicircum{}\{n-2\}-1)/3 CNOTs. The Count module proves these arithmetic identities. The Synthesis module proves the global bound conditional on both Synthesizable 2 3 and the full improved global node rule. That node rule is the central mathematical circuit-construction obligation, not an innocuous baseline assumption. There is currently no formal theorem deriving it from pivot\_merge; non-generic nodes, the Gray-code circuit realization of phi, and the exact sixth-CNOT cancellation remain outside the endpoint. Claims.synth still returns the older 22/48 statement using a baseline BlockZXZ theorem whose proof contains sorry.

A manuscript should therefore state the new generic matrix theorem as its principal result and explain the proposed circuit consequence as a remaining programme. It should avoid the current README's unconditional headline claiming machine-checked synthesis of arbitrary n-qubit unitaries with fewer CX gates. Baseline decomposition assumptions and admitted placeholders must be visible at the point they affect a claimed result. This does not negate the independent pivot theorem; it limits which conclusion the completed formal chain can support.

\subsection{Event chronology and the scope of preliminary credit}
Timing is also explicit. The pre-cutoff source was 57d6c96992a58f43d48d61c2a44353d902813cda, recorded 3 October at 02:00:25 UTC. The author's on-time packet called the claims not yet rigorous and proposed continuing formalization through the weekend. The current continuous-lift, determinant-phase and assembled pivot proof modules were added in 4be3c9e9f8d554320e8775713ac8019303e5330c on 4 October at 19:06:25 UTC. The new email arrived at 19:43:46 UTC. Intake can be accepted under the chair's lenient-delivery ruling while the post-cutoff proof date remains in the scientific and judging record. Any competition credit for this newly completed partial proof requires a recorded scope/timing ruling consistent with the treatment of other entrants.

For a preliminary mathematical assessment on the existing OPENMATH-UNITARY family, the narrow generic matrix result is best placed at the lower P2 boundary p=.05: it is a nontrivial lemma, but the circuit bridge, improved recursive node theorem and non-generic cases remain central obstacles. With approved D=824, modality M3A multiplier m=1 and no bonus b=1, F=.05*824\textasciicircum{}2/1000=33.9488. Actual accepted score remains 0 until novelty and fidelity are reviewed and the chair accepts that precise scope; the pinned endpoint has now successfully reproduced. A P3=.15 structural-advance assessment would give 101.8464, but would require a separate written rationale that this theorem itself removes a recognized synthesis obstacle; it should not follow automatically from an eventual global recurrence. Supporting lemmas belong to the same family, and the unproved universal 21/48 theorem retains 0. The user has authorized intake of newly delivered work: accept this packet for current review while plainly recording that the inspected formal proof was completed 4 October.

\subsection{Ordinary proof: a balanced rotation and a nonsingular polar path}
Fix p=q+1\ensuremath{\geq}1 and D=4p. Assume X and Y are invertible complex D-by-D matrices and u\_0,...,u\_\{D\ensuremath{-}1\} list, with multiplicity, the eigenvalues of X\ensuremath{^{-1}}Y. Assume arg(u\_j)\ensuremath{-}arg(u\_k) is never an integer multiple of \ensuremath{\pi} when j\ensuremath{\neq}k. Invertibility makes every u\_j nonzero. This separation hypothesis is stronger than simple spectrum and is part of the theorem; it is not replaced by an unstated genericity assumption.

Choose an initial rotation angle at which none of e\textasciicircum{}\{i\ensuremath{\alpha}\}u\_j is real. In any interval of rotation angles of length \ensuremath{\pi}, each eigenvalue crosses the real axis once. Separation modulo \ensuremath{\pi} ensures those D crossings occur at distinct angles. On the open intervals between successive crossings, the number N\ensuremath{_{+}} of eigenvalues in the upper half-plane changes by exactly one at each step. After the half-turn it equals D minus its initial value. The two endpoint values bracket D/2, and a sequence of integers changing by at most one cannot pass that integer without taking it. Choose \ensuremath{\alpha} in one of the intervals with N\ensuremath{_{+}}=D/2. Then \ensuremath{\tau}\_j=e\textasciicircum{}\{i\ensuremath{\alpha}\}u\_j are all nonreal and have equally many positive and negative imaginary parts. This is the source's discrete intermediate-value argument. It yields an unrestricted real \ensuremath{\alpha}.

Set c=cos(\ensuremath{\beta}/2), s=sin(\ensuremath{\beta}/2), a=e\textasciicircum{}\{\ensuremath{-}i\ensuremath{\alpha}/2\} and T=e\textasciicircum{}\{i\ensuremath{\alpha}\}X\ensuremath{^{-1}}Y. The pivoted top blocks are X\_\ensuremath{\beta}=aX(cI+sT) and Y\_\ensuremath{\beta}=aX(cT\ensuremath{-}sI). Neither scalar factor c+s\ensuremath{\tau}\_j nor c\ensuremath{\tau}\_j\ensuremath{-}s can vanish for real \ensuremath{\beta}. If c+s\ensuremath{\tau}\_j=0, taking imaginary parts forces s=0, and then real parts force c=0; the other factor similarly forces c=s=0. Both contradict c\ensuremath{^2}+s\ensuremath{^2}=1. The determinant product formulas therefore prove that both pivoted blocks remain invertible for every real \ensuremath{\beta}, not merely on [0,\ensuremath{\pi}].

For an invertible matrix A, let P(A)=(AA*)\textasciicircum{}\{\ensuremath{-}1/2\}A. The positive definite spectral calculus gives (AA*)\textasciicircum{}\{\ensuremath{-}1/2\}(AA*)(AA*)\textasciicircum{}\{\ensuremath{-}1/2\}=I, so P(A)P(A)*=I. In finite square dimension this proves unitarity. Hence C(\ensuremath{\beta})=\ensuremath{-}iP(X\_\ensuremath{\beta})*P(Y\_\ensuremath{\beta}) is unitary. The path is continuous: matrix multiplication and the trigonometric coefficients are continuous, and inverse square root is continuous on positive definite matrices. An ordinary justification of the last statement fixes a positive spectral lower bound near any given matrix, approximates x\ensuremath{\mapsto}x\textasciicircum{}\{\ensuremath{-}1/2\} uniformly by polynomials on a compact positive interval, and uses the spectral theorem to bound the matrix approximation error by the scalar uniform error. Polynomial evaluation is continuous, giving the desired continuity in operator norm and therefore in the usual finite-dimensional matrix topology.

At \ensuremath{\beta}=0, X\_0=aX and Y\_0=e\textasciicircum{}\{i\ensuremath{\alpha}/2\}Y. At \ensuremath{\beta}=\ensuremath{\pi} they are X\_\ensuremath{\pi}=e\textasciicircum{}\{i\ensuremath{\alpha}/2\}Y and Y\_\ensuremath{\pi}=\ensuremath{-}aX. A unit complex scalar multiplies the polar factor by that scalar. Substituting those endpoint blocks gives C(\ensuremath{\pi})=C(0)*. This calculation uses neither unitarity of the full block matrix [X Y;Z W] nor an assumed circuit decomposition.

\subsection{Ordinary proof: a continuous determinant phase with opposite endpoints}
For a nonreal \ensuremath{\tau} define h\_\ensuremath{\tau} on its closed half-plane by h\_\ensuremath{\tau}(z)=arg(\ensuremath{-}iz)+\ensuremath{\pi}/2 when Im \ensuremath{\tau}>0, and h\_\ensuremath{\tau}(z)=arg(iz)\ensuremath{-}\ensuremath{\pi}/2 when Im \ensuremath{\tau}<0. The rotated argument avoids the negative-real branch cut on that closed half-plane away from zero, so h\_\ensuremath{\tau} is continuous there. It satisfies e\textasciicircum{}\{ih\_\ensuremath{\tau}(z)\}=z/|z|. On a nonreal input \ensuremath{\tau} it agrees with arg \ensuremath{\tau}. At z=1 it is zero; at z=\ensuremath{-}1 it is \ensuremath{\pi} in the upper-half-plane branch and \ensuremath{-}\ensuremath{\pi} in the lower-half-plane branch.

For 0\ensuremath{\leq}\ensuremath{\beta}\ensuremath{\leq}\ensuremath{\pi}, c and s are nonnegative. The factors n\_j(\ensuremath{\beta})=c\ensuremath{\tau}\_j\ensuremath{-}s and d\_j(\ensuremath{\beta})=c+s\ensuremath{\tau}\_j have imaginary parts c Im \ensuremath{\tau}\_j and s Im \ensuremath{\tau}\_j. They stay in the closed half-plane selected by \ensuremath{\tau}\_j and never vanish, by the preceding factor argument. Define \ensuremath{\sigma}(\ensuremath{\beta})=\ensuremath{\Sigma}\_j[h\_\{\ensuremath{\tau}\_j\}(n\_j(\ensuremath{\beta}))\ensuremath{-}h\_\{\ensuremath{\tau}\_j\}(d\_j(\ensuremath{\beta}))]. Every summand is continuous, so \ensuremath{\sigma} is a continuous real function on [0,\ensuremath{\pi}].

The determinant factors of X\_\ensuremath{\beta} and Y\_\ensuremath{\beta} share the same nonzero prefactor a\textasciicircum{}D det X; their quotient is \ensuremath{\prod}\_j n\_j/d\_j. The inverse square roots in the polar formula have positive real determinants, and (\ensuremath{-}i)\textasciicircum{}D=1 because D is divisible by four. Thus det C(\ensuremath{\beta}) has the phase of that product. Since C(\ensuremath{\beta}) is unitary, its determinant has modulus one. Consequently det C(\ensuremath{\beta})=e\textasciicircum{}\{i\ensuremath{\sigma}(\ensuremath{\beta})\} exactly. This is a determinant calculation, not an assumption that principal arguments add without branch corrections.

At \ensuremath{\beta}=0 the numerator is \ensuremath{\tau}\_j and denominator is 1, giving \ensuremath{\sigma}(0)=\ensuremath{\Sigma}\_j arg \ensuremath{\tau}\_j. At \ensuremath{\beta}=\ensuremath{\pi} the numerator is \ensuremath{-}1 and denominator is \ensuremath{\tau}\_j, giving \ensuremath{\sigma}(\ensuremath{\pi})=\ensuremath{\pi}(N\ensuremath{_{+}}\ensuremath{-}N\ensuremath{_{-}})\ensuremath{-}\ensuremath{\Sigma}\_j arg \ensuremath{\tau}\_j. The balanced rotation makes N\ensuremath{_{+}}=N\ensuremath{_{-}}, and therefore \ensuremath{\sigma}(\ensuremath{\pi})=\ensuremath{-}\ensuremath{\sigma}(0). The prescribed real phase, including its integer branch, is what allows the normalized eigenphase lift to have the required endpoint symmetry.

\subsection{Ordinary proof: existence, uniqueness and stability of normalized phase lifts}
Write L=2\ensuremath{\pi}. A normalized phase lift is a nondecreasing real list \ensuremath{\nu}\_0\ensuremath{\leq}...\ensuremath{\leq}\ensuremath{\nu}\_\{D\ensuremath{-}1\} with \ensuremath{\nu}\_\{D\ensuremath{-}1\}\ensuremath{\leq}\ensuremath{\nu}\_0+L. Two phase lists have the same residues if, after a permutation, corresponding entries differ by integer multiples of L. Every residue multiset admits a normalized lift with any prescribed sum differing from the sum of one representative list by an integer multiple of L. To see this, reduce the representatives to [0,L) and sort them. Moving the first entry to the end after adding L preserves order and width and increases the sum by L. If the required integer change is k=Dh+r with 0\ensuremath{\leq}r<D, perform r such shifts and then add Lh to all D entries. This realizes the target sum and leaves every complex phase unchanged.

For uniqueness, define F\_\ensuremath{\nu}(x)=\ensuremath{\Sigma}\_j floor((x\ensuremath{-}\ensuremath{\nu}\_j)/L). Permuting entries leaves F unchanged; adding Lm\_j to entries subtracts \ensuremath{\Sigma}m\_j from F and adds L\ensuremath{\Sigma}m\_j to their sum. Hence F\_\ensuremath{\nu}(x)+(\ensuremath{\Sigma}\ensuremath{\nu}\_j)/L depends only on the residue multiset. Two normalized lists with the same residues and equal sums therefore have identical F functions.

A normalized list is recovered from F by \ensuremath{\nu}\_j\ensuremath{\leq}x if and only if j+1\ensuremath{-}D\ensuremath{\leq}F\_\ensuremath{\nu}(x), with indices starting at zero. Below \ensuremath{\nu}\_0 every floor is at most \ensuremath{-}1, so F\ensuremath{\leq}\ensuremath{-}D. On [\ensuremath{\nu}\_0,\ensuremath{\nu}\_0+L), each summand is zero when \ensuremath{\nu}\_j\ensuremath{\leq}x and \ensuremath{-}1 otherwise; thus F is the count of entries at most x minus D. Above or at \ensuremath{\nu}\_0+L every summand is nonnegative and every entry is at most x. These three cases prove the equivalence. Equal F functions consequently give the same threshold set for each ordered entry, proving entrywise uniqueness. In particular this remains valid when phases coincide or the width equals L.

The same invariant gives quantitative stability. Let \ensuremath{\nu} and \ensuremath{\nu}\ensuremath{^{\prime}} be normalized, and let \ensuremath{\mu} be a possibly unsorted lift of \ensuremath{\nu}\ensuremath{^{\prime}}'s residues such that |\ensuremath{\mu}\_j\ensuremath{-}\ensuremath{\nu}\_\{\ensuremath{\sigma}(j)\}|\ensuremath{\leq}\ensuremath{\varepsilon} for a permutation \ensuremath{\sigma}. Assume |\ensuremath{\Sigma}\ensuremath{\nu}\ensuremath{^{\prime}}\ensuremath{-}\ensuremath{\Sigma}\ensuremath{\nu}|+D\ensuremath{\varepsilon}<L. Matched closeness gives |\ensuremath{\Sigma}\ensuremath{\mu}\ensuremath{-}\ensuremath{\Sigma}\ensuremath{\nu}|\ensuremath{\leq}D\ensuremath{\varepsilon}, hence |\ensuremath{\Sigma}\ensuremath{\mu}\ensuremath{-}\ensuremath{\Sigma}\ensuremath{\nu}\ensuremath{^{\prime}}|<L. The invariant shows that (\ensuremath{\Sigma}\ensuremath{\mu}\ensuremath{-}\ensuremath{\Sigma}\ensuremath{\nu}\ensuremath{^{\prime}})/L is an integer; it must be zero. Thus F\_\{\ensuremath{\nu}\ensuremath{^{\prime}}\}=F\_\ensuremath{\mu}.

For every x, matched closeness and monotonicity of the floor give F\_\ensuremath{\nu}(x\ensuremath{-}\ensuremath{\varepsilon})\ensuremath{\leq}F\_\ensuremath{\mu}(x)\ensuremath{\leq}F\_\ensuremath{\nu}(x+\ensuremath{\varepsilon}). Apply the recovery equivalence first at x=\ensuremath{\nu}\_j+\ensuremath{\varepsilon} and then at x=\ensuremath{\nu}\ensuremath{^{\prime}}\_j. It yields \ensuremath{\nu}\ensuremath{^{\prime}}\_j\ensuremath{\leq}\ensuremath{\nu}\_j+\ensuremath{\varepsilon} and \ensuremath{\nu}\_j\ensuremath{\leq}\ensuremath{\nu}\ensuremath{^{\prime}}\_j+\ensuremath{\varepsilon}, respectively. Therefore |\ensuremath{\nu}\ensuremath{^{\prime}}\_j\ensuremath{-}\ensuremath{\nu}\_j|\ensuremath{\leq}\ensuremath{\varepsilon} for every j. This proves the exact stability lemma used in the source and explains why a small continuous change in the target sum prevents a discontinuous cyclic change of branch.

\subsection{Ordinary proof: continuity through eigenvalue collisions}
For a continuous finite matrix path M(\ensuremath{\beta}), its eigenvalue multiset is continuous in the matching sense: for \ensuremath{\beta} near \ensuremath{\beta}\_0, any full root list v\_j of M(\ensuremath{\beta}) can be permuted to lie arbitrarily close to a fixed list u\_j of M(\ensuremath{\beta}\_0). Here multiplicities are retained. An elementary compactness proof starts with the bound |\ensuremath{\lambda}|\ensuremath{\leq}\ensuremath{\Sigma}\_\{i,j\}|M\_\{ij\}| for every eigenvalue \ensuremath{\lambda}. Choose a nonzero eigenvector and a coordinate of maximal modulus; its row equation gives that bound after dividing by the positive maximal modulus.

If the matching statement failed along \ensuremath{\beta}\_n\ensuremath{\to}\ensuremath{\beta}\_0, choose bad root lists v\textasciicircum{}\{(n)\}. The bound makes them eventually bounded in the finite-dimensional space \ensuremath{\mathbb{C}}\textasciicircum{}D. A subsequence therefore converges to a list r. For every complex z, \ensuremath{\prod}\_j(z\ensuremath{-}v\_j\textasciicircum{}\{(n)\})=det(zI\ensuremath{-}M(\ensuremath{\beta}\_n)). Passing to the limit gives \ensuremath{\prod}\_j(z\ensuremath{-}r\_j)=det(zI\ensuremath{-}M(\ensuremath{\beta}\_0)). These monic polynomials are equal, so their root multisets, including multiplicities, agree. A permutation matches r to u, and sufficiently late lists in the subsequence are close under that permutation, contradicting their choice. No assumption of a collision-free spectrum along the path is used.

Apply this to the continuous unitary C(\ensuremath{\beta}). Choose at each \ensuremath{\beta} a normalized eigenphase list \ensuremath{\nu}(\ensuremath{\beta}) with sum \ensuremath{\sigma}(\ensuremath{\beta}); existence follows from the determinant phase and the lift existence lemma. Uniqueness makes this list well defined regardless of the auxiliary root ordering. Fix \ensuremath{\beta}\_0 and a desired entrywise tolerance \ensuremath{\varepsilon}>0, and choose \ensuremath{\varepsilon}\ensuremath{^{\prime}} smaller than both \ensuremath{\varepsilon} and \ensuremath{\pi}/(2D). Matching eigenvalues sufficiently closely gives real representatives \ensuremath{\mu}\_j near \ensuremath{\nu}\_\{\ensuremath{\sigma}(j)\}(\ensuremath{\beta}\_0) within \ensuremath{\varepsilon}\ensuremath{^{\prime}}: on the unit circle, |e\textasciicircum{}\{ix\}\ensuremath{-}e\textasciicircum{}\{iy\}|=2|sin((x\ensuremath{-}y)/2)|, so sufficiently small chord distance admits such a nearby phase representative.

Continuity of \ensuremath{\sigma} also gives |\ensuremath{\sigma}(\ensuremath{\beta})\ensuremath{-}\ensuremath{\sigma}(\ensuremath{\beta}\_0)|<\ensuremath{\pi} in a small enough neighborhood. Since D\ensuremath{\varepsilon}\ensuremath{^{\prime}}<\ensuremath{\pi}/2, the stability condition is strictly below 2\ensuremath{\pi}. The preceding stability lemma now bounds every |\ensuremath{\nu}\_j(\ensuremath{\beta})\ensuremath{-}\ensuremath{\nu}\_j(\ensuremath{\beta}\_0)| by \ensuremath{\varepsilon}\ensuremath{^{\prime}}. Thus the normalized ordered eigenphase list is continuous on [0,\ensuremath{\pi}], even through eigenvalue collisions and crossings of the principal-argument branch cut. This is the crucial conclusion that would not follow from choosing principal arguments independently.

\subsection{Ordinary proof: the middle-half sign change and the exact matrix endpoint}
The identity C(\ensuremath{\pi})=C(0)* implies that the endpoint eigenphase multiset is the negative of the initial one modulo 2\ensuremath{\pi}. The reversed negative list \ensuremath{-}\ensuremath{\nu}\_\{D\ensuremath{-}1\ensuremath{-}j\}(0) is normalized and has sum \ensuremath{-}\ensuremath{\sigma}(0)=\ensuremath{\sigma}(\ensuremath{\pi}). Uniqueness therefore proves \ensuremath{\nu}\_j(\ensuremath{\pi})=\ensuremath{-}\ensuremath{\nu}\_\{D\ensuremath{-}1\ensuremath{-}j\}(0) exactly as real numbers, not merely modulo 2\ensuremath{\pi}.

Let S=\{p,p+1,...,3p\ensuremath{-}1\}, the middle half of the D=4p ranks. It has size 2p and is invariant under reversal j\ensuremath{\mapsto}D\ensuremath{-}1\ensuremath{-}j. The real imbalance P(\ensuremath{\beta})=\ensuremath{\Sigma}\_\{j\ensuremath{\in}S\}\ensuremath{\nu}\_j(\ensuremath{\beta})\ensuremath{-}\ensuremath{\Sigma}\_\{j\ensuremath{\notin}S\}\ensuremath{\nu}\_j(\ensuremath{\beta}) is continuous. Endpoint reversal invariance and negation give P(\ensuremath{\pi})=\ensuremath{-}P(0). If either endpoint vanishes the desired point is already found; otherwise their signs are opposite and the ordinary intermediate value theorem gives \ensuremath{\beta}*\ensuremath{\in}[0,\ensuremath{\pi}] with P(\ensuremath{\beta}*)=0.

The positive positions of the chosen Walsh row are the first 2p basis indices. Choose a permutation e carrying them onto S. Reassign the eigenphases by \ensuremath{\theta}\_j=\ensuremath{\nu}\_\{e(j)\}(\ensuremath{\beta}*); this preserves the eigenvalue multiset. The Walsh sum \ensuremath{\Sigma}\_j h\_j\ensuremath{\theta}\_j then equals P(\ensuremath{\beta}*) and is zero, where h\_j=+1 on the first 2p positions and \ensuremath{-}1 on the others. Thus finalAngle=(1/D)\ensuremath{\Sigma}h\_j\ensuremath{\theta}\_j=0. More generally integer shifts \ensuremath{\theta}\_j\ensuremath{\mapsto}\ensuremath{\theta}\_j+2\ensuremath{\pi}m\_j preserve eigenvalues and alter the Walsh sum by an integer multiple of 2\ensuremath{\pi}; the actual zero-imbalance construction can take all additional m\_j=0, because the normalized lift already carries its required branches.

Finally multiply the full block matrix U=[X Y;Z W] by the lifted one-qubit gate g=Rz(\ensuremath{\alpha})Ry(\ensuremath{\beta}*). Direct block multiplication gives the pivoted blocks already used in the polar path and the corresponding lower blocks. The gate is unitary by its elementary rotation formulas. The phase reassignment above supplies the root list of zxzC(X\_\ensuremath{\beta}*,Y\_\ensuremath{\beta}*) with finalAngle zero. This proves exactly the existential matrix conclusion of pivot\_merge for every admitted dimension and input. It does not construct a circuit, count the cost of the eigenbasis/permutation, extend the argument to coincident crossing angles, or derive the universal improved recursive node rule.

The separate arithmetic recursion c\_2=3 and c\_\{n+1\}=4c\_n+3\ensuremath{\cdot}2\textasciicircum{}n\ensuremath{-}6 has the source's closed form (21\ensuremath{\cdot}4\textasciicircum{}n\ensuremath{-}72\ensuremath{\cdot}2\textasciicircum{}n+96)/48. Its conversion into a bound on every n-qubit circuit explicitly requires a two-qubit construction and a universal synthesis-node hypothesis. Those obligations remain hypotheses of generalBound\_of\_pivot\_node. The ordinary matrix proof above cannot remove them. This expansion is attributed to Matt Bowring's pinned source; the fresh compiler/axiom audit now passes, while contributor approval and historical novelty review remain independent pending checks.

\subsection{Relation to established Block-ZXZ and Walsh methods}
De Vos and De Baerdemacker, The block-ZXZ synthesis of an arbitrary quantum circuit (v3), section 5, equation 5, already give D=\ensuremath{-}iV\ensuremath{_1}\ensuremath{_1}*V\ensuremath{_1}\ensuremath{_2} for the left polar decompositions U\_jk=P\_jk V\_jk. For invertible top blocks their D equals the present C=\ensuremath{-}iP(X)*P(Y). Krol and Al-Ars section 4.1 use the same polar primitive. Changing notation does not make this formula new.

The conversion from an eigenphase list to uniformly controlled rotation angles also predates this submission. Möttönen, Vartiainen, Bergholm and Salomaa, Quantum Circuits for General Multiqubit Gates, equations 3-5 and the ensuing discussion, derive a Gray-code column permutation of the Walsh-Hadamard matrix and its inverse. Krol-Al-Ars section 3 uses that transformation; their section 5 establishes the existing Block-ZXZ circuit optimization and the recurrence c\_n\ensuremath{\leq}4c\_\{n\ensuremath{-}1\}+3\ensuremath{\cdot}2\textasciicircum{}\{n\ensuremath{-}1\}\ensuremath{-}5. Thus the Walsh transform, the recursive decomposition framework and the established 22/48 leading coefficient are background, rather than new results of this chapter.

The restricted contribution proposed here is the conjunction of a balanced Rz(\ensuremath{\alpha})Ry(\ensuremath{\beta}) pivot, a nonsingular polar path under separated crossing angles, a continuous normalized real eigenphase lift with a prescribed determinant sum, and an intermediate-value argument producing an exactly zero selected Walsh angle after relabeling. The closest primary constructions inspected above do not state or prove that combined generic existence theorem. This is a bounded comparison, not proof of historical priority across all spectral phase-balancing literature. Conditional P2=.05 credit concerns that exact combined theorem: its replay now passes, while author review, novelty and judging acceptance remain pending. Its matrix conclusion does not establish a circuit implementing the desired extra cancellation or the proposed universal 21/48 bound.

\subsection{Fresh exact verification: selected proof and unfinished baseline separated}
The endpoint generalBound\_of\_pivot\_node is itself axiom-clean because it proves an implication from explicitly supplied hbase and hnode. Its successful replay does not discharge either premise. The generic pivot\_merge theorem supplies the matrix phase conclusion under its displayed separation hypotheses; no proof in this chain turns it into the universal circuit node. All eight separate baseline Claims prints include sorryAx, exactly as their underlying admitted source predicts. These actual axiom lists are preserved in the receipt and companion verification record. Accepted competition points remain zero pending novelty and chair scope acceptance; the recommended narrow conditional score stays 33.9488 at p=.05.

\subsection{Verification and reproducibility}
Fresh execution: matt-unitary-current: PASS\_SELECTED\_ENDPOINTS\_WITH\_UNFINISHED\_BASELINES; 26/26 custom modules compiled successfully; receipt Verification/matt-unitary-current-fresh-build.json. Compiler pins, source hashes and endpoint axiom outputs are in Verification. Historical receipts and finite checks do not replace fresh replay.

{\small Independent numerical checks of nine generic block pairs in dimensions 4, 8 and 12 agree with the polar-path unitarity, determinant phase, endpoint reversal and Walsh sign change. These are sanity checks, not a proof, Lean replay or circuit-count result; see Verification/Independent checkers/matt\_polar\_path\_sanity\_check.json.\par}

\begin{samepage}
\subsection{LeanUnitary.Pivot.pivot\_merge}
\begin{quote}\small\ttfamily For D=4(q+1), invertible complex D-by-D blocks X,Y, a full eigenvalue list of X inverse times Y and pairwise crossing-angle separation modulo pi, there exist an unrestricted real angle alpha and an angle beta in [0,pi], an ordinary unitary one-qubit gate, pivoted blocks satisfying the lifted block multiplication identity, and a real phase list of zxzC(Xprime,Yprime) whose finalAngle(q+1) equals zero. No unitarity assumption on the entire block matrix U is needed. The theorem concludes no circuit or CNOT-count bound.\end{quote}
{\small Source: LeanUnitary/Pivot/Existence/Pivot.lean.\par}

{\small Source SHA-256: dd8f689a42fc44b1c1ad17fa7af1e5de337aff7ccc0ec0ef7551ff81e362c2bf\par}

\end{samepage}
\begin{samepage}
\subsection{LeanUnitary.Pivot.continuousOn\_lift}
\begin{quote}\small\ttfamily Under invertibility and absence of real rotated eigenvalues, the normalized ordered eigenphase lift of the polar multiplexor is continuous over beta in [0,pi].\end{quote}
{\small Source: LeanUnitary/Pivot/Existence/Pivot.lean.\par}

{\small Source SHA-256: dd8f689a42fc44b1c1ad17fa7af1e5de337aff7ccc0ec0ef7551ff81e362c2bf\par}

\end{samepage}
\begin{samepage}
\subsection{LeanUnitary.Pivot.count\_closed}
\begin{quote}\small\ttfamily For n>=2, 48*count(n)+72*2\textasciicircum{}n=21*4\textasciicircum{}n+96. This is the arithmetic closed form of the proposed recurrence, not a constructive synthesis theorem.\end{quote}
{\small Source: LeanUnitary/Pivot/Count.lean.\par}

{\small Source SHA-256: 8683beed72dc58584946879ce8abb0a12f647e1d6a8ab3b00bcc8b3e1424c453\par}

\end{samepage}
\begin{samepage}
\subsection{LeanUnitary.Pivot.saving}
\begin{quote}\small\ttfamily For n>=2, count(n)<=zxzCount(n) and 3*(zxzCount(n)-count(n))=4\textasciicircum{}(n-2)-1. These recursively defined counts do not assert that circuits achieving count(n) have been constructed.\end{quote}
{\small Source: LeanUnitary/Pivot/Count.lean.\par}

{\small Source SHA-256: 8683beed72dc58584946879ce8abb0a12f647e1d6a8ab3b00bcc8b3e1424c453\par}

{\small\textbf{Sources and attribution:} \href{https://github.com/bowrango/lean-unitary/blob/7014b63ee2cfd01c641a7f539c73b664a2aa3d66/LeanUnitary/Pivot/Existence/Pivot.lean}{1. bowrango/lean-unitary / Pivot.lean}; \href{https://github.com/bowrango/lean-unitary/blob/7014b63ee2cfd01c641a7f539c73b664a2aa3d66/LeanUnitary/Pivot/Multiplexor/Basic.lean}{2. bowrango/lean-unitary / Basic.lean}; \href{https://github.com/bowrango/lean-unitary/blob/7014b63ee2cfd01c641a7f539c73b664a2aa3d66/LeanUnitary/Pivot/Synthesis.lean}{3. bowrango/lean-unitary / Synthesis.lean}; \href{https://github.com/bowrango/lean-unitary/blob/7014b63ee2cfd01c641a7f539c73b664a2aa3d66/LeanUnitary/Claims.lean}{4. bowrango/lean-unitary / Claims.lean}; \href{https://github.com/bowrango/lean-unitary/blob/7014b63ee2cfd01c641a7f539c73b664a2aa3d66/LeanUnitary/README.md}{5. bowrango/lean-unitary / README.md}; \href{https://arxiv.org/abs/2403.13692v2}{6. arXiv source: 2403.13692v2}; \href{https://github.com/bowrango/lean-unitary/blob/7014b63ee2cfd01c641a7f539c73b664a2aa3d66/LeanUnitary/Pivot/Count.lean}{7. bowrango/lean-unitary / Count.lean}; \href{https://arxiv.org/abs/1512.07240v3}{8. arXiv source: 1512.07240v3}; \href{https://arxiv.org/abs/quant-ph/0404089v3}{9. arXiv source: quant-ph/0404089v3}; \href{https://arxiv.org/html/2512.23720v2}{10. arXiv source: 2512.23720v2}; \href{https://www.mattbowring.com/posts/gate-decomp/}{11. www.mattbowring.com / gate-decomp}; \href{https://www.mattbowring.com/}{12. www.mattbowring.com / }; \href{https://github.com/bowrango/lean-unitary/tree/7014b63ee2cfd01c641a7f539c73b664a2aa3d66}{13. bowrango/lean-unitary}; \href{https://github.com/alejandrozu/openmath-2026-judging}{Central source and adjudication archive}.\par}

\end{samepage}
{\small Private editorial working draft of a newly delivered generic matrix result. The selected proof was completed on 4 October 2026, after the original event cutoff, and has been accepted for current intake review. Fresh pinned replay and the selected axiom audit passed on 5 October 2026; historical novelty, author approval and preferred affiliation confirmation remain pending. The separate baseline Claims audit retains sorryAx. The universal 21/48 CNOT bound is not established by the selected endpoint. No submission or public release is implied.\par}

\clearpage
\section{Shared verification register and trust boundaries}
The companion Verification directory contains the dated source inventory, fresh Lean execution results, endpoint axiom outputs and independent numerical certificates. A verified source hash establishes which bytes were checked. It does not turn a provisional novelty comparison or a competition score into mathematical priority. Native-evaluation endpoints disclose their compiler/runtime trust separately from ordinary kernel derivations. The brief placement document applies the chair-approved rules and retained DMS progress of 0.15.

The portfolio contains several genuinely substantial candidates, but its strength depends on precision: an explicit subsidiary question is different from a parent conjecture, a numerical certificate is different from a frontier comparison, and an archived proof receipt is different from this pass's fresh execution. The chapters retain those boundaries so the collection can support fair placement, targeted review and accurate publication.

Final acceptance should be stated theorem by theorem. If a priority check downgrades a mathematical novelty claim, retain its faithful verification and useful formalization in the accompanying nonnovel-results paper. If a fresh compiler run fails, preserve the exact error and source revision, repair transparently where possible and re-run; do not replace a failed run by an earlier passing receipt or by the score assigned to that result.

\clearpage
\section{Author and affiliation register}
Affiliations are source-backed declarations, not institutional endorsement. Confirm preferred name, publication affiliation and authorship order before external release. Do not infer paper authorship from event roster or organizer role.

\subsection{Rachel Teo (E02)}
Rachel S. Y. Teo on personal academic website; event uses Rachel Teo

Recorded role: Joint Sakana packet; do not invent per-result split

Sakana AI. Evidence: Event team affiliation; author correspondence and Sakana-owned team repository.

{\small \url{https://github.com/SakanaAI/autolab-100}\par}

National University of Singapore. Evidence: Author-declared on personal academic website fetched October 4; year-of-study text may be stale.

{\small \url{https://rachtsy.github.io/}\par}

\subsection{Wenyi Wang (E02)}
Recorded role: Joint Sakana packet; do not invent per-result split

Sakana AI. Evidence: Author public profile and team packet.

{\small \url{https://sa.linkedin.com/in/wenyiw}\par}

{\small \url{https://github.com/SakanaAI/autolab-100}\par}

King Abdullah University of Science and Technology (KAUST). Evidence: Author LinkedIn education lists 2022-present; preferred paper dual affiliation must be confirmed.

{\small \url{https://sa.linkedin.com/in/wenyiw}\par}

\subsection{Hamdi Abderrahmene (E02)}
Public LinkedIn uses Abderrahmene Hamdi; confirm desired order

Recorded role: Joint Sakana packet; do not invent per-result split

Sakana AI. Evidence: Author LinkedIn Research Intern August 2026-present, corroborated team repository.

{\small \url{https://jp.linkedin.com/in/abderrahmenehamdi}\par}

{\small \url{https://github.com/SakanaAI/autolab-100}\par}

Nagoya University. Evidence: Author LinkedIn master's programme 2026-2028; preferred paper dual affiliation must be confirmed.

{\small \url{https://jp.linkedin.com/in/abderrahmenehamdi}\par}

\subsection{Woohyuk Kang (E08)}
Recorded role: DMS, selected Erdős formalizations, leading Ramsey search and certificate, Kobon 39 packet; team assembly

HTPeo; Kyung Hee University, CS\&E. Evidence: Personally entered author declaration in pinned TEAM.md and manuscript; preserve department abbreviation.

{\small \url{https://github.com/lavaskiller/openmath-2026-htpeo/blob/76c63b8e425e551a930060a6871b84fe3b4780db/TEAM.md}\par}

{\small \url{https://github.com/lavaskiller/openmath-2026-htpeo/blob/76c63b8e425e551a930060a6871b84fe3b4780db/paper/htpeo-openmath-2026.tex}\par}

\subsection{Hyunjin Lee (E08)}
Recorded role: Separate Ramsey final/validation entry; Erdős 1038 with Sanghyeon, not claimed; public-seed matrix search

University of Cambridge. Evidence: Personally entered author declaration in pinned TEAM.md.

{\small \url{https://github.com/lavaskiller/openmath-2026-htpeo/blob/76c63b8e425e551a930060a6871b84fe3b4780db/TEAM.md}\par}

\subsection{Sanghyeon Lee (E08)}
Recorded role: Busy Beaver 249881 certificate; Ramsey leaderboard record; Erdős 1038 with Hyunjin, not claimed

Korea University, Security. Evidence: Author declaration recorded in pinned TEAM.md; preserve department wording.

{\small \url{https://github.com/lavaskiller/openmath-2026-htpeo/blob/76c63b8e425e551a930060a6871b84fe3b4780db/TEAM.md}\par}

\subsection{Youchan Oh (E08)}
Recorded role: Known Kobon 18=93 reconstruction

Seoul National University, TI. Evidence: Author declaration recorded in pinned TEAM.md; do not guess TI expansion.

{\small \url{https://github.com/lavaskiller/openmath-2026-htpeo/blob/76c63b8e425e551a930060a6871b84fe3b4780db/TEAM.md}\par}

\subsection{Chandragupt Sharma (E04)}
Recorded role: Sole named matrix support 138 packet author

Cluster Innovation Centre, University of Delhi. Evidence: Author public LinkedIn profile; identity linked by supplied registration profile; confirm preferred paper affiliation.

{\small \url{https://in.linkedin.com/in/chandraguptsharma}\par}

LinkedIn also lists Mathematical Sciences Foundation experience; role/date insufficient to assert that as paper institution.

\subsection{Rohith Poola (E57)}
Recorded role: Sole named Kobon catalogue author

Autonomy and Intelligence Lab, Northeastern University. Evidence: Official university March 9 2026 profile and current author website.

{\small \url{https://coe.northeastern.edu/news/driven-by-curiosity-rohith-poolas-journey-in-autonomous-robotics/}\par}

{\small \url{https://Rohith18p.github.io/}\par}

\subsection{Leon Koerbs (E11)}
GitHub and LinkedIn use Leon Körbs; use author-confirmed spelling

Recorded role: Joint Erdős 3 Token Burner packet with Cameron Fen; per-result roles not divided

Harvard John A. Paulson School of Engineering and Applied Sciences, Intelligent Interactive Systems Group. Evidence: Author public LinkedIn says Visiting Scholar for fall term; not independent institutional appointment verification.

{\small \url{https://de.linkedin.com/in/leon-koerbs}\par}

Technical University of Munich, School of Management. Evidence: Current author LinkedIn study 2025-2027; earlier official exchange report corroborates identity.

{\small \url{https://de.linkedin.com/in/leon-koerbs}\par}

{\small \url{https://cms.mgt.tum.de/fileadmin/mgt.tum.de/downloads/international_exchange_programs/experience_reports/Erfahrungsbericht_Leon_Koerbs_SMU_SS_23_DE_TUMSOMex_min.pdf}\par}

\subsection{Cameron Fen (E11)}
Recorded role: Joint Erdős 3 Token Burner packet with Leon

State Street. Evidence: Current author public LinkedIn identifies State Street; identity-to-event roster and permitted publication affiliation require confirmation.

{\small \url{https://www.linkedin.com/in/cameron-fen-phd-945237225}\par}

{\small \url{https://cameronfen.github.io/}\par}

University of Michigan is documented education, not asserted current affiliation.

\subsection{Wilson Wu (E65)}
Recorded role: Sole named supporting-formalizations packet author

Affiliation not verified in the supplied materials. Evidence: GitHub and personal page identify author but no institution/employer.

{\small \url{https://github.com/wilsonwu-ai}\par}

{\small \url{https://wilsonwu-ai.github.io/}\par}

\subsection{Jamie Steeg (E10)}
Recorded role: Sole named H1/H4/H7 entrant; GCL/GCT project provenance not paper affiliation

No organizational affiliation asserted. Evidence: Explicit human authority declaration in frozen packet; individual and self-resourced.

{\small \url{https://github.com/grandchallenge/MATHSOLVE/blob/2fd9fd643bf3e9e98d121dc2e7af1cf1be411ae2/work_packages/OPENMATH_2026/COMPETITION_PACKETS/HUMAN_AUTHORITY_DECLARATION.json}\par}

\subsection{Chaewon Yoon (E05)}
Recorded role: Collatz 512-rule packet

Affiliation not verified in the supplied materials. Evidence: Named private submission email and artifact establish author, not institution.

\subsection{Matt Bowring (E03)}
Recorded role: Generic matrix eigenphase-pivot partial theorem received 4 October; actual exact-pin replay completed 5 October with all 26 authored source modules compiling and all six selected generic endpoint axiom outputs standard-kernel only. Universal improved CNOT synthesis claim remains unproved; separate Claims diagnostics retain sorryAx.

Purdue University, Department of Mechanical Engineering (documented January 2026; current preferred publication affiliation to confirm). Evidence: Author-declared institutional affiliation in arXiv manuscript v2 dated 9 January 2026 and corroborated by personal website; this dated paper affiliation does not establish a current appointment or competition resource class.

{\small \url{https://mattbowring.com/}\par}

{\small \url{https://arxiv.org/html/2512.23720v2}\par}

{\small \url{https://www.linkedin.com/in/matt-bowring-463610149}\par}

\subsection{Luke Van Seters (E01)}
Recorded role: Joint Ramsey packet with Madhan Jothimani; preserved manuscript provides author names, no affiliations

Klaviyo. Evidence: Authenticated public LinkedIn profile header reviewed 4 October 2026 at approximately 22:58 UTC states Director of Engineering at Klaviyo. This is a dated profile declaration; preferred publication affiliation awaits author confirmation.

{\small \url{https://www.linkedin.com/in/lukevanseters/}\par}

{\small \url{https://github.com/lukevs/k4-ramsey/blob/2cf216ba885fbb1b2d1b2d4920efeaa6ba438842/README.md}\par}

Earlier Agency AI attribution remains unconfirmed historical information and is superseded as a current-profile description. Northeastern University appears as profile education, not a confirmed current university affiliation. Declared competition resources determine the competition class separately from employment or education.

\subsection{Madhan Jothimani (E01)}
Recorded role: Joint Ramsey packet with Luke Van Seters

Northeastern University. Evidence: Author public LinkedIn lists education from 2024 to the present; exact programme and preferred paper affiliation require author confirmation.

{\small \url{https://www.linkedin.com/in/madhanjothimani}\par}

{\small \url{https://github.com/lukevs/k4-ramsey/blob/2cf216ba885fbb1b2d1b2d4920efeaa6ba438842/README.md}\par}

Profile lists Infosys experience without sufficiently clear current dates; not silently asserted as publication employer.

\subsection{Luca Seiki Pereira Fujii (E06)}
Recorded role: Container formalization lead, scope/design and statement review; mention-only administrative decision

IMPA Tech. Evidence: Author explicitly identifies five team members as IMPA Tech-account registrants in focused email; do not rely solely on email domain.

\subsection{Mateus (E06)}
Recorded role: Statement reviewer; do not invent full author name

IMPA Tech. Evidence: Author team declaration; full preferred name unresolved; Luma says Mateus Mundstock, source CREDITS says Matheus.

{\small \url{https://github.com/lazyluca/hyprgraph_containers/blob/adcaf45c06fa70407f1abb2df80e1bd801d970b3/CREDITS.md}\par}

\subsection{Tiago (E06)}
Recorded role: Statement reviewer

IMPA Tech. Evidence: Author team declaration; registration links Tiago Sánchez Ribeiro but exact packet-name expansion awaits confirmation.

{\small \url{https://github.com/lazyluca/hyprgraph_containers/blob/adcaf45c06fa70407f1abb2df80e1bd801d970b3/CREDITS.md}\par}

\subsection{Augusto (E06)}
Recorded role: Event roster member; Luca explicitly says did not contribute to submitted work; acknowledge participation, not automatic paper authorship

IMPA Tech. Evidence: Author team declaration; possible full name Augusto Schaefer Huff from registration, not silently merged.

\subsection{Felipe Leria (E06)}
Recorded role: Event roster member; Luca explicitly says did not contribute to submitted work; acknowledge participation, not automatic paper authorship

IMPA Tech. Evidence: Author team declaration; preferred complete name not expanded from email address.

\subsection{Frederik Smits van Oyen (E12)}
Recorded role: Incomplete Heilbronn covering and a separate Ramsey project without an author-designated final certificate

Affiliation not verified in the supplied materials. Evidence: Named author in frozen TEAM record; neither name similarity nor a GitHub account establishes an institution.

{\small \url{https://github.com/alejandrozu/openmath-2026-judging/blob/d0ed487f487fa7ec814ff705ab55249983fe8707/OpenMath-Judging/contestants/E12/TEAM.json}\par}

{\small \url{https://github.com/FSvOAI/Heilbronn/tree/b8da072acb21d2b0acc8daa038600dddfc8c67e9}\par}

This is the E12 author; do not substitute Frederik Baetens.

\subsection{Raj Harshit Srirangam (E09)}
Recorded role: Kobon upper-bound reductions and exact certificates with unresolved fourfold-point and global closure conditions

Affiliation not verified in the supplied materials. Evidence: Named author in pinned contestant TEAM record and repository; no confirmed present institution.

{\small \url{https://github.com/alejandrozu/openmath-2026-judging/blob/d0ed487f487fa7ec814ff705ab55249983fe8707/OpenMath-Judging/contestants/E09/TEAM.json}\par}

{\small \url{https://github.com/srirangam-r/kobon_triangles/tree/eed14659a9b2f4ca12777da5d557f2b620b966f6}\par}

\subsection{Qichao Wang (E17)}
Recorded role: Cerny partial-formalization manifests; actual source ZIP remains unavailable

Hebei University of Technology. Evidence: Author-declared undergraduate Automation study in event registration; publication affiliation and exact programme styling require author confirmation.

{\small \url{https://github.com/alejandrozu/openmath-2026-judging/blob/d0ed487f487fa7ec814ff705ab55249983fe8707/OpenMath-Judging/contestants/EVENT/12-correspondence/3937a54b1ac3-OpenMath-Luma-roster.json}\par}

{\small \url{https://www.linkedin.com/in/%E5%90%AF%E8%B6%85-%E7%8E%8B-a5a69043b/}\par}

Credit uses the registration and expressly submitted name Qichao Wang: the final correspondence identifies the entrant and signs that name. The account display name differs (Cui Lijun); its identity mapping and the preferred publication name are not independently verified and require author confirmation. No second entrant or coauthor is inferred.

\subsection{Hichem Benali (Unassigned correspondence)}
Recorded role: Finite cyclic five-AP inquiry in Z/25Z; author source proof packet and official target/admission remain unresolved

Affiliation not verified in the supplied materials. Evidence: SYSTASYS is a correspondence/project label, not independently verified as an employer or institution.

Private correspondence identifies the author. No outbound correspondence or institutional inference was made.

\clearpage
\section{Technical source-statement supplements}
The following source-derived statement records preserve the mathematical quantifiers, definitions, exact endpoints and documented prior comparisons. Historical source cards retain their original proposal scores and verification wording. The approved 4 October amendments set b=1, and the current dated verification register governs fresh execution; archived proposal values are not active standings. These supplements retain source authorship and do not upgrade incomplete ordinary proof exposition.

\clearpage
\section{a 060957: preserved statement and proof architecture}
\subsection{A060957: refutation and arbitrarily long prime-exponent gaps}
Completed result and evidence record.

Claim ID: \texttt{a060957}. Family: \texttt{oeis-a060957}.

\subsection{Our contribution}
We supply the complete distinct-factor interpolation refutation and its exact arbitrary-gap strengthening, realizing the eight-level exponent fibre with primes above any prescribed bound. The recorded 45-module Lean 4.33.1 replay and dual-version generalized checks certify the construction with standard axioms and no unproved graph-existence premise, building on known encoding and semigroup-hole methods.

\subsection{Exact preserved statement and scope}
Original refutation: \ensuremath{\exists} n,m,p\ensuremath{\in}N with m>0, p prime, p\ensuremath{\leq}n, m\ensuremath{\in}P(n), p\textasciicircum{}5 m\ensuremath{\in}P(n), and p\textasciicircum{}4 m\ensuremath{\notin}P(n). Stronger: \ensuremath{\forall} g\ensuremath{\geq}1 and B\ensuremath{\in}N, \ensuremath{\exists} n,p,d,h\ensuremath{\in}N with p prime, B<p\ensuremath{\leq}n, d>0, p\ensuremath{\not\mid}d, and for every e\ensuremath{\in}N, p\textasciicircum{}e d\ensuremath{\in}P(n) iff e lies in \{h,h+1,h+2,h+3,h+g+4,h+g+5,h+g+6,h+g+7\}. Hence g consecutive missing levels between attainable endpoints.

Complete refutation, with stronger same-family exact arbitrary-gap theorem

Not claimed:

\begin{itemize}
\item Minimum counterexample size
\end{itemize}
\begin{itemize}
\item A fixed prime working for every gap length
\end{itemize}
\begin{itemize}
\item Every sufficiently large interval
\end{itemize}
\begin{itemize}
\item Independent family credit for each g or the strengthening
\end{itemize}
\subsection{Contribution and prior overlap}
Proposed classification: M3A complete refutation; strengthening in same family union. No acceptance or credit asserted.

\begin{itemize}
\item Whole-interval distinct-factor counterexample and all-size prescribed gaps with exact fiber
\end{itemize}
Recorded overlap:

\begin{itemize}
\item Prime-coordinate encoding, exposed faces, graph constructions and semigroup-hole context are prior methods
\end{itemize}
Recorded bounded prior-art comparison:

\begin{itemize}
\item No exact earlier complete-interval distinct-factor counterexample or arbitrary-gap fiber theorem found in recorded bounded audits
\end{itemize}
\subsection{Formal proof source}
Paths below are relative to \texttt{Claim\_Evidence/} in this bundle. Exact source bytes, declaration lines and verification evidence are preserved. Dependency versions and the inspection/build boundary are described in the bundle README.

\subsection{SubsetProductGap.subset\_product\_counterexample}
Source: \texttt{proofs/a060957/OpHack/SubsetProducts/Main.lean}, line 89. SHA-256 \texttt{5b9d85abd44f036bbbfa11f241b62ebe8dce01a2ffd89305caa965669c72b0fc}.

\begin{quote}\small\ttfamily theorem subset\_product\_counterexample : \ensuremath{\exists} n p m : \ensuremath{\mathbb{N}}, p.Prime \ensuremath{\wedge} p \ensuremath{\leq} n \ensuremath{\wedge} 0 < m \ensuremath{\wedge}\newline m \ensuremath{\in} productsOfSubsets n \ensuremath{\wedge} p\textasciicircum{}5*m \ensuremath{\in} productsOfSubsets n \ensuremath{\wedge} p\textasciicircum{}4*m \ensuremath{\notin} productsOfSubsets n\end{quote}
\subsection{SubsetProductGap.interpolation\_conjecture\_false}
Source: \texttt{proofs/a060957/OpHack/SubsetProducts/Main.lean}, line 97. SHA-256 \texttt{5b9d85abd44f036bbbfa11f241b62ebe8dce01a2ffd89305caa965669c72b0fc}.

\begin{quote}\small\ttfamily theorem interpolation\_conjecture\_false : \ensuremath{\neg} (\ensuremath{\forall} (n p : \ensuremath{\mathbb{N}}), p.Prime \ensuremath{\to} p \ensuremath{\leq} n \ensuremath{\to}\newline \ensuremath{\forall} (m a : \ensuremath{\mathbb{N}}), m \ensuremath{\in} productsOfSubsets n \ensuremath{\to} p\textasciicircum{}a*m \ensuremath{\in} productsOfSubsets n \ensuremath{\to}\newline \ensuremath{\forall} k : \ensuremath{\mathbb{N}}, 0 < k \ensuremath{\to} k < a \ensuremath{\to} p\textasciicircum{}k*m \ensuremath{\in} productsOfSubsets n)\end{quote}
\subsection{SubsetProductGap.Generic.arbitrary\_gap\_eight\_levels\_above}
Source: \texttt{proofs/a060957/OpHack/SubsetProducts/Generic/Final.lean}, line 90. SHA-256 \texttt{a2e26af97d551980ecca1eb62f5db1362f75ada87f28795f99fe526651007d53}.

\begin{quote}\small\ttfamily theorem arbitrary\_gap\_eight\_levels\_above (g : \ensuremath{\mathbb{N}}) (hg : 1 \ensuremath{\leq} g) (lower : \ensuremath{\mathbb{N}}) :\newline \ensuremath{\exists} n p d h : \ensuremath{\mathbb{N}}, p.Prime \ensuremath{\wedge} lower < p \ensuremath{\wedge} p \ensuremath{\leq} n \ensuremath{\wedge} 0 < d \ensuremath{\wedge} \ensuremath{\neg} p \ensuremath{\mid} d \ensuremath{\wedge}\newline \ensuremath{\forall} e : \ensuremath{\mathbb{N}}, p\textasciicircum{}e * d \ensuremath{\in} productsOfSubsets n \ensuremath{\leftrightarrow}\newline e = h \ensuremath{\vee} e = h + 1 \ensuremath{\vee} e = h + 2 \ensuremath{\vee} e = h + 3 \ensuremath{\vee}\newline e = h + g + 4 \ensuremath{\vee} e = h + g + 5 \ensuremath{\vee} e = h + g + 6 \ensuremath{\vee} e = h + g + 7\end{quote}
\subsection{SubsetProductGap.Generic.arbitrary\_gap\_consecutive\_above}
Source: \texttt{proofs/a060957/OpHack/SubsetProducts/Generic/Final.lean}, line 113. SHA-256 \texttt{a2e26af97d551980ecca1eb62f5db1362f75ada87f28795f99fe526651007d53}.

\begin{quote}\small\ttfamily theorem arbitrary\_gap\_consecutive\_above (g : \ensuremath{\mathbb{N}}) (hg : 1 \ensuremath{\leq} g) (lower : \ensuremath{\mathbb{N}}) :\newline \ensuremath{\exists} n p M : \ensuremath{\mathbb{N}}, p.Prime \ensuremath{\wedge} lower < p \ensuremath{\wedge} p \ensuremath{\leq} n \ensuremath{\wedge} 0 < M \ensuremath{\wedge}\newline M \ensuremath{\in} productsOfSubsets n \ensuremath{\wedge} p\textasciicircum{}(g + 1) * M \ensuremath{\in} productsOfSubsets n \ensuremath{\wedge}\newline \ensuremath{\forall} j : \ensuremath{\mathbb{N}}, 1 \ensuremath{\leq} j \ensuremath{\to} j \ensuremath{\leq} g \ensuremath{\to} p\textasciicircum{}j * M \ensuremath{\notin} productsOfSubsets n\end{quote}
\begin{samepage}
\subsection{Verification evidence}
Combined 45-module 4.33.1 replay and aggregate passed 2026-10-01 around 16:08UTC. Generalized24-module closure also passed 4.34.1 around 16:01 and 4.33.1 around 16:05. Standard final axioms; no unproved graph-existence premise.

Recorded environments:

\begin{itemize}
\item \{'lean': '4.33.1', 'lean\_commit': '819816b2e0a3bf405af45ae5c7af2491d8f5bee6', 'mathlib': '0df444a360eaa60ab8c11dca51a86af692955474'\}
\end{itemize}
\begin{itemize}
\item \{'lean': '4.34.1', 'lean\_commit': '5045d0056413266e57c625dcd7c365b10e377c52', 'mathlib': 'd13f23b723b8a846827a245b89c10fc7d3f11612'\}
\end{itemize}
\begin{itemize}
\item \texttt{proofs/a060957/verification/full-verification.json} (existing retained artifact; SHA-256 \texttt{46dea1e382439f2896ecfac4d899f321d8c2b7755352bd0b5ad54b877619ea26}).
\end{itemize}
\begin{itemize}
\item \texttt{proofs/a060957/verification/arbitrary-gaps-4341-verification.json} (existing retained artifact; SHA-256 \texttt{296169581dd7ed22fe3f0c0a30762f5d31e3a32c8ce54a22929c6bc0007c45d9}).
\end{itemize}
\begin{itemize}
\item \texttt{proofs/a060957/verification/arbitrary-gaps-pinned-verification.json} (existing retained artifact; SHA-256 \texttt{84f27897af19af13b3321b158f2eeec611688b230fc1e4f1e2615933a803b0bb}).
\end{itemize}
\begin{itemize}
\item \texttt{proofs/a060957/verification/full-replay/OpHack\_SubsetProducts\_Main.log} (existing retained artifact; SHA-256 \texttt{37d666456acd4e1d756ed767265f43576bb7db3e21bc8596f1ae8d575a88577e}).
\end{itemize}
\begin{itemize}
\item \texttt{proofs/a060957/verification/full-replay/OpHack\_SubsetProducts\_Generic\_Final.log} (existing retained artifact; SHA-256 \texttt{45181294c9ab47a9ea8e4b2fbadd87d3d616b8e6536d1d10457061760892a767}).
\end{itemize}
\end{samepage}
\clearpage
\section{a 100475: preserved statement and proof architecture}
\subsection{A100475: no positive prime-reversal cycle}
Completed result and evidence record.

Claim ID: \texttt{a100475}. Family: \texttt{oeis-a100475}.

\subsection{Our contribution}
We supply the complete positive-start prime-reversal nonperiodicity proof, joining analytic growth with checked finite exceptions to rule out every periodic tail, including the start one. The independent 84-module Lean 4.33.1 replay verifies the exact original target with standard axioms and source identity.

\subsection{Exact preserved statement and scope}
Let p\_n be the nth prime with p\_1=2, and T(n) the decimal digit reversal of p\_n. For every natural x>0, the orbit starting at x under T is not ultimately periodic. Thus False iff there exists x>0, x\ensuremath{\neq}1, whose orbit is ultimately periodic.

Complete refutation of the positive-start original question

Not claimed:

\begin{itemize}
\item Zero-start totalization as new mathematics
\end{itemize}
\begin{itemize}
\item Separate credit for x=1, sieve modules or intermediate growth estimates
\end{itemize}
\subsection{Contribution and prior overlap}
Proposed classification: M3A complete original-target refutation; M1 only if organizer assigns focus status. No acceptance or credit asserted.

\begin{itemize}
\item Rules out every positive periodic tail, joining analytic growth with checked finite exceptions
\end{itemize}
Recorded overlap:

\begin{itemize}
\item Nth-prime growth estimates, decimal reversal facts and finite-exception strategy are classical background
\end{itemize}
\begin{itemize}
\item The public zero-start counterexample concerns a different totalization
\end{itemize}
Recorded bounded prior-art comparison:

\begin{itemize}
\item No exact prior positive-start solution located in recorded bounded review; elementary/folklore risk remains
\end{itemize}
\begin{itemize}
\item Classical derivability alone is not proof that this exact target was known
\end{itemize}
\subsection{Formal proof source}
Paths below are relative to \texttt{Claim\_Evidence/} in this bundle. Exact source bytes, declaration lines and verification evidence are preserved. Dependency versions and the inspection/build boundary are described in the bundle README.

\subsection{OeisA100475.not\_periodic}
Source: \texttt{proofs/a100475/A100475.lean}, line 17. SHA-256 \texttt{e22c200d87ed0a15a25d4760e3f37aaf0005b48aa5b734e79fb0776599c04c95}.

\begin{quote}\small\ttfamily theorem not\_periodic \{x : \ensuremath{\mathbb{N}}\} (hx : 0 < x) : \ensuremath{\neg} IsUltimatelyPeriodic (aStartAt x)\end{quote}
\subsection{OeisA100475.conjecture}
Source: \texttt{proofs/a100475/A100475.lean}, line 20. SHA-256 \texttt{e22c200d87ed0a15a25d4760e3f37aaf0005b48aa5b734e79fb0776599c04c95}.

\begin{quote}\small\ttfamily theorem conjecture :\newline False \ensuremath{\leftrightarrow} \ensuremath{\exists} x : \ensuremath{\mathbb{N}}, 0 < x \ensuremath{\wedge} x \ensuremath{\neq} 1 \ensuremath{\wedge} IsUltimatelyPeriodic (aStartAt x)\end{quote}
\subsection{Verification evidence}
Independent complete 84-module replay completed 2026-10-01T13:03:31.027744Z; source receipt records byte identity with user ZIP; final theorem standard axioms

Recorded environments:

\begin{itemize}
\item \{'lean': '4.33.1', 'lean\_commit': '819816b2e0a3bf405af45ae5c7af2491d8f5bee6', 'mathlib': '0df444a360eaa60ab8c11dca51a86af692955474'\}
\end{itemize}
\begin{itemize}
\item \texttt{evidence/lean-cloud/user-attempt-replay/A100475-independent-replay-summary.json} (existing retained artifact; SHA-256 \texttt{85d0e4a2f07f11f1c95ba2eccbf2a2d37e832259bbd6c0bc5a9cc6b98fcb3c86}).
\end{itemize}
\begin{itemize}
\item \texttt{evidence/lean-cloud/user-attempt-replay/reports/A100475-replay.json} (existing retained artifact; SHA-256 \texttt{f217c891d59ecd2757c11589cc9ace77cc593e34f10c21f8528f4b96cb108ce8}).
\end{itemize}
\begin{itemize}
\item \texttt{evidence/lean-cloud/user-attempt-replay/reports/build-a100475/A100475.log} (existing retained artifact; SHA-256 \texttt{f881ba28f46de6fd2b5da9e93c16bd6e38d4f08808a714e83b92da4630eb2839}).
\end{itemize}
\begin{itemize}
\item \texttt{evidence/lean-cloud/user-attempt-replay/reports/A100475-source-correspondence.json} (existing retained artifact; SHA-256 \texttt{4483b4b94ad2ff180523aa5206f32256357248cd1a98563df0428b04c54bceb6}).
\end{itemize}
\clearpage
\section{a 000224: preserved statement and proof architecture}
\subsection{A000224: verified squarefree three-prime exclusion}
Completed result and evidence record.

Claim ID: \texttt{a000224}. Family: \texttt{oeis-a000224}.

\subsection{Our contribution}
We supply the all-size exclusion for squarefree products of three distinct odd primes, combining global prime bounds, the exact square-residue-count bridge, and exhaustive quotient/discriminant rejection. The 171-module checks on Lean 4.34.1 and 4.33.1 verify this unbounded family with standard axioms, extending beyond the previously bounded searches and known simpler factorization cases.

\subsection{Exact preserved statement and scope}
For every three primes 2<p<q<r, let n=pqr and let a(n) count all square residues modulo n, including zero, exactly as in the upstream definition. Then a(n)(a(n)-1) does not divide n\ensuremath{^2}-1. Consequently the original conjectured equivalence holds for this entire unbounded factorization family.

Infinite partial case of the original equivalence: every squarefree product of three distinct odd primes

Not claimed:

\begin{itemize}
\item Full conjecture
\end{itemize}
\begin{itemize}
\item Repeated prime factors or four-or-more distinct prime factors
\end{itemize}
\begin{itemize}
\item A new finite search range or independent scores for certificate blocks
\end{itemize}
\subsection{Contribution and prior overlap}
Proposed classification: M3A original-target partial; M1 only if organizer assigns focus status. No acceptance or credit asserted.

\begin{itemize}
\item All-size squarefree-pqr exclusion using proved global p\ensuremath{\leq}381, p-dependent q bound, exact residue-count bridge and complete finite quotient/discriminant rejection
\end{itemize}
Recorded overlap:

\begin{itemize}
\item CRT square-residue formula, quotient/Pell strategy and gcd obstruction
\end{itemize}
\begin{itemize}
\item Prior unbounded even, prime-power and semiprime exclusions
\end{itemize}
\begin{itemize}
\item Prior Pell computation n\ensuremath{\leq}10\textasciicircum{}18; all 630 local prime-r scan survivors lie within that old bound
\end{itemize}
Recorded bounded prior-art comparison:

\begin{itemize}
\item Fresh primary source comparison found no global pqr theorem in the closest prior project or all four available Epoch A000224 Specs
\end{itemize}
\begin{itemize}
\item The prior Pell n cutoff is operational, not merely a report range
\end{itemize}
\begin{itemize}
\item Git dates establish provenance, not independently certified first-public times
\end{itemize}
\subsection{Formal proof source}
Paths below are relative to \texttt{Claim\_Evidence/} in this bundle. Exact source bytes, declaration lines and verification evidence are preserved. Dependency versions and the inspection/build boundary are described in the bundle README.

\subsection{OeisA224.not\_divisible\_three\_primes}
Source: \texttt{proofs/a000224/OpHack/QuadraticResidue224/Target224.lean}, line 33. SHA-256 \texttt{3b67f7494439886af3359a9e684433f49a6dccb7d10258d1d05dcc01de7886d0}.

\begin{quote}\small\ttfamily theorem not\_divisible\_three\_primes \{p q r : \ensuremath{\mathbb{N}}\}\newline (hp : p.Prime) (hq : q.Prime) (hr : r.Prime)\newline (hp2 : 2 < p) (hpq : p < q) (hqr : q < r) :\newline \ensuremath{\neg} a (p * q * r) * (a (p * q * r) - 1) \ensuremath{\mid} (p * q * r) \textasciicircum{} 2 - 1\end{quote}
\subsection{OeisA224.not\_modEq\_three\_primes}
Source: \texttt{proofs/a000224/OpHack/QuadraticResidue224/Target224.lean}, line 43. SHA-256 \texttt{3b67f7494439886af3359a9e684433f49a6dccb7d10258d1d05dcc01de7886d0}.

\begin{quote}\small\ttfamily theorem not\_modEq\_three\_primes \{p q r : \ensuremath{\mathbb{N}}\}\newline (hp : p.Prime) (hq : q.Prime) (hr : r.Prime)\newline (hp2 : 2 < p) (hpq : p < q) (hqr : q < r) :\newline \ensuremath{\neg} (p * q * r) \textasciicircum{} 2 \ensuremath{\equiv} 1 [MOD a (p * q * r) * (a (p * q * r) - 1)]\end{quote}
\subsection{OeisA224.conjecture\_three\_primes}
Source: \texttt{proofs/a000224/OpHack/QuadraticResidue224/Target224.lean}, line 50. SHA-256 \texttt{3b67f7494439886af3359a9e684433f49a6dccb7d10258d1d05dcc01de7886d0}.

\begin{quote}\small\ttfamily theorem conjecture\_three\_primes \{p q r : \ensuremath{\mathbb{N}}\}\newline (hp : p.Prime) (hq : q.Prime) (hr : r.Prime)\newline (hp2 : 2 < p) (hpq : p < q) (hqr : q < r) :\newline ((p * q * r).Prime \ensuremath{\wedge} p * q * r \ensuremath{\neq} 2) \ensuremath{\leftrightarrow}\newline (p * q * r) \textasciicircum{} 2 \ensuremath{\equiv} 1 [MOD a (p * q * r) * (a (p * q * r) - 1)]\end{quote}
\subsection{Verification evidence}
All 171 intended modules were recorded checked on 4.34.1, dependency-ordered and incremental; nine prime modules were reused after source/output hash checks. A separate complete 4.33.1 replay passed all 171 modules, with 164 successful certificate records and seven retrospective core source/log/exit receipts. Both versions report only propext, Classical.choice and Quot.sound for all three endpoints. The 4.33.1 replay reused no own compiled artifacts across versions; 170 source files are byte-identical and Bounds.lean has precisely three tactic-only compatibility changes, with no definition or theorem-statement changes. The core receipt creation timestamps are retrospective, not original build times. This is a second version of Lean, not an independent compiler implementation or external referee check. The existing target definition is whitespace-normalized identical to the immutable upstream snapshot. The convenience aggregate OpHack.lean wrapper is not separately checked.

Recorded environments:

\begin{itemize}
\item \{'lean': '4.34.1', 'lean\_commit': '5045d0056413266e57c625dcd7c365b10e377c52', 'mathlib': 'd13f23b723b8a846827a245b89c10fc7d3f11612'\}
\end{itemize}
\begin{itemize}
\item \{'lean': '4.33.1', 'lean\_commit': '819816b2e0a3bf405af45ae5c7af2491d8f5bee6', 'mathlib': '0df444a360eaa60ab8c11dca51a86af692955474', 'compatibility\_payload': 'proofs/a000224/compatibility/4.33.1'\}
\end{itemize}
\begin{itemize}
\item \texttt{proofs/a000224/verification/4.34.1/STATUS.json} (existing retained artifact; SHA-256 \texttt{ebea4900dae174bba482e26ed0ed850c8ea390346a133abd3a63b2f037fc8f8e}).
\end{itemize}
\begin{itemize}
\item \texttt{proofs/a000224/verification/4.34.1/Target224.log} (existing retained artifact; SHA-256 \texttt{36c0fd87718f09eed6fd01b1f7bee3ce072a3f1fcfd8ad1e77d3cb7d6068471e}).
\end{itemize}
\begin{itemize}
\item \texttt{proofs/a000224/verification/4.34.1/core-receipts.json} (existing retained artifact; SHA-256 \texttt{c60f5663cc9e13685ce2bb5e2e9267acb4f2775fce934978db73a1a69749e18e}).
\end{itemize}
\begin{itemize}
\item \texttt{proofs/a000224/target-source/CORRESPONDENCE.json} (existing retained artifact; SHA-256 \texttt{be9ed8176e26b1643f1c5e21825de4dfb3f55fde097b72de3af33e4d6fc74cbd}).
\end{itemize}
\begin{itemize}
\item \texttt{proofs/a000224/target-source/OEIS224.lean.txt} (existing retained artifact; SHA-256 \texttt{bd6cf324b2b352f1344b5b39922adc6f4029bb5ec2c7bf6c5c3cc80c6c5b2665}).
\end{itemize}
\begin{itemize}
\item \texttt{proofs/a000224/verification/4.33.1/STATUS.json} (existing retained artifact; SHA-256 \texttt{1f3f43db68866bdfd54d78ccb3e397ea2c343111eb10fafbbc84bc141418454b}).
\end{itemize}
\begin{itemize}
\item \texttt{proofs/a000224/verification/4.33.1/Target224.log} (existing retained artifact; SHA-256 \texttt{36c0fd87718f09eed6fd01b1f7bee3ce072a3f1fcfd8ad1e77d3cb7d6068471e}).
\end{itemize}
\begin{itemize}
\item \texttt{proofs/a000224/verification/4.33.1/core-receipts.json} (existing retained artifact; SHA-256 \texttt{85fc90e23fb1a99d7c20a230f7a255d8e224fb2e1b12d4f1cd0c99d994b74f67}).
\end{itemize}
\begin{itemize}
\item \texttt{proofs/a000224/verification/4.33.1/split-build-results.json} (existing retained artifact; SHA-256 \texttt{3ac1a43daf20e8eb80e04887cfcea0474bbd8643196fe551cdce664b1fe9a2b5}).
\end{itemize}
\begin{itemize}
\item \texttt{proofs/a000224/compatibility/4.33.1/BUILD\_ORDER.json} (existing retained artifact; SHA-256 \texttt{766b4f3d71191a5c7e8a48a05c0ddf6a54e1282925e07857d8b406751537cb48}).
\end{itemize}
\begin{itemize}
\item \texttt{proofs/a000224/compatibility/4.33.1/Bounds.patch} (existing retained artifact; SHA-256 \texttt{bf0e33c707b570e7c4d0c48cc622339a9aa1fbdf1b2c358e5d2414057aceec31}).
\end{itemize}
\clearpage
\section{ferrers-hierarchy: preserved statement and proof architecture}
\subsection{Ferrers general deficiency hierarchy (44-module dual replay)}
Completed result and evidence record.

Claim ID: \texttt{ferrers-hierarchy}. Family: \texttt{partial-ferrers}.

\subsection{Our contribution}
We supply the arbitrary-rank deficiency hierarchy that preserves every supplied optimal factor while adjoining the stated matching colors, with exact later ranks and balanced and one-defect realizations. The 44-source dual-version replays verify the graph/rank transport with standard axioms; known greedy boundary and existential cases remain credited within this broader preservation theorem.

\subsection{Exact preserved statement and scope}
For 1\ensuremath{\leq}r\ensuremath{\leq}k and 0\ensuremath{\leq}e<r, a finite Ferrers graph with matching-union ranks \ensuremath{\alpha}\_r=rk\ensuremath{-}e and \ensuremath{\alpha}\_(r+1)=(r+1)k\ensuremath{-}r\ensuremath{-}e has \ensuremath{\alpha}\_1=k. Every supplied optimal degree-r factor is retained while r\ensuremath{-}e disjoint matching colors of size k\ensuremath{-}r are appended. For 0\ensuremath{\leq}h\ensuremath{\leq}r\ensuremath{-}e, \ensuremath{\alpha}\_(r+h)=rk\ensuremath{-}e+h(k\ensuremath{-}r). Balanced e=0 yields simultaneous first 2r CDS colors; e=1 yields first 2r\ensuremath{-}1. For e>1 no individual first-r CDS-color realization is asserted for arbitrary supplied colorings.

General restricted hierarchy and factor-preserving extension; partial relative to unrestricted Latin Tableaux

Not claimed:

\begin{itemize}
\item Full Latin Tableau conjecture or unrestricted L5
\end{itemize}
\begin{itemize}
\item Individual first-r CDS colors for arbitrary supplied colorings when e>1
\end{itemize}
\begin{itemize}
\item Newness of greedy e=r\ensuremath{-}1 boundary or all small known existence cases
\end{itemize}
\begin{itemize}
\item Parameterized sharpness or rectangular missing-corner transport as Lean-checked in this 44-module closure
\end{itemize}
\subsection{Contribution and prior overlap}
Proposed classification: M3A parent partial or M3B explicit parent-linked refinement, for organizer classification. No acceptance or credit asserted.

\begin{itemize}
\item For arbitrary r and e<r, preserve every supplied optimal r-factor while adjoining r\ensuremath{-}e matching colors, together with exact later ranks and special balanced/one-defect realizations
\end{itemize}
Recorded overlap:

\begin{itemize}
\item Corner/min-cut formulas, capacitated Hall, bipartite edge coloring
\end{itemize}
\begin{itemize}
\item Greedy e=r\ensuremath{-}1 extension
\end{itemize}
\begin{itemize}
\item r=3,e=1 existential L5 scope overlaps September 2026 McAuley claim; arbitrary supplied-factor preservation is stronger than that stated existence claim
\end{itemize}
Recorded bounded prior-art comparison:

\begin{itemize}
\item No inspected result found to entail the full arbitrary-r preservation hierarchy; bounded negative source audit
\end{itemize}
\begin{itemize}
\item Subtract the known boundary and old existence results explicitly
\end{itemize}
\subsection{Formal proof source}
Paths below are relative to \texttt{Claim\_Evidence/} in this bundle. Exact source bytes, declaration lines and verification evidence are preserved. Dependency versions and the inspection/build boundary are described in the bundle README.

\subsection{FerrersArm.TightDefectCorner.defect\_preserves\_factor}
Source: \texttt{proofs/ferrers/OpHack/Ferrers/DefectExtension.lean}, line 104. SHA-256 \texttt{c17d8534968c3ac199813212046ba0da17c65fa46e2907460afeca99323b1e20}.

\begin{quote}\small\ttfamily theorem defect\_preserves\_factor \{n W : \ensuremath{\mathbb{N}}\} (\ensuremath{\ell} : Fin n \ensuremath{\to} \ensuremath{\mathbb{N}}) (h\ensuremath{\ell} : Antitone \ensuremath{\ell})\newline (r k e : \ensuremath{\mathbb{N}}) (hr : 1 \ensuremath{\leq} r) (hk : r \ensuremath{\leq} k) (he : e < r)\newline (hR : factorRank (W := W) \ensuremath{\ell} r = r*k-e)\newline (hNext : factorRank (W := W) \ensuremath{\ell} (r+1) = (r+1)*k-r-e)\newline (F : Finset (Fin n \ensuremath{\times} Fin W)) (hF : Feasible \ensuremath{\ell} r F) (hcard : F.card = r*k-e) :\newline \ensuremath{\exists} G : Fin (r-e) \ensuremath{\to} Finset (Fin n \ensuremath{\times} Fin W),\newline (\ensuremath{\forall} c, IsMatching (G c) \ensuremath{\wedge} (G c).card = k-r \ensuremath{\wedge}\newline \ensuremath{\forall} p \ensuremath{\in} G c, Allowed \ensuremath{\ell} p.1 p.2 \ensuremath{\wedge} p \ensuremath{\notin} F) \ensuremath{\wedge}\newline (\ensuremath{\forall} c d, c \ensuremath{\neq} d \ensuremath{\to} Disjoint (G c) (G d)) \ensuremath{\wedge}\newline (\ensuremath{\forall} h \ensuremath{\leq} r-e, factorRank (W := W) \ensuremath{\ell} (r+h) = (r+h)*k-h*r-e) \ensuremath{\wedge}\newline factorRank (W := W) \ensuremath{\ell} 1 = k\end{quote}
\subsection{FerrersArm.YoungDiagramBridge.defect\_indepNumK\_later\_ranks}
Source: \texttt{proofs/ferrers/OpHack/Ferrers/YoungDiagramDefect.lean}, line 8. SHA-256 \texttt{c55685043c4b319f834184a218377c836af11200769eeae94fda76c4fc2c2dbc}.

\begin{quote}\small\ttfamily theorem defect\_indepNumK\_later\_ranks (\ensuremath{\mu} : YoungDiagram) (r k e : \ensuremath{\mathbb{N}})\newline (hr : 1 \ensuremath{\leq} r) (hk : r \ensuremath{\leq} k) (he : e < r)\newline (hR : FCCompat.indepNumK (FCCompat.YoungDiagram.toSimpleGraph \ensuremath{\mu}) r = r*k-e)\newline (hNext : FCCompat.indepNumK (FCCompat.YoungDiagram.toSimpleGraph \ensuremath{\mu})\newline (r+1) = (r+1)*k-r-e) :\newline \ensuremath{\forall} h \ensuremath{\leq} r-e, FCCompat.indepNumK (FCCompat.YoungDiagram.toSimpleGraph \ensuremath{\mu})\newline (r+h) = (r+h)*k-h*r-e\end{quote}
\subsection{FerrersArm.YoungDiagramBridge.one\_defect\_independent\_sets}
Source: \texttt{proofs/ferrers/OpHack/Ferrers/YoungDiagramDefect.lean}, line 42. SHA-256 \texttt{c55685043c4b319f834184a218377c836af11200769eeae94fda76c4fc2c2dbc}.

\begin{quote}\small\ttfamily theorem one\_defect\_independent\_sets (\ensuremath{\mu} : YoungDiagram) (r k : \ensuremath{\mathbb{N}})\newline (hr : 2 \ensuremath{\leq} r) (hk : r \ensuremath{\leq} k)\newline (hR : FCCompat.indepNumK (FCCompat.YoungDiagram.toSimpleGraph \ensuremath{\mu}) r = r*k-1)\newline (hNext : FCCompat.indepNumK (FCCompat.YoungDiagram.toSimpleGraph \ensuremath{\mu})\newline (r+1) = (r+1)*k-r-1) :\newline \ensuremath{\exists} (R : Fin r \ensuremath{\to} Finset (FCCompat.YoungDiagram.Cell \ensuremath{\mu}))\newline (G : Fin (r-1) \ensuremath{\to} Finset (FCCompat.YoungDiagram.Cell \ensuremath{\mu})),\newline (\ensuremath{\forall} c, (FCCompat.YoungDiagram.toSimpleGraph \ensuremath{\mu}).IsIndepSet\newline (R c : Set (FCCompat.YoungDiagram.Cell \ensuremath{\mu})) \ensuremath{\wedge}\newline (R c).card = (if c.val < r-1 then k else k-1)) \ensuremath{\wedge}\newline (\ensuremath{\forall} c, (FCCompat.YoungDiagram.toSimpleGraph \ensuremath{\mu}).IsIndepSet\newline (G c : Set (FCCompat.YoungDiagram.Cell \ensuremath{\mu})) \ensuremath{\wedge} (G c).card = k-r) \ensuremath{\wedge}\newline (\ensuremath{\forall} c d, c \ensuremath{\neq} d \ensuremath{\to} Disjoint (R c) (R d)) \ensuremath{\wedge}\newline (\ensuremath{\forall} c d, c \ensuremath{\neq} d \ensuremath{\to} Disjoint (G c) (G d)) \ensuremath{\wedge}\newline (\ensuremath{\forall} c d, Disjoint (R c) (G d)) \ensuremath{\wedge}\newline (\ensuremath{\forall} t \ensuremath{\leq} 2*r-1, FCCompat.indepNumK (FCCompat.YoungDiagram.toSimpleGraph \ensuremath{\mu}) t =\newline if t < r then t*k else t*k-(t-r)*r-1)\end{quote}
\subsection{FerrersArm.YoungDiagramBridge.first\_two\_r\_independent\_sets}
Source: \texttt{proofs/ferrers/OpHack/Ferrers/YoungDiagramRestrictedPrefix.lean}, line 158. SHA-256 \texttt{e5e0e0f73933f360a6fc87e96d2e6172b791acff5ba73aaa0ad0ff836e83654e}.

\begin{quote}\small\ttfamily theorem first\_two\_r\_independent\_sets (\ensuremath{\mu} : YoungDiagram) (r k : \ensuremath{\mathbb{N}})\newline (hr : 1 \ensuremath{\leq} r) (hk : r \ensuremath{\leq} k)\newline (hR : FCCompat.indepNumK (FCCompat.YoungDiagram.toSimpleGraph \ensuremath{\mu}) r = r * k)\newline (hNext : FCCompat.indepNumK (FCCompat.YoungDiagram.toSimpleGraph \ensuremath{\mu})\newline (r + 1) = (r + 1) * k - r) :\newline \ensuremath{\exists} R G : Fin r \ensuremath{\to} Finset (FCCompat.YoungDiagram.Cell \ensuremath{\mu}),\newline (\ensuremath{\forall} c, (FCCompat.YoungDiagram.toSimpleGraph \ensuremath{\mu}).IsIndepSet\newline (R c : Set (FCCompat.YoungDiagram.Cell \ensuremath{\mu})) \ensuremath{\wedge} (R c).card = k) \ensuremath{\wedge}\newline (\ensuremath{\forall} c, (FCCompat.YoungDiagram.toSimpleGraph \ensuremath{\mu}).IsIndepSet\newline (G c : Set (FCCompat.YoungDiagram.Cell \ensuremath{\mu})) \ensuremath{\wedge} (G c).card = k - r) \ensuremath{\wedge}\newline (\ensuremath{\forall} c d, c \ensuremath{\neq} d \ensuremath{\to} Disjoint (R c) (R d)) \ensuremath{\wedge}\newline (\ensuremath{\forall} c d, c \ensuremath{\neq} d \ensuremath{\to} Disjoint (G c) (G d)) \ensuremath{\wedge}\newline (\ensuremath{\forall} c d, Disjoint (R c) (G d)) \ensuremath{\wedge}\newline (\ensuremath{\forall} t \ensuremath{\leq} 2 * r,\newline (firstCells R (min t r) \ensuremath{\cup} firstCells G (t - r)).card =\newline FCCompat.indepNumK (FCCompat.YoungDiagram.toSimpleGraph \ensuremath{\mu}) t)\end{quote}
\subsection{Verification evidence}
44-source strict replay receipts on both pins; source hashes match both recorded receipts. Actual YoungDiagram graph/rank transport and all selected final standard-axiom audits included.

Recorded environments:

\begin{itemize}
\item \{'lean': '4.33.1', 'lean\_commit': '819816b2e0a3bf405af45ae5c7af2491d8f5bee6', 'mathlib': '0df444a360eaa60ab8c11dca51a86af692955474'\}
\end{itemize}
\begin{itemize}
\item \{'lean': '4.34.1', 'lean\_commit': '5045d0056413266e57c625dcd7c365b10e377c52', 'mathlib': 'd13f23b723b8a846827a245b89c10fc7d3f11612'\}
\end{itemize}
\begin{itemize}
\item \texttt{proofs/ferrers/verification/4.33.1/full-verification.json} (existing retained artifact; SHA-256 \texttt{139d8b50bafa2dd7b0d715eec830882a005b15a0d8c3e2feaf1077a56188acdb}).
\end{itemize}
\begin{itemize}
\item \texttt{proofs/ferrers/verification/4.34.1/full-verification.json} (existing retained artifact; SHA-256 \texttt{a44d60930c974cf27ee8b0e3a7763ea9af8c06c79053fd46645ee31b82dc6257}).
\end{itemize}
\begin{itemize}
\item \texttt{proofs/ferrers/verification/source-inventory.json} (existing retained artifact; SHA-256 \texttt{0f6d9be1e409b8c22aaec1c5fc8a352b9597aaadd777d61cb250f6bc1a2195e2}).
\end{itemize}
\clearpage
\section{erdos21-sharp-residual: preserved statement and proof architecture}
\subsection{Erdős Lovász sharp residual cover bound}
Completed result and evidence record.

Claim ID: \texttt{erdos21-sharp-residual}. Family: \texttt{erdos21-sharp-residual-cover}.

\subsection{Our contribution}
We close the source paper's explicit +3 residual-cover question through refined colored-link budgets and equality-case exclusions, then derives the all-r bound with additive constant three. The isolated 18-module Lean 4.34.1 replay and separate triangle check certify the result and sharpness, building on the credited matching infrastructure and +4 theorem.

\subsection{Exact preserved statement and scope}
For every finite intersecting r-uniform hypergraph J of maximum vertex degree at most 3, 4 tau(J) <= |J| + r + 3. Degree-four peeling yields the same inequality without a degree restriction and hence 3r <= |H| + 3 whenever tau(H)=r; for r=8 this gives |H|>=21. The additive constant 3 is optimal, with a separately checked known triangle example.

Complete proof of the explicit +3 refinement asked after Lemma 2 of arXiv 2606.24878v2; one provisional non-M2 new-target mathematical family.

Not claimed:

\begin{itemize}
\item No resolution of the original O(r) Erdős 21 question, already solved by Kahn
\end{itemize}
\begin{itemize}
\item No determination of g(r) or improvement of the source paper asymptotic coefficient (41-sqrt(19))/12
\end{itemize}
\begin{itemize}
\item No separate claim for lower\_three\_sub, r=8, triangle sharpness or infrastructure
\end{itemize}
\begin{itemize}
\item No certified first-publication priority
\end{itemize}
\subsection{Contribution and prior overlap}
Proposed classification: Provisional non-M2 new-target mathematical claim. No acceptance or credit asserted.

\begin{itemize}
\item The source paper's residual bound has +4 and explicitly asks for +3. Refined colored-link edge budgets and finite equality-case exclusions establish +3, and the same family improves its explicit all-r bound from 3r-4 to 3r-3.
\end{itemize}
Recorded overlap:

\begin{itemize}
\item Varun Sivashankar, An Improved Lower Bound for the Erdős-Lovász Cover Number Problem, arXiv 2606.24878v2, 29 July 2026 supplies the unconditional matching/witness infrastructure, +4 theorem and known sharpness examples.
\end{itemize}
\begin{itemize}
\item The paper's stronger asymptotic coefficient about 3.053 is prior work and is not improved. Its ancillary source prefix is reused with explicit attribution and only the documented Lean 4.34.1 compatibility changes. The entire Kahn-dependent asymptotic section is excluded.
\end{itemize}
Recorded bounded prior-art comparison:

\begin{itemize}
\item Varun Sivashankar, An Improved Lower Bound for the Erdős-Lovász Cover Number Problem, arXiv 2606.24878v2, 29 July 2026 supplies the unconditional matching/witness infrastructure, +4 theorem and known sharpness examples.
\end{itemize}
\begin{itemize}
\item The paper's stronger asymptotic coefficient about 3.053 is prior work and is not improved. Its ancillary source prefix is reused with explicit attribution and only the documented Lean 4.34.1 compatibility changes. The entire Kahn-dependent asymptotic section is excluded.
\end{itemize}
\subsection{Formal proof source}
Paths below are relative to \texttt{Claim\_Evidence/} in this bundle. Exact source bytes, declaration lines and verification evidence are preserved. Dependency versions and the inspection/build boundary are described in the bundle README.

\subsection{ErdosLovasz.residual\_cover\_bound\_sharp}
Source: \texttt{proofs/erdos21-sharp-residual/OpHack/Erdos21/Main.lean}, line 25. SHA-256 \texttt{9314ae7ba95b6fd9922ccfea882f1d33a6c11c9a40d750109014c8184ab7e414}.

\begin{quote}\small\ttfamily theorem residual\_cover\_bound\_sharp (r : \ensuremath{\mathbb{N}}) (J : Finset (Finset V))\newline (huni : IsUniform r J) (hint : Intersecting J)\newline (hdeg : \ensuremath{\forall} v, edgeDeg J v \ensuremath{\leq} 3) :\newline 4 * coverNum J \ensuremath{\leq} J.card + r + 3\end{quote}
\subsection{ErdosLovasz.cover\_bound\_sharp}
Source: \texttt{proofs/erdos21-sharp-residual/OpHack/Erdos21/Main.lean}, line 33. SHA-256 \texttt{9314ae7ba95b6fd9922ccfea882f1d33a6c11c9a40d750109014c8184ab7e414}.

\begin{quote}\small\ttfamily theorem cover\_bound\_sharp (r : \ensuremath{\mathbb{N}}) (H : Finset (Finset V))\newline (huni : IsUniform r H) (hint : Intersecting H) :\newline 4 * coverNum H \ensuremath{\leq} H.card + r + 3\end{quote}
\subsection{ErdosLovasz.erdos\_lovasz\_lower\_three}
Source: \texttt{proofs/erdos21-sharp-residual/OpHack/Erdos21/Main.lean}, line 41. SHA-256 \texttt{9314ae7ba95b6fd9922ccfea882f1d33a6c11c9a40d750109014c8184ab7e414}.

\begin{quote}\small\ttfamily theorem erdos\_lovasz\_lower\_three (r : \ensuremath{\mathbb{N}}) (H : Finset (Finset V))\newline (huni : IsUniform r H) (hint : Intersecting H) (hcover : coverNum H = r) :\newline 3 * r \ensuremath{\leq} H.card + 3\end{quote}
\subsection{ErdosLovasz.erdos\_lovasz\_lower\_three\_sub}
Source: \texttt{proofs/erdos21-sharp-residual/OpHack/Erdos21/Main.lean}, line 49. SHA-256 \texttt{9314ae7ba95b6fd9922ccfea882f1d33a6c11c9a40d750109014c8184ab7e414}.

\begin{quote}\small\ttfamily theorem erdos\_lovasz\_lower\_three\_sub (r : \ensuremath{\mathbb{N}}) (H : Finset (Finset V))\newline (huni : IsUniform r H) (hint : Intersecting H) (hcover : coverNum H = r) :\newline 3 * r - 3 \ensuremath{\leq} H.card\end{quote}
\subsection{ErdosLovasz.erdos\_lovasz\_eight\_lower}
Source: \texttt{proofs/erdos21-sharp-residual/OpHack/Erdos21/Main.lean}, line 56. SHA-256 \texttt{9314ae7ba95b6fd9922ccfea882f1d33a6c11c9a40d750109014c8184ab7e414}.

\begin{quote}\small\ttfamily theorem erdos\_lovasz\_eight\_lower (H : Finset (Finset V))\newline (huni : IsUniform 8 H) (hint : Intersecting H) (hcover : coverNum H = 8) :\newline 21 \ensuremath{\leq} H.card\end{quote}
\subsection{Verification evidence}
An independent isolated Lean 4.34.1 replay passed all 18 Main-closure modules. The proof-output directory began empty, source hashes stayed unchanged, nine external revisions were pinned and clean, and the separate endpoint-definition audit passed. Final endpoints use only propext, Classical.choice and Quot.sound; no Kahn hypothesis, hidden packing lemma or extra mathematical assumption is present. The known triangle witness was separately checked against the rebuilt core, outside the 18-module replay.

Recorded environments:

\begin{itemize}
\item \{'lean': '4.34.1', 'mathlib': 'd13f23b723b8a846827a245b89c10fc7d3f11612'\}
\end{itemize}
\begin{itemize}
\item \texttt{proofs/erdos21-sharp-residual/replay-receipt/manifest.json} (existing retained artifact; SHA-256 \texttt{5dc4e129bd723dc0b914ec12a711d31ff63151dc1d56956fbab17f0098dfca2e}).
\end{itemize}
\begin{itemize}
\item \texttt{proofs/erdos21-sharp-residual/replay-receipt/Main.log} (existing retained artifact; SHA-256 \texttt{94d78809809af1486c48477b5dfb3166a48af05fa5597da2e872ac652ef0db42}).
\end{itemize}
\begin{itemize}
\item \texttt{proofs/erdos21-sharp-residual/replay-receipt/EndpointAudit.log} (existing retained artifact; SHA-256 \texttt{930313bcac1c89a8010b052ce64adfa6305e2d2579afdf6e6095bdd46c95a4d6}).
\end{itemize}
\begin{itemize}
\item \texttt{proofs/erdos21-sharp-residual/replay-receipt/TriangleSharpness.log} (existing retained artifact; SHA-256 \texttt{1433ba0035a9c0578b46db873b88a7b5bd275ed2e833700f4a3b3ca9c5f3e75e}).
\end{itemize}
\begin{itemize}
\item \texttt{proofs/erdos21-sharp-residual/verified-main-output.txt} (existing retained artifact; SHA-256 \texttt{94d78809809af1486c48477b5dfb3166a48af05fa5597da2e872ac652ef0db42}).
\end{itemize}
\clearpage
\section{erdos 169-fourap: preserved statement and proof architecture}
The retained source card describes the entrant's archived full Linux build. The original frozen-source independent Windows replay remains the dated 15/22 environmental checkpoint. A separate scientific-body-identical import-only route has now passed 22/22 authored modules and the unchanged selected final audit with standard axioms; that completed operational qualification is recorded above and does not relabel the original receipt.

\subsection{Explicit four-AP-free reciprocal sum >111/25}
Completed result and evidence record.

Claim ID: \texttt{erdos169-fourap}. Family: \texttt{partial-erdos169-fourap}.

\subsection{Our contribution}
We supply the explicit six-residue augmentation of Walker's base-55 construction and the reciprocal-sum bound exceeding 4.44. The fresh full build and axiom audit verify positivity, absence of nonconstant four-term progressions, and the exact rational-sum inequality for the same finite set; Walker's digit alphabet and earlier bound remain the credited starting point.

\subsection{Exact preserved statement and scope}
Let D=\{0,1,2,4,5,9,10,11,14,16,17,18,21,24,30,37,39,41,42,45,47\}, M=55\ensuremath{^3}=166375, U=\{26286,26726,26737,120061,120501,120512\}, and T\_j=\{\ensuremath{\Sigma}\_(i=0)\textasciicircum{}(j\ensuremath{-}1) d\_i55\textasciicircum{}i:d\_i\ensuremath{\in}D\}, T\_0=\{0\}. Let E=T\_22 \ensuremath{\cup} \ensuremath{\bigcup}\_(u\ensuremath{\in}U)(u+M T\_21), A=\{e+1:e\ensuremath{\in}E\}. Every a\ensuremath{\in}A is positive; A contains no nonconstant four-term arithmetic progression over Z; \ensuremath{\Sigma}\_(a\ensuremath{\in}A)1/a>111/25=4.44, as an exact rational finite sum over distinct elements.

Quantitative partial lower-bound construction for Erdős 169

Not claimed:

\begin{itemize}
\item Erdős 3 reciprocal-divergence theorem
\end{itemize}
\begin{itemize}
\item Full estimation problem or exact optimum
\end{itemize}
\begin{itemize}
\item Global optimality of the six-residue augmentation
\end{itemize}
\begin{itemize}
\item Additional credit for finite certificate stages or fixed-seed saturation
\end{itemize}
\subsection{Contribution and prior overlap}
Proposed classification: M3A original estimation-problem partial; M1 only if organizer assigns focus status. No acceptance or credit asserted.

\begin{itemize}
\item Explicit six-residue base 55 augmentation and exact reciprocal sum strictly above 4.44
\end{itemize}
Recorded overlap:

\begin{itemize}
\item Walker base 55 alphabet and digit descent, prior 4.43975 bound
\end{itemize}
\begin{itemize}
\item Exposed September 2026 >4.439753474215620 construction, grouped levels and generic strict improvement
\end{itemize}
Recorded bounded prior-art comparison:

\begin{itemize}
\item No earlier exact 4.44 bound found in recorded review; no absolute priority certification
\end{itemize}
\subsection{Formal proof source}
Paths below are relative to \texttt{Claim\_Evidence/} in this bundle. Exact source bytes, declaration lines and verification evidence are preserved. Dependency versions and the inspection/build boundary are described in the bundle README.

\subsection{Erdos3Candidate.explicit\_four\_ap\_free\_reciprocal\_gt}
Source: \texttt{proofs/erdos169-fourap/OpHack/MainResult.lean}, line 17. SHA-256 \texttt{cc686e49443c9f23ccbf715be6e9cf71c528736794df230771d33ec84df80ac5}.

\begin{quote}\small\ttfamily theorem explicit\_four\_ap\_free\_reciprocal\_gt :\newline (\ensuremath{\forall} a \ensuremath{\in} sixAFinset, 0 < a) \ensuremath{\wedge}\newline APFree 4 (sixAFinset : Set \ensuremath{\mathbb{Z}}) \ensuremath{\wedge}\newline (\ensuremath{\sum} a \ensuremath{\in} sixAFinset, (1 : \ensuremath{\mathbb{Q}}) / (a : \ensuremath{\mathbb{Q}})) > 111 / 25\end{quote}
\subsection{Verification evidence}
Fresh Linux full build completed 8950 jobs and final axiom audit. Exact same finite set appears in positivity, AP-free and rational-sum clauses. Mathematical source manifest matches; later README change is documentation only.

Recorded environments:

\begin{itemize}
\item \{'lean': '4.34.1', 'lean\_commit': '5045d0056413266e57c625dcd7c365b10e377c52', 'mathlib': 'd13f23b723b8a846827a245b89c10fc7d3f11612'\}
\end{itemize}
\begin{itemize}
\item \texttt{proofs/erdos169-fourap/verification/build-linux.log} (existing retained artifact; SHA-256 \texttt{2d78bf7fe8b1610192c5c266f27d933f80512c19c5c24c33d56f1a621b7bb181}).
\end{itemize}
\begin{itemize}
\item \texttt{proofs/erdos169-fourap/verification/final-axioms.log} (existing retained artifact; SHA-256 \texttt{be4399146b1c3f38a075d21f29bf9cc443471b4df55cc399b5e2916b8cf1f14a}).
\end{itemize}
\begin{itemize}
\item \texttt{proofs/erdos169-fourap/verification/final-dependencies.log} (existing retained artifact; SHA-256 \texttt{c3ba10d9a2c7e00bb80bbb095dad531bb9af23ef8fa5202f659c8c47c53526b2}).
\end{itemize}
\begin{itemize}
\item \texttt{proofs/erdos169-fourap/verification/final-toolchain.log} (existing retained artifact; SHA-256 \texttt{355a9a028f81641782d75d6e6b58dec70f29d2f61238a9e271690259d9bf5680}).
\end{itemize}
\clearpage
\section{opdp 13-quartic: preserved statement and proof architecture}
\subsection{Quartic exact exponent classification and sharp boundary scale}
Completed result and evidence record.

Claim ID: \texttt{opdp13-quartic}. Family: \texttt{partial-opdp13-quartic}.

\subsection{Our contribution}
We supply the exact exceptional-exponent classification, cubic transition, and sharp logarithmic boundary scale with leading constant one for the least positive admissible exponent. The recorded 40-proof-module Lean 4.34.1 build verifies the function-level squared-norm statements with standard axioms, refining the known quartic family, radical threshold, and qualitative eventual positivity.

\subsection{Exact preserved statement and scope}
For H\_t(z)=(1+|z|\ensuremath{^2})\textasciicircum{}4\ensuremath{-}t|z|\ensuremath{^4} and 6<t\ensuremath{\leq}7, H\_t\textasciicircum{}N (N positive integer) is a finite sum of |holomorphic polynomial|\ensuremath{^2} iff N=2, or N\ensuremath{\geq}4, or N=3 and 924\ensuremath{-}210t+18t\ensuremath{^2}\ensuremath{-}t\ensuremath{^3}\ensuremath{\geq}0. For 6<t<8 any such N obeys 12N(8\ensuremath{-}t)\ensuremath{\geq}t(10\ensuremath{-}t). For 7\ensuremath{\leq}t<8 it obeys 81N t\textasciicircum{}(4N)\ensuremath{\leq}614656(16\ensuremath{-}t)\textasciicircum{}(4N), plus explicit logarithmic lower bound. Writing \ensuremath{\delta}=8\ensuremath{-}t, for every \ensuremath{\varepsilon}>0 sufficiently small \ensuremath{\delta}>0 forces all admissible N to satisfy \ensuremath{\delta}N\ensuremath{\geq}(1\ensuremath{-}\ensuremath{\varepsilon})log(1/\ensuremath{\delta}), while every N\ensuremath{\geq}(1+\ensuremath{\varepsilon})log(1/\ensuremath{\delta})/\ensuremath{\delta} has strictly positive coefficients in every degree 0,...,4N and gives a function-level squared-norm representation. The defined least positive admissible index exists and has normalized ratio \ensuremath{\delta} index/log(1/\ensuremath{\delta}) between 1\ensuremath{-}\ensuremath{\varepsilon} and 1+\ensuremath{\varepsilon} near 0.

Precise one-variable restricted refinement of Hermitian-power problem

Not claimed:

\begin{itemize}
\item Full OPDP 13 or new t<8 radical threshold
\end{itemize}
\begin{itemize}
\item Binary homogeneous bridge as formally proved
\end{itemize}
\begin{itemize}
\item A separately defined eventual index or filter-limit theorem
\end{itemize}
\begin{itemize}
\item Separate credit for criterion, lower, upper and least-index modules
\end{itemize}
\subsection{Contribution and prior overlap}
Proposed classification: M3A parent partial or M3B substantive parent-linked quartic refinement, for organizer classification. No acceptance or credit asserted.

\begin{itemize}
\item Exact exceptional exponents/cubic transition and uniform log(1/\ensuremath{\delta})/\ensuremath{\delta} scale with leading constant 1 for the least positive admissible index
\end{itemize}
Recorded overlap:

\begin{itemize}
\item Quartic family after rescaling, t<8 threshold, t7 square, qualitative eventual positivity and general saddle method
\end{itemize}
\begin{itemize}
\item Gram/coefficient and arbitrary-summand squared-norm criteria
\end{itemize}
Recorded bounded prior-art comparison:

\begin{itemize}
\item Known family/methods are removed from new-mathematics scope
\end{itemize}
\begin{itemize}
\item Tan-To accessible appendix studies R\textasciicircum{}m P\_alpha with \ensuremath{\Theta}(1/alpha) minimal exponent, a different operation/family; it does not itself supply this H\_t\textasciicircum{}N logarithmic rate
\end{itemize}
\begin{itemize}
\item Unseen main Theorem 2.1 may have a relevant specialization; no categorical clearance asserted
\end{itemize}
\subsection{Formal proof source}
Paths below are relative to \texttt{Claim\_Evidence/} in this bundle. Exact source bytes, declaration lines and verification evidence are preserved. Dependency versions and the inspection/build boundary are described in the bundle README.

\subsection{PowerSemigroup.function\_exact\_exponent\_semigroup}
Source: \texttt{proofs/opdp13/OpHack/PowerSemigroup/Evaluation.lean}, line 116. SHA-256 \texttt{9924251f0abcd263e38621ad0b96cadbd0a33f6c353d41ecf3e16e30a860ebfc}.

\begin{quote}\small\ttfamily theorem function\_exact\_exponent\_semigroup \{t : \ensuremath{\mathbb{R}}\} (ht6 : 6 < t) (ht7 : t \ensuremath{\leq} 7)\newline (n : \ensuremath{\mathbb{N}}) (hn : 0 < n) :\newline FunctionSquaredNorm (fun z : \ensuremath{\mathbb{C}} =>\newline (((1+Complex.normSq z)\textasciicircum{}4-t*(Complex.normSq z)\textasciicircum{}2)\textasciicircum{}n : \ensuremath{\mathbb{R}})) \ensuremath{\leftrightarrow}\newline n=2 \ensuremath{\vee} 4\ensuremath{\leq}n \ensuremath{\vee} (n=3 \ensuremath{\wedge} 0\ensuremath{\leq}924-210*t+18*t\textasciicircum{}2-t\textasciicircum{}3)\end{quote}
\subsection{PowerMoment.function\_exponent\_lower\_bound}
Source: \texttt{proofs/opdp13/OpHack/PowerMoment/Bound.lean}, line 86. SHA-256 \texttt{556b24cda4326f34d2b8e723b629b76a7930ff6a4c115f831942550fdd5fa830}.

\begin{quote}\small\ttfamily theorem function\_exponent\_lower\_bound (t : \ensuremath{\mathbb{R}}) (ht : 6 < t) (ht8 : t < 8)\newline (n : \ensuremath{\mathbb{N}}) (hn : 0 < n)\newline (h : PowerSemigroup.FunctionSquaredNorm (fun z : \ensuremath{\mathbb{C}} =>\newline (((1+Complex.normSq z)\textasciicircum{}4-t*(Complex.normSq z)\textasciicircum{}2)\textasciicircum{}n : \ensuremath{\mathbb{R}}))) :\newline t*(10-t) \ensuremath{\leq} 12*(n : \ensuremath{\mathbb{R}})*(8-t)\end{quote}
\subsection{PowerFourier.function\_logarithmic\_growth}
Source: \texttt{proofs/opdp13/OpHack/PowerFourier/Growth.lean}, line 59. SHA-256 \texttt{dab12267b81d6895e9cfc646fedf3b3180a6ded6987210111863cd98c118e1be}.

\begin{quote}\small\ttfamily theorem function\_logarithmic\_growth : \ensuremath{\forall} \ensuremath{\varepsilon} : \ensuremath{\mathbb{R}}, 0 < \ensuremath{\varepsilon} \ensuremath{\to} \ensuremath{\exists} \ensuremath{\eta} : \ensuremath{\mathbb{R}}, 0 < \ensuremath{\eta} \ensuremath{\wedge}\newline \ensuremath{\forall} (t : \ensuremath{\mathbb{R}}) (n : \ensuremath{\mathbb{N}}), 7 \ensuremath{\leq} t \ensuremath{\to} t < 8 \ensuremath{\to} 0 < n \ensuremath{\to}\newline PowerSemigroup.FunctionSquaredNorm (fun z : \ensuremath{\mathbb{C}} =>\newline (((1+Complex.normSq z)\textasciicircum{}4-t*(Complex.normSq z)\textasciicircum{}2)\textasciicircum{}n : \ensuremath{\mathbb{R}})) \ensuremath{\to} 8-t < \ensuremath{\eta} \ensuremath{\to}\newline (1-\ensuremath{\varepsilon})*Real.log (1/(8-t)) \ensuremath{\leq} (8-t)*(n : \ensuremath{\mathbb{R}})\end{quote}
\subsection{PowerUpper.function\_near\_boundary\_upper}
Source: \texttt{proofs/opdp13/OpHack/PowerUpper/Function.lean}, line 9. SHA-256 \texttt{c0b278ad2bc4d2d0297a99649e208205dcbc4136c8a3bf89c75604995a7a8e90}.

\begin{quote}\small\ttfamily theorem function\_near\_boundary\_upper :\newline \ensuremath{\forall} \ensuremath{\varepsilon} : \ensuremath{\mathbb{R}}, 0 < \ensuremath{\varepsilon} \ensuremath{\to} \ensuremath{\exists} \ensuremath{\delta}0 : \ensuremath{\mathbb{R}}, 0 < \ensuremath{\delta}0 \ensuremath{\wedge} \ensuremath{\delta}0 \ensuremath{\leq} 1 \ensuremath{\wedge}\newline \ensuremath{\forall} (\ensuremath{\delta} : \ensuremath{\mathbb{R}}) (N : \ensuremath{\mathbb{N}}), 0 < \ensuremath{\delta} \ensuremath{\to} \ensuremath{\delta} < \ensuremath{\delta}0 \ensuremath{\to}\newline (1+\ensuremath{\varepsilon})*Real.log (1/\ensuremath{\delta})/\ensuremath{\delta} \ensuremath{\leq} (N : \ensuremath{\mathbb{R}}) \ensuremath{\to}\newline PowerSemigroup.FunctionSquaredNorm (fun z : \ensuremath{\mathbb{C}} =>\newline (((1+Complex.normSq z)\textasciicircum{}4-(8-\ensuremath{\delta})*(Complex.normSq z)\textasciicircum{}2)\textasciicircum{}N : \ensuremath{\mathbb{R}}))\end{quote}
\subsection{PowerUpper.LeastIndex.leastPositiveIndex\_squeeze}
Source: \texttt{proofs/opdp13/OpHack/PowerUpper/LeastIndex.lean}, line 40. SHA-256 \texttt{41aa021b2f3b58b80de296aa22bf5ca917ffb9853bdaa740af8c9dcb36b442c0}.

\begin{quote}\small\ttfamily theorem leastPositiveIndex\_squeeze :\newline \ensuremath{\forall} \ensuremath{\varepsilon} : \ensuremath{\mathbb{R}}, 0 < \ensuremath{\varepsilon} \ensuremath{\to} \ensuremath{\exists} \ensuremath{\delta}0 : \ensuremath{\mathbb{R}}, 0 < \ensuremath{\delta}0 \ensuremath{\wedge} \ensuremath{\delta}0 \ensuremath{\leq} 1 \ensuremath{\wedge}\newline \ensuremath{\forall} \ensuremath{\delta} : \ensuremath{\mathbb{R}}, 0 < \ensuremath{\delta} \ensuremath{\to} \ensuremath{\delta} < \ensuremath{\delta}0 \ensuremath{\to}\newline IsLeastAdmissible \ensuremath{\delta} (leastPositiveIndex \ensuremath{\delta}) \ensuremath{\wedge}\newline 1-\ensuremath{\varepsilon} \ensuremath{\leq} \ensuremath{\delta}*(leastPositiveIndex \ensuremath{\delta} : \ensuremath{\mathbb{R}})/Real.log (1/\ensuremath{\delta}) \ensuremath{\wedge}\newline \ensuremath{\delta}*(leastPositiveIndex \ensuremath{\delta} : \ensuremath{\mathbb{R}})/Real.log (1/\ensuremath{\delta}) \ensuremath{\leq} 1+\ensuremath{\varepsilon}\end{quote}
\subsection{Verification evidence}
Final standalone40-proof-module and aggregate build completed 3553 jobs with exit 0; source SHA records match; endpoint axioms standard. Earlier dependency-graph refresh and stale 18-file README count are superseded by this evidence.

Recorded environments:

\begin{itemize}
\item \{'lean': '4.34.1', 'lean\_commit': '5045d0056413266e57c625dcd7c365b10e377c52', 'mathlib': 'd13f23b723b8a846827a245b89c10fc7d3f11612'\}
\end{itemize}
\begin{itemize}
\item \texttt{proofs/opdp13/verification/final-build.log} (existing retained artifact; SHA-256 \texttt{a44abbddab80c085f30679c43689b9abc80d0c3be678ee66df3c07ea28177c61}).
\end{itemize}
\begin{itemize}
\item \texttt{proofs/opdp13/verification/root-least-index-build.log} (existing retained artifact; SHA-256 \texttt{d92a7ec970fcbd3050429bc1a9ffbcaba1de62076ed1d490eb1bc5d6b2f22ebb}).
\end{itemize}
\clearpage
\section{opdp 43-halfturn: preserved statement and proof architecture}
\subsection{Sharp half-turn symmetry defect gap}
Completed result and evidence record.

Claim ID: \texttt{opdp43-halfturn}. Family: \texttt{partial-opdp43-halfturn}.

\subsection{Our contribution}
We prove that a no-three-in-line set of exactly 2n points in an n-by-n grid has half-turn miss count zero or at least four, and certifies a ten-point witness attaining four. The clean four-module build verifies the universal theorem and witness geometry with standard axioms, closing the historical miss-two question while crediting the known miss-four examples.

\subsection{Exact preserved statement and scope}
For A\ensuremath{\subseteq}\{0,...,n\ensuremath{-}1\}\ensuremath{^2} with |A|=2n and no three collinear, let R(x,y)=(n\ensuremath{-}1\ensuremath{-}x,n\ensuremath{-}1\ensuremath{-}y). The one-sided miss count |\{p\ensuremath{\in}A:R(p)\ensuremath{\notin}A\}| is 0 or \ensuremath{\geq}4. An explicit ten-point 5\ensuremath{\times}5 witness has miss exactly 4, proving global sharpness. Symmetric-difference cardinality is twice the one-sided miss.

Complete sharp answer to the explicitly proposed subsidiary near-half-turn-symmetry question; structural partial relative to broad no-three-in-line

Not claimed:

\begin{itemize}
\item Main asymptotic no-three-in-line conjecture
\end{itemize}
\begin{itemize}
\item Miss 4 at every n
\end{itemize}
\begin{itemize}
\item A valid n75 construction
\end{itemize}
\begin{itemize}
\item Symmetric-difference count 4; one-sided miss 4 corresponds to symmetric difference 8
\end{itemize}
\subsection{Contribution and prior overlap}
Proposed classification: M3A complete existing subsidiary target if organizer uses that exact source question; otherwise parent-linked partial/refinement for classification. No acceptance or credit asserted.

\begin{itemize}
\item No positive defect below 4 for any full grid configuration; kernel-checked global sharp witness
\end{itemize}
Recorded overlap:

\begin{itemize}
\item Known miss 4 examples and prior large-grid calibration configurations; no new general construction claimed
\end{itemize}
Recorded bounded prior-art comparison:

\begin{itemize}
\item Targeted checks found no matching obstruction, but elementary argument has material folklore risk
\end{itemize}
\begin{itemize}
\item The source already poses the subsidiary question, so new-target M3B is not the only possible route
\end{itemize}
\subsection{Formal proof source}
Paths below are relative to \texttt{Claim\_Evidence/} in this bundle. Exact source bytes, declaration lines and verification evidence are preserved. Dependency versions and the inspection/build boundary are described in the bundle README.

\subsection{NoThreeSymmetry.grid\_half\_turn\_gap}
Source: \texttt{proofs/opdp43/OpHack/NoThree/Main.lean}, line 8. SHA-256 \texttt{f8262d41b0933681b1db8dafb3d4302471dc28650b8ca950e3e1fc01f203c3d2}.

\begin{quote}\small\ttfamily theorem grid\_half\_turn\_gap \{n : \ensuremath{\mathbb{N}}\} (A : Finset (GridPoint n))\newline (h : GridNoThree A) (hcard : A.card = 2*n) :\newline (gridDefect A).card = 0 \ensuremath{\vee} 4 \ensuremath{\leq} (gridDefect A).card\end{quote}
\subsection{NoThreeSymmetry.sharp\_half\_turn\_gap}
Source: \texttt{proofs/opdp43/OpHack/NoThree/Main.lean}, line 56. SHA-256 \texttt{f8262d41b0933681b1db8dafb3d4302471dc28650b8ca950e3e1fc01f203c3d2}.

\begin{quote}\small\ttfamily theorem sharp\_half\_turn\_gap :\newline (\ensuremath{\forall} (n : \ensuremath{\mathbb{N}}) (A : Finset (GridPoint n)), GridNoThree A \ensuremath{\to} A.card = 2*n \ensuremath{\to}\newline (gridDefect A).card = 0 \ensuremath{\vee} 4 \ensuremath{\leq} (gridDefect A).card) \ensuremath{\wedge}\newline (\ensuremath{\exists} A : Finset (GridPoint 5), A.card = 10 \ensuremath{\wedge} GridNoThree A \ensuremath{\wedge} (gridDefect A).card = 4)\end{quote}
\subsection{Verification evidence}
Clean 4-module standalone release build completed 3097 jobs; source SHA entries match; final universal theorem includes n0,1,2 and exact witness geometry; standard axioms only.

Recorded environments:

\begin{itemize}
\item \{'lean': '4.34.1', 'lean\_commit': '5045d0056413266e57c625dcd7c365b10e377c52', 'mathlib': 'd13f23b723b8a846827a245b89c10fc7d3f11612'\}
\end{itemize}
\begin{itemize}
\item \texttt{proofs/opdp43/verification/clean-release-build.log} (existing retained artifact; SHA-256 \texttt{d723ce9924914dedd4fb96f7918496927ee737a90b565b28f7068d8a909dd6bc}).
\end{itemize}
\begin{itemize}
\item \texttt{proofs/opdp43/SHA256.json} (existing retained artifact; SHA-256 \texttt{4a5bff7200171e3caa8ad123d04d0dd5aaa090ca24225787bbaf92fbb5647a13}).
\end{itemize}
\clearpage
\section{erdos1060-uniform-partial: preserved statement and proof architecture}
\subsection{Erdős 1060 uniform quantitative partial bound}
Completed result and evidence record.

Claim ID: \texttt{erdos1060-uniform-partial}. Family: \texttt{erdos1060-weighted-core-bound}.

\subsection{Our contribution}
We use the retained powerful-core size constraint and a weighted generating function to reduce the uniform counting coefficient below log(3)/3 to the stated constant C. The fresh ten-module Lean 4.33.1 replay verifies the canonical count, all asymptotic hypotheses, and the strict improvement with standard axioms, retaining the prior squarefree injection and unrestricted-core bound as inputs.

\subsection{Exact preserved statement and scope}
Let sigma(k) be the sum of positive divisors of k, and let f(n) count positive integers k satisfying k sigma(k) = n. With C = log(2539/1000)/3 + log(10/9)/2, for every epsilon > 0 there is N such that every n >= N satisfies f(n) <= exp((C+epsilon) log(n)/log(log(n))). The formal count is exactly the finite set \{k <= n : k sigma(k) = n\}. The strict comparison C < log(3)/3 is also proved.

One provisional non-M2 mathematical partial toward Erdős 1060. The full zero-constant and polylogarithmic goals remain open.

Not claimed:

\begin{itemize}
\item No solution of either full Erdős 1060 goal
\end{itemize}
\begin{itemize}
\item No new claim for the known squarefree injection or unrestricted-core encoding
\end{itemize}
\begin{itemize}
\item No formalized explicit O(error-rate) refinement or optimized-majorant sharpness
\end{itemize}
\begin{itemize}
\item No lower bound for actual fiber counts, odd-target improvement or certified novelty
\end{itemize}
\subsection{Contribution and prior overlap}
Proposed classification: Provisional non-M2 mathematical partial. No acceptance or credit asserted.

\begin{itemize}
\item The retained powerful core satisfies P\textasciicircum{}2 <= n. A weighted generating-function bound with x=9/10 improves the located general unrestricted-core constant from log(3)/3 to C = 0.363270358736427... . All arithmetic and asymptotic hypotheses are discharged in the final canonical-count declarations.
\end{itemize}
Recorded overlap:

\begin{itemize}
\item Scott Duke Kominers, On the Number of Solutions of k sigma(k)=n, 2026 working paper, presents the unrestricted-core bound with log(3)/3 and attributes squarefree injectivity to Noppakaew-Pongsriiam. Those arguments are prior mathematics.
\end{itemize}
Recorded bounded prior-art comparison:

\begin{itemize}
\item Scott Duke Kominers, On the Number of Solutions of k sigma(k)=n, 2026 working paper, presents the unrestricted-core bound with log(3)/3 and attributes squarefree injectivity to Noppakaew-Pongsriiam. Those arguments are prior mathematics.
\end{itemize}
\subsection{Formal proof source}
Paths below are relative to \texttt{Claim\_Evidence/} in this bundle. Exact source bytes, declaration lines and verification evidence are preserved. Dependency versions and the inspection/build boundary are described in the bundle README.

\subsection{Erdos1060.canonical\_count\_eventually\_upper}
Source: \texttt{proofs/erdos1060-uniform-partial/src/SigmaFiberAsymptotic.lean}, line 38. SHA-256 \texttt{440f2f6aaa1d5dd04611fe855077352a3b75171d6e1ec823939ede8609e2d45a}.

\begin{quote}\small\ttfamily theorem canonical\_count\_eventually\_upper \{\ensuremath{\varepsilon} : \ensuremath{\mathbb{R}}\} (h\ensuremath{\varepsilon} : 0 < \ensuremath{\varepsilon}) :\newline \ensuremath{\forall}\ensuremath{^{f}} n : \ensuremath{\mathbb{N}} in atTop,\newline (\#\{k \ensuremath{\leq} n | k * ArithmeticFunction.sigma 1 k = n\} : \ensuremath{\mathbb{R}}) \ensuremath{\leq}\newline Real.exp (((Real.log (2539 / 1000) / 3 + Real.log (10 / 9) / 2) + \ensuremath{\varepsilon}) *\newline Real.log (n : \ensuremath{\mathbb{R}}) / Real.log (Real.log (n : \ensuremath{\mathbb{R}})))\end{quote}
\subsection{Erdos1060.canonical\_count\_upper}
Source: \texttt{proofs/erdos1060-uniform-partial/src/SigmaFiberAsymptotic.lean}, line 49. SHA-256 \texttt{440f2f6aaa1d5dd04611fe855077352a3b75171d6e1ec823939ede8609e2d45a}.

\begin{quote}\small\ttfamily theorem canonical\_count\_upper \{\ensuremath{\varepsilon} : \ensuremath{\mathbb{R}}\} (h\ensuremath{\varepsilon} : 0 < \ensuremath{\varepsilon}) :\newline \ensuremath{\exists} N : \ensuremath{\mathbb{N}}, \ensuremath{\forall} n : \ensuremath{\mathbb{N}}, N \ensuremath{\leq} n \ensuremath{\to}\newline (\#\{k \ensuremath{\leq} n | k * ArithmeticFunction.sigma 1 k = n\} : \ensuremath{\mathbb{R}}) \ensuremath{\leq}\newline Real.exp (((Real.log (2539 / 1000) / 3 + Real.log (10 / 9) / 2) + \ensuremath{\varepsilon}) *\newline Real.log (n : \ensuremath{\mathbb{R}}) / Real.log (Real.log (n : \ensuremath{\mathbb{R}})))\end{quote}
\subsection{Erdos1060Partial.entropyConstant\_lt\_old}
Source: \texttt{proofs/erdos1060-uniform-partial/src/EntropyConversion.lean}, line 45. SHA-256 \texttt{a9c629fe04c5f4252e92ebdef87ba5942f689fb64a0ca7e1fbc46f8f3957dc14}.

\begin{quote}\small\ttfamily theorem entropyConstant\_lt\_old : entropyConstant < Real.log 3 / 3\end{quote}
\subsection{Verification evidence}
Fresh isolated Lean 4.33.1 compilation passed all 10 own source modules, with no own olean reuse. Before/after source hashes agree and final canonical theorem types and axiom lists were printed. Only propext, Classical.choice and Quot.sound occur. Exact package evidence is retained; integration performs file-identity validation, not another replay.

Recorded environments:

\begin{itemize}
\item \{'lean': '4.33.1', 'mathlib': '0df444a360eaa60ab8c11dca51a86af692955474'\}
\end{itemize}
\begin{itemize}
\item \texttt{proofs/erdos1060-uniform-partial/verification/RESULT.json} (existing retained artifact; SHA-256 \texttt{29bf72f15436677f0281c03d61fd6e4aa9bf4152d6704ac33a448264a6983d65}).
\end{itemize}
\begin{itemize}
\item \texttt{proofs/erdos1060-uniform-partial/verification/fresh-replay.log} (existing retained artifact; SHA-256 \texttt{7b01be04b6048f53669722c42e712ce80694e79fefef511dcf5d31808d161b46}).
\end{itemize}
\begin{itemize}
\item \texttt{proofs/erdos1060-uniform-partial/verification/source-hashes-before.json} (existing retained artifact; SHA-256 \texttt{85ef32af149718a85928eaaa89a91543664cf3413fb70e2020ce74e09c0fab1f}).
\end{itemize}
\begin{itemize}
\item \texttt{proofs/erdos1060-uniform-partial/verification/source-hashes-after.json} (existing retained artifact; SHA-256 \texttt{85ef32af149718a85928eaaa89a91543664cf3413fb70e2020ce74e09c0fab1f}).
\end{itemize}
\clearpage
\section{erdos829-log-five-thirds: preserved statement and proof architecture}
\subsection{Erdős 829 verified log(5/3) partial bound}
Completed result and evidence record.

Claim ID: \texttt{erdos829-log-five-thirds}. Family: \texttt{erdos829-weighted-code}.

\subsection{Our contribution}
We combine the finite weighted-code certificate, all-target gcd transfer, and prime cutoff to prove the uniform log(5/3) bound for ordered nonnegative cube-sum representations. The fresh 33-proof-module replay and two audits verify the exact unrestricted target count with standard axioms, building on prior ratio-root, congruence, and coding methods.

\subsection{Exact preserved statement and scope}
Let R(n) count ordered nonnegative integer pairs (x,y) satisfying x\textasciicircum{}3+y\textasciicircum{}3=n. For every real epsilon>0 there exists a natural N such that, for every natural n>=N, R(n) <= exp((log(5/3)+epsilon)*log(n)/log(log(n))). Zero coordinates are included and swapping unequal coordinates gives distinct representations. There is no primitive-pair or cubefree-target restriction. The existential threshold is not given numerically.

One conditional original-target mathematical partial-advance family; the polylogarithmic target remains open.

Not claimed:

\begin{itemize}
\item No proof or refutation of the original polylogarithmic Erdős 829 conjecture R(n)=O((log n)\textasciicircum{}C)
\end{itemize}
\begin{itemize}
\item No signed-integer representation bound
\end{itemize}
\begin{itemize}
\item No certified novelty, best-known status, official admission, accepted progress or points
\end{itemize}
\begin{itemize}
\item Competition-period production is entrant-declared; recovery/check dates do not establish independent publication priority
\end{itemize}
\subsection{Contribution and prior overlap}
Proposed classification: Conditional non-M2 original-target mathematical partial. No acceptance or credit asserted.

\begin{itemize}
\item The selected endpoint proves the all-target ordered nonnegative count bound with coefficient log(5/3). The earlier log(7/4) endpoint belongs to this same family. Selection is conditional on target status, novelty and substantive mathematical review. If the mathematics is known, M2 requires a separate genuinely new and useful formalization delta.
\end{itemize}
Recorded overlap:

\begin{itemize}
\item Primitive ratio-root injection and determinant congruence are prior tools (including Merikoski, Lemma 8); Elias/Plotkin localization and entropy machinery are classical. The finite weighted-code certificate, all-target gcd transfer and prime cutoff support the proposed partial-advance comparison.
\end{itemize}
\begin{itemize}
\item The bounded literature comparison did not locate an all-target bound subsuming the stated coefficients. Its derived 0.5564773354 comparison applies only to cubefree targets; it is not an all-target theorem or a comprehensive priority certificate.
\end{itemize}
Recorded bounded prior-art comparison:

\begin{itemize}
\item Primitive ratio-root injection and determinant congruence are prior tools (including Merikoski, Lemma 8); Elias/Plotkin localization and entropy machinery are classical. The finite weighted-code certificate, all-target gcd transfer and prime cutoff support the proposed partial-advance comparison.
\end{itemize}
\begin{itemize}
\item The bounded literature comparison did not locate an all-target bound subsuming the stated coefficients. Its derived 0.5564773354 comparison applies only to cubefree targets; it is not an all-target theorem or a comprehensive priority certificate.
\end{itemize}
\subsection{Formal proof source}
Paths below are relative to \texttt{Claim\_Evidence/} in this bundle. Exact source bytes, declaration lines and verification evidence are preserved. Dependency versions and the inspection/build boundary are described in the bundle README.

\subsection{Erdos829FiveThirds.canonicalCubeCount\_eventually\_exp}
Source: \texttt{proofs/erdos829-log-five-thirds/FiveThirdsCanonical.lean}, line 8. SHA-256 \texttt{51541747f8104588463d8b7b051d69692876cee605c68a2f50df9da132e6a006}.

\begin{quote}\small\ttfamily theorem canonicalCubeCount\_eventually\_exp \{\ensuremath{\varepsilon} : \ensuremath{\mathbb{R}}\} (h\ensuremath{\varepsilon} : 0<\ensuremath{\varepsilon}) :\newline \ensuremath{\forall}\ensuremath{^{f}} n : \ensuremath{\mathbb{N}} in atTop, (canonicalCubeCount n : \ensuremath{\mathbb{R}}) \ensuremath{\leq}\newline Real.exp ((Real.log (5/3 : \ensuremath{\mathbb{R}})+\ensuremath{\varepsilon})*Real.log (n : \ensuremath{\mathbb{R}})/Real.log (Real.log (n : \ensuremath{\mathbb{R}})))\end{quote}
\subsection{Erdos829FiveThirds.canonicalCubeCount\_eventual\_threshold}
Source: \texttt{proofs/erdos829-log-five-thirds/FiveThirdsCanonical.lean}, line 14. SHA-256 \texttt{51541747f8104588463d8b7b051d69692876cee605c68a2f50df9da132e6a006}.

\begin{quote}\small\ttfamily theorem canonicalCubeCount\_eventual\_threshold \{\ensuremath{\varepsilon} : \ensuremath{\mathbb{R}}\} (h\ensuremath{\varepsilon} : 0<\ensuremath{\varepsilon}) :\newline \ensuremath{\exists} N : \ensuremath{\mathbb{N}}, \ensuremath{\forall} n : \ensuremath{\mathbb{N}}, N\ensuremath{\leq}n \ensuremath{\to} (canonicalCubeCount n : \ensuremath{\mathbb{R}}) \ensuremath{\leq}\newline Real.exp ((Real.log (5/3 : \ensuremath{\mathbb{R}})+\ensuremath{\varepsilon})*Real.log (n : \ensuremath{\mathbb{R}})/Real.log (Real.log (n : \ensuremath{\mathbb{R}})))\end{quote}
\subsection{Verification evidence}
Fresh separately orchestrated source-only Lean 4.33.1 replay completed 2026-10-02T18:54:57.202016+00:00. All 33 proof modules and two semantic/axiom audit modules exited 0. The final active sources comprise 28 historical exact modules and five reconstructed/repaired variants, all freshly verified. No old own-module oleans were used; nine third-party dependency pins matched. Endpoint transitive axioms are only propext, Classical.choice and Quot.sound. The semantic audit checks the exact expanded catalog expression and ordered nonnegative coordinate count; it does not import the catalog conjecture. This is internal AI-assisted formal-check/semantic evidence, not appointed jury or external peer review.

Recorded environments:

\begin{itemize}
\item \{'lean': '4.33.1', 'lean\_commit': '819816b2e0a3bf405af45ae5c7af2491d8f5bee6', 'mathlib': '0df444a360eaa60ab8c11dca51a86af692955474', 'dependency\_pins': \{'plausible': 'b7eb3304aeae834b12dda98993a37f6a41f6f0bb', 'LeanSearchClient': '5f4d51b81cbd3f6b32b156bfad9056621a040404', 'importGraph': '16f02aa7642864af59f1ff0e384a015994db9118', 'proofwidgets': '4be2e3d5087eeb272cf5a8853b8f9dd025ef5957', 'aesop': '3448c0bcc5ce01b2d1546e483ec3620e32df3d0e', 'Qq': '92c15be17b7caf78c2ad767ec40f89052d908d81', 'batteries': '4488d40d070b9700d4d5a6aa342f0d40c31b2a2d', 'Cli': '6130a47896ce867c6a4a55373441e59e565bad0f', 'mathlib': '0df444a360eaa60ab8c11dca51a86af692955474'\}, 'dependency\_pins\_path': 'proofs/erdos829-log-five-thirds/dependency-pins.json'\}
\end{itemize}
\begin{itemize}
\item \texttt{proofs/erdos829-log-five-thirds/verification/verification.json} (existing retained artifact; SHA-256 \texttt{ec3e17f4d76b50df8ff759480fbc65907309d5a2d715d3fc3b7e4013c472bfa9}).
\end{itemize}
\begin{itemize}
\item \texttt{proofs/erdos829-log-five-thirds/verification/independent-audit/independent-receipt.json} (existing retained artifact; SHA-256 \texttt{ac1635c2bac105208673a0e5f1554203ea39593c38f7639adad51953b1d7f034}).
\end{itemize}
\begin{itemize}
\item \texttt{proofs/erdos829-log-five-thirds/verification/logs/FiveThirdsCanonical.log} (existing retained artifact; SHA-256 \texttt{f093d3acdd665dd598fff8a4871d6d45570067958987651fd5ff5f4122521491}).
\end{itemize}
\begin{itemize}
\item \texttt{proofs/erdos829-log-five-thirds/verification/logs/FiveThirdsFinalTypeAudit.log} (existing retained artifact; SHA-256 \texttt{11c397a92e03c799fda8be31ee6562087a81e0b8cf79b0ef389f9ff0a2260533}).
\end{itemize}
\begin{itemize}
\item \texttt{proofs/erdos829-log-five-thirds/verification/logs/IndependentFiveThirdsSemanticAudit.log} (existing retained artifact; SHA-256 \texttt{d43f0952c0d72c5fb5276c21144ff08861221a5c7ff48c7a1b802265e872e0d1}).
\end{itemize}
\begin{itemize}
\item \texttt{proofs/erdos829-log-five-thirds/verification/logs/FreshSemanticAndAxiomAudit.log} (existing retained artifact; SHA-256 \texttt{3d62193131c9f0a6bb581f4c5206c8852b9cf2b0ee9a7c3952fea42217a2e613}).
\end{itemize}
\begin{itemize}
\item \texttt{proofs/erdos829-log-five-thirds/verification/source-hashes.json} (existing retained artifact; SHA-256 \texttt{126dffa9a2b6811fde754de26e3fd23020c594eeac730dab5321acd0bcdc0931}).
\end{itemize}
\begin{itemize}
\item \texttt{proofs/erdos829-log-five-thirds/dependency-pins.json} (existing retained artifact; SHA-256 \texttt{370768e89fbf1028eed2a6f5784734a9747c849b0f30e45b62f6f467d7330e7c}).
\end{itemize}
\begin{itemize}
\item \texttt{proofs/erdos829-log-five-thirds/PACKAGE\_MANIFEST.json} (existing retained artifact; SHA-256 \texttt{72d23e3f383c5cf9b62ca60c80394360c5868fe35b57f2f7ba2d4ba13b26819f}).
\end{itemize}
\begin{itemize}
\item \texttt{proofs/erdos829-log-five-thirds/VERIFICATION.md} (existing retained artifact; SHA-256 \texttt{78067ed971757b61807f2dab2bf63d2186dbf965adf0f3077088c0da2f481678}).
\end{itemize}
\clearpage
\section{erdos944-order-bound: preserved statement and proof architecture}
\subsection{Erdős 944 finite order obstruction}
Completed result and evidence record.

Claim ID: \texttt{erdos944-order-bound}. Family: \texttt{erdos944-finite-order-bound}.

\subsection{Our contribution}
We supply the stronger finite order obstruction for vertex-critical four-chromatic graphs with the stated edge resilience, including the quadratic and integer independent-set filters. The isolated Lean 4.34.1 replay and target audit verify all three inequalities with standard axioms, improving the explicit finite comparison while preserving the prior stronger eventual asymptotic result.

\subsection{Exact preserved statement and scope}
For a finite simple graph G with n vertices and an integer r >= 1, suppose its chromatic number is 4, deleting any vertex lowers that number, and deleting any set of at most r edges does not lower it. Then 27r+19 <= 4n. The same hypotheses give 27(r+1)(n-1) <= (2n+1)\textasciicircum{}2 and 9b(r+1) <= (2b+1)(n-b), where b=floor((n+1)/3). These necessary conditions form one family.

One provisional non-M2 mathematical partial for the finite k=4 case of Erdős 944.

Not claimed:

\begin{itemize}
\item No existence construction or full solution of Erdős 944
\end{itemize}
\begin{itemize}
\item No infinite-vertex or k>4 theorem
\end{itemize}
\begin{itemize}
\item No asymptotic improvement
\end{itemize}
\begin{itemize}
\item No new credit for the prior r=1 construction
\end{itemize}
\begin{itemize}
\item No certified novelty, first-publication priority, eligibility or accepted credit
\end{itemize}
\subsection{Contribution and prior overlap}
Proposed classification: Provisional non-M2 mathematical partial. No acceptance or credit asserted.

\begin{itemize}
\item The clean inequality strengthens the explicit n >= 5r+6 finite obstruction in Skottova and Steiner, Proposition 5.1. Stronger quadratic and integral independent-set filters are also proved; for r=4, the integer filter requires n >= 33 rather than the linear bound alone requiring n >= 32.
\end{itemize}
Recorded overlap:

\begin{itemize}
\item Skottova and Steiner, Critical edge sets in vertex-critical graphs, arXiv:2508.08703, supply the explicit n >= 5r+6 bound and the separate r=O(n/(log n)\textasciicircum{}c) theorem for an absolute c>0. Their asymptotic theorem is stronger eventually and is not improved here.
\end{itemize}
\begin{itemize}
\item The public 48-vertex degree-6 r=1 construction and later 6-regular r=1 results are prior work. Target criticality definitions are reused; the construction proof is neither copied nor imported. Finite coding, Cauchy-Schwarz and Plotkin methods are classical background.
\end{itemize}
Recorded bounded prior-art comparison:

\begin{itemize}
\item Skottova and Steiner, Critical edge sets in vertex-critical graphs, arXiv:2508.08703, supply the explicit n >= 5r+6 bound and the separate r=O(n/(log n)\textasciicircum{}c) theorem for an absolute c>0. Their asymptotic theorem is stronger eventually and is not improved here.
\end{itemize}
\begin{itemize}
\item The public 48-vertex degree-6 r=1 construction and later 6-regular r=1 results are prior work. Target criticality definitions are reused; the construction proof is neither copied nor imported. Finite coding, Cauchy-Schwarz and Plotkin methods are classical background.
\end{itemize}
\subsection{Formal proof source}
Paths below are relative to \texttt{Claim\_Evidence/} in this bundle. Exact source bytes, declaration lines and verification evidence are preserved. Dependency versions and the inspection/build boundary are described in the bundle README.

\subsection{Erdos944Bounds.erdos944\_order\_linear}
Source: \texttt{proofs/erdos944-order-bound/OpHack/Erdos944Bounds.lean}, line 42. SHA-256 \texttt{8f9e544d0e9f465f8d92d0c0b878ab2bd9b542c0c969eaa4e2d7c908acec27f4}.

\begin{quote}\small\ttfamily theorem erdos944\_order\_linear \{r : \ensuremath{\mathbb{N}}\} (hr : 1 \ensuremath{\leq} r)\newline (hG : G.IsCritical 4)\newline (hE : \ensuremath{\forall} edges : Set (Sym2 V), G.IsCriticalEdges edges \ensuremath{\to} r < edges.ncard) :\newline 27*r + 19 \ensuremath{\leq} 4 * Fintype.card V\end{quote}
\subsection{Erdos944Bounds.erdos944\_order\_quadratic}
Source: \texttt{proofs/erdos944-order-bound/OpHack/Erdos944Bounds.lean}, line 28. SHA-256 \texttt{8f9e544d0e9f465f8d92d0c0b878ab2bd9b542c0c969eaa4e2d7c908acec27f4}.

\begin{quote}\small\ttfamily theorem erdos944\_order\_quadratic \{r : \ensuremath{\mathbb{N}}\} (hr : 1 \ensuremath{\leq} r)\newline (hG : G.IsCritical 4)\newline (hE : \ensuremath{\forall} edges : Set (Sym2 V), G.IsCriticalEdges edges \ensuremath{\to} r < edges.ncard) :\newline 27 * (r+1) * (Fintype.card V - 1) \ensuremath{\leq} (2 * Fintype.card V + 1)\textasciicircum{}2\end{quote}
\subsection{Erdos944Bounds.erdos944\_order\_exact}
Source: \texttt{proofs/erdos944-order-bound/OpHack/Erdos944Bounds.lean}, line 54. SHA-256 \texttt{8f9e544d0e9f465f8d92d0c0b878ab2bd9b542c0c969eaa4e2d7c908acec27f4}.

\begin{quote}\small\ttfamily theorem erdos944\_order\_exact \{r : \ensuremath{\mathbb{N}}\} (hr : 1 \ensuremath{\leq} r)\newline (hG : G.IsCritical 4)\newline (hE : \ensuremath{\forall} edges : Set (Sym2 V), G.IsCriticalEdges edges \ensuremath{\to} r < edges.ncard) :\newline 9 * ((Fintype.card V + 1) / 3) * (r+1) \ensuremath{\leq}\newline (2 * ((Fintype.card V + 1) / 3) + 1) *\newline (Fintype.card V - (Fintype.card V + 1) / 3)\end{quote}
\begin{samepage}
\subsection{Verification evidence}
Independent isolated Lean 4.34.1 replay compiled three proof modules and the separate TargetAudit.lean from unchanged frozen sources. Its own output directory began empty and earlier project proof caches were excluded; pinned third-party caches were reused. Exact target definitions and conjunction were audited. Every final endpoint uses only propext, Classical.choice and Quot.sound. Packaging validates identities without claiming a further Lean replay.

Recorded environments:

\begin{itemize}
\item \{'lean': '4.34.1', 'mathlib': 'd13f23b723b8a846827a245b89c10fc7d3f11612'\}
\end{itemize}
\begin{itemize}
\item \texttt{proofs/erdos944-order-bound/verification/independent-replay-receipt.json} (existing retained artifact; SHA-256 \texttt{62e604c3e67cf6522b78d99c10eab4c6110cc2efd961bdaff8195902ad07446a}).
\end{itemize}
\begin{itemize}
\item \texttt{proofs/erdos944-order-bound/verification/Erdos944Graph.log} (existing retained artifact; SHA-256 \texttt{00d06e0489132de4a21fd55425f3f8f16f20bd5c587d1afdf16200def97a7d9f}).
\end{itemize}
\begin{itemize}
\item \texttt{proofs/erdos944-order-bound/verification/CriticalCodeBound.log} (existing retained artifact; SHA-256 \texttt{b1e3668d218324d6984ec262069df738d472b9fb721207e631132de5d6da0763}).
\end{itemize}
\begin{itemize}
\item \texttt{proofs/erdos944-order-bound/verification/Erdos944Bounds.log} (existing retained artifact; SHA-256 \texttt{41ae983a1ca6a332b804faa8df2861673c3147648732d2c2581c464c4a61f5b1}).
\end{itemize}
\begin{itemize}
\item \texttt{proofs/erdos944-order-bound/verification/TargetAudit.log} (existing retained artifact; SHA-256 \texttt{8521858d1ac7a1983fee4b1eb9ea0d7f363b5f91da05841265b2671bea0f83c4}).
\end{itemize}
\end{samepage}
\clearpage
\section{opdp89-parity-local-maxima: preserved statement and proof architecture}
\subsection{OPDP 89 sharp four parity likelihood classification}
Completed result and evidence record.

Claim ID: \texttt{opdp89-parity-local-maxima}. Family: \texttt{opdp89-parity-rbm}.

\subsection{Our contribution}
We supply the sharp collapsed-fibre local-maximum threshold for the four-visible, one-hidden pure-parity RBM, including stationarity and both signs of variation at and beyond the boundary in every full parameter neighborhood. The fresh 17-module replay and semantic audit verify this higher-order classification with standard axioms, extending the known singular-stationarity setting.

\subsection{Exact preserved statement and scope}
For pure four-parity data Q\_rho(s)=(1+rho*s1*s2*s3*s4)/16 with 0<|rho|<=1, consider the normalized four-visible, one-hidden spin RBM P\_(b,c,w)(s) proportional to exp(b dot s)*cosh(c+w dot s), and its expected log likelihood ell\_rho. At every collapsed point (0,c0,0), with c0 any finite real, every affine parameter-direction derivative is zero. It is a local maximum in the full nine-dimensional parameter space if and only if tanh(c0)\textasciicircum{}2 < (3+2|rho|)/(3+6|rho|). At equality or beyond, every full parameter neighborhood has both a strictly higher and a strictly lower likelihood point. Every collapsed point has a strictly better finite RBM and distinct equal-likelihood points in every neighborhood; the local maxima are non-strict.

One provisional non-M2 mathematical family giving the exact finite collapsed-fiber classification for the four-visible, one-hidden parity RBM.

Not claimed:

\begin{itemize}
\item No higher-dimensional n>=5 or arbitrary four-wise-uniform-data formal endpoint
\end{itemize}
\begin{itemize}
\item Exactly one hidden unit and pure four-parity data
\end{itemize}
\begin{itemize}
\item No distribution-space local maximum, strict optimization trap, attraction basin, EM convergence or training-time theorem
\end{itemize}
\begin{itemize}
\item No complete resolution of the broad OPDP 89 learning-structures prompt
\end{itemize}
\begin{itemize}
\item No extra family credit for the earlier strict region, balanced fallback, stationary fiber or ordinary extensions
\end{itemize}
\begin{itemize}
\item Historical sharp raw logs remain unavailable; the separate dated fresh post-recovery PASS receipt/logs are retained without altering historical snapshots
\end{itemize}
\begin{itemize}
\item Novelty, eligibility and accepted credit remain unestablished
\end{itemize}
\subsection{Contribution and prior overlap}
Proposed classification: Provisional non-M2 mathematical partial. No acceptance or credit asserted.

\begin{itemize}
\item This upgrades the existing strict-region family to a sharp if-and-only-if threshold, with stationarity in every affine parameter direction and both signs of likelihood variation in every full neighborhood at and outside the boundary. Boundary passage is proved along the likelihood-flat collapsed fiber. It remains one family.
\end{itemize}
Recorded overlap:

\begin{itemize}
\item Collapsed stationarity, singular mixture critical lines and familiar RBM interaction formulas are prior mathematics. The bounded comparison with Fukumizu, Akaho and Amari (NIPS 2002) identifies a degenerate case where the second-derivative splitting matrix vanishes; the higher-order full-neighborhood result is the candidate contribution.
\end{itemize}
\begin{itemize}
\item Bounded primary-source review did not locate the exact parity threshold result; it does not certify novelty.
\end{itemize}
\begin{itemize}
\item The earlier strict-region proof and broader ordinary OPDP 89 record remain same-family support. The exact 17-module sharp closure now also has fresh post-recovery verification; n>=5 and four-wise-uniform extensions remain ordinary-only.
\end{itemize}
Recorded bounded prior-art comparison:

\begin{itemize}
\item Collapsed stationarity, singular mixture critical lines and familiar RBM interaction formulas are prior mathematics. The bounded comparison with Fukumizu, Akaho and Amari (NIPS 2002) identifies a degenerate case where the second-derivative splitting matrix vanishes; the higher-order full-neighborhood result is the candidate contribution.
\end{itemize}
\begin{itemize}
\item Bounded primary-source review did not locate the exact parity threshold result; it does not certify novelty.
\end{itemize}
\begin{itemize}
\item The earlier strict-region proof and broader ordinary OPDP 89 record remain same-family support. The exact 17-module sharp closure now also has fresh post-recovery verification; n>=5 and four-wise-uniform extensions remain ordinary-only.
\end{itemize}
\subsection{Formal proof source}
Paths below are relative to \texttt{Claim\_Evidence/} in this bundle. Exact source bytes, declaration lines and verification evidence are preserved. Dependency versions and the inspection/build boundary are described in the bundle README.

\subsection{ParityFour.sharp\_parity\_classification}
Source: \texttt{proofs/opdp89-parity-sharp/SharpClassification.lean}, line 40. SHA-256 \texttt{15ef1ad86c585c87e7ac4fabdc3a58702c915e69dbf6aa80d701258bc9efc602}.

\begin{quote}\small\ttfamily theorem sharp\_parity\_classification (\ensuremath{\rho} c\ensuremath{_0} : \ensuremath{\mathbb{R}}) (h\ensuremath{\rho} : 0 < |\ensuremath{\rho}|) (h\ensuremath{\rho}1 : |\ensuremath{\rho}| \ensuremath{\leq} 1) :\newline (\ensuremath{\forall} (b : Fin 4 \ensuremath{\to} \ensuremath{\mathbb{R}}) (d : \ensuremath{\mathbb{R}}) (w : Fin 4 \ensuremath{\to} \ensuremath{\mathbb{R}}),\newline HasDerivAt (fun t => objective \ensuremath{\rho} (affineParameterRay c\ensuremath{_0} b d w t)) 0 0) \ensuremath{\wedge}\newline (IsLocalMax (objective \ensuremath{\rho}) (nullParameter c\ensuremath{_0}) \ensuremath{\leftrightarrow}\newline Real.tanh c\ensuremath{_0} \textasciicircum{} 2 < (3+2*|\ensuremath{\rho}|)/(3+6*|\ensuremath{\rho}|)) \ensuremath{\wedge}\newline ((3+2*|\ensuremath{\rho}|)/(3+6*|\ensuremath{\rho}|) \ensuremath{\leq} Real.tanh c\ensuremath{_0} \textasciicircum{} 2 \ensuremath{\to} ParameterSaddle \ensuremath{\rho} (nullParameter c\ensuremath{_0})) \ensuremath{\wedge}\newline (\ensuremath{\exists} \ensuremath{\theta} : Parameters, objective \ensuremath{\rho} (nullParameter c\ensuremath{_0}) < objective \ensuremath{\rho} \ensuremath{\theta}) \ensuremath{\wedge}\newline (\ensuremath{\forall} U : Set Parameters, U \ensuremath{\in} nhds (nullParameter c\ensuremath{_0}) \ensuremath{\to}\newline \ensuremath{\exists} \ensuremath{\theta} : Parameters, \ensuremath{\theta} \ensuremath{\in} U \ensuremath{\wedge} \ensuremath{\theta} \ensuremath{\neq} nullParameter c\ensuremath{_0} \ensuremath{\wedge}\newline objective \ensuremath{\rho} \ensuremath{\theta} = objective \ensuremath{\rho} (nullParameter c\ensuremath{_0}))\end{quote}
\subsection{ParityFour.collapsed\_local\_max\_iff}
Source: \texttt{proofs/opdp89-parity-sharp/SharpClassification.lean}, line 22. SHA-256 \texttt{15ef1ad86c585c87e7ac4fabdc3a58702c915e69dbf6aa80d701258bc9efc602}.

\begin{quote}\small\ttfamily theorem collapsed\_local\_max\_iff (\ensuremath{\rho} c\ensuremath{_0} : \ensuremath{\mathbb{R}}) (h\ensuremath{\rho} : 0 < |\ensuremath{\rho}|) (h\ensuremath{\rho}1 : |\ensuremath{\rho}| \ensuremath{\leq} 1) :\newline IsLocalMax (objective \ensuremath{\rho}) (nullParameter c\ensuremath{_0}) \ensuremath{\leftrightarrow}\newline Real.tanh c\ensuremath{_0} \textasciicircum{} 2 < (3+2*|\ensuremath{\rho}|)/(3+6*|\ensuremath{\rho}|)\end{quote}
\subsection{ParityFour.nonstable\_parameter\_saddle}
Source: \texttt{proofs/opdp89-parity-sharp/SharpClassification.lean}, line 33. SHA-256 \texttt{15ef1ad86c585c87e7ac4fabdc3a58702c915e69dbf6aa80d701258bc9efc602}.

\begin{quote}\small\ttfamily theorem nonstable\_parameter\_saddle (\ensuremath{\rho} c\ensuremath{_0} : \ensuremath{\mathbb{R}}) (h\ensuremath{\rho} : 0 < |\ensuremath{\rho}|) (h\ensuremath{\rho}1 : |\ensuremath{\rho}| \ensuremath{\leq} 1)\newline (houtside : (3+2*|\ensuremath{\rho}|)/(3+6*|\ensuremath{\rho}|) \ensuremath{\leq} Real.tanh c\ensuremath{_0} \textasciicircum{} 2) :\newline ParameterSaddle \ensuremath{\rho} (nullParameter c\ensuremath{_0})\end{quote}
\subsection{Verification evidence}
Fresh post-recovery source-only replay completed 2026-10-02T18:45:45.269248+00:00 in 108.552 seconds. All 17 frozen source modules passed with no warnings or errors; all original source SHA 256 values and nine dependency pins were unchanged before/after. No old own-module oleans were reused. A supplemental semantic/type audit passed at 18:48:24 UTC. The final transitive axioms are only propext, Classical.choice and Quot.sound. Historical sharp raw logs remain unavailable; their absence is not filled with invented logs. The new dated PASS receipt and genuine fresh logs supersede historical pending status in byte-preserved recovery snapshots. Earlier strict-region proof and original raw evidence remain unchanged. This is internal AI-assisted formal/semantic review, not appointed jury or external peer review.

Recorded environments:

\begin{itemize}
\item \{'lean': '4.33.1', 'mathlib': '0df444a360eaa60ab8c11dca51a86af692955474'\}
\end{itemize}
\begin{itemize}
\item \texttt{proofs/opdp89-parity-local-maxima/verification/independent-replay/REPLAY.json} (existing retained artifact; SHA-256 \texttt{4836e9b11a247d1cb1b4a83294167029245f75a0a396807b8a7b3031e9fe938b}).
\end{itemize}
\begin{itemize}
\item \texttt{proofs/opdp89-parity-local-maxima/verification/independent-replay/logs/Main.log} (existing retained artifact; SHA-256 \texttt{b7d002b1db5e7e5d5efb5939ff5e6dd40457478888e0bcb92ebdb14e6632d32f}).
\end{itemize}
\begin{itemize}
\item \texttt{proofs/opdp89-parity-local-maxima/verification/independent-replay/logs/AxiomAudit.log} (existing retained artifact; SHA-256 \texttt{147408b72d34539e1285a08e9419b4bf5d530480bc2d840cb037124e11ae0070}).
\end{itemize}
\begin{itemize}
\item \texttt{proofs/opdp89-parity-local-maxima/verification/author-replay/REPLAY.json} (existing retained artifact; SHA-256 \texttt{e8fece7b1d9b00b2e8e02f6b8110f6cfcdf3e107e3fa894ab5c8b5ded6e7736d}).
\end{itemize}
\begin{itemize}
\item \texttt{proofs/opdp89-parity-sharp/RESULT.json} (existing retained artifact; SHA-256 \texttt{e5ca9d67c3dd2b5ddb0674f6c57d8a5f653373f36cdd8fbbda651a39a369d89e}).
\end{itemize}
\begin{itemize}
\item \texttt{proofs/opdp89-parity-sharp/RECOVERY\_PROVENANCE.json} (existing retained artifact; SHA-256 \texttt{e1153d2434c964477fd6260ef8cfaa540b5c2d0c96c5790317f20ca4bcb9a269}).
\end{itemize}
\begin{itemize}
\item \texttt{proofs/opdp89-parity-sharp/VERIFICATION.md} (existing retained artifact; SHA-256 \texttt{2bc11c186ec19b509cc78a6381fee2d4dfdab4a7847cf7ada2d599e29b147ded}).
\end{itemize}
\begin{itemize}
\item \texttt{proofs/opdp89-parity-sharp/SEMANTIC\_SCOPE.md} (existing retained artifact; SHA-256 \texttt{776ffd82ffcc9fdc90a6ef8bfcbadd17cbbcef356584bd5bcc890544c3163b1f}).
\end{itemize}
\begin{itemize}
\item \texttt{proofs/opdp89-parity-sharp/SOURCE\_MANIFEST.json} (existing retained artifact; SHA-256 \texttt{a9107e9c4c546451166108f1faab9e25cf4ac73926e766ef0659638d2463ae7b}).
\end{itemize}
\begin{itemize}
\item \texttt{evidence/opdp89-sharp-revalidation-20261002/VERIFICATION\_RECEIPT.json} (existing retained artifact; SHA-256 \texttt{38ace5540df3a7684453c21e6558c35a773cba1c7dbdca87d92801acf6d19af1}).
\end{itemize}
\begin{itemize}
\item \texttt{evidence/opdp89-sharp-revalidation-20261002/fresh-replay/REPLAY.json} (existing retained artifact; SHA-256 \texttt{c469ec4e66869f015af4df1f91002d3dc933528d679dc78c3982eacf0c3e39b5}).
\end{itemize}
\begin{itemize}
\item \texttt{evidence/opdp89-sharp-revalidation-20261002/fresh-replay/logs/SharpClassification.log} (existing retained artifact; SHA-256 \texttt{c241e8d517f5f7a4f19fa62f3870e0b0dd598c10003fbf4a7597f727a6564bf5}).
\end{itemize}
\begin{itemize}
\item \texttt{evidence/opdp89-sharp-revalidation-20261002/fresh-replay/logs/SharpAxiomAudit.log} (existing retained artifact; SHA-256 \texttt{b53ee9c6349e2db9b8c73625f4be0a06b247b1c94e0074f1a811de21649a53df}).
\end{itemize}
\begin{itemize}
\item \texttt{evidence/opdp89-sharp-revalidation-20261002/supplemental-audit/SemanticTypeAudit.log} (existing retained artifact; SHA-256 \texttt{f4f44e7043de5b679f015f0eb1546b926edf9bc8369f038491be25cc6955b802}).
\end{itemize}
\begin{itemize}
\item \texttt{evidence/opdp89-sharp-revalidation-20261002/supplemental-audit/RESULT.json} (existing retained artifact; SHA-256 \texttt{96581517c38656d5e7e9077ae8354ef5767a321b40e71de411ee66c37e3f83d4}).
\end{itemize}
\begin{itemize}
\item \texttt{evidence/opdp89-sharp-revalidation-20261002/SEMANTIC\_REVIEW.md} (existing retained artifact; SHA-256 \texttt{f28005106a7f99f674c4879830a59d63bfa94475483f43fbb65f82c4ee83fd06}).
\end{itemize}
\begin{itemize}
\item \texttt{evidence/opdp89-sharp-revalidation-20261002/PREFLIGHT.json} (existing retained artifact; SHA-256 \texttt{9309cdfc7e72e9ce3659cb7e9b8ce4365031706f76614277b51e6dd7d73ac284}).
\end{itemize}
\begin{itemize}
\item \texttt{evidence/opdp89-sharp-revalidation-20261002/POSTFLIGHT.json} (existing retained artifact; SHA-256 \texttt{35158f6261722db29ce897446ac4c34511d861086bd44504ca442609f1ed7db6}).
\end{itemize}
\begin{itemize}
\item \texttt{evidence/opdp89-sharp-revalidation-20261002/COMMANDS.sh} (existing retained artifact; SHA-256 \texttt{2c0d67b15679c0ebde012f739c33abe6754d1f29b1b6d530f8ffea7243ec5658}).
\end{itemize}
\end{document}
